This commentary was written by an AI model (ChatGPT, Pro mode) from page-by-page audits of the published edition of Federigo Enriques's \emph{Le superficie algebriche} (1949) and of the scans of its print. Some of its sentences were revised for the reader by another AI model (Claude Opus 5.5, Anthropic). Part 1 brings the Introduzione and Capitoli I--VI into one continuous account; Part 2 covers Capitoli VII--XI. The published text and its scans control statements and quotations, and the modern arguments below consolidate the necessary repairs and examples.

\chapter*{Introduzione}
\markboth{Introduzione}{Introduzione}
\label{enriques-introduzione}

\section*{Editorial setting and conventions}

The purpose of this commentary is to follow Enriques's mathematical development, section by section, while making explicit which arguments he completes, which he sketches, which he states without proof, and which he assumes from earlier work. These descriptions concern the argument at the cited place: a true theorem may have a defective local proof, and a modern replacement does not identify a historical repairer unless a reference establishes that attribution. All page references are to the book's original pagination.

Unless a singular local model is explicitly being discussed, \(X\) denotes a smooth, connected, projective surface over \(\mathbb C\), and \(Y\) a possibly singular projective image. Purely algebraic arguments often work over any algebraically closed field of characteristic zero; this does not authorize a positive-characteristic extension of separability, generic smoothness, Picard smoothness, vanishing, or classification results. Normalization separates branches and is finite; resolution may further require a birational modification. For the ordinary projection models of the Introduzione the normalization is already smooth. Enriques's instruction to separate the sheets over the double curve (p.~8) governs subsequent intrinsic statements.

We use \(D\sim E\) for linear equivalence and \(D\equiv E\) only for numerical equivalence. Enriques often writes \(\equiv\) for the former. The canonical divisor is always \(K_X\); the letter \(K\) in Chapter I's sums and auxiliary constructions remains an arbitrary divisor. Set \[
q(X)=h^1(X,\mathcal O_X),\qquad
p_g(X)=h^0(X,\mathcal O_X(K_X)),
\] \[
\chi(\mathcal O_X)=1-q(X)+p_g(X),
\] \[
p_a(X)=\chi(\mathcal O_X)-1=p_g(X)-q(X),\qquad
P_m(X)=h^0(X,\mathcal O_X(mK_X)).
\] For a divisor class, the same letter \(p_a\) has the different, explicitly indicated argument \[
p_a(D)=1+\frac{D^2+K_X\cdot D}{2}.
\] For an effective Cartier divisor this is \(1-\chi(\mathcal O_D)\), including reducible and nonreduced divisors. It is not generally the genus of the normalization. On a specified smooth model the relative linear genus is \(p^{(1)}=K_X^2+1\); the invariant version uses the smooth minimal model and is distinguished as \(\bar p^{(1)}\) when both occur, as in Chapter II, Section 6. Every formula involving this number must retain its model convention. A ``pure'' canonical system means the canonical system on the minimal model, pulled back when appropriate, and does not mean that every fixed canonical component has been discarded.

A linear system is \(\mathbb P(V)\), for a nonzero subspace \(V\subset H^0(X,L)\); the complete system is \(|L|=\mathbb P H^0(X,L)\), with projective dimension \(-1\) if it is empty. Completeness is independent of generation at every point. With an assigned base cluster resolved by \(\pi:X'\to X\), its moving class is \[
M=\pi^*D-\sum_i m_iE_i^*,
\] where \(E_i^*\) are total exceptional transforms; fixed divisorial parts are removed separately. The effective degree is \(M^2\); the virtual degree is the square of the specified divisor class on its specified model. Neither is the projective order \(H\cdot D\) of a curve. Effective genus refers to the normalized general integral member; virtual genus refers to \(p_a(D)\), with the same assigned-point convention.

Enriques's \(S_r\) means \(\mathbb P^r\), and \(\infty^r\) means parameter dimension \(r\). An ``involution'' usually means the finite generic equivalence relation defined by a map and need not be a group action. ``Simple'' means generically birational. A ``fundamental curve'' is contracted by the specified system. ``Normal singularities'' are ordinary projection singularities, not normal local rings. His ``normal surface'' on pp.~25--26 refers to completeness of hyperplane sections on the chosen smooth or normalized model. The qualifications to these translations are developed in their respective sections.

Four distinctions will recur: the rational map versus the polynomials representing it; strict versus total transform; a complete intrinsic series versus the subseries actually cut by an auxiliary system; and a parameter count versus a proved genericity or transversality statement. ``General'' means in a suitable nonempty Zariski-open set of the indicated parameter space.

\subsection*{The prefaces and the irregularity problem}

The unnumbered title and copyright leaves contain no mathematical theorem. Castelnuovo's preface occupies the leaf corresponding to p.~V through p.~VIII. His principal mathematical assertion is that an irregular surface has continuous algebraic families of curves not contained in one linear system, and that this is characteristic of irregularity. He explicitly distinguishes the truth of the result from the inadequacy of Enriques's 1904 algebraic proof. He mentions Poincaré's transcendental work and points the reader to Chapter IX, while saying that he does not see how the route proposed there can yet be made rigorous. This is \textbf{an explicit acknowledgment of an unresolved proof problem in the edition's historical setting}, not a hidden error discovered in Chapter I.

Castelnuovo also reports that the manuscript was complete but the final proof corrections were unfinished when Enriques died; Pompilj and Franchetta undertook the editorial work, whose difficult parts could have required years rather than months (pp.~V--VI). The prefaces trace the exposition to the Enriques--Campedelli lecture volumes of 1932 and 1934 and to revision during racial persecution and war. Enriques explicitly records uncertainty about the publication of some Madrid submissions and incomplete access to the recent literature (p.~X). These statements locate the witnesses historically; they do not certify or discredit an unrelated proof. Castelnuovo recommends the algebraic and transcendental literature of Picard--Simart, Lefschetz, Zariski, Hodge, Poincaré, and Severi (p.~VII), without supplying a complete bibliography there.

The additional caution about Chapters VII, VIII, and X concerns detailed classifications and existence questions. It supplies no particular theorem from those chapters for us to adjudicate. Castelnuovo's remarks about intuition, formalism, and the future of mathematics (pp.~VII--VIII) are historiographical judgments, not mathematical claims with proofs. His preface is dated January 1949. Enriques's own preface (pp.~IX--X) says the text attained its definitive form by April 1942 and was retained with limited bibliographic additions; it is dated September 1945. These dates explain why ``1949'' must not be read as a claim that every result then known had been incorporated (Enriques, 1949, pp.~V--X).

Enriques identifies the four volumes of Enriques--Chisini and Conforto's \emph{Le superficie razionali} as prerequisites. The opening chapters cite them repeatedly for curve theory, infinitely near points, Lüroth's theorem, degeneration, and rational-surface examples. Their citations are preserved at the relevant places below. They are not silently replaced by an assertion that these opening chapters proves all of those results.

The contents, pp.~XI--XV, are also informative about scope. Canonical systems begin in Chapter II, numerical genus and surface Riemann--Roch in Chapter IV, and the principal irregularity discussion in Chapter IX. Consequently the cohomological apparatus used here is \textbf{modern explanatory context} for Chapter I and its forward references; it is not falsely attributed to these opening pages. The first-order Picard calculation is also concrete: transition functions of the form \(1+\varepsilon a_{ij}\) identify its tangent space with \(H^1(\mathcal O_X)\). In positive characteristic its tangent dimension need not equal the dimension of its reduced variety (Liedtke, 2008, Introduction and Section 2). The contents' expression \(p^{(1)}=6,5\) lists the two integral cases 6 and 5; it is not a fractional genus 6.5. The printed spelling ``Disgressione'' in the Chapter IX contents is a textual issue, not a mathematical one.

\subsection*{A precise form of the preface's assertion}

For a smooth projective complex surface, \(q(X)>0\) if and only if there exists a flat family of effective Cartier divisors over a connected reduced algebraic parameter scheme whose members are not all linearly equivalent. A stronger useful form is this: for a sufficiently positive divisor \(D\), the algebraically equivalent divisors near \(\lvert D\rvert\) form a smooth family of dimension

\[
\dim\lvert D\rvert+q(X).
\]

The proof uses the Picard variety, which is a smooth projective group of dimension \(q(X)\) in characteristic zero. Choose a Poincaré bundle \(\mathcal P\) on \(X\times\operatorname{Pic}^0(X)\). After tensoring with a sufficiently high power \(A\) of an ample bundle on \(X\), relative Serre vanishing gives vanishing of higher cohomology for all twists \(A\otimes\alpha\), uniformly in \(\alpha\in\operatorname{Pic}^0(X)\). Global generation, and then very ampleness, can be achieved uniformly as well. The pushforward of the twisted Poincaré bundle is therefore a vector bundle of constant rank. Its projective bundle of nonzero sections parametrizes effective divisors, with fiber \(\lvert A\otimes\alpha\rvert\). Distinct \(\alpha\)'s give distinct linear-equivalence classes. If \(q=0\), any map from a connected reduced parameter space to the discrete Picard scheme is constant, so a connected family stays inside one complete linear system. This proves both directions, including the existence assertion omitted by a mere tangent-space count.

The distinction between tangent directions and actual families matters. For a smooth divisor \(C\), the sequence

\[
0\longrightarrow\mathcal O_X\longrightarrow\mathcal O_X(C)
\longrightarrow N_{C/X}\longrightarrow0
\]

gives

\[
0\longrightarrow H^0(\mathcal O_X(C))/\mathbb C
\longrightarrow H^0(N_{C/X})
\longrightarrow H^1(\mathcal O_X)
\longrightarrow H^1(\mathcal O_X(C)).
\]

Thus the deficiency of the characteristic series of the \textbf{complete linear system} is

\[
\delta(C)=\dim\ker\{H^1(\mathcal O_X)\longrightarrow H^1(\mathcal O_X(C))\}\le q.
\]

For \(H^1(\mathcal O_X(C))=0\), it equals \(q\). But \(H^0(N_{C/X})\) is the tangent space to the divisor scheme: without a smoothness or unobstructedness argument, its elements need not integrate into actual nearby curves. That is the precise modern reason one cannot infer the full completeness theorem simply from infinitesimal intersections.

For attribution, Grothendieck's 1962 Picard-scheme exposés provide the algebraic framework; In Exposé 236, Proposition 2.10(iii) identifies the tangent dimension and establishes smoothness in characteristic zero. The earlier completeness discussion is in \textbf{Exposé 221, §§5.4--5.6}, as identified by Kleiman (2005, Remark 5.18). Mumford's \emph{Lectures on curves on an algebraic surface} (1966, pp.~157--159) treats the semiregular obstruction argument, and Kleiman's ``Completeness of the characteristic system'' (1973, pp.~304--310) gives a further algebraic treatment. Kleiman's 2005 exposition, Theorem 5.11, Corollary 5.14, and Remark 5.18, verifies the precise formulation and these attributions. These are replacements for identified logical steps, not a claim that each author corrected this particular sentence of the 1949 book. In positive characteristic the Picard scheme can be nonreduced, so its dimension and its tangent dimension \(h^1(\mathcal O_X)\) must be separated. The arguments here remain in characteristic zero.

The two prefaces are therefore candid about the work's scope. Over \(\mathbb C\) the irregularity assertion is true; the historical problem is the adequacy of the proposed proof. It is not a presently open theorem. The positive family just constructed also separates actual divisors from arbitrary first-order normal sections, an issue that returns in Chapter I, Section 13.

\section{Richiamo di nozioni elementari}
\label{enriques-intro-1}

\subsection{Equations, smoothness, and multiplicity}

An integral surface in affine three-space is a hypersurface \(f(x,y,z)=0\); its projective closure is defined by a homogeneous equation \(F(x_0,x_1,x_2,x_3)=0\). Irreducibility is assumed. Enriques also assumes the projective notions of degree, class, and polar surfaces, citing Enriques--Chisini, Book III, volume II (1918). In modern terms, degree is the intersection number with a general line. The class, in the nondefective situation, is the degree of the dual surface. A first polar has equation \(\sum a_i F_{x_i}=0\). Dual-defective cases require separate conventions for ``class''; no such classification is attempted here.

The line-intersection definition of a simple point must use a \textbf{general} line through the point, with intersection multiplicities counted. Tangent lines and lines contained in the surface are excluded. If the local expansion is

\[
f=f_m+f_{m+1}+\cdots,\qquad f_m\ne0,
\]

where \(f_i\) is homogeneous of degree \(i\) in local coordinates, then \(m=\operatorname{mult}_pX\). On a general line through \(p\), the local intersection multiplicity is \(m\); a degree-\(n\) hypersurface therefore has \(n-m\) residual intersections. Smoothness is equivalent to \(m=1\), or equivalently \(df(p)\ne0\). These are correct elementary facts; the book does not prove Bézout or the Jacobian criterion anew.

At a smooth point, choose coordinates so that \(f_z(p)\ne0\). The implicit function theorem gives \(z=g(x,y)\). Its Taylor polynomial through any prescribed order represents the corresponding jet of the surface. After placing the tangent plane at \(z=0\), the quadratic approximation is \(z=g_2(x,y)\), Enriques's osculating paraboloid. Algebraically the completed local ring is \(\widehat{\mathcal O}_{X,p}\simeq\mathbb C[[u,v]]\). This is a statement about local analytic or formal coordinates, not a claim that every smooth surface is rational or has a Zariski-open subset isomorphic to a plane. Enriques's finite-order formulation already avoids that erroneous global inference.

Consequently, a curve germ on a smooth surface is a plane-curve germ in suitable local coordinates. Its embedded singularity can be described by a succession of blowups at proper and infinitely near points. Enriques explicitly refers to Enriques--Chisini, Book IV, volume II, for the omitted theory. An infinitely near point is a point on an exceptional divisor of an iterated blowup; an ``infinitesimal curve'' is not an extra reduced curve lying in the original surface.

\subsection{Ordinary and special multiple points}

At an ordinary point of an \(r\)-fold curve, the surface has \(r\) smooth local branches with distinct tangent planes. This describes an ordinary local configuration, rather than every multiplicity-\(r\) curve. Enriques immediately records the exceptions: branches may be tangent or merge into a higher-order branch, as along a cuspidal double curve. At an isolated multiplicity-\(r\) point, the projectivized tangent cone is \(V(f_r)\subset\mathbb P^2\). A line direction lies on that cone precisely when its intersection multiplicity is at least \(r+1\). Irreducibility or smoothness of the cone is an additional ordinary-case hypothesis. Reducible cones and cones supported on a repeated plane are explicitly allowed in the source.

One sentence requires a more substantial qualification. Page 2 says that a uniplanar double point has three distinct infinitely near double points. A doubled tangent plane alone does not imply this. Write, after a change of coordinates,

\[
f=z^2+g_3(x,y)+\text{terms involving }z\text{ or of higher order}.
\]

In the blowup chart \(x=t,\ y=tv,\ z=tw\), the strict transform begins

\[
w^2+t\,g_3(1,v)+\cdots=0.
\]

On the exceptional line, singular points occur at the zeros of the binary cubic \(g_3\). If that cubic has three distinct zeros, each yields an ordinary double point: locally the leading equation is \(w^2+c\,t(v-v_0)=0\), with \(c\ne0\). This proves the three-point assertion for the ordinary uniplanar case. Without that condition, the conclusion fails. For example, \(z^2+x^3+y^4=0\) has a doubled tangent plane, but its first blowup has only one singular point over the origin: in the \(y\)-chart its equation is \(w^2+t u^3+t^2=0\), singular on the exceptional line only at \(u=0\). This calculation supplies the qualification without a claim of historical priority. The test germ \(z^2+xy(x-y)=0\) gives the three simple roots explicitly. Chapter III, pp.~83--84, supplies the ordinary qualification that p.~2 leaves implicit. The scan confirms the unqualified wording; it is not a transcription error.

The tacnode description is likewise a statement about general plane sections. A useful local test case is \(z^2+x^4+y^4=0\). A general plane section through the origin has a double point followed by a double infinitely near point. After blowing up, the strict transform is singular along the exceptional line: \(w^2+t^2(1+v^4)=0\) in one chart. This explains Enriques's ``infinitesimal double line.'' Special plane sections can degenerate further and should not be included in the generic formulation.

\section{Trasformazioni razionali}
\label{enriques-intro-2}

\subsection{The rational-map formulas and generic degree}

The displayed formulas

\[
(X,Y,Z)=\left(\frac{\varphi_1}{\varphi_0},
\frac{\varphi_2}{\varphi_0},\frac{\varphi_3}{\varphi_0}\right)
\]

describe a rational map. Projectively, the \(\varphi_i\) are homogeneous forms of a common degree, after homogenization and cancellation of common divisorial factors. ``Nondegenerate'' must mean dominant when Enriques counts a finite number of inverse images of a general point of the target three-space. The induced extension of function fields is then finite; in characteristic zero its degree is the number of distinct points of a general fiber. His ``involution of order \(n\)'' means the partition by these general fibers, not necessarily the orbit partition of a group of order \(n\). No group action is assumed or proved.

A surface chosen generally relative to this finite correspondence is not preserved by a nontrivial subcorrespondence. Its map to the image is therefore generically one-to-one. A clean formulation is that the off-diagonal relation in the fiber product meets \(f\times f\) in dimension at most one; hence it cannot identify the generic point of \(f\) with another point of \(f\). If the surface is instead saturated by all or part of the correspondence, the restriction can have degree \(m>1\). The image as a reduced variety remains a surface; the pushforward cycle has multiplicity \(m\). That is the modern meaning of ``a multiple surface.'' The page-3 phrase ``la trasformata di \(F\)'' should read ``di \(f\)'' in this sentence. The capital \(F\) is printed in the scan, so this is a source slip, not a transcription error.

Enriques correctly distinguishes a birational map between surfaces from a Cremona transformation of their ambient three-spaces. A birational surface map need not extend to a birational map of the ambient spaces.

\subsection{Base loci, blowups, and branch separation}

The common zero set of the homogeneous coordinate sections is the base locus of the chosen presentation. One must first remove common factors; a common zero of a nonminimal presentation can be removable. At a genuine base point the map is resolved by a modification of the source. Over a smooth surface point, a point blowup introduces \(\mathbb P^1\); a succession of blowups introduces a configuration of rational curves. Its image on another birational model may be one curve, several curves, or a contracted configuration. Thus the introductory singular ``exceptional curve'' is preliminary terminology, refined later in Chapter I.

For an ordinary multiple curve with distinct branches, normalization separates the generic branches. Over a general multiplicity-\(r\) point, there are \(r\) distinct points of the normalization. Blowing up an isolated ordinary hypersurface point exposes its projectivized tangent cone as the exceptional curve; if this cone is smooth, the first blowup resolves the singularity along it. These assertions explain the source's ``generally,'' but are not true for an arbitrary collection of forms merely because they vanish at the center. The forms must generate an appropriate ideal and separate the directions or branches in question.

The phrase on p.~4 about reducing ``from \(r\) to \(r-1\) the neighborhood'' does not define a decreasing invariant. A blowup makes first-neighbor points proper points of the new surface; it need not lower the multiplicity. The preceding \(z^2+x^3+y^4\) example still has a double point after its first blowup. The charitable reading is that one level of the infinitely near structure has been exposed. It cannot be used as a termination proof for resolution. A second test germ, \(z^2+x^5+y^5=0\), likewise remains double after its first blowup: in the \(x\)-chart its strict transform is \(w^2+t^3(1+v^5)=0\).

\subsection{The five equations and what their count establishes}

For two points \(p=(x,y,z)\), \(q=(x',y',z')\) in the chart \(\varphi_0\ne0\), equality of their images is expressed by

\[
\varphi_i(p)\varphi_0(q)-\varphi_i(q)\varphi_0(p)=0
\quad(i=1,2,3),\qquad f(p)=f(q)=0.
\]

The transcription of these equations agrees with the scan. They are five equations in six affine variables. The diagonal \(p=q\) is a two-dimensional component, and the expected dimension of the residual off-diagonal locus is one. Enriques also gives the ambient version in a footnote: the off-diagonal finite-correspondence relation is three-dimensional before intersecting the two copies of the surface.

This is an expected-dimension argument, not a theorem that every rational map creates a double curve. It requires proper intersection, exclusion of the diagonal and base-locus components, and generic finiteness. For an ambient birational map there may be no off-diagonal relation at all on its isomorphism locus. In scheme language, one works with the fiber product on the domain of definition, removes the diagonal there, and then takes closure; cross-multiplication without this restriction can add spurious components. General projection theorems justify the expected count in the setting subsequently used by Enriques.

The last observations are possibilities: special maps may create higher multiple curves or contract curves to isolated singularities. Contracting a curve does not invariably produce a singular point; the exceptional curve of a point blowup contracts to a smooth point, as Chapter I will emphasize. Reading p.~5 as a universal claim to the contrary would create an artificial contradiction.

\section[Superficie negli iperspazi]{Superficie negli iperspazi: scioglimento delle singolarità}
\label{enriques-intro-3}

\subsection{Higher-dimensional embeddings}

With \(r+1\) sections one obtains a rational map to \(\mathbb P^r\). The image is the closure of its graph-theoretic image, equivalently an irreducible component of a suitable algebraic intersection after elimination. Enriques's count predicts that a surface projected generally to \(\mathbb P^4\) has only finitely many identifications, whereas projection to \(\mathbb P^5\) has none. Point separation alone, however, does not establish smoothness: tangent directions must also be separated.

Here is the exact embedding argument. Embed a smooth projective surface \(X\) in \(\mathbb P^N\). The closure of the union of its secant lines, including limiting tangent lines, has dimension at most five: pairs of points give at most four parameters and the point on the line gives one more. A general center \(\Lambda\simeq\mathbb P^{N-6}\) avoids this set because

\[
\dim\Lambda+\dim\operatorname{Sec}(X)\le N-1.
\]

Projection from \(\Lambda\) separates every pair of points and every nonzero tangent direction. Equivalently, it separates length-two subschemes. The resulting proper morphism is a closed immersion into \(\mathbb P^5\). This proves the five-dimensional embedding assertion once a smooth projective model exists. It also explains why the tangent variety, not just the set of distinct-point secants, must be included.

\subsection{A genuine qualification to Levi's graph remark}

Page 6 says that the surface of pairs \((x,y,z,X,Y,Z)\) can have only singularities corresponding to singularities common to the two surfaces. Used to separate branches identified by one projection while the other projection is a local isomorphism, this is correct. As a statement about the closure of an arbitrary birational graph it is false.

An explicit counterexample is the birational map between two affine planes

\[
X=\frac{x^2}{y},\qquad Y=\frac{x^3}{y^2},
\qquad
x=\frac{X^2}{Y},\qquad y=\frac{X^3}{Y^2}.
\]

Its graph closure has equations

\[
yX=x^2,\qquad yY=xX,\qquad xY=X^2.
\]

This is the affine cone over the twisted cubic, parametrized by\[ (y,x,X,Y)=(s^3,s^2t,st^2,t^3). \] It is a two-dimensional irreducible variety whose three defining equations have no linear parts at the origin. The Zariski tangent space there has dimension four, so the origin is singular. Both endpoint planes are smooth. If desired, insert \(z=Z=0\) to place the example in exactly the six-coordinate ambient space used by Enriques. Normalizing the graph does not automatically remove this particular defect: the twisted-cubic cone is already normal.

The repair is precise: a graph is smooth at points where one of its projections is a local isomorphism to a smooth surface. Over simultaneous indeterminacy, one must resolve the graph, or principalize the base ideal on a smooth model. No claim of historical priority is made for this counterexample and its calculation. The source's resolution theorem does not depend on accepting the graph assertion in this unqualified form, because the theorem is explicitly imported from independent references.

A complementary graph example has a different, equally valid normal singularity. Put \[
(x,y)\dashrightarrow(u,v)=\left(\frac{x^2}{y},\frac{x^3}{y}\right),
\qquad (u,v)\dashrightarrow(x,y)=\left(\frac vu,\frac{v^2}{u^3}\right).
\] Its graph closure is \(v=xu,\ uy=x^2\), a normal quadratic cone with a singular vertex. Both affine endpoint planes remain smooth. Thus both examples distinguish graph closure, normalization, and a common smooth resolution; they should not be mistaken for conflicting formulas for a single map.

\subsection{Resolution: what is cited and what a modern proof uses}

Enriques expressly declines to prove resolution. His references on p.~6 are B. Levi (1897), O. Chisini (1921), G. Albanese (1924), R. J. Walker (1935), and O. Zariski (1939), with an additional pointer to \emph{Annals of Mathematics} (1941). He credits Levi with the first rigorous treatment. This commentary does not independently certify that historical priority or the completeness of every early cited proof. Walker's and Zariski's papers are already part of Enriques's reference list; they should not be presented as later repairs newly unknown to him.

The modern input is the surface-resolution theorem: an integral projective surface over \(\mathbb C\) has a projective birational modification \(\pi:X\to Y\) with \(X\) smooth. Lipman's theorem applies more generally to excellent surfaces. A precise, checked modern statement and proof is The Stacks Project, Theorem 54.14.5, Tag 0BGP; its local construction and termination are developed in §§54.7--54.14. Normalize first, reduce to finitely many normal surface singularities, resolve those by normalized blowups, and glue the local modifications. Excellence ensures the finiteness and formal-normality hypotheses. This supplies a genuine termination theorem rather than an assumed decrease of multiplicity. Composing a projective resolution with the preceding secant-avoidance embedding proves the exact theorem on p.~6 (Lipman, 1978; The Stacks Project Authors, 2026). The printed theorem reads: ``Una superficie \(f\) dotata di singolarità qualsiansi si può sempre trasformare birazionalmente in un'altra priva di singolarità, appartenente ad uno spazio a cinque dimensioni.''\footnote{The print reads ``qualsiansi.'' Enlarging the scan confirms that spelling in the italic theorem as well as in the next paragraph. The edition prints ``qualsiasi,'' and its apparatus asserts that the \(n\) is absent. For an exact quotation the scan controls; the spelling has no mathematical effect.}

This is a proof of Enriques's formulation \textbf{using} a stated modern resolution theorem. It is not a purported new derivation of Lipman's entire 57-page argument. The nontrivial vanishing, boundedness, and local termination results remain explicitly cited dependencies.

\section{Superficie dotate di singolarità normali}
\label{enriques-intro-4}

\subsection{Correct modern statement and terminology}

For a smooth nondegenerate projective surface in characteristic zero, a general linear projection to \(\mathbb P^3\) is finite and birational onto an image with only ordinary singularities. The local analytic models at singular points are

\[
xy=0,\qquad xyz=0,\qquad x^2-y^2z=0.
\]

They are respectively an ordinary point of the double curve, an ordinary triple point, and a pinch point. The reduced double curve is smooth at a pinch point; ``cuspidal point'' on p.~8 describes the surface's folding behavior, not a cusp of that reduced curve. Its only singular points in the strict ordinary model are ordinary triple points. The normalization is smooth. These statements and the stronger formulation for a surface in \(\mathbb P^4\) with only ordinary improper double points are recorded in Ciliberto and Flamini (2011, Definition 2.1, Remark 2.2, Theorem 2.5).

``Normal singularities'' in Enriques's terminology must not be confused with normality of a variety. A hypersurface with a double curve is not normal. An ``improper double point'' in \(\mathbb P^4\) is two smooth surface branches with independent tangent planes, rather than an intrinsic normal surface node such as \(xy-z^2=0\).

\subsection{What the dimension counts prove}

For a nondegenerate surface in \(\mathbb P^5\), the secant-line family has dimension four. Its point-line incidence has dimension five. Hence a general point lies on finitely many secants if the secant variety fills \(\mathbb P^5\), and on none if it does not. ``Finite'' must allow zero: the Veronese surface is the important exceptional example. Avoiding the tangent variety makes the projection to \(\mathbb P^4\) an immersion, with only finitely many pairs identified in general. The two tangent planes at each such pair must also meet transversely; that assertion uses genericity beyond a bare count.

For projection to \(\mathbb P^3\) from a line \(L\subset\mathbb P^5\), secant lines meeting \(L\) form a one-dimensional family: meeting a fixed line has codimension three in \(G(1,5)\). Trisecant planes containing \(L\) are expected to be finite: planes spanned by three surface points form a six-dimensional family, and containment of \(L\) imposes six conditions in \(G(2,5)\). These counts explain the double curve and isolated triple points.

Enriques invokes the trisecant lemma for space curves, citing Enriques--Chisini, Book III, chapter IV, §43, p.~289. His reduction is useful: if every general secant of the surface were trisecant, a suitable nonplanar hyperplane-section curve would have the same impossible property. Likewise, if all planes through three general surface points contained a fourth, inner projection of an appropriate hyperplane-section curve would produce a nonplanar curve whose general chord was trisecant. The bad family of quadrisecant planes is consequently proper in the six-dimensional trisecant-plane family; adjoining a line in such a plane gives dimension at most seven, smaller than \(\dim G(1,5)=8\), so a general center line avoids it.

These arguments control numbers of distinct points in fibers. They do not, on their own, classify nonreduced fibers, prove transverse intersections of all branches, or show that ramification points are simple crosscaps. Those are the missing local and incidence statements behind the sentence ``rigorously demonstrated'' on p.~8. This is a precise limitation of the displayed proof, not a counterexample to the ordinary-projection theorem.

\subsection{A complete existence proof after a sufficiently ample re-embedding}

The existential theorem on p.~8 can be justified without assuming arbitrary embeddings have all the required jet properties. The following is a modern reconstruction. Its hypotheses are deliberately sufficient rather than optimal.

Start with a smooth projective resolution \(X\). Choose a sufficiently high power \(L\) of a very ample line bundle so that sections separate all the finite jets used below, simultaneously at up to four distinct points. Such an \(L\) exists: the relevant zero-dimensional schemes have bounded length, their Hilbert schemes are projective, and relative Serre vanishing gives uniform surjectivity of the evaluation maps. A general quadruple of sections defines a morphism \(f:X\to\mathbb P^3\). Four simultaneous zeros impose four independent conditions on a two-dimensional source, so the general quadruple has no base point.

Surjectivity of the jet evaluations means that the corresponding universal evaluation maps are smooth. Thus ordinary codimension calculations in jet space apply to a general choice of sections, by generic smoothness and the fiber-dimension theorem. There are only finitely many kinds of bad strata needed here:

\begin{enumerate}
\item
  For \(k\) distinct source points to have a common image imposes \(3(k-1)\) conditions on a configuration space of dimension \(2k\). The expected dimension is \(3-k\). Double fibers therefore form a curve, triple fibers are isolated, and fibers with four distinct points do not occur.
\item
  At one point, the derivative is a \(3\times2\) matrix. Rank at most one has codimension two, so ramification points are isolated. Rank zero has codimension six and does not occur.
\item
  At a double fiber of immersion points, equality of the two tangent planes is an additional codimension-two condition. It therefore does not occur in the one-dimensional double-fiber locus. At an isolated triple fiber, dependence of the three normal directions is an additional condition and is absent generally.
\item
  A rank-one point sharing its image with a second point imposes three coincidence conditions and two rank conditions on a pair space of dimension four. It is absent. The failure of the crosscap condition on the two-jet at a rank-one point is one additional condition and is absent as well.
\end{enumerate}

To identify the remaining ramification germ, choose coordinates with

\[
f(u,v)=(u,g(u,v),h(u,v)),\qquad dg(0)=dh(0)=0.
\]

The generic two-jet condition is that \((g_{vv},h_{vv})\) and \((g_{uv},h_{uv})\) are independent. A target linear change makes \(g_{vv}\ne0\). The analytic Morse lemma with parameter \(u\) changes \(g\) to \(v^2\), after removing a function of \(u\). Write the remaining analytic function uniquely as \(A(u,v^2)+vB(u,v^2)\). Remove \(A\) by a target coordinate change. Independence gives \(B_u(0)\ne0\), so \(x=B(u,v^2)\) is a permissible new coordinate. The map becomes

\[
(x,v)\longmapsto(x,v^2,xv),
\]

whose image has equation \(z^2=x^2y\), the pinch-point model. Two or three mutually transverse immersion branches give \(xy=0\) or \(xyz=0\) by the analytic implicit function theorem.

Finally, the map is finite: a contracted curve would have degree zero against \(f^*\mathcal O_{\mathbb P^3}(1)=L\), contradicting ampleness. A proper morphism with finite fibers is finite. The one-dimensional coincidence locus also shows that the map is generically one-to-one, hence birational in characteristic zero. This proves the existence of the required ordinary model. The stronger theorem for every fixed smooth embedding is a genuine additional generic-projection theorem, not a consequence of assuming that every very ample bundle separates arbitrary higher jets.

\subsection{The apparent nodes and the polar calculation}

Page 8 permits ``incroci,'' or nodes, of the double curve. Pages 10--11 then seek to rule out precisely the associated tacnodal configurations for a general center. The final modern ordinary model does not include these extra nodes. It is more accurate to describe p.~8 as an overbroad interim list, subsequently narrowed, than to claim that Enriques's final theorem predicts such nodes.

The pinch-point assertion about polars can be checked directly. For

\[
F=x^2-y^2z,
\quad
\nu(u,v)=(uv,u,v^2),
\]

the polar \(aF_x+bF_y+cF_z\) pulls back to

\[
u(2av-2bv^2-cu).
\]

The factor \(u\) is the fixed conductor/double-curve component. The residual polar curves have equations \(2av-2bv^2-cu=0\): they all pass through \((0,0)\), and a general member has nonzero linear term. The pinch point is therefore a simple base point of the residual polar system, exactly as claimed. At a general double point, normalization instead gives two distinct smooth points, one for each branch. This justifies the bookkeeping convention used in the rest of the book.

\section{Nota: irreducibility of the double curve}
\label{enriques-intro-5}

\subsection{The theorem and its dependencies}

The final statement is: if \(X\subset\mathbb P^5\) is smooth, irreducible, and nondegenerate, then the reduced double curve of its general projection to \(\mathbb P^3\) is irreducible, except when \(X=v_2(\mathbb P^2)\) is the Veronese surface. That exceptional image is the Steiner quartic, with three concurrent double lines. The intermediate statement for a surface in \(\mathbb P^4\) with at most ordinary improper double points is also a classical Franchetta theorem. Enriques cites Franchetta (1941), with a 1947 addendum. Ciliberto and Flamini (2011, Theorem 2.6) state the modern theorem, while Lopez and Ran (2003; checked preprint pp.~2 and 7) explain the simpler incidence argument in the \(\mathbb P^5\) case and credit it to Mumford.

The existential refinement at the start of §5 remains true even for a rational surface: one may change its embedding. The exception is a particular embedding by plane conics, not an obstruction attached to the birational class of \(\mathbb P^2\). A sufficiently ample re-embedding avoids that exceptional situation.

\subsection{A rigorous replacement for the tangent-plane lemma}

The lemma on p.~9 says that a three-dimensional family of pairs \((p,q)\) with tangent planes meeting in a line cannot have its joining secants fill \(\mathbb P^4\). Its conclusion has a short differential proof that avoids the compressed infinitesimal integration in the text.

Let \(T\) be an irreducible family of such pairs, with distinct smooth points general on it. On local affine cones choose lifts \(\widehat p,\widehat q\). The secant evaluation is

\[
\Psi:T\times\mathbb P^1\dashrightarrow\mathbb P^4,
\qquad(p,q,[a:b])\mapsto[a\widehat p+b\widehat q].
\]

Every derivative of \(\Psi\) lies in the projective span \(\langle T_pX,T_qX\rangle\). If the tangent planes meet in a line, this span is a \(\mathbb P^3\). Thus the generic differential rank is at most three. A dominant map to \(\mathbb P^4\) in characteristic zero has generic differential rank four. The alleged filling is impossible. This proves exactly the needed lemma, including the case where the global surface has isolated singularities but the selected points are smooth. It does not rely on classifying all developable plane families.

Enriques's proof fixes \(p\), considers a curve of possible \(q\)'s, and argues that the tangent planes along that curve lie in a fixed hyperplane. It then obtains planar contact curves, a ruling or a developable family of planes, and non-filling. The potentially delicate step is treating an infinitesimal containment as a fixed hyperplane along the entire curve. The differential-rank proof supplies the missing justification for the desired conclusion without claiming that every subsidiary classification sentence has been proved by the brief printed argument.

On the ordinary locus, however, the fixed-hyperplane step itself has a valid short justification. With \(\pi=\mathbb P(\Pi)=T_pX\), the map \[
q\longmapsto\langle\pi,q\rangle\in\mathbb P(V/\Pi)\simeq\mathbb P^1
\] has zero derivative along a contact component outside \(\pi\) when \(\langle\pi,T_qX\rangle\) is a hyperplane. Characteristic zero makes the map constant.\footnote{The preceding paragraph calls this step potentially delicate; the differential explanation just given supports it. Inspection of the printed neighboring-tangent-plane argument on p. 9 supports the qualified reconciliation: the step is justified off the fixed tangent plane, while components lying in that plane and the subsequent developable classification still require treatment. The direct secant-rank proof above establishes the needed lemma without those subsidiary cases.}

\subsection{Connectedness and linearity (pp. 9–10)}

The apparent double divisors \(C_O=\sum K\) on the normalization vary with the center \(O\). A rationally parametrized family of divisor classes is constant in the Picard variety: a rational map from projective space to an abelian variety is constant. One proof restricts to general lines; a morphism \(\mathbb P^1\to A\) is constant because invariant one-forms span the cotangent bundle of \(A\), while \(H^0(\mathbb P^1,\Omega^1)=0\). Resolving the rational map gives the same conclusion (Milne, 2008, I, Propositions 3.9--3.10). Thus the divisors lie in one complete linear system. Enriques attributes this principle to earlier work and cites Enriques--Chisini, Book V, chapter I, §9, pp.~78 and 485.

If the general member of the relevant linear system is integral, every member has connected support. One can restrict to a pencil through the proposed member and a general integral member, then use the proper universal family and Stein factorization: its generic fiber is geometrically connected, and the finite birational factor over the normal base \(\mathbb P^1\) is trivial. Enriques calls this the degeneration principle, citing Book V, chapter III, §36, p.~405. If the moving general member is instead reducible, Bertini's second theorem says that it is composed with a pencil, after removing fixed components. Distinct fibers have no moving intersections. The curves obtained by fixing one point \(A\) and varying its two secant partners provide the variable intersections that exclude this case. The distinction between moving intersections and fixed base points is essential to this argument.

This discussion is on a smooth normalization or resolution. Working directly with Cartier divisors on the singular \(\mathbb P^4\) image without controlling their pullbacks would be an additional gap. The source's convention that ordinary improper double points stand for two distinct smooth points supplies the intended interpretation.

\subsection{The local case analysis and where its proof is compressed}

Distinct irreducible components of the reduced double curve can meet at an ordinary triple point, or, in the interim singularity list, at a node of that curve. Enriques treats the second possibility first.

For a node arising from two distinct source points, the corresponding image branches have a common tangent plane. Their two source tangent planes then lie in a common \(\mathbb P^3\), meeting in a line. If this occurred for a general center, the associated bad secants would fill \(\mathbb P^4\), contradicting the preceding lemma. The dimension/dominance statement must be formulated on an incidence variety; saying merely ``there are many pairs'' is insufficient.

For an infinitesimal pair, Enriques interprets the image as tacnodal and the corresponding hyperplane section as having two successive double points. He separately lists the multiplicity distributions \((2,2)\), \((3,1)\), and \((4,0)\), rather than silently ignoring the degenerations. The first two are related to a degenerate Gauss map or developability. The argument does not explicitly establish the necessary dominance of the family of higher-contact hyperplanes. The modern ordinary-projection theorem, or the sufficiently ample jet argument above for the existential result, excludes these bad strata directly.

There is also a small numerical error on p.~11. The scan reads that a developable surface in \(\mathbb P^4\) has only \(\infty^1\), rather than \(\infty^3\), tangent hyperplanes of this kind. For a nonplanar developable surface, its tangent planes form a one-dimensional family and each tangent plane lies in a one-dimensional family of hyperplanes; the generic total is two-dimensional. A cone over a nonplanar curve makes this explicit. Thus \(\infty^2\) is the appropriate count, and it still contradicts the required \(\infty^3\). This is a source-level numerical slip with no effect on the stated contradiction, not an OCR correction.

The last subcase uses the assertion that a space curve whose osculating plane has contact at least four at every general point is planar. In characteristic zero this follows from a Wronskian: for a homogeneous local lift \(v(t)\),

\[
\det(v,v',v'',v''')\equiv0
\]

forces the four coordinate functions to be linearly dependent over the constants. The surface conclusion then uses that a general hyperplane section of a nondegenerate integral surface in \(\mathbb P^4\) spans its hyperplane \(\mathbb P^3\). The link from the postulated family of quadruple hyperplane contacts to this pointwise condition is stated very tersely in the text and still requires the incidence and jet argument. The Wronskian observation by itself is not the missing global proof.

At a triple point, the partitions of the three double-curve branches among irreducible components are \(1+1+1\) and \(1+2\). For three preimages \(p_1,p_2,p_3\), the corresponding apparent double curves have the pair incidences \((p_2,p_3)\), \((p_1,p_3)\), and \((p_1,p_2)\). In the \(1+2\) case one component passes once through \(p_2,p_3\), while the other passes once through those two points and twice through \(p_1\). These multiplicities agree with the normalization of the local model \(xyz=0\): on each branch the conductor has two crossing branches.

The center is then varied along the trisecant line. If an apparent component stays fixed, the book deduces that it is a plane conic. The geometric reasoning can be made explicit: for a general point \(p\) of that component, its partners obtained as the center varies lie in the plane spanned by \(p\) and the center line. If the component is irreducible and the partners vary, it must lie in that plane. Projection from a general point of the line is generically two-to-one on this plane curve, so its degree is two. The next inference, from a two-dimensional conic family to the Veronese surface, invokes a classification result and retains the smooth nondegenerate \(\mathbb P^5\) lifting hypothesis. The bare assertion ``a surface with \(\infty^2\) conics is Veronese'' is false: a cubic scroll in \(\mathbb P^4\) also has a two-dimensional family of conics. That scroll does not satisfy the retained lifting setup and does not contradict Franchetta's theorem.

If the component moves, Enriques asserts that some member passes through the third preimage, producing tangent branches and a forbidden filling family of secants. Properness makes an incidence specialization plausible, but the text does not show that the required specialization occurs at an admissible center, rather than only at one of the points where projection is undefined. Nor does it explicitly establish the dominance needed after specialization. These are genuine missing steps in the printed argument. They are not evidence that the theorem is false, and discovery of these particular textual gaps is not attributed to a later author without a verified source.

\subsection[Modern proof in projective five-space]{Modern proof for the smooth \(\mathbb P^5\) case}

Here is a proof with its deep input made explicit. Let \(X\subset\mathbb P^5\) be smooth and nondegenerate, and first suppose \(\operatorname{Sec}(X)=\mathbb P^5\). Blow up the diagonal in \(X\times X\). The resulting smooth fourfold parametrizes ordered distinct pairs together with limiting tangent directions. Over it take the projective-line bundle of their secant lines. Its total space \(I\) is smooth and irreducible of dimension five, and evaluation gives a proper dominant generically finite map

\[
e:I\longrightarrow\mathbb P^5.
\]

For a general line \(L\), the curve \(e^{-1}(L)\) is irreducible by successive applications of Bertini's irreducibility theorem to the pullback hyperplane system. At each step before the final curve the image has dimension at least two, so the system is not composed with a pencil. Its image in the first factor \(X\) is the apparent double divisor; away from the finite exceptional configurations, each first point has a unique partner. The image of an irreducible curve is irreducible, and its further image in \(\mathbb P^3\) is the double curve. This is the Mumford--Bertini proof described by Lopez and Ran and located by them in Moishezon (1977, p.~115).

The remaining input is Severi's classification of smooth defective surfaces in \(\mathbb P^5\): the only smooth nondegenerate surface with \(\dim\operatorname{Sec}(X)<5\) is the Veronese surface of conics. This is the exact extra theorem needed to identify the exception; it is not proved by the preceding Bertini argument. Ciliberto and Flamini give its statement as Theorem 2.4 and refer to Severi and to Mezzetti--Portelli, Theorem 3. The latter full text was not independently consulted, so that theorem-number pointer is reported through the checked source rather than presented as a directly read proof.

The exception itself is completely concrete. Represent the Veronese surface by rank-one symmetric \(3\times3\) matrices. Its secant variety is the cubic hypersurface of matrices of rank at most two. A general center line meets that cubic in three rank-two points. Each has a conic as entry locus on the Veronese surface, and each entry conic maps two-to-one to a line in the image. The image has degree four and the three double lines are concurrent at an ordinary triple point. In suitable coordinates a representative Steiner quartic is

\[
x^2y^2+y^2z^2+z^2x^2-xyzw=0.
\]

The explicit projection \[
[u:v:w]\longmapsto[uv:uw:vw:u^2+v^2+w^2]
\] verifies this equation by substitution. Its three coordinate lines through \([0:0:0:1]\) are double lines. At the common point the lowest-degree term is \(-xyz\), giving the ordinary triple point. This verifies the asserted geometry without mistaking the three conics upstairs for the three lines downstairs.

For the existential statement, the defective-surface classification can even be avoided: a sufficiently positive embedding separates first jets at two points, so two general affine tangent spaces are independent and the secant variety has dimension five. A general projection to \(\mathbb P^5\) preserves the embedding and makes that secant variety fill \(\mathbb P^5\). The preceding incidence proof then gives an irreducible double curve. This also clarifies why changing the embedding removes the Veronese exception at the level of birational models.

\subsection{Scope of the reconstruction}

The local examples, the two singular graph closures, and the contrast between the conic embedding of \(\mathbb P^2\) and a higher Veronese embedding provide concrete checks on the distinctions used throughout the Introduzione. Two further questions concern the exposition rather than an unsolved theorem: the exact proper-incidence reconstruction of the specialization on pp.~12--13, including inadmissible boundary centers, and which hypotheses Franchetta's original proofs use in the \(\mathbb P^4\) version. The original papers were not inspected in full. Determining which arguments survive for a specified positive-characteristic embedding is a separate proposed investigation, since the differential-rank and Wronskian arguments require separability. No positive-characteristic analogue is asserted here.

\chapter{Sistemi lineari di curve}\label{enriques-I}

The chapter develops linear systems before canonical systems have been defined. We continue to distinguish a fixed divisor, an assigned proper or infinitely near base cluster, and the remaining moving system; the surface image may be nonnormal even though the intrinsic calculations take place on \(X\). The arbitrary auxiliary curve \(K\) is not \(K_X\).

\section{Fasci lineari}
\label{enriques-I-1}

\textbf{Definitions and formulas.} A nonconstant \(t\in\mathbb C(X)\), written \(\varphi/\psi\), gives the rational map \(t:X\dashrightarrow\mathbb P^1\). Its level divisors are the members of the pencil generated by \(\varphi\) and \(\psi\), after both are regarded as sections of one line bundle. The fibers at \(0\) and \(\infty\) are the zero and pole divisors. Enriques states without proof that every parameter value gives a nonempty projective level divisor. Resolve \(t:X\dashrightarrow\mathbb P^1\); the resulting nonconstant proper morphism is surjective. This does not imply that a nonconstant function takes every value on a specified affine open. If a chosen representation has a common effective component \(K\), it can either be retained as a fixed component of every member or removed. A common component in that representation is not an actual simultaneous positive and negative component of \(\operatorname{div}(t)\): principal divisors have already undergone this cancellation.

For \[
\tau=\frac{\alpha t+\beta}{\gamma t+\delta},
\qquad \alpha\delta-\beta\gamma\ne0,
\] the pencil is unchanged and its parameter is changed by an automorphism of \(\mathbb P^1\). Enriques leaves the nonzero determinant implicit in the word ``substitution''; a singular matrix would give a constant and must be excluded. Three distinct chosen members, labelled \(0,\infty,1\), determine the parameter uniquely. Equal zero and pole divisors imply \(t/\tau\in\mathbb C^\times\): on a normal projective model this ratio has no zero or pole and is therefore a global unit. Projectivity and normality are the modern hypotheses behind this familiar argument. Equality only of supports, without multiplicities, would not suffice.

For an irreducible affine equation \(F=0\), equality of two rational expressions means equality in the function field, hence \[
\frac{\varphi}{\psi}=\frac{\varphi_1}{\psi_1}
\iff \varphi\psi_1-\psi\varphi_1\in(F)
\iff \varphi\psi_1-\psi\varphi_1=AF.
\] The equality is asserted where the fractions are defined, then extended as a function-field identity. It is not pointwise evaluation at their indeterminacy points. The displayed denominator on p.~16 is indeed \(\psi_1\), as shown by magnifying the scan; a visually tempting \(\varphi_1\) is not an error in the source. The printed punctuation following the equality is a semicolon.

All members have the same projective degree because they are linearly equivalent. The apparent drop of affine degree for a special value of the parameter means that a component has moved into the hyperplane at infinity. For instance, if \(\deg\varphi<a=\deg\psi\), homogenizing \(\varphi\) to degree \(a\) introduces the factor \(x_4^{a-\deg\varphi}\). This is an actual component with that multiplicity in the homogeneous zero divisor. Enriques both states this phenomenon and explains its repair by homogenization.

The ambient equations \(\varphi-t\psi=0\) give a pencil of surfaces cutting the given pencil on \(F\); replacing the ambient equations modulo \(F\) need not change the induced pencil. Their ambient common curve \(\Gamma\) usually meets \(F\) in finitely many points. After removing fixed divisorial parts, these intersections form the base scheme. A base point has multiplicity \(i\) when a general member has order \(i\) there. Enriques's argument that a genuine moving base point is an essential indeterminacy point is sound: distinct level curves approach it with distinct constant \(t\)-values. Algebraically, after cancelling a common local divisorial factor on the smooth model, the ideal \((\varphi,\psi)\) is not generated by one invertible local section at a genuine base point. A removable \(0/0\) introduced merely by multiplying numerator and denominator by the same factor is different and was already excluded by the cancellation convention.

\textbf{Proof status.} This section consists chiefly of definitions and short proofs by rational-function algebra. It does not purport to prove finite resolution of base ideals. The projective reformulation and the possibility of different ambient pencils are proved directly. A modern reference for the correspondence between generating sections and maps to projective space is the Stacks Project, Lemma 27.13.1, Tag 01ND; the elementary divisor calculations here are supplied explicitly.

\section{Fasci irrazionali}
\label{enriques-I-2}

An algebraic one-parameter family of moving curves has index one if exactly one member contains a general point of the surface. On normalizing its parameter curve \(B\), the incidence correspondence defines \(f:X\dashrightarrow B\). When \(B\simeq\mathbb P^1\), the coordinate on \(B\) pulls back to a rational function on \(X\), so the family is linear. Enriques's ``single-valued, hence rational'' is justified by the algebraicity of this incidence construction: a generically one-to-one correspondence between integral varieties in characteristic zero induces a birational map. It would not be a valid inference for an arbitrary single-valued analytic function.

An irrational pencil means \(g(B)>0\). This genus is the genus of the parameter curve, not the genus of its fibers. The generators of a ruled surface over \(B\) are an example: the fibers are rational lines although the pencil has genus \(g(B)\). For an ordinary spatial ruled model, a general plane section meets each generator once and is birational to \(B\).

The cone construction on p.~17 projects a curve \(B\) of genus \(p\) from a vertex line in \(\mathbb P^4\); intersecting the resulting threefold scroll/cone with a hypersurface gives a surface fibered over \(B\). The phrase ``fascio irrazionale di curve di genere \(p\)'' is potentially misleading. The construction determines the \textbf{base genus} \(p\). A smooth plane fiber cut by a degree-\(d\) hypersurface has genus \((d-1)(d-2)/2\), independently of \(p\). The scan really contains the phrase; the mathematical repair is to distinguish base and fiber genus, not to invent a different transcription. The construction is an example sketch, with generality and irreducibility of the chosen intersection implicit.

\textbf{Base-point rationality theorem.} An index-one pencil of moving integral curves with a genuine base point at a smooth point of \(X\) has rational parameter curve. The word ``moving'' is essential: adding a fixed curve to every fiber of an irrational fibration would otherwise manufacture ``base points'' without making its base rational.

Enriques first considers distinct moving tangent directions. At a simple base point the pencil is parametrized by those directions; at multiplicity \(i\), it parametrizes degree-\(i\) groups on the exceptional \(\mathbb P^1\), and Lüroth's theorem gives rationality. He explicitly notices the defect in this first argument when all tangents are fixed, then passes to the first moving infinitely near point. The preliminary result cited is Enriques--Chisini, Book II, chapter I, §3, volume I, p.~167.

A precise completion is short. Resolve the rational map to \(B\) by finitely many point blowups. Every component of the exceptional fiber over the smooth base point is a rational curve. Since the strict transform of every general member meets that exceptional fiber, at least one of its finitely many components dominates \(B\). Thus there is a nonconstant \(\mathbb P^1\to B\). Lüroth, or the Riemann--Hurwitz formula, implies \(g(B)=0\). This supplies exactly the finiteness and the first moving direction used implicitly on p.~18. It also proves that an irrational pencil on a smooth projective surface has no genuine base points. It does not make such a statement at an arbitrary singular point of a singular model.

\section{Sistemi lineari}
\label{enriques-I-3}

Restricting an ambient vector space \(W=\langle\varphi_0,\ldots,\varphi_R\rangle\) of sections gives a quotient \(V\). If the projectivized kernel, consisting of hypersurfaces containing \(F\), has dimension \(h\), then \[
r=\dim\mathbb P(V)=R-h-1.
\] This is elementary rank-nullity; if the kernel is zero, its projective dimension is conventionally \(-1\). Fixed components are common divisorial factors, and the base scheme is the common zero scheme after accounting for them. The scan reads ``delle \(C\)'' on p.~18. The edition likewise reads ``delle.''

An intrinsic description is \[
V=\langle1,t_1,\ldots,t_r\rangle\subset
L(D):=\{f\in\mathbb C(X):\operatorname{div}(f)+D\ge0\}\cup\{0\},
\] where \(D\) bounds the poles. The functions \(1,t_1,\ldots,t_r\), not merely the \(t_i\), must be linearly independent. Enriques's displayed relation in the homogeneous \(\varphi_i\) says precisely this. ``Same pole curve'' must mean that the poles are bounded by a common divisor: cancellation can reduce the actual pole divisor of a sum. The equation \(t_0=\sum_i\lambda_it_i\), with the level value allowed to vary, produces \(\sum_{i=0}^r\lambda_i\varphi_i=0\). Choosing a different denominator corresponds to a change of basis in \(V\); no member of the system is intrinsically the polar member.

Exactly one member passes through \(r\) general points. Indeed the evaluation functionals at suitably chosen points span \(V^\vee\), since the rational map defined by \(V\) has nondegenerate image. A general \(r\)-tuple gives \(r\) independent linear equations in \(r+1\) coefficients and hence a one-dimensional solution space. This property holds for reducible linear systems too; its converse needs the extra hypothesis discussed in Section 4.

The system is irreducible if its general member is integral. Its effective genus is the genus of the normalization of that member, equivalently the genus of the smooth general strict transform after resolving the fixed singularities. Its effective degree is the number \(n\) of intersections of two general members outside the common base cluster, counted with intersection multiplicity. On the resolved model \(n=M^2\). For a pencil all pairwise intersections are base intersections, so \(n=0\).

For an irreducible system, restricting its sections to a general strict transform \(C\) gives the characteristic series \[
\operatorname{im}\!\left(V\longrightarrow H^0(C,\mathcal O_C(M))\right),
\qquad \deg\mathcal O_C(M)=n.
\] The kernel is the line generated by the section defining \(C\). Therefore its projective dimension is \(r-1\), yielding a \(g_n^{r-1}\). The curve inequality \(h^0(C,A)\le \deg A+1\) for an effective line bundle gives \[
r-1\le n,\qquad r\le n+1.
\] One proof of the curve inequality is to impose vanishing at \(n+1\) distinct points: a nonzero section of degree \(n\) cannot vanish at all of them. Enriques gives this restriction argument, invoking the familiar curve-theoretic inequality rather than proving it. The hypothesis of an integral general member is indispensable: a higher-dimensional system of sums of fibers can have \(n=0\). The text itself discusses that case immediately afterwards, so its inequality should not be exported to all reducible systems.

The characteristic series need not be complete. Even for a complete \(|D|\), the exact sequence \[
0\longrightarrow\mathcal O_X
\longrightarrow\mathcal O_X(D)
\longrightarrow\mathcal O_C(D)
\longrightarrow0
\] shows that the cokernel of restriction injects into \(H^1(X,\mathcal O_X)\), with the subsequent map to \(H^1(X,\mathcal O_X(D))\). Thus its defect is at most \(q=h^1(\mathcal O_X)\), and equals \(q\) when \(H^1(\mathcal O_X(D))=0\). This is modern context for the characteristic series, not a result that Enriques has already established on p.~19. It links this chapter to the irregularity problem mentioned in the front matter and to the virtual characteristic series on p.~40.

For a system without fixed components, \(n=0\) means that its resolved morphism has one-dimensional image. To see this, a surface image would give \(M^2=(\deg f)\deg f(X)>0\), whereas a curve image gives a pullback from a curve and square zero. Stein factorization gives \(X\to B\to f(X)\), and the members are sums of fibers. This is the precise modern proof of Enriques's final argument on p.~19. It does not force \(B\) to be rational when the members are reducible.

Enriques attributes the characteristic series to Segre and Castelnuovo around 1890, specifically citing Castelnuovo's 1891 ``Ricerche generali sopra i sistemi lineari di curve piane,'' reprinted in the 1937 selected papers, pp.~137 ff.~This historical citation is recorded as his citation; its entire 1891 text has not been rechecked here.

\section{Proprietà caratteristica dei sistemi lineari}
\label{enriques-I-4}

\textbf{Correct theorem.} Let \(T\) be a complete irreducible \(r\)-dimensional algebraic family of effective divisors on a smooth projective complex surface, with integral general member and \(r>1\). Suppose exactly one member contains \(r\) general points of the surface. Then the family is a projective linear subsystem of a complete linear system. For a family initially defined only on an open parameter set, its closure is the linear system; deleting arbitrary members obviously need not leave the entire projective space.

The statement at the end of p.~19 does not include irreducibility, and p.~20 calls it a simplifying supposition. It is an essential restriction. If \(f:X\to B\) is an irrational fibration, the family \[
\left\{f^*(b_1+\cdots+b_r):(b_1+\cdots+b_r)\in\operatorname{Sym}^rB\right\}
\] has dimension \(r\), and exactly one member contains \(r\) general points of \(X\). It is not a linear system: the degree-\(r\) divisors on \(B\) vary in \(\operatorname{Pic}^r(B)\), rather than in one linear-equivalence class. For example, on \(X=B\times\mathbb P^1\), restriction to a section of \(f\) would identify these varying classes if the family were linear.

This is not a newly discovered failure of Enriques's 1893 theorem. The cited original paper explicitly gives the same sums-of-fibers counterexample and then states the theorem for irreducible curves (Enriques, 1893, pp.~4--5). Thus the defensible repair is to restore Enriques's own original hypothesis. There is no reason to credit a later author with repairing a theorem already correctly formulated in 1893.\footnote{The sums-of-fibers construction is not an independent counterexample: the original 1893 paper already contains the same counterexample and the integral-member hypothesis on pp. 4--5. The 1949 wording is unqualified on p. 19, and p. 20 says ``Pongasi per semplicità di discorso che le curve del sistema \(|C|\) siano irriducibili.'' The correction therefore concerns the formulation of the book, while restoring Enriques's earlier formulation.}

\textbf{What p.~20 proves.} For \(r=2\), Enriques forms a quotient surface from the groups of \(n\) common intersections, reducing to an incidence family of degree one. Pencils through two chosen points on that quotient are rational by Section 2. He identifies them with plane pencils and infers that the quotient curves are lines. He then sketches the induction, representing nets through fixed points as stars of planes. This is a substantial geometric proof sketch. It suppresses the construction of the algebraic quotient, the handling of exceptional/boundary incidence, and compatibility of the projective identifications. Those details should not be replaced by the bare assertion that an arbitrary set-theoretic incidence plane is an algebraic projective plane.

\textbf{A modern completion, supplied here.} First remove common fixed base conditions on a smooth model, and work on the dense open where the family and its incidence have their expected dimensions.

For a two-dimensional family, the members through a general point \(x\) form an index-one pencil: through \(x\) and a second general point there is exactly one member. There is only one positive-dimensional component of this pencil, since two such components, each containing moving integral curves through \(x\), would both sweep the surface and violate uniqueness. By Section 2 the pencil is linear. Two general members of the original family meet away from fixed base points. Otherwise all distinct members on the resolved model would be disjoint, incompatible with the positive-dimensional family of members through a general point. An intersection point may be chosen in the good open locus: requiring that an intersection lie on a prescribed proper curve has incidence dimension at most \(2r-1\), rather than \(2r\). The two members therefore lie in one of the linear pencils and are linearly equivalent.

For \(r>2\), impose \(r-2\) general points. The resulting two-dimensional families satisfy the preceding case. Pairs of members belonging to these nets dominate \(T\times T\): their parameter count is \(2(r-2)+4=2r\), and the fiber over a pair of distinct integral curves is finite, consisting of choices among their finitely many moving intersections. Hence two general original members are linearly equivalent. The family lies in some \(|D|\).

Finally, the condition of containing a point \(x\) is a hyperplane in \(|D|\). Intersecting the closure of \(T\) with \(r\) general point-incidence hyperplanes gives one reduced point, because the corresponding incidence map is generically one-to-one and separable. Boundary families of smaller dimension cannot contain \(r\) general surface points. Thus the projective degree of \(T\) is one; an integral projective variety of degree one is a linear subspace. This proves the asserted linearity. The argument uses only the base-point rationality theorem, standard incidence dimensions, and the elementary degree-one criterion, and avoids attributing an unverified modern proof to a named author.

The complementary Picard-map completion reaches the same conclusion. At a general member impose \(r\) general points at which the incidence projection is unramified, and free each point condition in turn. Each resulting one-parameter family has genuine smooth base points and hence rational parameter; its divisor-class map to \(\operatorname{Pic}(X)\) is constant. The resulting \(r\) tangent directions span the parameter tangent space, so the differential of the class map vanishes. In characteristic zero the class map is constant. The preceding degree-one argument then gives the linear subsystem. In the sums-of-fibers counterexample, the fixed points belong to a retained fixed sum of fibers and are not genuine base points of the moving integral pencil, so this reasoning does not apply.

\section{Estensione dei teoremi di Bertini}
\label{enriques-I-5}

\textbf{Smoothness theorem.} On a smooth variety in characteristic zero, a general member of a linear system has no singularities away from its base locus and fixed components. On a singular variety, the singular locus of the ambient variety must also be excluded. On the ordinary spatial model \(F\), apparent double points caused by identifying two distinct points of the normalization are not singularities of the intrinsic curve on \(X\). This distinction is required to reconcile the initial convention of p.~15 with the statement on p.~21.

Enriques's infinitesimal argument can be made rigorous directly. Locally write a pencil as \(u+t v=0\), and suppose a singular point \(x(t)\) of its general member moves algebraically with \(t\). Differentiating \[
u(x(t))+t\,v(x(t))=0
\] gives \(v(x(t))=0\), because the differential in the surface direction vanishes at a singular point. Therefore also \(u(x(t))=0\). Varying the pencil direction through a general member puts the point in every member. A fixed divisorial locus is treated as a fixed component, and after its removal the remaining locus is the base locus. To justify differentiating, one works at the generic point of a component of the singular incidence and uses separability in characteristic zero.

There is also a clean sheaf-theoretic proof. Over the complement \(U\) of the base locus, evaluation \(V\otimes\mathcal O_U\to L|_U\) is surjective. Its kernel is a vector bundle, and the universal incidence \[
\mathcal I=\{(x,[s]):s(x)=0\}\subset U\times\mathbb P(V)
\] is a projective bundle over the smooth \(U\). It is smooth. The generic fiber over \(\mathbb P(V)\) is regular, hence geometrically regular in characteristic zero; spreading out smoothness gives a dense open set of smooth members. This proves the theorem in its proper scope. Zariski's 1944 paper gives the intrinsic development; modern embedded and incidence forms are treated in the Stacks Project, Tags 0FD4 and 0G4C. Characteristic zero is substantive: a linear span of \(p\)th powers in characteristic \(p\) has nonreduced general members even without base points.

\textbf{Reducibility theorem.} If a linear system on a smooth projective surface has no fixed components and its general member is reducible, it is composed with a pencil, whose parameter curve need not be rational. After resolving the base ideal and taking Stein factorization, the system has the form \[
f:X\longrightarrow B,\qquad
C=f^*(b_1+\cdots+b_d),
\] with the degree-\(d\) point divisors varying in a linear series on \(B\). The letter \(r\) used on p.~21 for the number of components is not necessarily the earlier projective dimension \(r\); modern notation \(d\) avoids that collision.

A proof on the resolved surface makes the logical step precise. If the associated map has surface image, its line bundle \(M\) is nef and \(M^2>0\). The general divisor is smooth by the first theorem. It must be connected. Indeed, if a smooth member decomposed as disjoint effective divisors \(A+B\), then \[
A\cdot B=0,\qquad A^2=M\cdot A\ge0,\qquad
B^2=M\cdot B\ge0.
\] At least one square is positive; the Hodge index theorem would make the other square strictly negative unless that divisor were numerically zero. A nonzero effective divisor cannot be numerically zero because it meets an ample divisor positively. This is a contradiction. Smoothness and connectedness imply irreducibility. Consequently a reducible general member forces a curve image. Stein factorization gives connected fibers over a normal projective curve \(B\); generic smoothness makes the general connected fiber integral. Pulling back the hyperplane divisor on \(B\) gives the asserted sums. This is a complete modern reduction to Bertini, Stein factorization, and the Hodge index theorem, not a claim that Enriques proves those later tools here.

Adding back fixed components gives exactly the three alternatives on p.~21: a fixed part plus an irreducible moving part; a sum of several moving fibers; or that latter sum plus a fixed part. Systems of dimension greater than one and effective degree zero fall in the fiber-composed alternatives. The first theorem concerns a \textbf{general} member; special singular or reducible members are allowed and ubiquitous.

\textbf{Historical assessment.} Enriques explicitly refers to Enriques--Chisini, Book II, chapter I, §5, volume I, p.~180, and cites his 1893 and 1896 papers. His component-motion argument does leave the algebraicity of the families of components implicit. Here a named repair is verifiable: van der Waerden (1937, especially p.~709) supplies this missing algebraic-family step using Chow coordinates; Zariski (1941) gives another rigorous proof. Kleiman (1998, §5, Theorem 5.3) identifies this precise gap in the Enriques--Campedelli treatment and reconstructs the argument. This supports a specific diagnosis, not the assertion that the entire Italian proof was invalid. The derivative proof of the first theorem is itself repairable and essentially sound.

The citation ``Bertini \ldots{} 1880'' on p.~20 is a manuscript date: the paper appeared in 1882. The p.~21 reference to Noether, volume VIII, ``1874'' is bibliographically inaccurate for the cited volume-8 paper, which appeared in 1875, pp.~495--533. Moreover, the historical assertion that Noether had already observed the reducibility theorem is stronger than what Kleiman's examination supports: he reports that Noether did not clearly formulate or prove a particular case. Accordingly the ledger records Enriques's attribution as an attribution, not as independently established priority.

\section{Superficie immagini di sistemi lineari}
\label{enriques-I-6}

For \(V=\langle\varphi_0,\ldots,\varphi_r\rangle\), the map is \[
\phi_V:X\dashrightarrow\mathbb P(V^\vee),\qquad
x\longmapsto[\varphi_0(x):\cdots:\varphi_r(x)].
\] The members through \(x\) form a hyperplane of \(V\), equivalently the star of hyperplanes through \(\phi_V(x)\). This proves the intrinsic construction described on p.~22 and its coordinate formula on p.~23.

When the general member is integral and \(r\ge2\), the image is a surface, apart from the pencil case already separated. On a resolution let \(M=\phi_V^*\mathcal O(1)\), and let \(m=[\mathbb C(X):\mathbb C(Y)]\) be the generic degree onto the reduced image \(Y\). Then \[
n=M^2=m\deg Y,\qquad \deg Y=\frac n m.
\] This follows by pulling back two general hyperplanes: they meet \(Y\) in \(\deg Y\) general points, each with \(m\) preimages. Thus the opening assertion ``image of order \(n\)'' on p.~22 is valid in the birational case, and the later paragraph explicitly corrects it to \(n/m\) for a composite system. It should not be extracted without that qualification.

Enriques calls the general fibers of such a degree-\(m\) map an ``involution.'' This does not mean an automorphism of order two, or even a quotient by a group: a finite function-field extension need not be Galois. A ``simple'' system is one for which \(m=1\), hence one giving a birational map onto its image. His dimension count for simplicity when \(r\ge3\) is a genericity heuristic, not a proof that every such system is simple; the text itself immediately gives counterexamples. A pullback of a very ample system along a fixed finite cover gives nonsimple systems of arbitrarily large dimension. A general projection of an embedding to \(\mathbb P^3\) is birational, but a specified subsystem can factor through a multiple cover.

For a net, \(r=2\), the nondegenerate surface image is \(\mathbb P^2\), so the degree of the map is \(n\). This is the ``multiple plane.'' The pencil of members through a general point has all \(n\) points of that point's map fiber as base points: on the resolved moving surface, two independent section equations cut precisely that generic fiber. The branch locus is more precisely the discriminant locus of the finite map obtained after resolution and Stein factorization. At a general simple branch point two points of the fiber coalesce. Arbitrary nets can have higher ramification, so ``a double point in the group'' is not a blanket theorem of simple branching.

\textbf{A genuine printed error.} At the bottom of p.~22, let \(F\subset\mathbb P^3\) have degree \(m\), and take degree-\(a\) cones with a fixed vertex outside \(F\). Projection from the vertex has degree \(m\); the cone system is its composition with the degree-\(a\) Veronese map of \(\mathbb P^2\). Therefore \[
M=aH,\qquad M^2=ma^2,\qquad
\deg(\text{image})=a^2.
\] Enriques uses \(n\) for the cone order and the scan prints \(mn\), without a square. The correct effective degree is \(mn^2\). This is not an OCR error: the transcription is faithful at this point. The calculation also confirms that the map belongs to an involution of degree \(m\), as claimed.

The asserted relation between birational geometry with a chosen simple system and projective geometry of its image is exact in the following sense. A birational map identifying two systems identifies their spaces of sections up to scalars and hence gives a projective equivalence of the image models; conversely a projective equivalence of birational image models lifts to the birational identification. A subspace \(W\subset V\) gives the linear projection \[
\mathbb P(V^\vee)\dashrightarrow\mathbb P(W^\vee)
\] with center \(\mathbb P(\ker(V^\vee\to W^\vee))\). The intersection of this center with the image records the base scheme of \(W\); a fixed component records a divisorial part of that intersection. Enriques sketches these correspondences correctly. On a singular or nonnormal image the full scheme structure, not just the set of intersection points, is needed.

\textbf{One ordinary base point.} Blow up a smooth point \(O\) of multiplicity \(i\). If there is no fixed tangent, the leading homogeneous forms give a base-point-free subsystem of \[
H^0(E,\mathcal O_E(i)),\qquad E\simeq\mathbb P^1.
\] Consequently \(M\cdot E=i\). What always follows is \[
\deg(E\longrightarrow\phi(E))\cdot\deg\phi(E)=i.
\] The stronger claim on p.~23 that the image curve itself has degree \(i\) requires degree one on \(E\). It is valid for the standard sufficiently separating model, but does not follow just from global birationality and the absence of fixed tangents.

\textbf{Explicit counterexample.} On \(\mathbb P^2\), use \[
[x:y:z]\longmapsto[x^2z:y^2z:x^3:y^3].
\] The only base point is \(O=[0:0:1]\), of multiplicity two, with no common tangent factor. On a dense open the inverse affine coordinates are \(x/z=X_2/X_0\) and \(y/z=X_3/X_1\), so the map is birational. On the exceptional curve of the blowup it becomes \[
[u:v]\longmapsto[u^2:v^2:0:0],
\] a degree-two map onto a line. The image degree is one, not \(i=2\). The moving class is \(3H-2E\), whose square is five; its image has equation \(X_0^3X_3^2-X_1^3X_2^2=0\) and is nonnormal along the image line. This example shows exactly which exceptional separation condition is missing. The base ideal, inverse coordinates, restriction to \(E\), image equation, and intersection number are all explicit.

\textbf{One additional simple infinitely near point.} If the base cluster is \(O\) of multiplicity \(i\) followed by one simple point \(O_1\), and no further base points, use total exceptional classes \(E_0^*,E_1\): \[
M=\pi^*D-iE_0^*-E_1,\qquad
E_0'=E_0^*-E_1,\qquad
M\cdot E_0'=i-1,\quad M\cdot E_1=1.
\] This proves the degree splitting \(i-1+1=i\) asserted in the standard birational model. As above, intersection degrees must be divided by mapping degrees when a component is mapped multiply. If \(i=1\), \(E_0'\) is a \((-2)\)-curve with zero \(M\)-degree. In the normal model contracting just that curve, its image is the ordinary \(A_1\) double point, analytically \(ab=c^2\), and the image of \(E_1\) is a line through it. This is the quadric-cone model invoked by Enriques. Normality and the exclusion of additional contracted components are needed for this exact singularity type.

\textbf{The note on p.~24.} The asserted correspondence between exceptional components and ``branches'' is too loose to serve as a modern counting theorem. Exceptional divisors on a smooth resolution are indexed by blowup centers, with their incidence and multiplicities determined by proximity; they are not in general counted by the analytic branches of one plane-curve germ. An ordinary multiple point has several analytic branches but one exceptional component after its first blowup; a unibranch singularity such as \(y^2=x^5\) needs repeated blowups even just to make its strict transform smooth. Further contractions can alter the number of components visible on a singular image. The correct invariant data are the cluster of infinitely near points, its proximity relations, and the total versus strict exceptional divisors. Enriques explicitly cites that theory and returns to it in §§11--12, so the comment is best treated as an imprecise preview rather than substituted for his later formulas.

The cited paper is Barber and Zariski (1935), ``Reducible exceptional curves of the first kind,'' pp.~119--141. The scan's ``E. Barber'' is a wrong initial; the author is S. F. Barber. The journal is American Journal of Mathematics, whereas the scan really prints ``Mathematik.'' The abbreviated Du Val reference has been identified against the journal issue as Du Val (1936), ``Reducible exceptional curves,'' pp.~285--289. It must not be confused with his 1934 papers on isolated surface singularities.

\section{Curve equivalenti e sistemi lineari completi}
\label{enriques-I-7}

Two effective divisors \(C,K\) are linearly equivalent if \(C-K=\operatorname{div}(t)\) for some nonzero rational function. Equivalently they are members of one linear system; distinct members span a pencil and a single member gives dimension zero. Equal degree in a chosen embedding is a consequence, not a sufficient condition. The identification with zero and pole curves is literal when the effective divisors have no common components; otherwise their common part is cancelled, or retained as a fixed part of the pencil.

Reflexivity uses \(t=1\), symmetry uses \(1/t\), and transitivity uses products or quotients. With Enriques's choices \[
\operatorname{div}(t)=K-C,\qquad
\operatorname{div}(\tau)=K-L,
\] one gets \(\operatorname{div}(\tau/t)=C-L\). Thus the quotient on pp.~24--25 has the right sign. There is no hidden sign error to correct.

\textbf{Completeness theorem.} The totality of effective divisors linearly equivalent to \(C\) is \[
|C|=\mathbb P H^0(X,\mathcal O_X(C)).
\] A section represented by \(f\) corresponds to the effective divisor \(C+\operatorname{div}(f)\). Two nonzero sections give the same divisor exactly when their ratio is a nonzero constant. Conversely every linearly equivalent effective divisor has this form. The finiteness of \(H^0\) on a projective variety makes this a finite-dimensional projective linear system. This proves both existence and uniqueness of the complete system and shows that it can be recovered from any one member. Enriques gives the rational-function construction and uniqueness argument; modern sheaf cohomology supplies a uniform finiteness statement even outside the irreducible moving situation.

The enlargement procedure on p.~25 is valid after fixing the divisor class and the prescribed base conditions: any strictly increasing chain of subspaces of \(H^0\) must terminate, and their unique maximal member is \(H^0\) itself. The earlier bound \(r\le n+1\) justifies termination only in its irreducible moving setting. It should not be applied to a reducible degree-zero system. This is a limitation of the suggested proof, not a failure of the completeness theorem.

\textbf{Completeness does not imply base-point-freeness.} Enriques actually says ``in general,'' then introduces the convention ``virtually'' without base points. These words must not be promoted to a universal theorem. For a concrete counterexample, blow up eight general points of \(\mathbb P^2\). Cubics through them form a complete pencil on the blowup and still have a ninth base point, by Bézout. Its general member is an integral cubic. Completeness means that all sections of the given line bundle have been used; it does not say that they generate the line bundle at every point.

A complete system with prescribed base cluster \((p_j,m_j)\) is \[
\mathbb P H^0\!\left(Y,\pi^*\mathcal O_X(C)
                 \otimes\mathcal O_Y(-\sum m_jE_j^*)\right)
\] on the blown-up surface. Equivalently it is a complete space of sections satisfying the corresponding ideal conditions on \(X\). The exceptional divisors appearing in the total pullback must be subtracted to obtain the strict/moving transform. After further blowing up any remaining base points, the moving series is base-point-free. Enriques correctly stresses this subtraction. His maximality statement applies to an irreducible complete relative system containing a subsystem of dimension at least two. ``Same degree'' concerns the moving degree with those base conditions retained; removing imposed conditions can increase that degree.

\textbf{What ``normal surface'' means here.} Fix a smooth or normalized model and its hyperplane line bundle \(L\). Enriques calls an image normal when its hyperplane sections give the complete intrinsic system \(H^0(L)\), so that it cannot be lifted to a larger projective space by adjoining missing sections while preserving the degree. If the embedded surface is already smooth, this is modern linear normality. If the image is nonnormal, it refers to sections on the normalization and is stronger than ambient restriction onto \(H^0(F,\mathcal O_F(1))\). In fact every hypersurface in \(\mathbb P^3\) has surjective restriction of linear forms, whereas Enriques's quintic example still lifts from \(\mathbb P^3\) to \(\mathbb P^5\). Calling both notions simply ``linear normality'' without naming the model would conceal that difference. Neither intrinsic completeness nor linear normality is the same as projective normality, which also concerns higher multiplication maps and the homogeneous coordinate ring.\footnote{Enriques's usage should not be abbreviated as ``linearly normal'': it concerns completeness on the normalization, which must be distinguished from ambient linear normality of the nonnormal spatial image. The printed lifting example on p. 26 settles the intended meaning: all spatial hypersurfaces have complete ambient linear sections, yet Enriques's quintic lifts to its complete six-section system on the smooth model.}

The uniqueness of the final complete image in pp.~25--26 is consequently a uniqueness for the fixed polarized birational model: any two bases of \(H^0(L)\) differ by an invertible linear transformation. All chains of adjoining sections end in the same projective model up to projective equivalence. It is not a general uniqueness theorem for all projective embeddings of the underlying surface.

\textbf{The quintic example, reconstructed.} Let \[
S=\operatorname{Bl}_{A_1,\ldots,A_4}\mathbb P^2,
\quad D=3H-E_1-\cdots-E_4=-K_S,
\] with the four points in general position. Cubics form a vector space of dimension ten; the four independent point conditions leave \(h^0(D)=6\). Also \[
D^2=9-4=5,\qquad
p_a(D)=1+\frac{D(D+K_S)}2=1.
\] The complete system embeds the degree-five del Pezzo surface in \(\mathbb P^5\). A four-dimensional vector subspace \(V\subset H^0(D)\) gives a general quintic image in \(\mathbb P^3\). Intermediate five-dimensional subspaces \(W\) containing \(V\) are parametrized by \(\mathbb P(H^0(D)/V)\simeq\mathbb P^1\), explaining Enriques's \(\infty^1\) family of lifts to \(\mathbb P^4\). The final six-dimensional system is unique.

The intermediate models are generally projectively inequivalent. One must read ``birationally distinct'' here as ``not carried to one another with their specified linear systems''; all the surfaces themselves are rational and hence birational. The polarized automorphism group of the smooth degree-five del Pezzo surface is finite, namely \(\mathfrak S_5\), so it cannot identify a one-parameter family into a single projective orbit. Special coincidences need not be ruled out by this argument. The anticanonical model and this automorphism group are verified in Bauer and Catanese (2021, §2).

Enriques additionally specifies a degree-five double curve lying on a quadric cone and one triple point at its vertex. These details are stated without a calculation. Here is a modern reconstruction for a general projection with ordinary singularities. The conductor divisor on \(S\) is \[
\Delta=(5-4)D-K_S=2D,
\] so \(D\Delta=10\); since it maps two-to-one onto the double curve \(B\), \(\deg B=5\). Adjunction gives \(p_a(\Delta)=6\). The conductor exact sequence gives \[
\chi(\mathcal O_F)=\chi(\mathcal O_S)+\chi(\mathcal O_B)
                 -\chi(\mathcal O_\Delta).
\] Here the quintic hypersurface has \(\chi(\mathcal O_F)=5\), the rational smooth model has \(\chi(\mathcal O_S)=1\), and \(\chi(\mathcal O_\Delta)=-5\). Hence \(p_a(B)=2\).

Let \(t\) be the number of ordinary triple points and \(b\) the number of pinch points. The rank-one degeneracy locus of \(df:T_S\to f^*T_{\mathbb P^3}\) has, by the Thom--Porteous formula, \[
b=c_2(f^*T_{\mathbb P^3}-T_S)
 =6D^2+4D K_S+K_S^2-c_2(S)
 =30-20+5-7=8.
\] An ordinary triple point contributes three nodes to \(\Delta\), but contributes delta invariant two to the spatial curve \(B\), whose three tangents are not coplanar. Normalizing and applying Riemann--Hurwitz to the two-to-one conductor cover yields \[
-5+3t=2(-1+2t)-\frac b2.
\] Thus \(t=1\). These equations also work componentwise, using Euler characteristics if a conductor curve is disconnected. Here Franchetta's theorem gives irreducible \(B\); its normalization has genus \(2-2t=0\), and the conductor normalization has genus \(6-3t=3\). Its double cover is connected, since a disconnected degree-two cover of a connected smooth curve would split into degree-one covers and have no branch points. The ordinary-projection input is also recorded by Friedman, Lehman, Leyenson, and Teicher (2012, Section 2). The unique quadric adjoint is explained by \[
H^0(\mathbb P^3,\mathcal I_B(2))
\simeq H^0(S,\mathcal O_S(2D-\Delta))
\simeq\mathbb C.
\] At the triple point its gradient must vanish because the three tangent directions of \(B\) span the ambient tangent space. The quadric cannot split into two planes: its pullback is the reduced divisor \(\Delta\), so no component of \(B\) can lie in both planes. Each plane pullback would then map two-to-one onto the double-curve components it supports, forcing its hyperplane degree \(D^2=5\) to be even. A double plane would likewise give a nonreduced pullback. Thus the unique quadric is irreducible and singular, hence a quadric cone, and the triple point is its vertex. This reconstruction uses the ordinary-projection hypotheses, the conductor computation in Section 8, and the standard degeneracy-locus theorem (Fulton, 1998, §14.4). This reconstruction is not a proof present on p.~26.

\section[Somma e differenza di sistemi lineari]{Somma e differenza di sistemi lineari: teorema del resto}
\label{enriques-I-8}

\textbf{Addition and the distinction between two sums.} If \(C_1\sim C_2\) and \(K_1\sim K_2\), multiplication of their rational functions gives \[
C_1+K_1\sim C_2+K_2.
\] If positive and negative components overlap, the principal divisor cancels them; the displayed equivalence still holds. The complete sum is \(|C+K|\). The smallest linear system containing every sum of a member of \((L,V)\) and a member of \((N,W)\) is instead \[
\mathbb P\,\operatorname{im}
\left(V\otimes W\longrightarrow H^0(X,L\otimes N)\right).
\] This is Enriques's minimal sum. The multiplication map need not be onto, even if the two starting systems are complete relative to their assigned conditions.

If the two positive-dimensional systems have integral general members and are not the same pencil, their minimal sum already has a surface image. Indeed its map is the Segre composite of their two maps; a curve image would force both systems to come from the same one-dimensional fibration. Bertini then makes its general member irreducible. The complete sum, which contains it, is also generically irreducible. In particular multiples of a non-pencil irreducible moving system are irreducible in this sense. The pencil exception means that the argument gives no such guarantee. It must not be read as a claim that every unrestricted multiple of every pencil remains reducible: base-point conditions may disappear when the complete divisor class is enlarged.

The sentence that assigned base points remain base points in the sum is correct only with its adjective ``assigned'': prescribed vanishing conditions are retained in taking the relative complete sum. It is false for the unrestricted complete divisor class. For example, pencils of lines through two different points of \(\mathbb P^2\) both have class \(H\), but \(|2H|\) is the base-point-free complete system of conics. In the relative interpretation the imposed conditions persist, and their multiplicities add at any common base point.

\textbf{The quadric example and its degree restriction.} Put \(O=[0:0:0:1]\) on a smooth degree-\(d\) surface \(F\subset\mathbb P^3\). Hyperplanes through \(O\) give a three-dimensional vector space. Their pairwise products span the six-dimensional space of quadratic forms in \(x_0,x_1,x_2\), namely cones with vertex \(O\). A quadric section has multiplicity at least two at \(O\) on \(F\) precisely when its value and its two tangent derivatives vanish: three independent linear conditions in the ten-dimensional space of quadrics leave dimension seven. For \(d\ge3\), restriction of quadrics to \(F\) is injective, so Enriques's projective dimensions are \[
\mathbb P^5\subsetneq\mathbb P^6.
\] This proves the claimed failure of the minimal double to be complete relative to the double base point.

The printed example says merely ``degree \(n\)'' and needs \(n\ge3\). If \(F\) is a smooth quadric, its own equation lies in the seven-dimensional tangency space and vanishes on \(F\), leaving only six independent restrictions. The six cone quadrics restrict injectively and already fill the system; both projective dimensions are five. For a plane, the starting hyperplane-section system itself drops dimension. These are genuine omitted low-degree qualifications, not transcription errors.

\textbf{Residual systems and the exact algebra.} If an effective member of \(|C|\) is \(K+L\), the residual system to \(K\) is \[
|C-K|=\mathbb P H^0(X,\mathcal O_X(C-K)).
\] Multiplication by the canonical section of \(\mathcal O_X(K)\) identifies it with the subspace of \(|C|\) consisting of divisors containing \(K\). If \(K'\sim K\), then \(\mathcal O_X(C-K')\simeq\mathcal O_X(C-K)\), so the \textbf{complete residual system} is the same. Equivalently, every effective \(L\sim C-K\) satisfies \(L+K'\sim C\). This proves the residual theorem on pp.~27--28 and the identities \(|C-K|=|L|\), \(|L+K|=|C|\).

Completeness is essential. For an arbitrary incomplete subsystem \(V\subset H^0(C)\), the residual vector space is the inverse image of \(V\) under multiplication by the particular section \(s_K\); replacing \(K\) by a linearly equivalent divisor can change it. The italic formulation on p.~28 uses ``a system \(|C|\),'' but its immediately preceding construction explicitly assumes the complete system. The modern statement must retain that assumption.

\textbf{The projective complement, asserted for later proof on p.~28.} This is distinct from the intrinsic residual theorem. For a smooth hypersurface \(F\) of degree \(d\), \[
0\longrightarrow\mathcal O_{\mathbb P^3}(e-d)
\longrightarrow\mathcal O_{\mathbb P^3}(e)
\longrightarrow\mathcal O_F(e)\longrightarrow0.
\] Since \(H^1(\mathbb P^3,\mathcal O(k))=0\) for every integer \(k\), all sections of \(\mathcal O_F(e)\) are cut by degree-\(e\) ambient surfaces. This is completeness of the ambient cut, not merely an asymptotic assertion.

For an ordinary singular hypersurface, let \(\nu:X\to F\) be its smooth normalization, \(B\) its reduced double curve, and \(\Delta\) the conductor divisor on \(X\). The conductor ideal on \(F\) is \[
\mathfrak c=
\mathcal{H}om_{\mathcal O_F}(\nu_*\mathcal O_X,\mathcal O_F)
=\mathcal I_{B/F}=\nu_*\mathcal O_X(-\Delta).
\] The equality with the reduced double-curve ideal can be checked in the ordinary local models \(xy=0\), \(xyz=0\), and \(x^2=y^2z\). At the pinch point, normalize by \(y=u,\ x=uv,\ z=v^2\); the conductor is \((x,y)\) downstairs and \((u)\) upstairs. At a triple point it is the ideal \((xy,xz,yz)\); on each normalization branch it cuts the two intersecting conductor curves.

Writing \(H=\nu^*\mathcal O_F(1)\), restriction gives \[
0\longrightarrow\mathcal O_{\mathbb P^3}(e-d)
\longrightarrow\mathcal I_B(e)
\longrightarrow\nu_*\mathcal O_X(eH-\Delta)
\longrightarrow0.
\] The same vanishing therefore proves that adjoint surfaces of degree \(e\), passing through \(B\), cut the complete residual system \(|eH-\Delta|\). Finite duality and hypersurface adjunction identify \[
K_X\sim(d-4)H-\Delta,\qquad
eH-\Delta\sim K_X+(e-d+4)H.
\] Thus this is also the modern canonical meaning of ``adjoint.'' No gap in the truth of the p.~28 completeness assertion appears under the ordinary-singularity convention. For arbitrary singularities one must use the conductor or adjoint ideal; passage through only the reduced singular locus is not sufficient. The required standard facts are the Stacks Project, Lemma 30.8.1, Tag 01XS, and finite duality in Remark 48.20.12, Tag 0AU5. The local conductor checks and exact-sequence argument have been given here.

Noether's Restsatz consequently has the two parts Enriques explicitly distinguishes: intrinsic independence of the complete residual system from the representative of \(K\), and completeness of the projectively cut adjoint system. Once both are available, every \(C'\sim C\) is residual to every \(K'\sim K\) in the appropriate complete adjoint class. The book cites Noether, Mathematische Annalen VIII, and, for the curve analogue, Brill--Noether VII and Enriques--Chisini, Book V; p.~28 postpones the projective proof rather than pretending it has just proved it.

\textbf{Subtraction can destroy irreducibility and introduce base loci.} If \(f:X\to\mathbb P^1\) has connected fibers \(L\) and no base points, then \[
H^0(X,\mathcal O_X(2L))\simeq H^0(\mathbb P^1,\mathcal O(2)).
\] A general member of \(|2L|\) is the sum of two fibers. Choose a moving system \(K\) such that \(|K+2L|\) is irreducible; subtracting \(K\) leaves this reducible system. On \(\mathbb P^1\times\mathbb P^1\), \(|(1,2)|-|(1,0)|=|(0,2)|\) makes the example explicit. On \(\operatorname{Bl}_p\mathbb P^2\), \(|H|-|(H-E)|=|E|\) gives an irreducible starting system whose residual has a fixed exceptional component. On the blowup of eight general plane points, taking \(C=nA\), with \(A\) ample and \(n\) large, and choosing an effective \(K\sim nA+K_X\), gives residual \(|-K_X|\), the cubic pencil with a ninth base point described above. Blowing up a newly created base point turns its contribution in the total pullback into a fixed exceptional component, exactly as on p.~29.

\textbf{Virtual divisors.} If \(H^0(X,\mathcal O_X(C-K))=0\), the effective residual system is empty, but the divisor class \(C-K\) remains meaningful in \(\operatorname{Pic}(X)\). Formal differences of effective classes obey \[
C-K\sim L-M
\iff C+M\sim K+L.
\] Addition is associative in the divisor class group: \[
D+(C-K)=D+C-K.
\] If \(A\) is ample, \(\mathcal O_X(mA+C-K)\) has sections, and is indeed globally generated, for \(m\gg0\). This supplies the precise meaning of ``sufficiently ample'' in p.~29; it follows from the ample-twist theorem and Serre vanishing, not from a numerical assertion about an arbitrary large self-intersection. Conversely every line bundle on a projective variety is a difference of effective divisors after such a twist. The analogy with negative numbers is the group completion of the monoid of effective divisor classes.

Enriques also proposes the narrower notion of a noneffective class having an effective positive multiple. That is a separate condition, not the definition of all virtual classes. For example, \(-H\) on \(\mathbb P^2\) has no effective positive multiple. By contrast, on \(E\times\mathbb P^1\) let \(M\ne\mathcal O_E\) be an order-\(a\) torsion line bundle and put \[
L=p_E^*M\otimes p_{\mathbb P^1}^*\mathcal O(1).
\] Then \(H^0(L)=0\) because a nontrivial degree-zero line bundle on an elliptic curve has no section, but \[
L^{\otimes a}\simeq p_{\mathbb P^1}^*\mathcal O(a),
\qquad h^0(L^{\otimes a})=a+1.
\] A simpler torsion example keeps the zero effective divisor: any nontrivial torsion line bundle on a projective integral variety has no section (an effective nonzero divisor cannot have a trivial positive multiple), whereas its torsion power is trivial. The canonical class of a complex Enriques surface is a later example. The displayed product construction has the additional advantage of giving effective nonzero multiples. This is an explicit noneffective virtual class with effective nonzero multiples. It is an example supplied here, not attributed to Enriques.

The historical paragraph continuing onto p.~30 credits the idea to Enriques's canonical-class work of 1893--1896, the terminology ``virtual curve'' to Severi (1905), and the topological interpretation to Lefschetz. Its mathematically precise topological map is \[
\operatorname{Pic}(X)\longrightarrow H^2(X,\mathbb Z),
\qquad L\longmapsto c_1(L),
\] or the Poincaré-dual two-cycle class on the underlying compact oriented four-manifold. This does not identify virtual divisors with all two-cycles and does not recover linear equivalence: classes in \(\operatorname{Pic}^0(X)\) have zero first Chern class, and rational cohomology also forgets torsion. The historical priority claims are retained as statements made by the book; the specific 1905 Severi paper and an exact Lefschetz passage were not independently read.

\section{Superficie immagine del sistema somma}
\label{enriques-I-9}

Throughout this section, work on a smooth projective complex surface \(X\), after resolving the assigned base points of the systems. The symbol \(K\) in Enriques's \(|C+K|\) denotes an auxiliary system; it is \textbf{not} the canonical divisor, which will be written \(K_X\). If \(V\subset H^0(X,L)\) and \(W\subset H^0(X,M)\) have dimensions \(r+1\) and \(s+1\), the minimum sum is the projectivization of

\[
U=\operatorname{im}\bigl(V\otimes W\longrightarrow H^0(X,L\otimes M)\bigr).
\]

Its map is the product of the two rational maps followed by the Segre embedding, with any linear dependencies among the products removed. The complete sum uses all of \(H^0(X,L\otimes M)\). The cokernel of the multiplication map measures the difference; asking which positivity conditions make it vanish leads to projective-normality and syzygy questions, rather than following from completeness alone. Enriques's assertion that the image is generally a surface and \(R=\dim U-1>r+s\) has a general-position qualification in the printed text. It must not be universalized: two copies of one pencil give a conic image and \(R=2=r+s\). The construction in the ensuing paragraphs concerns two systems whose joint map has two-dimensional image. A common finite correspondence can still make that map generically many-to-one (Enriques, 1949, p.~30).

The strict dimension inequality can be justified in the intended case, with both systems positive-dimensional and with integral general members, as in the preceding discussion. The family of sums of one integral member from each system has dimension \(r+s\): the decomposition of a general divisor into its two integral components is unique, up to the finite interchange when appropriate. These sums are reducible. The minimum sum has an irreducible general member when its joint map has two-dimensional image, by the Bertini argument of §8. Thus the \(r+s\)-dimensional family of decomposable members is a proper closed subset of \(\mathbb P(U)\), forcing \(R>r+s\). This proves the bound rather than assuming injectivity of the entire multiplication map.

Here is the linear algebra behind the construction, including the step that Enriques leaves implicit. Choose \(c_0\in V\), \(k_0\in W\) with divisors \(C_0,K_0\) having no common component. The subspaces \(k_0V,c_0W\subset U\) meet precisely in the line spanned by \(c_0k_0\): a section divisible by both divisors is divisible by their sum, and the quotient is a global regular function on the connected projective surface. The subsystem

\[
U_0=k_0V+c_0W,\qquad \dim U_0=r+s+1,
\]

therefore gives a map to \(\mathbb P^{r+s}\). The two centers lie in the common hyperplane represented by the product section \(c_0k_0\). It is the projection from the common \(\mathbb P^{R-r-s-1}\) of the two centers in \(\mathbb P^R\). The fixed \(K_0\) lies in the center \(\mathbb P^{R-r-1}\), and fixed \(C_0\) in \(\mathbb P^{R-s-1}\). After projection they lie in \(\mathbb P^{s-1}\) and \(\mathbb P^{r-1}\), respectively. On the open set \(c_0k_0\ne0\), the coordinates recover all ratios \(c_i/c_0\) and \(k_j/k_0\). Thus the projected map generates exactly the joint function field of the two original maps. It is birational precisely when that field is \(\mathbb C(X)\). This proves the assertion about systems not belonging to a common involution, with ``involution'' read as their common generic finite equivalence relation, not as an automorphism of order two.

Two small \textbf{printed errors} occur in this otherwise sound argument. The ambient space in the second paragraph is printed \(S_r\), although it must be \(S_R\). The following paragraph interchanges \(C\) and \(K\) when naming the curves in the two centers. Both discrepancies are present in the scan; neither should be attributed to the transcription. The dimensions following the projection and the specialization \(r=2,s=1\) confirm the intended assignment.

For the net--pencil case, choose coordinates

\[
[c_0k_0:c_1k_0:c_2k_0:c_0k_1]\in\mathbb P^3.
\]

The curve \(K_0\) goes to \(O=[0:0:0:1]\), while \(C_0\) maps to the line \(a=\{x_0=x_3=0\}\), which does not contain \(O\). Planes through \(O\) cut the net and planes through \(a\) cut the pencil. Write \(C^2=n\), \(C\cdot K=i\), and \(K^2=0\) on the model where the pencil is a morphism. Assuming the chosen \(C_0,K_0\) meet transversely in \(i\) points, blow those points up, with exceptional divisors \(E_1,\ldots,E_i\). The moving divisor of the four-dimensional vector space above is

\[
D=\pi^*(C+K)-\sum_{j=1}^{i}E_j,\qquad
D^2=n+i.
\]

Consequently a birational image has degree \(n+i\). Each \(E_j\) has \(D\cdot E_j=1\), so maps to a line. These are the \(i\) lines in the plane \(Oa=\{x_0=0\}\) left after removing the multiple line \(a\). The net contributes multiplicity \(i\) at \(O\); a general plane through \(a\) has residual pencil member of degree \(i\), so \(a\) has multiplicity \(n\). These are generic statements about the prescribed model. Nontransverse intersections must be replaced by their scheme lengths and infinitely near contributions; if the joint map has degree \(d>1\), the basic identity is \(D^2=d\deg\Phi\), rather than \(\deg\Phi=D^2\). Enriques explicitly imposes birationality before his degree computation (pp.~30--31).

Finally, products of a plane through \(O\) with one through \(a\) span the quadrics through \(O\cup a\). Their ambient vector space has dimension \(10-3-1=6\), hence projective dimension five. Restricting them to the surface and removing fixed parts recovers the minimum sum. This does not assert that the restricted system always has dimension five: when the image surface itself is a quadric, its equation gives a relation. This distinction prevents an artificial ``dimension error'' in Enriques's last paragraph.

\textbf{Assessment.} Enriques gives a genuine geometric construction, including the birationality condition and its numerical consequences. The modern calculation above fills in kernels, dimensions, and base schemes; it does not replace a false theorem. The two letter-level source slips change the labels, not the construction just proved.

\section{Caratteri virtuali}
\label{enriques-I-10}

\subsection{Addition and the intrinsic definitions}

For Cartier divisors \(A,B\) on \(X\), the modern statements are

\[
(A+B)^2=A^2+B^2+2A\cdot B,\qquad
p_a(A+B)=p_a(A)+p_a(B)+A\cdot B-1,
\]

where

\[
p_a(D)=1+\frac{D^2+D\cdot K_X}{2}.
\]

For an effective divisor this equals \(1-\chi(\mathcal O_D)\); for a formal divisor class the same numerical expression defines its virtual arithmetic genus. Enriques writes these as \(n=n_1+n_2+2i\) and \(\pi=\pi_1+\pi_2+i-1\), with degrees and genera evaluated after resolving the assigned base cluster. If the general curves are smooth there, their geometric genera equal these arithmetic genera. The first formula is bilinearity; Enriques assumes the second from Noether and cites Enriques--Chisini, book V, vol.~III, p.~399 (Enriques, 1949, p.~31).

A proof valid also for nonreduced divisors uses

\[
0\longrightarrow\mathcal O_X(-D)\longrightarrow\mathcal O_X
\longrightarrow\mathcal O_D\longrightarrow0
\]

and surface Riemann--Roch. The latter says

\[
\chi(\mathcal O_X(D))=\chi(\mathcal O_X)
+\frac{D\cdot(D-K_X)}2.
\]

Subtracting Euler characteristics gives the displayed genus formula; expanding its quadratic expression gives the addition formula. This explains precisely the ``continuity'' used in the text: Euler characteristic is constant in a flat proper family, whereas normalization genus can drop. The required intersection and Riemann--Roch framework is given in Beauville (1996, chap.~I); an independently accessible proof is Vakil (2002, class 6, §2). No claim is made that these modern expositions were the first historical repairs.

Enriques defines the virtual degree of an isolated curve \(C\) by choosing \(|K|\) such that \(|K|\) and \(|C+K|\) have irreducible general members, then setting

\[
n=N-m-2i,\qquad N=(C+K)^2, m=K^2, i=C\cdot K.
\]

Existence of such an auxiliary system follows by taking an ample \(H\) and \(m\gg0\), so that both \(mH-C\) and \(mH\) are very ample. Its independence is not merely asserted: p.~32 computes the degree of \(C+K_1+K_2\) in two orders and cancels terms to obtain \(n_1=n_2\). This is a complete algebraic verification once the addition formula and the availability of sufficiently positive auxiliaries are granted. The analogous genus calculation, only indicated in the text, is

\[
p_a(C+K_1+K_2)
=p_a(C)+p_a(K_1)+p_a(K_2)
+C K_1+C K_2+K_1K_2-2.
\]

It is visibly symmetric in \(K_1,K_2\). The same argument applies to differences \(C-K\), regardless of whether their linear systems contain an effective divisor. Serre vanishing and very-ampleness after sufficiently high twisting supply the auxiliary systems (Stacks Project, Tag 0B5U).

Multiple components cause no failure. For an integral \(C\) and integer \(m>0\), for example,

\[
(mC)^2=m^2C^2,\qquad
p_a(mC)=1+m(p_a(C)-1)+\frac{m(m-1)}2C^2.
\]

This is the scheme-theoretic genus, not the genus of the reduced support. It can be negative. The assertion on p.~32 is therefore stronger and more useful than a statement only about smooth moving curves.

\subsection{Cohomological justification of the virtual formulas}

For a smooth curve \(C\subset X\), the conormal sequence

\[
0\to\mathcal O_C(-C)\to\Omega_X^1|_C\to\Omega_C^1\to0
\]

gives, after taking determinants,

\[
\omega_C\simeq(\omega_X\otimes\mathcal O_X(C))|_C,\qquad
2g(C)-2=C\cdot(C+K_X).
\]

Enriques does not prove surface Riemann--Roch in Chapter I; its dedicated treatment is Chapter IV. The complete reduction to curve Riemann--Roch is given in Chapter IV, Section 11. Its short form writes \(D=A-B\) with \(A,B\) smooth very ample divisors, obtained by a sufficiently positive twist and Bertini. The two restriction sequences then give
\[
\begin{aligned}
\chi(\mathcal O_X(D))-\chi(\mathcal O_X)
 &=\chi(\mathcal O_A(A))-\chi(\mathcal O_B(A))\\
 &=\frac{D\cdot(D-K_X)}2.
\end{aligned}
\]
Thus the virtual character formulas use an Euler-characteristic identity, without an assumption of vanishing.

Serre duality \(h^2(\mathcal O_X(D))=h^0(\mathcal O_X(K_X-D))\) rewrites this as

\[
h^0(D)-h^1(D)+h^0(K_X-D)
=1+p_a(X)+\tfrac12D\cdot(D-K_X).
\]

This identity is not the assertion that \(h^1(D)=0\). Treating an expected dimension as an actual dimension without a vanishing theorem is a logical step requiring separate justification. The standard textbook locations are Hartshorne (1977), Chapter V, §§1 and 3, with surface Riemann--Roch at Theorem V.1.6; the underlying curve formula, divisor sequences, and projective cohomology can also be checked in the separately identified Stacks Project references in the bibliography. The algebra displayed here is supplied in full so that the check does not depend on an inaccessible book passage.

For any effective Cartier divisor \(D\), possibly reducible or nonreduced, the sequence \(0\to\mathcal O_X(-D)\to\mathcal O_X\to\mathcal O_D\to0\) now yields

\[
\boxed{p_a(D):=1-\chi(\mathcal O_D)=1+\tfrac12D\cdot(D+K_X).}
\]

This is Enriques's virtual genus on an unmarked smooth model. In particular,

\[
\begin{aligned}
(A+B)^2&=A^2+2A\cdot B+B^2,\\
p_a(A+B)&=p_a(A)+p_a(B)+A\cdot B-1,\\
p_a(sD)&=1+\tfrac12(s^2D^2+sD\cdot K_X),\\
p_a(A-B)&=p_a(A)+p_a(B)-1-A\cdot B-B\cdot K_X.
\end{aligned}
\]

The last formula follows by expansion and avoids the tempting, incorrect idea that arithmetic genus is additive under formal subtraction. For a reduced integral curve, normalization gives \(p_a(C)=g(\widetilde C)+\sum_P\delta_P\). For a disconnected divisor, arithmetic genus may be negative. For the zero divisor the displayed virtual-genus polynomial gives 1, consistently with \(1-\chi(\mathcal O_{\varnothing})=1\). Thus a virtual genus is a numerical function, not an assertion that a smooth curve of that genus exists.

\subsection{Birational behavior and assigned points}

Let \(\pi:Y\to X\) blow up a smooth point \(p\in X\), with exceptional divisor \(E\), and let a curve \(C\) have multiplicity \(m\) at \(p\). Then

\[
\pi^*C=\widetilde C+mE,\quad E^2=-1,\quad
K_Y=\pi^*K_X+E,\quad\pi^*C\cdot E=0,
\]

and hence

\[
\widetilde C^{\,2}=C^2-m^2,\qquad
p_a(\widetilde C)=p_a(C)-\binom m2.
\]

The canonical coefficient can also be checked in the chart \(x=u,\ y=uv\), where \(dx\wedge dy=u\,du\wedge dv\). Total transform preserves both \(C^2\) and \(p_a(C)\). Strict transform does not. Enriques's \(\bar C+o\) on pp.~32--33 is the multiplicity-one illustration; the general expression requires \(mE\). His discussion is not a claim that self-intersection is an invariant of a bare curve under arbitrary birational maps. It is about a divisor together with the decision whether particular proper or infinitely near points are assigned base points.

More precisely, for an assigned cluster with multiplicities \(m_i\), let \(E_i^*\) denote the total transforms of the exceptional divisors. The corresponding class is

\[
D'=\pi^*D-\sum_i m_iE_i^*.
\]

Because \(E_i^*\cdot E_j^*=-\delta_{ij}\), one obtains

\[
(D')^2=D^2-\sum_i m_i^2,\qquad
p_a(D')=p_a(D)-\sum_i\binom{m_i}{2}.
\]

These equations also explain the unusual but deliberate permission on p.~33 to assign a multiplicity smaller than the curve's actual multiplicity. If the actual multiplicity is \(a>m\), the virtual transform includes \((a-m)E\); it need not equal the reduced strict transform. The appropriate modern setting is the divisor class on the specified blowup, or the language of virtual multiplicities and clusters (Casas-Alvero, 2000, chaps. 3--4).

Two qualifications are essential. First, \(\binom m2\) is the \textbf{arithmetic genus drop of one blowup}. The normalization-genus defect of an arbitrary multiplicity-\(m\) singularity may be larger, because additional singularities remain infinitely near. For \(y^2=x^5\), the first drop is one and the total delta invariant is two. Second, \(m^2\) is the contribution removed by the first blowup from the intersection of two multiplicity-\(m\) curves. Their entire local intersection number equals \(m^2\) only if their tangent cones have no common direction. For instance \(y=0\) and \(y=x^2\) have multiplicities one but local intersection two; one unit remains infinitely near. Enriques's infinitesimal conventions can accommodate both facts, but the last two sentences of p.~33 are not safe as unrestricted formulas about normalization genus or total local intersection.

\section{Curve eccezionali}
\label{enriques-I-11}

\subsection{Definitions and examples}

An irreducible exceptional curve of the first kind is, in the modern smooth-surface formulation, a smooth rational curve \(E\simeq\mathbb P^1\) with \(E^2=-1\). Equivalently, it is the entire exceptional fiber of the blowup of a smooth point, or a curve contractible by a birational morphism to a smooth point with no other exceptional component in that local fiber. The qualification ``entire'' matters: a \((-2)\)-component of a larger exceptional fiber need not itself be contractible to a smooth point. Enriques explains this distinction explicitly in §12. His words about the neighborhood of a smooth point are stronger than simply ``a curve mapped to some point'' (pp.~33--35).

His second kind refers to curves made exceptional through a birational transformation with assigned base points on the curve. For a smooth rational curve with self-intersection \(\nu\ge0\), blowing up \(\nu+1\) smooth points on it gives a \((-1)\)-curve, which can then be contracted. The distinction depends on the model, not solely on the birational class of the surface. A safely formulated version of the converse requires keeping track of the assigned singularity cluster and of the entire exceptional fiber. Resolving an arbitrary rational map can produce several exceptional components, so ``resolve the base points'' alone is not a proof that every named component is a \((-1)\)-curve.

All the elementary examples on p.~34 fit this interpretation. On a smooth cubic surface the 27 lines are \((-1)\)-curves. Its familiar presentation as the blowup of six general plane points gives their classes: the six \(E_i\), fifteen \(H-E_i-E_j\), and six \(2H-\sum_{j\ne i}E_j\). Each has square \(-1\), arithmetic genus zero, and anticanonical degree one; the anticanonical embedding therefore makes it a line. This uses the classical cubic-surface representation assumed in Enriques's example. Its plane representation may be chosen so that any prescribed one of the 27 lines is one of the exceptional divisors over the six base points; that choice is imported from the classical representation theorem. In the plane, a degree-\(d\) curve has square \(d^2>0\), so there is no first-kind curve. A line is contracted by the standard quadratic map after assigning two of its base points on the line; its strict transform has square \(1-1-1=-1\). A smooth conic is contracted by a homaloidal cubic net of type \((3;2,1,1,1,1)\), whose five base points lie on it: the net meets the conic in \(6=2+1+1+1+1\) assigned intersections, and the conic's strict transform has square \(4-5=-1\). The double point here is a double base point of the \emph{net}, not a double point of the conic. Generic placement is required for this to be the stated homaloidal net.

The intersection count \(C\cdot\omega=\sum i_j\) in the transformation discussion on p.~34 likewise presupposes multiplicity-one branches of \(\omega\) at the assigned centers. For a singular curve the resolved formula is \(\sum i_j\operatorname{mult}_{p_j}\omega_j\), including the infinitely near centers and strict transforms. Adding a positive auxiliary system should be understood as constructing a projective model of the resolved graph: a sufficiently ample auxiliary realizes the required morphisms. Multiplying sections alone on the original surface does not erase their original base ideal.

The formulas \(\varrho=0,\nu=-1\) for the first kind are correctly obtained by Enriques from \(C=E+C_1\), \(E\cdot C_1=1\), and the sum formulas. Explicitly his equations are \(n=(n-1)+\nu+2\) and \(\pi=\pi+\varrho+1-1\), giving \(\nu=-1\) and \(\varrho=0\). In modern terms, \(N_{E/X}\simeq\mathcal O_{\mathbb P^1}(-1)\), so \(E^2=-1\); adjunction gives \(p_a(E)=0\). Conversely, an integral curve of arithmetic genus zero is smooth rational: its normalization sequence gives \(p_a(E)=g(\widetilde E)+\sum\delta_p\), with nonnegative summands. Thus Enriques's assertion that a first-kind curve cannot have a double point is correct. It does not depend on an unproved general smoothing principle.

For the second kind, \(\varrho=0\) on p.~35 must be read with the assigned-point convention of §10 when the original curve is singular. It is \textbf{not} a true assertion that every birationally contractible integral plane curve has arithmetic genus zero. A nodal plane cubic has \(p_a=1\); blow up its node to get a smooth rational curve of square \(9-4=5\), then blow up six smooth points on its strict transform to obtain a \((-1)\)-curve and contract it. The resulting rational map contracts the original cubic to a smooth point. Its adjusted genus is \(1-\binom22=0\), exactly as the assigned-cluster convention predicts. The stated numerical examples \(\nu=0,1,4\) for a ruling fiber, a line, and a conic are their self-intersections on the indicated models.

\subsection{The converse: what is and is not proved on p. 36}

Enriques anticipates Castelnuovo's criterion. Starting with a sufficiently general system \(|C|\), he adds \(E\) until \(C+iE\), where \(i=C\cdot E\), has zero intersection with \(E\). The first addition has \((C+E)\cdot E=i-1\). The calculations

\[
(C+E)^2=n+2i-1,\qquad
p_a(C+E)=\pi+i-1,
\]

are correct, and \(n-2\pi\) increases by one. The crucial assertion is that the addition actually produces new sections and eventually a base-point-free contracting system. Enriques explicitly says that its effectiveness will be justified later. Thus these pages offer a sketch with an acknowledged dependency, not a completed proof of contraction.

A modern completion makes the positivity requirement explicit. Choose a very ample \(H\), sufficiently positive that \(H^1(X,\mathcal O_X(H))=0\), and put \(i=H\cdot E>0\). For \(1\le j\le i\), use

\[
0\to\mathcal O_X(H+(j-1)E)\to\mathcal O_X(H+jE)
\to\mathcal O_{\mathbb P^1}(i-j)\to0.
\]

Inductively, \(H^1(H+jE)=0\), the restriction map is onto, and

\[
h^0(H+jE)-h^0(H+(j-1)E)=i-j+1.
\]

The old sections generate away from \(E\); the lifted sections of \(\mathcal O_{\mathbb P^1}(i-j)\) generate on \(E\). Hence \(|H+iE|\) is generated everywhere and restricts trivially to \(E\). Its map contracts \(E\) and is birational and quasi-finite off \(E\). Take its Stein factorization \(b:X\to Y\). The remaining issue is that \(b(E)\) is smooth, which does not follow just from \(E\cdot(H+iE)=0\).

For that last step, successive infinitesimal neighborhoods of \(E\) have quotients \(\mathcal O_E(-mE)\simeq\mathcal O_{\mathbb P^1}(m)\). Their first cohomology vanishes. The associated graded ring at the contracted point is generated by the two sections of \(\mathcal O_{\mathbb P^1}(1)\), with degree-\(m\) piece of dimension \(m+1\); the formal-functions argument identifies the completed local ring with \(\mathbb C[[u,v]]\). Thus the point is regular and \(b\) is its blowup. This supplies precisely the step missing from a purely numerical contraction argument. Full scheme-theoretic versions of the construction and smoothness argument are Stacks Project, Tags 0C2M and 0C2L. The name Castelnuovo's criterion is standard; this commentary does not attribute the first complete historical proof to the modern online exposition.

\subsection{Birational invariants and the chosen model}

For a smooth projective surface, \(p_g=h^0(K_X)\), \(q=h^1(\mathcal O_X)\), and \(P_m=h^0(mK_X)\) are birational invariants. One can check this for a blowup and then use factorization of birational maps between smooth surfaces. Indeed \(\pi_*\mathcal O_Y=\mathcal O_X\) and \(R^1\pi_*\mathcal O_Y=0\), while any pluricanonical form on \(Y\) is a pluricanonical form away from the exceptional point and extends across that codimension-two point on smooth \(X\). Pullback gives the reverse inclusion. Therefore \(p_a(X)=p_g-q\) is invariant too. By contrast \(K_Y^2=K_X^2-1\); an expression involving \(K_X^2\), such as \(p_a(K_X)=K_X^2+1\), requires its chosen-model or minimal-model convention.

\subsection{Numerical inequality, ruled surfaces, and termination}

For a smooth general member \(D\) of a moving system, adjunction rewrites Enriques's inequality as

\[
N>2\Pi-2\quad\Longleftrightarrow\quad K_X\cdot D<0.
\]

After resolving its base points, \(D\) is nef. An effective pluricanonical divisor would have nonnegative intersection with it. Hence all plurigenera \(P_m(X)\) vanish. Over \(\mathbb C\), the classification identifies these surfaces as birationally ruled, including rational surfaces. Conversely, sufficiently positive sections on a ruled model have square larger than \(2g-2\). Enriques announces that classification implication for later chapters; it is not proved within pp.~1--40. The modern classification is summarized in Beauville (1996, chaps. III, V, and the classification chapters). The stated bound \(n\le2\pi-2\) on a nonruled surface is consequently correct with the same resolved-system convention.

There is also a clean proof of the claim about second-kind curves in its smooth rational formulation. If \(E\simeq\mathbb P^1\) and \(E^2=\nu\ge0\), then \(E\) is nef and \(K_X\cdot E=-\nu-2<0\). Again no pluricanonical divisor can be effective. Alternatively, Enriques's repeated addition is made effective by starting with a sufficiently positive \(H\) as above; restriction degrees now never decrease. Direct calculation gives

\[
(H+jE)^2-2p_a(H+jE)+2
=H^2-2p_a(H)+2+j(\nu+2),
\]

which eventually becomes positive. This justifies his iterative argument on pp.~36--37, without treating arithmetic growth alone as a guarantee of sections. Equivalently, the normal bundle of a smooth rational curve with nonnegative square has vanishing \(H^1\) and positive-dimensional sections; its unobstructed deformations sweep an open subset. The resulting uniruledness implies ruledness over \(\mathbb C\).

For finitely many first-kind curves \(E_j\), one must preserve the original numerical hypothesis when choosing the positive starting divisor. On the smooth model under discussion, the moving divisor \(C\) is nef. Put \(G=\sum_jE_j-K_X\). The source's hypothesis is exactly \(C\cdot G>0\), by adjunction. Choose an ample divisor \(A\) and then \(m\gg0\) so that \((mC+A)\cdot G>0\). Since \(C\) is nef, \(mC+A\) is ample. A sufficiently large multiple \(H=t(mC+A)\) is very ample, has \(H^1(X,\mathcal O_X(H))=0\), and still satisfies \(H\cdot G>0\). Thus replacement by \(H\) retains the needed inequality; an arbitrary positive replacement would not justify that step.

Now set \(i_j=H\cdot E_j\). If the curves are disjoint, simultaneous contraction gives

\[
D=H+\sum_j i_jE_j,\quad
D^2=H^2+\sum_j i_j^2,\quad
p_a(D)=p_a(H)+\sum_j\binom{i_j}{2}.
\]

Thus the displayed hypothesis \(\sum i_j>2\pi-2-n\) implies the desired inequality. Intersecting \((-1)\)-curves cannot be simultaneously interpreted as disjoint blowup fibers. Nevertheless the \emph{existence inequality} need not be discarded: starting with sufficiently positive \(H\), add each \(E_j\) exactly \(i_j\) times. At each stage the restriction degree is at least \(i_j-m\ge0\), augmented by nonnegative intersections with previously added curves, so the same cohomology argument gives a generated system with irreducible general member. Cross terms change its square and genus, but cancel in

\[
-K_X\cdot D=-K_X\cdot H+\sum_j i_j.
\]

Finally, a succession of actual \((-1)\)-contractions terminates because each lowers the Picard number by one. This yields a minimal smooth model without any numerical circularity. Uniqueness for a nonruled surface is a further theorem, not a consequence of termination alone. On a nonruled surface the second-kind phenomenon above is absent, so this is the interpretation of Enriques's model ``without exceptional curves.'' It may still contain rational \((-2)\)-curves. Such curves can contract to singular points and are not exceptions to the smooth-point assertion.

\section{Nota sulle curve eccezionali riducibili}
\label{enriques-I-12}

The correct modern objects are \textbf{total exceptional divisors}, distinguished from the strict transforms of their components. Blow up a sequence of proper or infinitely near points \(p_0=O,p_1=O_1,\ldots,p_s=O_s\). Denote the final strict transforms by \(E_0,\ldots,E_s\), and by \(E_i^*\) the total transform of the exceptional divisor created at step \(i\). At creation each is a \((-1)\)-curve. The total-transform classes always satisfy

\[
E_i^*\cdot E_j^*=-\delta_{ij},\qquad
p_a(E_i^*)=0.
\]

For distinct proper points these cycles are disjoint reduced curves. Here ``connected'' refers to intersections of components; Chapter I does not define numerical \(1\)-connectedness, which would require \(A\cdot B\ge1\) for every decomposition into nonzero effective divisors. Nor does it prove a converse for arbitrary connected divisors with square \(-1\) and arithmetic genus zero. The actual cycles carry the stronger structure of successive blowups and proximity relations. Specializing centers to infinitely near points makes their supports connected and can make the cycles nonreduced. This supplies a rigorous replacement for the ``principle of continuity'': use pullbacks in a sequence of blowups, or a flat family of blowups when a specialization is explicitly desired. Numerical invariance follows immediately from pullback, rather than from an unspecified degeneration.

For a free chain along a smooth branch, the source's formulas become

\[
\Omega_i=E_i^*=E_i+E_{i+1}+\cdots+E_s,\qquad 0\le i\le s.
\]

The graph is a chain with \(s+1\) vertices and \(s\) edges. One has \(E_j^2=-2\) for \(j<s\), \(E_s^2=-1\), adjacent intersections one, and all other intersections zero. Therefore

\[
\Omega_0^2=s(-2)-1+2s=-1,\qquad
p_a(\Omega_0)=\sum_{j=0}^s0+s-s=0.
\]

The scan on p.~38 instead prints \(s-1\) in both calculations and stops the sequence of blowups at \(F_s\), despite listing \(O,O_1,\ldots,O_s\). These are consistent indexing slips, not evidence against the claimed characters. Under the displayed indexing, the last surface must be \(F_{s+1}\) and the edge count must be \(s\). The same shift affects the sentence locating the creation of \(\omega_i\) on \(F_i\).

The degrees \(i-i_1,i_1-i_2,\ldots,i_s\) on p.~37 are \textbf{projective orders with respect to the transforming system}, not self-intersections. On the blown-up surface that system has class \(L=\pi^*A-\sum m_iE_i^*\). For a free chain,

\[
L\cdot E_j=m_j-m_{j+1},\quad j<s,\qquad
L\cdot E_s=m_s.
\]

These are degrees of the restricted line bundles. As in §6, if the induced map on \(E_j\) has degree \(d_j\) onto its image, the precise equality is \(L\cdot E_j=d_j\deg\varphi(E_j)\). The printed projective orders therefore require degree one on each component; sufficiently strong separation of the relevant jets supplies that condition. Positivity of the differences alone does not. Subject to this mapping-degree qualification, the formulas give the stated telescoping degree sum. They also prevent a mistaken identification of these positive orders with the negative virtual degrees computed on p.~38.

The first satellite example is equally exact. Blow up \(p_0\), then \(p_1\in E_0\), then the intersection of the two resulting exceptional components. On the final surface,

\[
E_0^2=-3,\quad E_1^2=-2,\quad E_2^2=-1,\qquad
E_0E_1=0, E_0E_2=E_1E_2=1,
\]

and

\[
\Omega_0=E_0+E_1+2E_2,\quad
\Omega_1=E_1+E_2,\quad\Omega_2=E_2.
\]

For assigned multiplicities \(m_0,m_1,m_2\), the same example gives \[
L\cdot E_0=m_0-m_1-m_2,\qquad
L\cdot E_1=m_1-m_2,\qquad L\cdot E_2=m_2.
\] In particular the first order drops from \(m_0-m_1\) to \(m_0-m_1-m_2\) and the new component has restricted degree \(m_2\), with the same mapping-degree qualification as above. The factor two occurs because the last center has multiplicity two on the previous total exceptional divisor. Directly,

\[
\Omega_0^2=-3-2-4+4+4=-1,\qquad
K_Y\Omega_0=1+0-2=-1,\qquad p_a(\Omega_0)=0.
\]

The source correctly prints \(\omega+\omega_1+2\omega_2\) on p.~38 but incorrectly says on p.~39 that the transformed \(\omega_1\) must be counted twice. The latter must be \(\omega_2\). This is a scan-confirmed source slip.

In general, write \(p_j\succ p_i\) when \(p_j\) lies on the current strict transform of the exceptional divisor created at \(p_i\). Then

\[
E_i=E_i^*-\sum_{p_j\succ p_i}E_j^*.
\]

If \(t_i\) later centers are proximate to \(p_i\), this gives \(E_i^2=-1-t_i\). For the degree of the restricted line bundle, and hence for projective order under the same degree-one qualification, it gives the proximity expression

\[
L\cdot E_i=m_i-\sum_{p_j\succ p_i}m_j.
\]

The triangular proximity matrix is invertible over the integers; its inverse expresses every total transform as an effective weighted sum of strict components. This provides a reproducible calculation with Gram matrix \(-I\), and tests free and satellite configurations alike. The final intersection graph must be updated at every blowup: blowing up a smooth point of one component appends an edge; blowing up an intersection removes that edge and inserts the new component between its two endpoints. This is the precise version of Enriques's statement that \(\omega_i\) may meet \(\omega_{i+t}\) rather than its chronological successor. Merely knowing that two centers were consecutive does not determine the final adjacency.

The coefficients in \(\Omega_0=\sum h_iE_i\) have a simple recursion. The initial coefficient is one; a free blowup inherits the coefficient of the component through its center, while a satellite blowup receives the sum of the two coefficients through its center. To interpret Enriques's minimum-order rule, choose a curvette transverse to \(E_i\) away from the other exceptional components. Its image is a branch through the specified chain, with multiplicity \(h_i\) at \(O\). Any branch passing through that chain has at least this multiplicity; this follows by intersecting its strict transform with the total transform of a general line through \(O\). Thus the coefficient equals the minimum possible order of such a branch. Apply the same construction starting at \(p_i\) to obtain every \(\Omega_i\).

Enriques attributes the general analysis to Barber and Zariski. The verified citation is \textbf{Barber, S. F., and Zariski, O. (1935), ``Reducible exceptional curves of the first kind,'' American Journal of Mathematics, 57(1), 119--141, DOI 10.2307/2372025}. The printed ``E. Barber'' and ``Ann. Journ. of Math.'' are bibliographic errors. The paper's identity and bibliographic data were verified; its full proof was not consulted, so the recurrence above is supplied directly rather than represented as a quotation or page-specific paraphrase of that paper. Casas-Alvero (2000, chaps. 3--4) is a verified modern reference for the infinitely near and virtual-multiplicity framework.

There is a further mixed transcription/source issue in the footnotes. After ``Enriques'' the enlarged scan reads \textbf{(2)}, whereas \textbf{(3)} is the intended reference. Note (2) explains how to regard an infinitely near point as proper; note (3) is the Enriques--Chisini reference. The printed (2) is thus a cross-reference slip; the edition retains it, with an uncertainty note. A faithful edition should print the source's (2), flag it, and identify (3) as the intended reference. The superscript in \(\omega_i^{h_i}\) really is \(h_i\). The odd ``con-le'' in the final paragraph is present in the scan.

\textbf{Relative canonical interpretation.} The final sum on p.~39 has a natural interpretation not yet given there:

\[
\Omega_0+\cdots+\Omega_s=K_Y-\pi^*K_X.
\]

It is the relative canonical divisor of the composite blowup. This follows by applying the one-point canonical formula at each step. In the satellite example it is \(E_0+2E_1+4E_2\). Its square is \(-(s+1)\) and its arithmetic genus is \(-s\), as expected for a flat limit of \(s+1\) disjoint rational curves. Connected nonreduced support therefore does not force arithmetic genus to be nonnegative. This is a useful preparatory example for the canonical systems of the next chapter; it is an explanatory extension, not a claim of a new result.

\section{Serie caratteristica virtuale}
\label{enriques-I-13}

For a smooth effective curve \(C\subset X\), its intrinsic virtual characteristic series is

\[
|N_{C/X}|=|\mathcal O_C(C)|,\qquad \deg N_{C/X}=C^2.
\]

The characteristic series obtained by deforming \(C\) \emph{within its complete linear system} is the projectivized image of

\[
H^0(X,\mathcal O_X(C))/\mathbb C\longrightarrow H^0(C,N_{C/X}).
\]

The two need not be equal on an irregular surface. The defining sequence is

\[
0\to\mathcal O_X\to\mathcal O_X(C)\to N_{C/X}\to0.
\]

Its boundary \(\delta:H^0(N_{C/X})\to H^1(\mathcal O_X)\) measures the missing directions. Thus the defect in vector-space dimension is \(\dim\operatorname{im}\delta\le q\), with equality if \(H^1(X,\mathcal O_X(C))=0\). This makes the relation to the front matter's irregularity discussion precise.

Enriques constructs the virtual series using an auxiliary \(|K|\), restricting \(|C+K|\) to \(C\), imposing the divisor \(C\cap K_0\) for a fixed \(K_0\in|K|\), and removing it. The words ``passes through the intersections'' must include this removal: the unreduced degree would be \(C^2+C\cdot K\). After removal, the underlying line bundle is intrinsically \(N_{C/X}\). His final sentence about choosing \(K\) sufficiently large is justified by

\[
0\to\mathcal O_X(K)\to\mathcal O_X(C+K)
\to\mathcal O_C(C+K)\to0.
\]

For \(K=mH\), \(m\gg0\), Serre vanishing gives \(H^1(X,\mathcal O_X(K))=0\), so restriction is surjective and the residual series is the entire \(|N_{C/X}|\). This is a correct modern completion of the existence claim (Stacks Project, Tag 0B5U).

There is, however, a real qualification to the preceding independence assertion. Arbitrary auxiliary systems satisfying only the two stated irreducibility conditions can produce different \textbf{subseries}. If \(k\in H^0(X,\mathcal O_X(K))\) defines \(K_0\), a diagram of the two exact sequences shows that the vector space of residual sections obtained from this \(K_0\) is

\[
V_K=\{s\in H^0(N_{C/X}):
 k_*\delta(s)=0\text{ in }H^1(X,\mathcal O_X(K))\}.
\]

It equals all of \(H^0(N_{C/X})\) when \(H^1(K)=0\), but this need not happen for a small auxiliary pencil. Consequently the line bundle, its complete series, or the eventual sufficiently positive construction is independent; a literal assertion about every actually cut subseries is too strong. Comparing three sum systems does not by itself prove that all their restriction maps are surjective.\footnote{Independence is proved above after explicitly choosing a sufficiently ample auxiliary; the counterexample below tests the preceding, unrestricted sentence. On p. 40 Enriques first says ``Questa serie invero riesce indipendente dalla scelta di \(|K|\)'' and only subsequently says that \(|K|\) may be chosen to cut a complete series. The scan confirms that order. Thus the positive-auxiliary proof and the counterexample concern different assertions. The correct invariant is the completed intrinsic series; literal equality of every cut subseries is false.}

\textbf{Worked example.} Let \(E\) be an elliptic curve, \(f:E\to\mathbb P^1\) a map of degree two, \(A=f^*\mathcal O_{\mathbb P^1}(1)\), and \(X=E\times\mathbb P^1\). Let \(C\) be the graph of \(f\), and let \(|K|\) be the pencil \(E\times\{t\}\). Then

\[
\mathcal O_X(C)=p_1^*A\otimes p_2^*\mathcal O(1),\quad
C^2=4,\quad N_{C/X}=A^2,
\]

\[
h^0(X,\mathcal O_X(C))=2\cdot2=4,\qquad
h^0(C,N_{C/X})=4.
\]

The ordinary characteristic image has vector dimension three after quotienting by the section defining \(C\). Multiplication \(H^1(\mathcal O_X)\to H^1(\mathcal O_X(K))\) by \(k\) is injective, by the product cohomology formula. Therefore this auxiliary pencil still yields vector dimension three, whereas a sufficiently ample auxiliary system yields four. Both \(|K|\) and \(|C+K|\) have irreducible general members; the latter system is base-point-free with a two-dimensional image, so its general member is smooth and irreducible by Bertini. For an explicit double cover \(E:y^2=q_4(x)\), the three pullback sections \(1,x,x^2\) lie in \(H^0(A^2)\), while \(y\) supplies its fourth direction. This verifies the defect without any historical attribution.

The hypothesis \(n=C^2>0\) on p.~40 is not needed to define the normal bundle. Nor does positive degree alone ensure a nonempty complete linear series on a curve of high genus. The modern definition allows an empty projectivized space, degree zero, and negative degree. For Cartier divisors that are singular or nonreduced, \(\mathcal O_C(C)\) and its normal sheaf still make sense; interpreting a linear series via moving effective point divisors requires the usual additional care about sections which vanish on a component.

\begin{example}[Positive square with empty virtual characteristic series]
Choose a smooth curve \(B\) of genus \(g\ge2\) and a line bundle \(N\) of degree \(1\le d<g\) with \(H^0(B,N)=0\). On the ruled surface \(\mathbb P(\mathcal O_B\oplus N)\), using the quotient convention, the section corresponding to the quotient onto \(N\) has normal bundle \(N\), square \(d>0\), and no infinitesimal normal section. This example gives the precise reason that positive virtual degree does not guarantee an effective characteristic series.
\end{example}

Finally, first-order embedded deformations of a Cartier divisor have tangent space \(H^0(N_{C/X})\). A family \(C_t\) differentiates to such a section, and the characteristic divisor is its zero divisor when defined. This proves the local meaning of the final paragraph's ``infinitely near curves.'' Locally a deformation over \(\mathbb C[\varepsilon]/(\varepsilon^2)\) is described by an \(\mathcal O_X\)-linear map \(\mathcal I_C\to\mathcal O_C\), factoring through \(\mathcal I_C/\mathcal I_C^2\); gluing identifies these with \[
H^0\bigl(\mathcal{H}om(\mathcal I_C/\mathcal I_C^2,\mathcal O_C)\bigr)
=H^0(N_{C/X}).
\] The normal bundle and divisor sequence are The Stacks Project Authors (2026), Tags 0C4S and 0B3P; the local deformation interpretation and its gluing obstructions are treated by Hartshorne (2010, Chapters 1--2). It does \textbf{not} assert that every tangent vector integrates to an actual smooth family, or that every family has complete characteristic series. Obstructions and semiregularity are separate questions; the precise modern treatment is the normal-sheaf and Picard-scheme framework discussed in Kleiman (2005, Remark 5.18), together with the sources cited there. Thus p.~40 introduces the object needed for the later completeness problem, rather than solving that problem.

The historical notes distinguish three attributions. Enriques cites his 1896 \emph{Introduzione}, §§15--16, for virtual characters; Severi's 1904 paper for the virtual characteristic series; and a Castelnuovo communication dated 25 December 1898, no. 4, for the implicit earlier appearance. The verified full Severi reference is \emph{Osservazioni sui sistemi continui di curve appartenenti ad una superficie algebrica}, \emph{Atti della Reale Accademia delle Scienze di Torino}, 39, 376--392. The specific letter and its text were not independently located, so that claim remains explicitly Enriques's attribution. His genus-two symmetric-product example concerns the surface of unordered pairs of points of a genus-two curve; it is birational to its Jacobian by the degree-two Abel map, with the canonical pencil contracted. The old adjective ``iperellittica'' here must not be identified with the modern bielliptic-surface class.

\subsection{Semiregularity and the later completeness theorem}

For an effective Cartier divisor \(C\), put \(N=\mathcal O_C(C)\) and consider the two connecting maps

\[
H^0(N)\xrightarrow{\delta}H^1(\mathcal O_X),\qquad
H^1(N)\xrightarrow{\sigma}H^2(\mathcal O_X).
\]

The first measures the difference between deformations within one linear system and deformations which change the line bundle. The second is the semiregularity map: \(C\) is \textbf{semiregular} when \(\sigma\) is injective. These are different assertions.

Set

\[
R=q+h^0(\mathcal O_X(C))-1-h^1(\mathcal O_X(C)).
\]

The long exact sequence itself gives

\[
h^0(N)=R+\dim\ker\sigma.
\]

To check this without an expected-dimension shortcut, write \(a\) for the rank of \(H^1(\mathcal O_X)\to H^1(\mathcal O_X(C))\). Exactness gives \(h^0(N)=h^0(C)-1+q-a\) and \(\dim\ker\sigma=h^1(C)-a\), which is the asserted identity.

The additional deformation theorem says, over \(\mathbb C\), that semiregularity implies smoothness of the divisor scheme at \([C]\). Obstruction classes lie in \(H^1(N)\) and have zero image under \(\sigma\); injectivity therefore kills them. The assertion about their image is a genuine theorem, not a consequence of the preceding dimension count: Mumford (1966, pp.~157--159) gives the obstruction calculation, and Kleiman (1973, theorem on p.~307; 2005, Remark 5.18) gives the algebraic formulation. Together with the tangent-space identification \(T_{[C]}\operatorname{Div}(X)=H^0(N)\), this yields the precise equivalence

\[
C\text{ semiregular}
\quad\Longleftrightarrow\quad
\operatorname{Div}(X)\text{ smooth of dimension }R\text{ at }[C].
\]

The dimension qualification cannot be omitted. In the sufficiently positive case actually used in the front-matter proof, \(H^1(\mathcal O_X(C))=0\); here the Abel map to the smooth Picard scheme is smooth, and the required family is already explicitly constructed as a projective bundle. Thus that argument does not assume the full semiregularity theorem in order to manufacture the family it needs. Grothendieck (1961, Exposé 221, Corollary 5.4 and §§5.5--5.6) directly treats the normal-sheaf tangent space and the two relevant vanishing cases. Page 40 of Enriques introduces the characteristic object; none of these later deformation-theoretic conclusions is claimed to be proved on that page.

The auxiliary formula \(V_K=\ker(k_*\delta)\) also leaves a concrete exercise: for a fixed \((X,C)\), determine the effective auxiliaries whose multiplication map annihilates \(\operatorname{im}\delta\). This is a problem about specified cohomology maps, not an asserted open theorem. The surviving historical ambiguity is whether Enriques intended an implicitly completed series or an unstated positivity condition in the earlier independence sentence; either reading must retain the distinction proved above. Similarly, general contraction to singular points requires more than Castelnuovo's smooth-point criterion: rational \((-2)\)-curves and the nonreduced cycles of Section 12 already mark that boundary.

\chapter{Sistemi covarianti e invarianti}

The divisor theory of Chapter I now acquires its intrinsic covariant class. Enriques constructs the Jacobian system, subtracts three copies of the original system to obtain the canonical class, and uses adjunction to define the geometric genus, linear genus, and plurigenera (pp.~41--65; p.~66 is blank). The governing identities are \[
J(D)\sim K_X+3D,\qquad D'=K_X+D.
\] They are identities of line-bundle classes, including virtual classes without effective representatives. In the notation used throughout Part 1, \(X\) is smooth and projective over \(\mathbb C\), \(K_X\) is its canonical class, and \(p_a(D)=1+(D^2+K_X\cdot D)/2\). Enriques's relative characters are \(p^{(1)}=1+K_X^2\) and \(p^{(2)}=K_X^2\); his surface genus is \(p_g=h^0(X,\omega_X)\). His ``degree'\,' in these formulas is self-intersection, with assigned base conditions resolved, rather than projective order. The distinctions between complete systems and their restricted series, between strict and total transforms, and between discrepancy and inherited fixed components remain essential.

\section[Curve jacobiane]{Curve jacobiane (pp.~41--42)}\label{commentary-II-1}

\subsection{Definition and Enriques's argument}

A net is a projective two-dimensional linear system. Its Jacobian is the locus where a member has a double point, equivalently where two members are tangent. At a point where the evaluation and first-derivative conditions have their expected rank, a linear combination of two tangent local equations has zero value and zero differential, proving the equivalence. This is a determinantal construction: its multiplicities, not just its reduced support, are essential.

For plane curves of degree \(n>1\), Enriques states that the Jacobian has degree \(3n-3\). At an ordinary assigned base point of multiplicity \(i\), its multiplicity is generically \(3i-1\). A simple base point gives a double point of the Jacobian; a common tangent there gives generically a triple point with that tangent. ``In general'' is essential: further infinitesimal conditions can raise multiplicity. The planar multiplicities are imported from Enriques and Chisini (1918, Book III, chapter I, Section 5); the surface extension is sketched by osculation. Transfer to a smooth surface is justified by finite-order osculation and rational local models. An exact local-coordinate calculation makes the required finite-order argument precise.

For a fixed irreducible system of projective dimension greater than two, Enriques proves linear equivalence of its defined Jacobian divisors. Two nets sharing a pencil span a projective three-dimensional subsystem. Nets through that pencil form a \(\mathbb{P}^1\); their Jacobians form a rational pencil of curves. Two nets sharing only one curve are compared through an auxiliary net sharing a pencil with each. Nets with no common member are linked through a third sharing one member with each. Transitivity finishes the proof. It uses the earlier result that a rational pencil is linear.

An implicit hypothesis must be retained in ``arbitrary nets'': each Jacobian under comparison must be a divisor, not an identically zero determinant. An irreducible net over \(\mathbb C\) gives a generically finite map, but a special subnet of a larger irreducible system may be composed with a pencil. For example \(\langle x^2,xy,y^2\rangle\subset H^0(\mathbb{P}^2,\mathcal O(2))\) has identically zero Jacobian. This identifies a genericity convention; it does not disprove equivalence of the defined Jacobian divisors.

\subsection{Modern proof by principal parts}

For \(L=\mathcal O_X(D)\), the first-principal-parts sequence is \[
0\longrightarrow\Omega_X^1\otimes L\longrightarrow P_X^1(L)
\longrightarrow L\longrightarrow0.
\] The middle bundle has rank three and determinant \(\omega_X\otimes L^3\). See the Stacks Project (n.d., Section 18.34, Lemma 18.34.6) for the principal-parts sequence. Three sections \(V\subset H^0(X,L)\) give \[
V\otimes\mathcal O_X\longrightarrow P_X^1(L).
\] Its determinant, when nonzero, defines a divisor in \(|K_X+3D|\). In local coordinates and a trivialization of \(L\), its equation is \[
\det\begin{pmatrix}
f_0&f_1&f_2\\
\partial_x f_0&\partial_x f_1&\partial_x f_2\\
\partial_y f_0&\partial_y f_1&\partial_y f_2
\end{pmatrix}=0.
\] Thus the result concerns a line-bundle class, not merely a numerical class. On \(\mathbb{P}^2\), \(K_X=-3H\), giving degree \(3n-3\).

At an \(i\)-fold assigned point, the naïve lowest determinant order is \(3i-2\). Euler's relation on the initial homogeneous degree-\(i\) terms makes that leading determinant vanish, so the generic order becomes \(3i-1\). Equivalently, on the blowup, \[
\widetilde D=\pi^*D-iE,\quad
K_{\widetilde X}=\pi^*K_X+E,\quad
K_{\widetilde X}+3\widetilde D=\pi^*(K_X+3D)-(3i-1)E.
\] Excess actual multiplicity does not change this virtual prescribed coefficient.

The same determinant explains the dimension extension stated without proof on p.~64. On a smooth variety of dimension \(d\), the first-principal-parts bundle has rank \(d+1\), with determinant \(\omega_X\otimes L^{d+1}\). A nonzero determinant of \(d+1\) sections has class \(K_X+(d+1)D\); a common fixed factor occurs \(d+1\) times. This calculation assumes that the determinant is nonzero. In positive characteristic inseparability can make it vanish identically, and the homogeneous proof below also divides by degrees. These qualifications delimit an extension of the argument; the commentary remains in characteristic zero.

\section[Sistema jacobiano]{Sistema jacobiano (pp.~42--48)}\label{commentary-II-2}

\subsection{Completion and infinitely near conditions}

Enriques defines the completion and then gives local, analytic, and projective proofs of the fixed-component rule. The complete Jacobian system contains the Jacobians already constructed, imposing multiplicity \(3i-1\) at each assigned \(i\)-fold base point. Without assigned base conditions, completion adds none. This does not imply that the completed system has no actual base points.

For successive simple points \(O,O_1\), with \(O_1\) specifying the common tangent, Enriques distinguishes an actual triple point at \(O\) plus a simple point at \(O_1\) from the virtual pair of double points. On the twice-blown-up surface, let \(E_1^*\) be the total transform of the first exceptional divisor and \(E_2\) the second. The intrinsic Jacobian class is \[
\pi^*(K_X+3D)-2E_1^*-2E_2.
\] An actual strict Jacobian with multiplicities \(3,1\) has class \[
\pi^*(K_X+3D)-3E_1^*-E_2.
\] Their difference is \(E_1^*-E_2\), the strict first exceptional curve. This exactly explains the virtual adjustment. It does not guarantee an actual enlargement with precisely two ordinary double points; Enriques explicitly allows the enlargement to fail.

Adding a fixed curve \(A\) to a net yields \[
J(A+D)=3A+J(D).
\] Enriques calls this fixed curve \(K\) temporarily, before introducing the canonical notation. The coefficient is exactly three unless the residual Jacobian already contains \(A\). Locally replace every \(f_i\) by \(af_i\) in the jet determinant. Row subtraction on \(a\ne0\) gives \(a^3J(f_0,f_1,f_2)\), an identity of regular expressions extending everywhere. This proves the statement without osculating surfaces.\footnote{The reference printed on p. 43 reads ``Libro XV, cap. IV, \S~39, vol. II, pag. 634 e seg.'' The scan confirms volume II; changing it to III would be an unsupported emendation. The unusual Book XV locator remains a bibliographic question, not a premise of the jet proof.}

\subsection{Analytic verification and the printed determinant error}

Let \(Y=(f=0)\subset\mathbb{P}^3\) be smooth of degree \(n\); let \(\varphi_i\) have degree \(m\) and \(\varphi\) degree \(s\). Enriques chooses an auxiliary curve \(L\), not containing the fixed component, such that \(C+L\) is cut by degree-\(m\) hypersurfaces. The net is represented by \[
\varphi(\lambda_1\varphi_1+\lambda_2\varphi_2+\lambda_3\varphi_3)=0.
\] Completeness of these residual systems is assumed from the completeness theorem for adjoints, expressly referred forward to Chapter III. For a smooth hypersurface, completeness itself follows from \(0\to\mathcal O_{\mathbb{P}^3}(m-n)\to\mathcal O_{\mathbb{P}^3}(m)\to\mathcal O_Y(m)\to0\) and the vanishing of the intervening \(H^1\); sufficiently positive twisting supplies the auxiliary residual divisor avoiding the chosen fixed component.

Write \([g_1,g_2,g_3,g_4]\) for the determinant with their four gradients as columns. The first determinant on p.~44 is \[
D=[\varphi\varphi_1,\varphi\varphi_2,\varphi\varphi_3,f].
\] In its bottom row, the third entry has misplaced parentheses. It must read \(\partial(\varphi\varphi_3)/\partial x_4\), in place of the printed \((\partial\varphi\varphi_3)/\partial x_4\). Multilinearity gives \[
D=\varphi^3 A+\varphi^2\Phi,\qquad A=[\varphi_1,\varphi_2,\varphi_3,f],
\] \[
\begin{aligned}
\Phi={}&\varphi_1[\varphi,\varphi_2,\varphi_3,f]\\
&+\varphi_2[\varphi_1,\varphi,\varphi_3,f]\\
&+\varphi_3[\varphi_1,\varphi_2,\varphi,f].
\end{aligned}
\] Terms with two columns proportional to \(d\varphi\) vanish. After the stated column exchanges, these three terms are the first-row expansion of \[
\Phi=\det\begin{pmatrix}
\varphi_1&\varphi_2&\varphi_3&0&0\\
\partial_1\varphi_1&\partial_1\varphi_2&\partial_1\varphi_3&\partial_1\varphi&\partial_1f\\
\partial_2\varphi_1&\partial_2\varphi_2&\partial_2\varphi_3&\partial_2\varphi&\partial_2f\\
\partial_3\varphi_1&\partial_3\varphi_2&\partial_3\varphi_3&\partial_3\varphi&\partial_3f\\
\partial_4\varphi_1&\partial_4\varphi_2&\partial_4\varphi_3&\partial_4\varphi&\partial_4f
\end{pmatrix}.
\] The plus in the last column of the fourth row corrects a printed minus.\footnote{The p. 45 scan has \(-\partial f/\partial x_3\) in that entry. Retaining it would contradict both the preceding determinant expansion and the following Euler operation. The correction is mathematical; the transcription faithfully preserves the misprint.}

Multiply the first row by \(m\), then subtract \(x_1,\ldots,x_4\) times the derivative rows. Euler's identity gives \[
(0,0,0,-s\varphi,-nf).
\] The resulting determinant is \(m\Phi\), vanishing when \(f=\varphi=0\). More precisely, with \(B=[\varphi_1,\varphi_2,\varphi_3,\varphi]\), expansion gives \[
m\Phi=s\varphi A-nfB,
\] \[
D=\frac{m+s}{m}\varphi^3A-\frac nm f\varphi^2B,\qquad
D|_Y=\frac{m+s}{m}\varphi^3A|_Y.
\] This exact factorization streamlines the analytic argument. It shows why a fourth copy occurs only when \(A|_Y\) already contains the fixed component. Division by \(m\) and the nonvanishing of \(m+s\) use characteristic zero.

The last \(K\) in the opening paragraph of p.~46 is unbarred in the scan; an overline in the earlier transcription has no mathematical significance. On pp.~46--47 Enriques instead differentiates at \(P\) with \(f(P)=\varphi(P)=0\). The first three differentiated-column determinants vanish because the first row is zero; the other two give \[
m(\partial_i\Phi)_P=s(\partial_i\varphi)_P A(P)-n(\partial_i f)_P B(P).
\] If the tangent planes coincide, write \[
(\partial_i\Phi)_P=\rho(\partial_i f)_P,\quad i=1,\ldots,4.
\] Then \[
(m\rho+nB(P))df_P=sA(P)d\varphi_P.
\] If \(A(P)\ne0\), \(df_P,d\varphi_P\) are proportional, contradicting the generic transversality of the simple fixed component. Otherwise \(A(P)=0\), precisely the residual-Jacobian alternative. Multiplicity is tested at its generic smooth point; special singular or tangency points do not invalidate that argument.

\subsection{The projective verification}

The last proof uses four degree-\(m\) forms: \[
\lambda_1\varphi_1+\lambda_2\varphi_2+\lambda_3\varphi_3+\lambda_4f\psi=0,
\qquad\deg\psi=m-n.
\] Their ambient Jacobian has degree \(4m-4\). Its restriction to \(Y\) contains \(\psi\cap Y\), of order \(n(m-n)\), besides the net Jacobian. Hence the latter has order \[
n(4m-4)-n(m-n)=n(3m+n-4).
\] The text chooses \(m>n\); the displayed construction only needs \(m\ge n\), which is achieved by enlarging the auxiliary residual choice. Multiplying by the fixed degree-\(s\) form gives four copies in the ambient Jacobian and removes one in the extraneous intersection, leaving three. The resulting order is \[
n\{3(m+s)+n-4\}.
\] Intrinsically \(K_Y=(n-4)H\), \(H^2=n\), and \((K_Y+3mH)H=n(3m+n-4)\), confirming the count.

\section[Teorema fondamentale]{Teorema fondamentale (pp.~48--49)}\label{commentary-II-3}

For systems without assigned base conditions, the fixed-component calculation in both orders gives \[
J(C+L)\sim J(C)+3L\sim3C+J(L),\qquad
J(C)-3C\sim J(L)-3L.
\] Enriques extends to a pencil or isolated curve by defining \[
J(C):=3C+J(L)-3L
\] using a sufficiently ample auxiliary \(L\). This may have no effective representative; independence of \(L\) follows from the proved relation.

If \(C\) has an assigned \(i\)-fold point \(O\) and \(L\) an assigned \(s\)-fold point \(O'\), distinct, let \(E,E'\) be their exceptional divisors on a common blowup. Then \[
\widetilde C=\pi^*C-iE,\quad \widetilde L=\pi^*L-sE',
\] \[
\widetilde J_C=\pi^*(K_X+3C)-(3i-1)E,\quad
\widetilde J_L=\pi^*(K_X+3L)-(3s-1)E'.
\] The equality is \[
3\widetilde C+\widetilde J_L+E
\sim3\widetilde L+\widetilde J_C+E'.
\] Here the symbols \(\widetilde J_C,\widetilde J_L\) mean transforms of the Jacobian classes built with each system's own assigned conditions. They are not the intrinsic Jacobian classes of the two transformed systems on the common blowup \(Y\): those are respectively \(K_Y+3\widetilde C=\widetilde J_C+E'\) and \(K_Y+3\widetilde L=\widetilde J_L+E\). This distinction explains the exceptional terms and avoids treating two different constructions as identical.

This is the displayed \(|3C+L_j+O|=|3L+C_j+O'|\) on p.~48. A point symbol records the exceptional correction.\footnote{The prose on pp. 48--49 associates the points with the two systems in the opposite order under ``rispettivamente.'' The scan confirms the wording slip. The displayed formula is correct, as the exceptional coefficients above show; exchanging its points would create an error.}

Define \[
K_X=J(C)-3C,\qquad C'=J(C)-2C=C+K_X.
\] For plane degree-\(m\) curves with prescribed \(i\)-fold points, adjoints have degree \(m-3\) and prescribed multiplicities \(i-1\), because \[
K_X+\widetilde C=(m-3)H-\sum(i-1)E_i.
\] The fundamental adjunction identity is \[
C'+L=L'+C.
\] It holds in the Picard group even when some systems are empty. Enriques also says that adjoints will generally be effective for sufficiently positive systems and leaves the relevant bounds to later theory. A precise modern sufficient condition is \(C=mH\) with \(H\) ample and \(m\gg0\): Serre vanishing and surface Riemann--Roch give \(h^0(K_X+mH)>0\). The passage should not be converted into a theorem about every divisor of large arithmetic genus without positivity hypotheses. Absence of assigned base conditions does not erase the actual common points of a complete system.

\section[Curve canoniche]{Curve canoniche (pp.~50--52)}\label{commentary-II-4}

\subsection{Birational covariance and the hypersurface description}

The covariance argument is proved from the fundamental relation; the projective adjoint description is a sketch whose completeness is explicitly deferred. Enriques compares canonical classes constructed from systems on birational surface models. Without fundamental points or exceptional curves they transform to the same system. Under a blowup, the transformed canonical system acquires the exceptional divisor once. This is covariance of divisor classes: an individual canonical divisor's strict transform need not be canonical on the blowup.

For \(Y_n\subset\mathbb{P}^3\) with ordinary singularities, the Jacobian of a plane-section net is cut by a polar of degree \(n-1\), removing the double curve. Subtracting three hyperplane sections suggests adjoints of degree \(n-4\). Enriques discusses forced cusp or pinch points of the polar construction: these belong to that subsystem, not to the prescribed conditions of its completion. Degree-\(n-4\) adjoints through the double curve cut canonical curves outside it.

That these hypersurfaces cut the \emph{complete} canonical system is expressly anticipated from Chapter III, as is completeness in other degrees. Polar counting alone does not prove it. It would be inaccurate to present this declared dependency as a concealed gap.

\subsection{Modern proof for ordinary singular models}

Let \(\nu:X\to Y\subset\mathbb{P}^3\) be the smooth normalization of a hypersurface with ordinary double-curve, triple-point, and pinch-point singularities. Let \(\mathfrak c\subset\mathcal O_Y\) be the conductor and \(\Delta\) its divisor upstairs. Finite duality gives \[
\nu_*\omega_X\simeq
\mathcal{H}om_F(\nu_*\mathcal O_X,\omega_Y)
\simeq\mathfrak c\otimes\mathcal O_Y(n-4),
\] equivalently \[
K_X\sim(n-4)H-\Delta.
\] For these ordinary singularities \(\mathfrak c\) is the ideal of the double curve with its proper scheme structure. Dolgachev (2016, Section 2, pp.~3--4) explicitly verifies this conductor formulation and the ordinary local models. This supplies the modern justification without assigning historical priority for the repair.

If \(\Gamma\subset\mathbb{P}^3\) denotes the conductor subscheme on \(Y\), then \[
0\to\mathcal O_{\mathbb{P}^3}(d-n)\to\mathcal I_\Gamma(d)
\to\mathfrak c(d)\to0.
\] Since \(H^1(\mathbb{P}^3,\mathcal O(k))=0\) for every \(k\), degree-\(d\) adjoints cut the complete system \[
\bigl|K_X+(d-n+4)H\bigr|.
\] Local duality identifies the correct ideal; projective-space cohomology proves completeness. At \(d=n-4\) this is \(|K_X|\).

Enriques also states the orders \(i-1\) along an ordinary \(i\)-fold curve and \(s-2\) at an ordinary isolated \(s\)-fold conical point, without their full local proofs. The isolated-point rule follows by blowing up the ambient smooth threefold: its canonical divisor gains \(2E\), while the strict surface subtracts \(sE\). Adjunction gives discrepancy \((2-s)E\), requiring the adjoint numerator to vanish \(s-2\) times. This assumes an ordinary point with smooth projectivized tangent cone. Multiplicity alone does not determine the adjoint ideal at an arbitrary singularity.

For a sequence of blowups \(f:X\to Y\), \[
K_X=f^*K_Y+K_{X/Y}.
\] Enriques's ``impure'' system includes the discrepancy divisor. His later precise ``pure'' system is the pullback of the canonical system on the contracted model. It is \textbf{not} obtained by removing every fixed component: base points of \(|K_Y|\) may force further exceptional components in \(f^*|K_Y|\), which he retains as pure.

\section[Proprietà delle curve canoniche]{Proprietà delle curve canoniche (pp.~52--53)}\label{commentary-II-5}

\subsection{The characteristic-series proof}

Let \(C\) be a smooth general member of a mobile system, of genus \(\pi\) and self-intersection \(n\). Its characteristic pencil has degree \(n\). Curve Riemann--Hurwitz gives ramification degree \(2n+2\pi-2\) and class \(K_C+2C|_C\). Enriques identifies this ramification divisor with the cut of the net Jacobian. Subtracting two characteristic divisors shows that adjoints cut canonical groups; subtracting three gives the characteristic residual: \[
(K_X+C)|_C\simeq K_C,\qquad K_X|_C\simeq K_C-C|_C.
\] Consequently \[
C\cdot C'=2\pi-2,\qquad C\cdot K_X=2\pi-2-n.
\] This proves the assertion in the stated smooth, mobile setting, using curve Riemann--Hurwitz as an earlier result.

For arbitrary curves Enriques adds a sufficiently positive auxiliary \(L\), so that \(C+L\) has a mobile irreducible representative, and subtracts the formulas for \(L\). If \(p_a(C)=\pi\), \(p_a(L)=\rho\), and \(CL=s\), then \[
p_a(C+L)=\pi+\rho+s-1,\qquad(C+L)^2=C^2+2s+L^2.
\] Cancellation yields the same formula for \(C\). For singular or reducible curves \(\pi\) is the arithmetic/virtual genus, not normalization genus.

\subsection{Modern adjunction and its limits}

For smooth \(C\), determinants in \[
0\to\mathcal O_C(-C)\to\Omega_X^1|_C\to\Omega_C^1\to0
\] give \(\omega_C\simeq(\omega_X(C))|_C\). The same identity using the dualizing sheaf holds for any effective Cartier divisor on a smooth surface. Riemann--Roch supplies the virtual version: \[
\chi(\mathcal O_X(D))=\chi(\mathcal O_X)+\frac{D(D-K_X)}2.
\] From \(0\to\mathcal O_X(-D)\to\mathcal O_X\to\mathcal O_D\to0\), \[
\chi(\mathcal O_D)=-\frac{D(D+K_X)}2,\qquad
p_a(D)=1+\frac{D(D+K_X)}2.
\] References are Hartshorne (1977), II, Proposition 8.20, and V, Theorem 1.6; Brosnan's course notes (n.d., Sections 3.1--3.2) give these statements and textbook cross-references.

The line-bundle identity does \textbf{not} imply that every canonical section on \(C\) lifts to an adjoint on \(X\). In \[
0\to\omega_X\to\omega_X(C)\to\omega_C\to0,
\] the restriction on global sections can have cokernel in \(H^1(X,\omega_X)\). Completeness of adjoint hypersurfaces on \(X\) and completeness of the canonical series they induce on \(C\) are different questions. Enriques says the cuts are canonical groups; Section 5 does not assert surjectivity. Attributing that stronger assertion to him would invent a mistake.

An abelian surface makes the distinction concrete. On an abelian surface, take a smooth sufficiently ample \(C\) with \(H^1(\mathcal O_X(C))=0\). Then \(K_X=0\), \(g(C)=1+C^2/2\), and \(h^0(\mathcal O_X(C))=g(C)-1\). Restriction to \(\omega_C\) has a one-dimensional kernel, hence image of dimension \(g(C)-2\), versus \(h^0(\omega_C)=g(C)\). The defect is two, the irregularity. This illustrates the exact-sequence obstruction without contradicting adjunction.

\section[Genere superficiale e lineare]{Genere superficiale e genere lineare (pp.~53--56)}\label{commentary-II-6}

\subsection{Definitions and geometric genus}

An absolute invariant survives every birational map of smooth models. A relative invariant survives transformations without fundamental points or exceptional curves but may change under blowups. Enriques defines \[
p_g(X)=h^0(X,\omega_X)=\dim|K_X|+1,
\] with \(p_g=0\) if the system is empty. The ``number of linearly independent curves'' is the dimension of a section space.

For a smooth degree-\(n\) hypersurface, \[
p_g=\binom{n-1}{3}=\frac{(n-1)(n-2)(n-3)}6,
\] obtained by counting degree-\(n-4\) forms. It vanishes for \(n=1,2,3\). This concerns a general, thus smooth, hypersurface, not every singular surface of that degree. Rational surfaces have \(p_g=0\), by birational invariance and \(K_{\mathbb{P}^2}=-3H\). Enriques qualifies his cubic rationality assertion to exclude cones.

For a ruled surface, a general fibre \(f\simeq\mathbb{P}^1\) has \(f^2=0\) and \(K_X f=-2\). An effective canonical divisor cannot have negative intersection with a general moving fibre. Hence \(p_g=0\); likewise \(P_m=0\) for every \(m>0\). Enriques states this in adjoint language: a curve cutting a fibre in \(n\) points has adjoint cutting it in \(n-2\).

\subsection{Linear genus, canonical degree, and blowups}

The relative linear genus and degree are \[
p^{(1)}=p_a(K_X),\qquad p^{(2)}=K_X^2.
\] As \(|K_X'|=|2K_X|\), adjunction gives \[
K_X(2K_X)=2p^{(1)}-2,\qquad
p^{(2)}=p^{(1)}-1,\qquad p^{(1)}=K_X^2+1.
\] These use the entire virtual class; the genus of its moving part need not obey the identity.

For a blowup at a smooth point, \[
K_{\widetilde X}=\pi^*K_X+E,\quad E^2=-1,\quad
\pi^*D\cdot E=0,\quad K_{\widetilde X}^2=K_X^2-1.
\] In a local chart \(x=u,y=uv\), the identity \(\pi^*(dx\wedge dy)=u\,du\wedge dv\) gives the single exceptional zero. Projection and the normal bundle \(\mathcal O_{\mathbb{P}^1}(-1)\) give the intersection formulas. Both relative invariants decrease by one. See Hartshorne (1977, V, Propositions 3.1--3.3) and Brosnan (n.d., Section 4.2) for these blowup formulas.

For \(f:X\to X_{\min}\) comprising \(\sigma\) blowups over a nonruled minimal model, \[
\bar p^{(1)}=p^{(1)}+\sigma,\quad
\bar p^{(2)}=p^{(2)}+\sigma,\quad
\bar p^{(2)}=\bar p^{(1)}-1.
\] Here \(\sigma\) counts blowup steps, including infinitely near points, or the rank of the exceptional lattice. It is not the number of currently visible irreducible \((-1)\)-curves: after two infinitely near blowups their strict exceptional curves have squares \(-2,-1\), although \(\sigma=2\).

The assertion that ruled surfaces lack a model free of exceptional curves uses the classical notion including second-kind exceptions. It must not become the false modern assertion that ruled surfaces have no models without \((-1)\)-curves. Minimal ruled surfaces exist. Enriques separately defines a maximal relative linear genus across a ruled birational class.

\subsection{The numerical examples}

Plane quartics have \(\pi=3\), \(C^2=n=16\), and adjoint lines have \(\pi'=0\). The genus-sum identity \[
\pi'=\pi+p^{(1)}+(2\pi-2-n)-1
\] gives \(p^{(1)}=10\), agreeing with \(K_{\mathbb{P}^2}^2=9\). Thus both the printed self-intersection 16 and the final subtraction of one are correct.

For a ruled surface with base genus \(g>0\) and hyperplane class \(H\), \(H^2=m\), Enriques obtains \[
p_a(2H)=2g+m-1,\quad (2H)^2=4m.
\] The adjoint \(K_X+2H\) is pulled back from a degree-\(m+2g-2\) divisor on the base, represented by that many fibres when effective. Thus \[
p_a(K_X+2H)=1-(m+2g-2)=-m-2g+3,
\] and substitution yields \[
p^{(1)}=9-8g=-8(g-1)+1,\qquad K_X^2=8(1-g).
\] Negative arithmetic genus of a disconnected fibre sum is legitimate.

For a direct check in the quotient convention \(X=\mathbb{P}_B(\mathcal E)\), put \(\xi=c_1(\mathcal O_X(1))\). Then \[
K_X=-2\xi+\pi^*(K_B+\det\mathcal E).
\] For \(H=\xi+\pi^*A\), the base degree of \(K_X+2H\) is \(2g-2+\deg\mathcal E+2\deg A=2g-2+H^2\), proving the asserted fibre count.

Maximality of 10 in the rational class and \(9-8g\) in an irrational ruled class requires classification of minimal rational and ruled models. The arithmetic does not prove maximality. Enriques attributes the invariant to Castelnuovo (1897) and refers forward. In modern classification (Beauville, 1996), minimal rational models have \(K_X^2=9\) on \(\mathbb{P}^2\) or \(8\) on Hirzebruch surfaces; minimal irrational ruled models have \(K_X^2=8(1-g)\). Blowups only lower it.

\section[Curve eccezionali e sistema canonico]{Le curve eccezionali come parti fisse del sistema canonico (pp.~56--60)}\label{commentary-II-7}

\subsection{Fixed components and disjointness}

The arguments using negative intersection are complete. For composite divisors the discrepancy rule is justified by continuity, and existence of the contracted model is expressly deferred. Assume \(p_g>0\), so an effective canonical divisor exists. For a \((-1)\)-curve \(E\simeq\mathbb{P}^1\), adjunction gives \(K_X E=-1\); thus every effective canonical divisor contains \(E\).

Enriques correctly proves that two such curves cannot meet. If \(EF=i\ge1\), remove their full multiplicities \(r,s\) from an effective \(K_X\), leaving \(R\) containing neither. Then \[
RE=-1+r-si,\qquad RF=-1+s-ri.
\] Both must be nonnegative, whereas their sum is \(-2+(r+s)(1-i)<0\). Finiteness follows because all such curves are components of any one effective canonical divisor. This concerns irreducible \((-1)\)-curves on the current model, not every prime component of an exceptional configuration.

A classical second-kind exceptional curve can, by the earlier argument, be represented on a birational model by a rational curve \(E\) with square \(\nu\ge0\). There \[
K_X E=-2-\nu,\qquad (K_X-sE)E=-2-(s+1)\nu<0
\] for every \(s\ge0\). No finite fixed multiplicity can suffice. Birational invariance of \(p_g\) preserves the effective canonical divisor on that model, giving the contradiction.

After removing one \(E\), the residual class has intersection zero with it, but may still contain it. If the blowdown point has canonical base multiplicity \(a\), then \[
|K_X|=f^*|K_Y|+E
\] has \(a+1\) copies of \(E\), of which \(a\) still belong to the pure pullback system. Enriques explicitly recognizes this on p.~57. ``Virtually absent'' describes an unassigned point, not an empty actual base scheme.

\subsection{Infinitely near points and discrepancy}

The text invokes continuity and cites Franchetta (1940) for composite exceptional multiplicities. Blowup induction gives an exact modern proof. If \(E_j^*\) is the total transform on \(X\) of the exceptional divisor born at step \(j\), then \[
K_X=f^*K_Y+\sum_{j=1}^{\sigma}E_j^*.
\] The total-transform classes are orthogonal: \(E_i^*E_j^*=-\delta_{ij}\). Each prime component's discrepancy coefficient is the sum of its multiplicities in these total transforms. This proves the stated rule without specialization. Additional fixed multiplicity inherited from \(|K_Y|\) is still possible.

For three points following a smooth branch the total exceptional curves are \[
\omega+\omega_1+\omega_2,\qquad\omega_1+\omega_2,\qquad\omega_2,
\] giving coefficients \(1,2,3\). For a satellite third point at the intersection of the first two prime components they are \[
\omega+\omega_1+2\omega_2,\qquad\omega_1+\omega_2,\qquad\omega_2,
\] giving \(1,2,4\). These are different proximity configurations.

The rigorous existence of a model free of exceptional curves is expressly deferred to surface Riemann--Roch and Castelnuovo--Enriques (1901). The modern completion uses Castelnuovo contraction: a smooth rational curve of square \(-1\) contracts to a smooth point on a projective surface. Each contraction lowers the Picard number, so the process terminates. The Stacks Project, Lemma 54.16.8 (Tag 0C2M), provides the algebraic construction: twist a sufficiently positive line bundle \(L\) by \((L\cdot E)E\), use the resulting system trivial on \(E\), and take its Stein factorization; the local contraction criterion ensures smoothness. This is a modern proof source, without a first-repair attribution.

For \(p_g>0\), the minimal model is nonruled and unique. Given a common resolution of two minimal models, remove any \((-1)\)-curve exceptional for both morphisms. A remaining first curve to be contracted for one side is a \((-1)\)-curve not exceptional for the other. Every \((-1)\)-curve contracted on the other side is disjoint from it by the proved disjointness theorem. It therefore remains a \((-1)\)-curve throughout those contractions, contradicting minimality on the other side. This uses standard elimination of indeterminacy and factorization into point blowups.

\subsection{Additivity and the irreducible criterion}

The function \[
D\longmapsto2p_a(D)-2-D^2=K_X D
\] is additive. This also follows by cancelling the cross term in the genus-sum formula, as the text notes. If integral \(C\) satisfies \[
C^2>2p_a(C)-2,
\] then \(K_X C<0\), so \(C\) occurs in every effective canonical divisor. After removing its full multiplicity \(s\), the residual intersection is \[
2p_a(C)-2-(s+1)C^2.
\] If \(C^2\ge0\), it is still negative, impossible. Thus \(C^2<0\). Since \(p_a(C)\ge0\), the original strict inequality forces \(p_a(C)=0,C^2=-1\). Arithmetic genus zero of an integral curve implies that its normalization genus and all singularity \(\delta\)-invariants are zero. Hence \(C\) is smooth rational, exactly a \((-1)\)-curve. This proof is substantially present in Enriques.

\subsection{A qualification of the general composite assertion on p.~59}

For effective divisors, define numerical 1-connectedness by \(A\cdot B\ge1\) for every expression \(D=A+B\) with \(A,B>0\), including shared prime components. For reduced curves this agrees with connected support; for nonreduced divisors it does not. On p.~59 Enriques glosses \emph{sconnesse} by ``aventi con esse un numero virtuale nullo di punti comuni'\,'; the exercise on p.~60 requires attention to decompositions with shared components. However, neither the reviewed Chapter I nor this passage states the fully quantified definition of numerical 1-connectedness. Numerical 1-connectedness is the precise sufficient hypothesis used in the repaired criterion (see also Alzati and Tortora, 2002, for the modern terminology). The exercise disproves a support-connected reading but does not uniquely identify the author's intended definition.

The preceding broader assertion on p.~59 deserves a separate qualification. Read literally as an assertion about effective summands, it says that every \(D\) with \(K_X D<0\) is obtained by adding exceptional curves orthogonal to an effective remainder \(A\) with \(K_X A\ge0\). This orthogonal-decomposition claim is too strong for arbitrary nonreduced effective divisors. The scan confirms the zero-intersection wording; it does not say nonpositive intersection.

An effective-divisor counterexample is obtained as follows. Blow up a smooth quartic at a point and then at a point on the first exceptional curve. On the final surface let \(F\) be the first strict exceptional curve and \(E\) the second. Then \[
F^2=-2,\quad E^2=-1,\quad FE=1,\quad K_X=F+2E.
\] For \(D=2F+E\), \[
D^2=-5,\qquad K_X D=-1,\qquad p_a(D)=-2,
\] so \(D^2>2p_a(D)-2\). An effective remainder \(A\le D\) with \(K_X A\ge0\) must be \(aF\), where \(a=0,1,2\). The only first-kind exceptional cycles supported there are \(E\) and \(F+E\): a cycle \(uF+vE\) with genus zero and square \(-1\) has \(K_X(uF+vE)=-v=-1\), hence \(v=1\), and the square then forces \(u=0\) or 1. If \(A=0\), \(D\) is not a positive sum of these cycles, because two copies of \(F\) require at least two of \(E\). If \(A=F\), the exceptional remainder \(F+E\) meets \(A\) in \(-1\), not zero. If \(A=2F\), the remainder \(E\) meets \(A\) in 2.

Thus the orthogonal decomposition fails \textbf{if its remainder is required effective}. Chapter I permits virtual divisor classes, so this is not a categorical refutation of every classical reading. Indeed the very same example has \[
D=-E+2(F+E),\qquad K_X(-E)=1,\qquad(-E)(F+E)=0.
\] The virtual remainder \(-E\) salvages the statement.

There is a precise general modern formulation. Contract to a minimal model \(f:X\to Y\) and use the orthogonal total-transform basis \(E_i^*\). Write \[
D=f^*D_0+\sum a_iE_i^*,\qquad
K_XD=K_YD_0-\sum a_i.
\] Here \(D_0=f_*D\) is effective or zero and \(K_Y\) is nef: the irreducible criterion proves nefness on a minimal surface with \(p_g>0\). Put \[
B=\sum_{a_i>0}a_iE_i^*,\qquad
A=f^*D_0+\sum_{a_i\le0}a_iE_i^*.
\] Then \(B\) is a positive sum of total exceptional divisors, \(A\cdot E_i^*=0\) for every positive coefficient in \(B\), and \[
K_XA=K_YD_0-\sum_{a_i\le0}a_i\ge0.
\] Moreover \(K_XD<0\) forces at least one positive \(a_i\). This proves the text's orthogonal-addition statement with a \textbf{possibly virtual remainder}, and explains why the coefficients cannot all be presumed nonnegative. This qualifies the sketched argument by allowing a virtual remainder, without attributing the reconstruction to a historical source. The example is not 1-connected and does not invalidate p.~60, whose effective criterion is proved next.

\subsection{Complete modern proof of the composite exceptional criterion}

\textbf{Statement.} If \(p_g(X)>0\), an effective 1-connected divisor \(D\) with \(p_a(D)=0,D^2=-1\) is a total exceptional divisor of a contraction to a smooth point. Conversely a total pullback of a \((-1)\)-curve under successive point blowups is 1-connected with these invariants.

For the forward implication, \(K_X D=-1\). Some prime component has negative canonical intersection, hence is a \((-1)\)-curve \(E\). If \(D\ne E\), 1-connectedness applied to \(E+(D-E)\) gives \(t=DE\ge0\). Contract \(E\), put \(D'=\pi_*D\), and obtain \[
D=\pi^*D'-tE,\qquad (D')^2=D^2+t^2,\qquad
K_YD'=K_XD-t<0.
\] The pushforward \(D'\) is again 1-connected. For \(D'=A'+B'\), let \(m_A,m_B\) be the \(E\)-coefficients of the total pullbacks. Since \(D\) contains \(E\), \(m_A+m_B-t\ge1\). Choose integers \(a,b\ge0\), \(a\le m_A,b\le m_B,a+b=t\). Then \[
D=(\pi^*A'-aE)+(\pi^*B'-bE)
\] is an effective decomposition. Its product is \(A'B'-ab\ge1\), so \(A'B'\ge1\).

Again \(K_YD'<0\) supplies a \((-1)\)-component. Picard number decreases at each step, so eventually the entire divisor is one \((-1)\)-curve; a multiple \(rE\), \(r>1\), is excluded by 1-connectedness. Initial and final squares are both \(-1\), while each step increases square by \(t^2\). Every \(t\) is therefore zero. The original \(D\) is the total pullback of that last exceptional curve.

Conversely total pullback preserves square and arithmetic genus and preserves 1-connectedness. Indeed, decompose \(\pi^*D=A+B\), write \(A=\pi^*A'+aE,B=\pi^*B'-aE\), and compute \(AB=A'B'+a^2\). If both pushforwards are nonzero this is at least one; if one is zero, its summand is a positive multiple of \(E\), giving \(a^2\ge1\). Induction from a reduced \((-1)\)-curve finishes the proof. This blowdown proof is a modern reconstruction; it is not attributed to Franchetta's original article.

\subsection{The quartic exercise}

Let \(S\subset\mathbb{P}^3\) be smooth quartic containing a line \(a\). With hyperplane class \(H\), \(K_S=0,H^2=4,Ha=1,a^2=-2\). The residual plane cubics through \(a\) have class \(F=H-a\), so \[
F^2=0,\qquad p_a(F)=1.
\] Blow up a smooth point of a smooth such cubic; let \(E\) be exceptional and \(C=\pi^*F-E\). Then \[
C^2=-1,\quad CE=1,\quad E^2=-1,\quad K_X=E.
\] The exercise divisor is \[
D=C+2E=\pi^*F+E,\qquad D^2=-1,\quad K_X D=-1,\quad p_a(D)=0.
\] It is not exceptional. Its support contains an elliptic curve, and the intersection matrix \[
\begin{pmatrix}-1&1\\1&-1\end{pmatrix}
\] is not negative definite. Most directly \(D=(C+E)+E\) with \((C+E)E=0\). Thus it is not 1-connected although its support is connected. This resolves Enriques's paradox and shows why merely topological connectedness would make the criterion false. Numerical 1-connectedness is the explicit sufficient replacement adopted here, not a definition quoted from Chapter I.

The footnote extends the criterion to \(p_g=0\), with an exception concerning a rational pencil and its fixed components, citing Franchetta (1941). This is assumed earlier work here. Its full original formulation and proof were not verified in that paper. The proof above retains \(p_g>0\), where all inputs have been established.

\section[Curve pluricanoniche]{Curve bicanoniche e pluricanoniche (pp.~60--62)}\label{commentary-II-8}

\subsection{Definitions and identities}

A class may lack effective representatives while a multiple has them. Since \(C'=C+K_X\), \[
|2K_X|=|2C'-2C|,\qquad |iK|=|iC'-iC|,
\qquad |iC'-iC|=|iL'-iL|.
\] If \(C''=(C')'=C+2K_X\), then \[
|(C+C')'|=|C+C''|=|2C'|,\qquad |2K_X|=|C''-C|.
\] Enriques proves these Picard-group identities; none presumes that a lower multiple is effective. He states birational invariance of the dimensions, whose blowup proof is given below.

The \(i\)-th plurigenus is \(P_i=h^0(X,\mathcal O_X(iK_X))\), with bigenus \(P_2\). Enriques defines these and states absolute birational invariance. The virtual genera and degrees are relative: \[
p_a(2K_X)=2p^{(1)}+p^{(2)}-1=3p^{(1)}-2,\qquad
(2K_X)^2=4p^{(2)}=4(p^{(1)}-1),
\] \[
p_a(iK)=1+\frac{i(i+1)}2(p^{(1)}-1),\qquad
(iK)^2=i^2(p^{(1)}-1).
\] Substitution in arithmetic adjunction proves all these. They concern entire divisor classes, not the genus of a moving part or the smoothness of a general representative.

\subsection{The tetrahedral sextic}

For a sufficiently general sextic with the six tetrahedron edges as double curves, the smooth normalization \(X\) has \(p_g=0\) and \(2K_X\sim0\). The unique effective bicanonical divisor is the zero divisor. Its plane-section class \(H\) has \(H^2=6,p_a(H)=4\); \(H'=H+K_X\) is cut by cubics through the edges. Although \(H'\not\sim H\), \(2H'\sim2H\). Smooth normalization with the expected ordinary singularities is required; this does not apply to every degenerate sextic containing the doubled edges.

These statements are announced and sketched here; their full conductor and monomial verification belongs with the developed example in Chapter III, Section 12 (Dolgachev, 2016, Section 2, pp.~3--4). The residual cuts by sextics double along the six edges represent \(|2H'|=|2H|\). Their completeness uses the adjoint/conductor description, while the prescribed double-edge conditions, the span of products of cubics, and the vertex conditions must be distinguished as in that later discussion.

Torsion excludes \((-1)\)-curves because their canonical intersection would be \(-1\). The general exceptional-curve extension below also excludes second-kind exceptions. After one blowup \(K_{\widetilde X}=\pi^*K_X+E\), so the unique bicanonical divisor becomes \(2E\), exactly as on p.~61.

\subsection{Birational invariance and positive plurigenera}

For one blowup, \[
iK_{\widetilde X}=\pi^*(iK_X)+iE.
\] Every effective divisor in this class contains \(E\) at least \(i\) times: its intersection with \(E\) is \(-i\), and remains negative after removing \(r<i\) copies. Removing \(iE\), then using \(\pi_*\mathcal O_{\widetilde X}=\mathcal O_X\), gives \[
H^0(\widetilde X,iK_{\widetilde X})\simeq H^0(X,iK_X).
\] A common resolution proves invariance for arbitrary smooth birational models. Hartshorne (1977, V, Propositions 3.3--3.4) and Brosnan (n.d., Section 4.2) supply the blowup and pushforward facts.

If some \(P_i>0\), replace the effective canonical divisor in Section 7 by an effective \(iK\). For two meeting \((-1)\)-curves the residual intersection equations begin with \(-i\), and their sum is again negative. For a rational curve with nonnegative square, subtraction again never terminates. Thus irreducible exceptional curves are finite, pairwise disjoint, and first-kind. Their forced multiplicity is \(i\) times the canonical discrepancy; more copies may come from pluricanonical base points downstairs. This is Enriques's ``in general \(i\) times,'' including composite cycles through the total-transform rule.

The same contraction theory supplies a nonruled minimal model even if \(p_g=0\) but \(P_i>0\), and the absolute linear genus is \(K_{X_{\min}}^2+1=p^{(1)}+\sigma\). This does not assert positive plurigenus for ruled surfaces: all their positive plurigenera vanish.

Canonical self-intersection alone loses torsion information. A smooth quartic has \(K_X\sim0\), whereas the tetrahedral sextic normalization has \(K_X\not\sim0\), \(2K_X\sim0\). Both have \(K_X^2=0\) and \(p^{(1)}=1\), but their geometric genera and plurigenera differ. This is why the line-bundle class and its section spaces must be retained before passing to numerical characters.

\section[Nota storica]{Nota storica (pp.~62--65)}\label{commentary-II-9}

The historical note restates the preceding mathematics and attributes its development. Its citations are reports of earlier work, not proofs in this chapter; a modern reconstruction does not establish priority.

Enriques attributes geometric genus to Clebsch's 1868 number of everywhere finite double integrals. Their integrands are holomorphic two-forms, a space of dimension \(h^0(\omega_X)\). In hypersurface language, adjoint numerators of degree \(n-4\) multiply residue two-forms, parallel to degree-\(n-3\) adjoints for plane curves. This explains the canonical divisors' birational covariance.

Noether is credited with an algebraic invariance proof and the linear genus. Enriques emphasizes that \(p^{(2)}=p^{(1)}-1\) holds for virtual canonical characters, not necessarily the genus and degree of a moving effective part. The footnote gives Mathematische Annalen volumes 2 and 8, dated 1869 and 1874. Bibliographic references usually date publication 1870 and 1875; earlier dates can reflect writing or submission. The original journal pages were not verified, so the date difference remains a bibliographic qualification rather than a demonstrated mathematical error.

The Ricerche and Introduzione are said to characterize canonical restrictions, subject to an explicit exception and exclusions for proper fundamental curves, then prove birational invariance by special transformations. The exception is not completely specified in this passage, which refers forward to Chapter III. This cannot be replaced by the unrestricted assertion that restriction to one moving curve determines a line bundle on the surface.

The Ricerche date 1893 is defensible: the official collected-works archive lists a Clausen offprint of 1893 and the journal publication in volume 44 in 1894. The archive also records the 1896 Introduzione. A blanket ``1893 to 1894'' correction would erase a real distinction.

The conceptual extension is the canonical class when \(p_g=0\): \[
C+D'=C'+D,\qquad K_X=C'-C=D'-D.
\] Its virtual numerical characters remain defined, and \(sK_X\) may be effective although \(K_X\) is not. Plurigenera supply additional \emph{dimensions}; their numerical genus and degree are already polynomials in \(K_X^2\).

The account also states its extension to arbitrary dimension, verified by the principal-parts calculation in Section 1, and locates the Jacobian presentation in the Bologna course of 1897--98, the 1899 published programme, and the 1901 Intorno ai fondamenti della geometria sopra le superficie algebriche. The official archive confirms the programme's 1897--98 title and the paper in Atti of Turin, 37, 19--40. The paper's opening recollection uses a different lecture-year wording; exact chronology should be checked against both originals before settling priority.

Enriques compares Segre's and Bertini's 1894 expositions, citing Segre Section 75 and Bertini Section 21(b),(g), and says their proof uses uniqueness of the canonical \(g_{2g-2}^{g-1}\) via curve Riemann--Roch. The mathematical uniqueness is immediate: if \(\deg L=2g-2,h^0(L)=g\), Riemann--Roch yields \[
h^0(K_C\otimes L^{-1})=1.
\] A degree-zero line bundle with a nonzero section is trivial; hence \(L\simeq K_C\). The effective-series formulation has \(g\ge1\); on \(\mathbb{P}^1\) the canonical class is the ineffective \(\mathcal O(-2)\). This proves the comparison's mathematics, but not the universal historical assertion that no other proof was known. The entire Segre and Bertini papers were not independently reread for that priority claim.

The closing references to Enriques--Chisini volume III and Enriques's Bologna memoir of 1914 point to variants of the method. The preference for simplicity is an expository judgment, not an additional theorem.

The locator for Enriques and Chisini on p.~43, the full original formulations of Franchetta's 1940 and 1941 exceptional-curve papers, and the entire priority comparison between Segre and Bertini were not verified. The mathematical replacements above therefore stand on their displayed calculations and modern references. They do not imply a first-repair attribution to an unexamined historical paper.

\chapter{Le superficie aggiunte}

Chapter III gives the projective meaning of the classes constructed in Chapter II: which hypersurfaces cut canonical, adjoint, and bicanonical systems, and which local conditions are required at singular points (pp.~67--100). Write \(Y\subset\mathbb P^3\) for a projective model, \(X\) for a smooth resolution, and \(H\) for the pullback of a plane section. Enriques's \(p\) and \(P\) become \(p_g\) and \(P_2\), both vector-space dimensions. For a divisor \(C\), one has \(C'=K_X+C\), \(C''=2K_X+C\), but \((2C)'=K_X+2C\). The characteristic series is the image of restriction, with underlying bundle \(\mathcal O_C(C)\), and need not be complete.

Three questions must be kept separate: completeness of an ambient system on \(X\); surjectivity of its restriction to a curve; and independence, in the chosen degree, of the local adjunction conditions. Enriques explicitly recognizes the third distinction on p.~74. Similarly, reduced exceptional support, fundamental cycle, and maximal-ideal fiber cycle need not coincide. These distinctions organize both the singularity calculations and the final equivalence criteria. The background identities are surface Riemann--Roch, Serre duality, and adjunction (Beauville, 1978, I.12--I.16 and III.20).

\section[Superficie aggiunte di ordine n-3]{\texorpdfstring{Superficie d'ordine \(n-3\) aggiunte ad una superficie d'ordine \(n\) priva di punti multipli propri (pp.~67--69)}{Superficie d'ordine n-3 aggiunte ad una superficie d'ordine n priva di punti multipli propri (pp.~67--69)}}\label{commentary-III-1}

\subsection{Singular curves, neutral groups, and proper singular points}

Enriques begins with a generic \(i\)-fold curve of \(Y\). On a model separating its \(i\) local sheets, a point on that curve is represented by a group of \(i\) points. A hyperplane containing the original point contains the entire group and imposes one linear condition, not \(i\) independent conditions. This is the meaning of a neutral group relative to \(|H|\). It is a statement about the image of the chosen linear system, not a claim that distinct points never impose independent conditions on a complete linear system.

Exceptional multiple curves are handled in the source by adjoining their infinitely near multiple curves, without a general resolution proof here; the modern adjoint ideal below supplies the invariant formulation. A multiple point is called improper when a general plane section constrained to pass through it has the same geometric genus as an unconstrained general section. It is proper when that genus decreases. Ordinary intersections of sheets, such as a triple point of both the surface and its double curve, are the first sort. The words ``ordinary'' and ``general'' matter: additional tangencies or singularities of the sheets can create proper singularity data that are invisible in the multiplicity alone.

If an improper point has been replaced by exceptional curves, each connected component of the appropriate exceptional cycle has arithmetic genus zero and meets the residual plane-section system once. The genus-of-a-sum identity \[
p_a(A+B)=p_a(A)+p_a(B)+A\cdot B-1
\] explains the absence of a genus loss. Enriques refers back to the theory of exceptional curves from the earlier chapters for contractibility. This is not an assertion that every reduced configuration of rational curves is a single \((-1)\)-curve. A connected configuration can be contracted by a succession of blowdowns, with the relevant multiplicities retained.

\subsection{The sweeping argument}

Assume \(Y\) is not ruled and its general plane section \(C\) has genus \(\pi>1\). Suppose a curve \(L\) on \(Y\) cuts a canonical divisor on every general \(C\). The plane adjunction theorem supplies a curve of degree \(n-3\) through the singularities of the plane curve \(C\), with multiplicity \(i-1\) at an ordinary \(i\)-fold point. That point absorbs \(i(i-1)\) intersections. The remaining intersection divisor has degree \(2\pi-2\). The plane adjoint is unique for the specified residual divisor because its degree is smaller than \(n\).

Enriques fixes a general line \(a\), takes this plane adjoint in each plane through \(a\), and sweeps out a surface of degree \(n-3+h\) having multiplicity \(h\) along \(a\). If \(h>0\), take \(A\in a\cap Y\) outside \(L\) and the multiple curve. A suitable plane through \(a\), tangent to the swept surface at \(A\), makes its residual degree-\(n-3\) curve pass through \(A\). The prescribed canonical and singularity intersections already account for the full Bézout number with the section of \(Y\); the additional point forces a common component. Since the residual adjoint has smaller degree, that section of \(Y\) must be reducible. Varying the generic axis gives a two-dimensional family of such reducible sections. The conclusion is conditional on an imported classification: exclusion of the ruled and Steiner alternatives uses the Kronecker--Castelnuovo classification cited by Enriques through Enriques--Conforto, Book II, Chapter II, Section 7, p.~273. That classification is assumed, not proved in these pages. The assertion that the plane curves vary algebraically also uses the algebraic dependence of the solutions of their linear equations on the plane; uniqueness on a dense open set supplies a rational family. With these inputs, the sweeping argument identifies an adjoint surface of degree \(n-3\).

The local condition along an ordinary \(i\)-fold curve is multiplicity \(i-1\). The elementary translational model in the footnote explains that condition at a general point of the multiple curve. It does not by itself settle the extra conditions at exceptional points of that curve or at isolated singularities; precisely those questions occupy the subsequent sections.

\subsection{A modern completeness theorem}

A modern replacement for the construction follows, valid without the ruled-surface exclusion. Let \(Y\) be a reduced degree-\(n\) hypersurface in \(\mathbb{P}^3\), and let \(\pi:X\to Y\) be a resolution. The adjoint ideal \(\operatorname{adj}(Y)\) is defined on an embedded resolution \(\mu:W\to\mathbb{P}^3\), with smooth strict transform \(\widetilde Y\), by \[
\operatorname{adj}(Y)
=\mu_*\mathcal O_W(K_{W/\mathbb{P}^3}-\mu^*Y+\widetilde Y).
\] Ein and Lazarsfeld (1997, Proposition 3.1, p.~254) prove the exact sequence \[
0\longrightarrow\mathcal O_{\mathbb{P}^3}(-4)
\longrightarrow\operatorname{adj}(Y)(n-4)
\longrightarrow\pi_*\omega_X\longrightarrow0.
\] Twisting by \(r\), and using \(H^1(\mathbb{P}^3,\mathcal O(r-4))=0\), gives a surjection \[
H^0(\mathbb{P}^3,\operatorname{adj}(Y)(n-4+r))
\twoheadrightarrow H^0(X,K_X+rH).
\tag{III.1}
\] Thus adjoint forms cut the complete system \(|K_X+rH|\); forms differing by the equation of \(Y\) times a form of degree \(r-4\) have the same section on \(Y\). The proof pushes forward the ordinary adjunction sequence on \(W\); the vanishing of \(R^1\mu_*\omega_W\) is the vanishing input. This is a complete modern replacement for the construction, without a claim of historical first repair.

If the normalization of \(Y\) is smooth, finite-morphism duality describes its canonical sheaf through the conductor. If that normalization still has proper singularities, one must pass on to its resolution: the conductor alone need not impose all adjunction conditions. This is exactly why the first sweeping construction has to be modified at proper multiple points.

The surjectivity in (III.1) concerns sections on \(X\). Restriction to \(C\) is governed by \[
0\to\omega_X\to\omega_X(C)\to\omega_C\to0,
\] and its cokernel can involve \(H^1(X,\omega_X)\). No claim that every canonical section on \(C\) extends to \(X\) should be read into Enriques's completeness statement.

\section[Superficie aggiunte]{Superficie aggiunte d'ordine qualsiasi (pp.~69--70)}\label{commentary-III-2}

On a plane section, an adjoint of degree \(n-3+r\) cuts \[
K_C+rH|_C.
\] For \(r\geq3\), adding the equation of \(C\) times a form of degree \(r-3\) leaves the residual divisor unchanged. The vector space of these additions has dimension \[
\binom{r-1}{2}=\frac{(r-1)(r-2)}2.
\] This explains the loss of uniqueness. For a fixed section, the representatives form an affine translate of this kernel. For a fixed residual divisor, they form an open part of the projectivization of the span of one representative and the kernel, of projective dimension \(\binom{r-1}{2}\), as on p.~70.\footnote{The text above describes an affine family of representatives, while the source records a projective dimension. These concern a fixed section and a fixed divisor respectively; their dimensions agree. The scan confirms the printed dimension \((r-1)(r-2)/2\).} Enriques suggests specifying points and tangencies along the axis to choose one representative; the independence of such conditions is not proved for arbitrary choices. On a dense open set, choosing a linear complement to the kernel makes the required rational selection. For negative \(r\), Enriques adds an auxiliary surface to reduce to a positive degree and then uses the forced common factor to subtract it. Modern twisting of the adjoint sequence gives the same conclusion directly, with an empty linear system permitted when there are no sections.

For the adjoint of \(rH\), however, the formula is \[
(rH)'=K_X+rH=H'+(r-1)H.
\] Consequently its restriction to a general \(H\) is \[
(K_X+rH)|_H=K_H+(r-1)H|_H.
\tag{III.2}
\] The printed statement on p.~70 uses the \(r\)-fold characteristic series alongside \(|(rC)'|\), an off-by-one error.\footnote{Both the scan and transcription have ``dello \(r\)-plo della serie caratteristica.'' The preceding displayed identity has \(C'+(r-1)C\), which gives the corrected restriction.} The same mismatch recurs in Section 3. It is not the earlier plane-curve formula that is wrong: that earlier formula uses degree \(n-3+r\), one higher than the degree \(n-4+r\) of adjoints representing \(K_X+rH\).

At \(r=0\), degree-\(n-4\) adjoint forms represent canonical sections. At \(r=1\), degree-\(n-3\) adjoints represent \(|K_X+H|\). For a smooth hypersurface this is the familiar identity \(K_Y=(n-4)H\). The remainder of the chapter explains how singularities change the ideal of admissible numerators.

There is a separate completeness warning on p.~70. The complete multiple \(|rH|\) can be larger than the subsystem cut by ambient degree-\(r\) hypersurfaces. The relevant comparison is between \(H^0(X,rH)\) and the image of \(H^0(\mathbb{P}^3,\mathcal O(r))\); equality for all degrees is a projective-normality issue. It must not be inferred from completeness of the adjoint systems established by the twisted adjoint sequence.

\section[Superficie rigate]{Superficie rigate (pp.~70--71)}\label{commentary-III-3}

Let \(\rho:X=\mathbb{P}_B(\mathcal E)\to B\) be a smooth ruled model over a curve of genus \(g\), let \(f\) be a ruling fiber, and suppose the chosen hyperplane class \(H\) is a unisection: \(Hf=1\). Then \(K_Xf=-2\). Hence \[
(K_X+H)f=-1,
\] so \(|K_X+H|\) is empty: an effective divisor cannot have negative intersection with a general moving fiber.

Nevertheless, a canonical divisor on \(B\), pulled back to \(X\), cuts a canonical divisor on every unisection \(C\simeq B\). For \(g>1\), this consists of \(2g-2\) generators; for \(g=1\) it is the zero divisor. Thus the property of cutting canonical groups on the plane sections does not characterize \(|K_X+H|\) for scrolls. This is a genuine exception to the characterization, not a failure of the completeness theorem (III.1).

Enriques's projective obstruction agrees with the intersection argument. Let \(\Delta\) be the conductor divisor on the smooth normalization of the ordinary scroll. Conductor adjunction gives \(K_X=(n-4)H-\Delta\); intersecting with a general ruling fiber gives \(-2=n-4-\Delta f\), hence \(\Delta f=n-2\). These are the \(n-2\) intersections of a generator with the double curve counted in the source. This is an added proof of a fact Enriques assumes (Dolgachev, 2016, Section 2, equation (1), gives the conductor formula). A degree-\(n-3\) surface cannot meet the generator at that many points without containing it. Requiring this for every generator forces the whole scroll as a component, which is impossible in the stated degree.

In the convention \(H=\mathcal O_{\mathbb{P}(\mathcal E)}(1)\), \(H^2=\deg\mathcal E=n\), \[
K_X=-2H+\rho^*(K_B+\det\mathcal E),
\qquad K_X+2H=\rho^*(K_B+\det\mathcal E).
\] This vertical system cuts the series \(K_C+H|_C\), of degree \(2g-2+n\). More generally, \((K_X+rH)f=r-2\). Therefore the vertical system that merely reproduces the restricted divisors agrees with an adjoint of \(rH\) only for \(r=2\). The printed \(r\)-fold characteristic series must again be replaced by \((r-1)\)-fold in this comparison. For higher \(r\), genuine adjoints have a horizontal part. For higher degree Enriques sketches adjunction on an auxiliary degree-\(r\) surface, then explicitly postpones the dimension argument for completeness. Equation (III.1) supplies that missing proof independently of the ruled-surface exclusion.

\section[Proprietà caratteristica degli aggiunti]{Proprietà caratteristica delle curve canoniche e delle curve aggiunte ad un sistema lineare (pp.~71--72)}\label{commentary-III-4}

The projective construction can be applied to a net represented by planes through a point \(O\) of multiplicity \(s\) on a surface of degree \(m=n+s\). In the scan, \(m>n\) is clear. The adjoint has degree \(m-3\), multiplicity \(i-1\) along an \(i\)-fold curve, and multiplicity \(s-1\) at \(O\); the latter includes the one additional order imposed by taking plane sections through \(O\). These are conditions for the adjoint to the chosen net, not the \(s-2\) condition for canonical adjoints of unrestricted hyperplanes in Section 6.

The two projective hypotheses on p.~72 are that the net has no one-dimensional family of reducible sections and that it has no unresolved proper multiple points, including infinitely near points over \(O\), or multiple lines through \(O\) along which the constructed adjoint has insufficient multiplicity.\footnote{The second hypothesis could be mistaken for noncompositeness of a general member, but it is the proper-singularity qualification. The scan settles the discrepancy: the printed second condition excludes ``punti multipli proprii'' and the specified ``rette multiple per \(O\).''} The theorem then applies to a system of dimension \(r>1\) whose reducible locus has dimension at most \(r-2\). The strict inequality is essential.\footnote{The earlier transcription has \(r\geq1\); the scan unequivocally has \(r>1\). The source statement, rather than the transcription, is used here.} A fundamental curve for a system of dimension \(r\) produces an \((r-1)\)-dimensional family of reducible members: sections vanishing at its contracted image acquire that divisor. This explains the invariant reducible-locus condition. A pencil by itself would be a false extension, since one can add fiber divisors without changing restriction to a general fiber.

Here is a modern justification of the characterization under the stated base-locus and completeness conventions. Choose a general pencil in the parameter space that misses the reducible locus. If \(L\) cuts \(K_C\), the line bundle \[
M=\mathcal O_X(L-K_X-C)
\] is trivial on the generic member of that pencil. Resolve the pencil to \(f:\widetilde X\to\mathbb{P}^1\). A trivialization on its generic fiber gives a rational section of the pullback of \(M\), whose divisor is vertical. Push down to \(X\). Exceptional curves of the resolution disappear; since all the original pencil members are integral, the remaining divisor is an integral sum of pencil members. They are linearly equivalent, so \(M\simeq\mathcal O_X(sC)\). Restriction to \(C\) gives \(sC^2=0\). A noncomposite system of dimension \(>1\) has \(C^2>0\), and therefore \(s=0\).

This proves \(L\sim K_X+C\). Replacing the cut canonical series by its residual to the characteristic series proves the corresponding characterization of canonical divisors. If the chosen system contracts proper curves, vertical correction terms can survive on \(X\); those are precisely the exceptions that subsequent sections analyze. This generic-fiber argument reconstructs the proof; it is not the argument written in the book. The same comparison treats adjoints of multiples, with restriction \(K_C+(r-1)C|_C\).

\section[Punti multipli isolati]{Punti multipli isolati (pp.~72--73)}\label{commentary-III-5}

With a proper multiple point present, a curve \(L\) cutting canonical groups on general hyperplane sections still comes from an ambient form satisfying the multiple-curve conditions. Its system can exceed \(|K_X+H|\) by an exceptional divisor \(\theta\), because those general sections do not meet the isolated singular point. The additional point conditions must be determined on a smooth model, rather than inferred from their restrictions alone.

A polar of \(Y\) has degree \(n-1\). At an ordinary \(i\)-fold curve its multiplicity is \(i-1\); at an ordinary isolated \(s\)-fold point it has multiplicity \(s-1\). The polar cut omits one exceptional copy from the complete Jacobian divisor at an isolated point. Adding that copy relaxes the numerator condition from \(s-1\) to \(s-2\), the canonical adjunction order. The text explains this through a reducible net member singular at a general point of the exceptional curve and through subsequent examples. This is a geometric sketch; the ambient blowup calculation below proves the coefficient.

The same order \(s-2\) is a lower bound at a proper \(s\)-fold point lying on the multiple curve, under the genericity stated by Enriques. The multiple-curve conditions and infinitely near singularities may force additional orders. The point multiplicity alone is not a complete classification.

\section[Aggiunte e subaggiunte]{Aggiunte e subaggiunte (pp.~73--75)}\label{commentary-III-6}

Enriques separates two ideals of conditions. A subadjoint satisfies the generic conditions along multiple curves and cuts the intended groups on general plane sections. An adjoint additionally satisfies the conditions at proper multiple points that are necessary for the resulting differential to extend over a smooth model. This distinction is essential: a general plane misses an isolated point and cannot detect those extra extension conditions.

The number of local linear conditions eventually stabilizes as the ambient degree grows. In modern language, an adjunction ideal of finite colength defines a zero-dimensional quotient, and Serre vanishing makes the global restriction map onto that quotient surjective after a sufficiently positive twist. In small degree the image can be smaller. Enriques explicitly warns of dependence between point conditions and conditions already imposed by the multiple curve or other singularities. The ``diminution of genus'' is therefore virtual unless independence is checked. It is not an error that later numerical losses are called virtual.

The numerical genus that he anticipates is \(p_a(X)=\chi(\mathcal O_X)-1=p_g-q\) on a smooth surface. For a singular projective hypersurface one must distinguish its arithmetic genus from that of its resolution; the resolution defect is what the local conditions measure in the examples. Surface Riemann--Roch packages these virtual counts, while \(H^1\) records the possible global dependence (Beauville, 1978, I.12--I.13).

\subsection{\texorpdfstring{The order \(s-2\) and its condition count}{The order s-2 and its condition count}}

Suppose the tangent cone of an isolated \(s\)-fold point is smooth. Blow up the point in the ambient smooth threefold. If \(E\) is the exceptional divisor, then \[
K_W=\mu^*K_{\mathbb{P}^3}+2E,\qquad
\widetilde Y=\mu^*Y-sE.
\] Adjunction on \(\widetilde Y\) gives \[
K_{\widetilde Y}=\pi^*K_Y+(2-s)E|_{\widetilde Y}.
\] Thus an ambient canonical numerator must vanish to order at least \(s-2\). The number of coefficients of degree \(<s-2\) in three local variables is \[
\ell\bigl(\mathcal O_{\mathbb A^3,O}/\mathfrak m_O^{s-2}\bigr)
=\binom{s}{3}=\frac{s(s-1)(s-2)}6.
\tag{III.3}
\] This proves the local count in the ordinary case. A singular tangent cone can require further infinitely near conditions, which are not determined by \(s\) alone.

An ordinary double point imposes no canonical condition. An ordinary triple point imposes one; a quadruple point imposes four. A smooth quartic has \(p_g=1\). A quartic with an ordinary triple point has \(p_g=0\), and is rational by projection from that point: a general line through a point of multiplicity \(n-1\) has one residual intersection. The same argument makes a quintic with an ordinary quadruple point rational. A nodal quartic retains \(p_g=1\).

\subsection{A double line and a triple point on a quintic}

For a general quintic with a double line \(a\), canonical adjoints are planes containing \(a\). Their vector dimension is two, so \(p_g=2\), and the residual canonical curves are plane cubics. An additional ordinary triple point \(O\notin a\) imposes the extra condition that the plane contain \(O\), leaving one canonical divisor and \(p_g=1\). The unique plane \(Oa\) has residual intersection three lines; these are exceptional curves on the smooth model.

Enriques sketches the birational construction using quadrics through \(a\) and \(O\), obtaining a degree-eight genus-five surface in \(\mathbb{P}^5\), and reversing the projection from a line through three points of a plane cubic. The following intersection calculation makes the construction transparent. Let \(S\) be a smooth K3 surface with a very ample divisor \(A\), \(A^2=8\), and an elliptic plane cubic \(E\), \(AE=3\), \(E^2=0\). Blow up three distinct collinear points \(P_i\) of \(E\), with exceptional curves \(e_i\), and set \[
H=A-e_1-e_2-e_3.
\] Then \(H^2=5\), \(K_X=e_1+e_2+e_3\), \(H K_X=3\), and \(p_a(H)=5\). The strict transform \(E'=E-\sum e_i\) has \(H E'=0\), \(E'^2=-3\), and genus one. Each \(e_i\) has image degree one. Because the three points are collinear in the degree-eight embedding, hyperplanes through them impose only two independent conditions on \(H^0(A)\), leaving four sections. Under the genericity assumptions in the construction this gives the quintic model, contracts \(E'\) to the triple point, and realizes the canonical divisor as the three exceptional lines.

The genus-five K3 surface containing this plane elliptic cubic is a special, trigonal model; one must not automatically describe it as a complete intersection of three quadrics. Beauville (1978, VIII, Exercise 11) explicitly distinguishes this situation. The genus and intersection computations above verify the claimed birational model numerically; the particular ordinary-singularity presentation still uses the generic projection assumptions of Enriques's construction.

\section[Punti multipli di prima specie]{Punti multipli di prima specie (pp.~76--77)}\label{commentary-III-7}

A proper singular point is of first species when, after removing its first exceptional cycle from the hyperplane system, no further proper fundamental curve remains. The plane-section test is then sufficient: an adjoint numerator, together with one line through the point in the section plane, must be a plane adjoint to the constrained section. Enriques proves this by Section 4 under the no-pencil-of-reducibles hypothesis. Where that hypothesis fails for plane sections, he directs the reader to use intersections with sufficiently high-degree surfaces instead.

A tacnode has general plane sections with two consecutive double points; an oscnode has three. The text obtains:

A first-species tacnode requires the canonical numerator to pass through the point, giving one virtual condition. A first-species oscnode requires passage through the point and the prescribed tangent plane, giving three virtual conditions.

Here ``have the prescribed tangent plane'' means that the linear term is a multiple of the linear equation of that plane. It gives three conditions in the local three-variable ring: one constant coefficient and two tangential linear coefficients.

For explicit modern checks, take \(z^2+f_4(x,y)+\cdots=0\), with \(f_4\) a product of four distinct linear factors. The weighted blowup of weights \((1,1,2)\) has discrepancy \(-1\); its exceptional curve is the genus-one double cover of \(\mathbb{P}^1\) branched at four points and has square \(-2\). The canonical ideal is \((x,y,z)\). For \(z^2+f_6(x,y)+\cdots=0\), weights \((1,1,3)\) give discrepancy \(-2\), an exceptional curve of genus two and square \(-2\), and ideal \((x^2,xy,y^2,z)\). These calculations prove the stated local counts for these ordinary models. They do not equate all singularities sharing a coarse plane-section description.

\subsection{The tacnodal quintic and its quartic model}

For the quintic with a double line and a tacnode off the line, the unique canonical plane through both contains, residually, a conic and a line tangent at the tacnode. The conic and line are exceptional in the map to a smooth quartic K3 surface containing a line \(a'\).

On that quartic let \(H\) be a hyperplane section and \(E=H-a'\) the elliptic cubic pencil. Then \[
\begin{aligned}
H^2&=4,\quad E^2=0,\quad HE=3,\\
D&=2H-a'=H+E,\quad D^2=10,\quad p_a(D)=6.
\end{aligned}
\] The system \(|D|\) has seven sections. At a general \(P\in E\), imposing a double point removes three independent first-jet conditions, leaving four sections. On \(E\), the remaining divisor has degree \(3-2=1\), so there is an automatic residual base point \(B\) with \(2P+B\sim H|_E\). For general \(P\), \(B\ne P\). Blowing up \(P,B\) gives \[
\widetilde D=D-2e_P-e_B,\quad
\widetilde D^2=5,\quad
K_X=e_P+e_B,\quad K_X\widetilde D=3.
\] The strict transform \(E'=E-e_P-e_B\) has square \(-2\), genus one, and zero intersection with \(\widetilde D\). The two exceptional curves have image degrees two and one. This checks the quintic, the tacnode, and the canonical conic-plus-line description.

The two references to \(C_4\) in the inverse construction must read \(C_7\), the degree-seven residual cuts by quadrics through \(a'\).\footnote{The p. 77 scan prints \(C_4\) twice, after correctly introducing \(C_7\). The wrong subscripts are printed errors, not transcription errors.} The dimension calculation and the automatic base point are correct. The more special case \(B=P\), namely a flex, is relevant to the correction in Exercise 2 below.

\section[Piani doppi]{Piani doppi (pp.~77--79)}\label{commentary-III-8}

A double plane has function field \(\mathbb{C}(x,y)(\sqrt f)\). Replacing \(f\) by \(f g^2\) does not change that field. A branch curve of odd affine degree must be completed by the line at infinity; in projective language the branch divisor has even degree \(2n\).

Enriques uses the ordinary projective closure of \(z^2=f_{2n}(x,y)\), of degree \(2n\). Its point at infinity on the \(z\)-axis has multiplicity \(2n-2\). A general plane through the axis gives a hyperelliptic curve of genus \(n-1\). The plane-curve arithmetic genus is \((2n-1)(n-1)\), so the total delta contribution at infinity is \(2(n-1)^2\). The initial multiplicity contributes \((n-1)(2n-3)\), leaving \(n-1\) double-point contributions infinitely near. This is the source's count. Its ambient canonical numerator factors as the infinity-plane equation to power \(n-1\), times a degree-\(n-3\) form; the factor accounts for exceptional intersections at infinity.

\subsection{Intrinsic double-cover formulas}

Let \(f:S\to Y\) be a smooth double cover with branch \(B\in|2L|\). Directly from the algebra \(\mathcal O_Y\oplus L^{-1}\) and the ramification formula, \[
f_*\mathcal O_S=\mathcal O_Y\oplus L^{-1},\qquad
K_S=f^*(K_Y+L).
\] The projection formula then gives \[
H^0(S,mK_S)=H^0(Y,mK_Y+mL)
\oplus H^0(Y,mK_Y+(m-1)L).
\tag{III.4}
\] These basic identities are stated in Calabri and Ferraro (2002, Section 3, equations (3.2)--(3.3)); their Section 14 applies the same decomposition to pluricanonical conditions. For a smooth branch curve of degree \(2n\) in \(\mathbb{P}^2\), they yield, by direct substitution, \[
p_g=\binom{n-1}{2},\quad q=0,\quad K_S^2=2(n-3)^2,
\] with \(h^0(\mathcal O(d))=0\) for \(d<0\). Canonical divisors are represented by degree-\(n-3\) plane curves because the second canonical summand is \(H^0(K_{\mathbb{P}^2})=0\). For bicanonical divisors the second summand can be nonzero; this is the point of the footnote on pp.~91--92.

\subsection{The branch blowup rule and unloading}

Blow up a point of multiplicity \(m\) in the current branch divisor, put \(r=\lfloor m/2\rfloor\), and let \(E\) be the new exceptional curve. Normalize the pulled-back double cover. Its branch divisor and half-branch class are \[
B'=\widetilde B+(m\bmod2)E,\qquad
L'=\sigma^*L-rE.
\] Thus \[
K_{Y'}+L'=\sigma^*(K_Y+L)-(r-1)E.
\tag{III.5}
\] An even point of multiplicity \(2i\) imposes multiplicity \(i-1\) on the canonical plane numerator. An odd point of multiplicity \(2i+1\) has the same first correction, but the exceptional curve is itself branch; consequently a following point must be tested against the transformed branch divisor including that exceptional component. This is why replacing an odd cluster by a list of independent ordinary multiplicities would be wrong.

For two successive odd points \([2i+1,2i+1]\), the initial formal multiplicities \((i-1,i)\) violate proximity. Subtracting the forced exceptional component, or ``unloading,'' replaces them by \((i,i-1)\). In particular a \([3,3]\) branch point imposes one ordinary base point on the degree-\(n-3\) canonical curves. A branch point of multiplicity \(2i\), or \(2i+1\) without further nonnegligible infinitely near data, contributes \(i(i-1)/2\) to the virtual genus loss.

For completeness, the Euler-characteristic calculation underlying the count is short. From \(f_*\mathcal O_S\) and Riemann--Roch, \[
\chi(\mathcal O_S)=2\chi(\mathcal O_Y)+\tfrac12 L(L+K_Y).
\] Substitute \(L'=\sigma^*L-rE\), \(K_{Y'}=\sigma^*K_Y+E\), \(E^2=-1\). The loss is \(r(r-1)/2\). Similarly \(K_X^2\) drops by \(2(r-1)^2\) before any exceptional \((-1)\)-curves of the covering surface are contracted. This is a local virtual calculation; global dependencies remain as in Section 5.

The historical paragraph credits Clebsch with rational double planes of branch degrees two and four and with sextics having a quadruple point; Geiser's cubic-surface projection supplies the quartic-branch example. It credits Noether (1878) with reduction to three birational types, related through De Paolis's transformations to Bertini's plane involutions; the third type has a sextic branch curve with a {[}3,3{]} point. Enriques's own 1896 memoir treats genus-one double planes. These are citations of earlier classifications, not classifications proved in this section. The Clebsch date is a printed bibliographic error:\footnote{The p. 79 scan reads ``Math. Annalen, Bd. III, 1861.'' The identified paper is dated 1870 by the publisher and often 1871 by the bound-volume convention, including the bibliography of Dolgachev (2012); neither supports 1861.} the identified paper is in volume 3, pp.~45--75 (Clebsch, 1870/1871). The other historical references are recorded as citations used by Enriques, not as papers whose full proofs have all been independently inspected in the commentary.

\section[Genere delle curve fondamentali]{Il genere delle curve fondamentali e l'influenza sulle aggiunte (pp.~80--82)}\label{commentary-III-9}

Let \(E=\theta\) be an irreducible first-species fundamental curve with \(p_a(E)=\rho\), \(E^2=-s<0\), and \(H E=0\). In this first-species single-cycle situation the source identifies \(s\) with the multiplicity of the contracted point. That identification uses the actual hyperplane/fiber cycle; it does not apply to an arbitrary reduced exceptional support. Write \(H=C_1+E\). Then \(C_1E=s\). Since \(C'=K_X+H\) and \(E'=K_X+E\), \[
C'=C_1+E',\qquad
E C'=E C_1+E E'=s+2\rho-2.
\] Enriques's adjunction identity here is correct.\footnote{The p. 80 scan has unprimed \(C_1\) in both lines. The transcription inserts a prime in the second identity and proposes a compensating correction to the first; neither change is warranted.}

If the subadjoint system extends \(|C'|\) by \(xE\), and \(E\) is not to remain a forced fixed component in the enlarged system, a necessary condition is \[
(C'+xE)E=2\rho-2-(x-1)s\geq0,
\qquad
x\leq1+\frac{2\rho-2}{s}.
\tag{III.6}
\] The preceding prose says the extension stops when intersection becomes ``zero or negative,'' whereas the displayed inequality correctly permits zero. Degree zero is the boundary case: a nontrivial degree-zero line bundle has no sections; the trivial one has one. Read the prose as a warning about that boundary, not as a strict inequality.

The exact sequence for adding the \(x\)-th copy is \[
0\to\mathcal O_X(C'+(x-1)E)
\to\mathcal O_X(C'+xE)
\to\mathcal O_E(K_E+(x-1)E)\to0.
\tag{III.7}
\] This gives two possible obstructions beyond the degree bound: the line bundle on \(E\) may have no section, and a section may fail to lift because of \(H^1\) on \(X\). Sufficient positivity in the external hyperplane twist removes the latter obstruction, not the former. Enriques explicitly demonstrates that the numerical bound need not be attained; there is no missing assertion of universal sufficiency.

For \(\rho=0\), a proper first-species exceptional curve has \(s\geq2\), so (III.6) forces \(x=0\). Such rational singularities impose no canonical correction; the source includes an ordinary conical double point and an ordinary triple point that is simple on the double curve. For \(\rho=1\), the first boundary bundle is \(K_E\simeq\mathcal O_E\), giving the familiar one-condition correction. More generally, a degree-\(s\) tangent cone with \(\binom{s-1}{2}-1\) ordinary double generators has normalization genus one. Its subadjoints can have order \(s-3\), whereas true adjoints have order \(s-2\); the extra exceptional copy expresses the single elliptic correction. The tacnode has \((\rho,s)=(1,2)\). The oscnode has \((2,2)\), permits \(x=2\), and contributes three canonical conditions, as already computed.

Enriques's counterexample to attaining the upper bound is particularly instructive. An ordinary sextuple tangent cone with six nodes has normalization of genus \(10-6=4\). The six nodes correspond to six branches of the surface's double curve. Here \(s=6\), so (III.6) permits \(x=2\). Let \(h\) be the plane-section divisor on the normalized sextic, and let \(\Delta\) be the twelve preimages of its six nodes. Then \[
K_E=3h-\Delta,\qquad
\mathcal O_E(E)=\mathcal O_E(-h),\qquad
K_E+E=2h-\Delta.
\] This last bundle has degree zero. It is trivial exactly when a conic passes through the six nodes. For general six nodes there is no such conic. Hence the last possible numerical enlargement does not occur: the subadjoint has order three at the point, whereas the adjoint has order four, so only one copy is added. This supplies a precise line-bundle explanation of the source's geometric example.

\section[Singolarità di specie superiore]{Singolarità puntuali di specie superiore (pp.~82--89)}\label{commentary-III-10}

\subsection{Successive fundamental cycles}

Enriques defines higher species inductively: after the first fundamental part is removed, a further fundamental part remains for the residual system. His descriptions of connectedness and monovalence concern the successive systems, not just the unweighted support of an exceptional divisor. The conditions on sections must therefore be computed with the actual multiplicities and proximity relations. A monovalent removal is the exceptional cycle forced by one point condition; adding the appropriate next cycle gives a bivalent removal. Independence of these successive conditions is secured by sufficiently positive jet separation, and must not be presumed for every small system.

On a modern resolution there are at least three relevant objects: the reduced exceptional divisor, the smallest positive exceptional cycle with nonpositive intersection with every exceptional component, and the fiber cycle defined by the divisorial part of the pulled-back maximal ideal. They can differ. Calabri and Ferraro (2002, Theorem 9.7 and Sections 10--11) explicitly compute the fiber and fundamental cycles for double points. Their distinction provides a precise framework for the source's successive-removal procedure; it does not make the historical word ``fundamental curve'' automatically synonymous with one modern cycle.

\subsection{A uniplanar double point and the cubic-surface calculation}

An ordinary uniplanar double point has a double tangent plane and three exceptional tangent directions. A local model is \[
z^2+xy(x+y)=0,
\] a rational double point of type \(D_4\). Its minimal resolution has a central rational curve \(E_0\) and three arms \(E_1,E_2,E_3\), all of square \(-2\), with \(E_0E_i=1\) and pairwise disjoint arms. The maximal-ideal cycle is \[
Z=2E_0+E_1+E_2+E_3,\qquad Z^2=-2,\qquad p_a(Z)=0.
\] The exceptional intersection matrix is negative definite, and \(K_XE_i=0\) for all four components by adjunction. The resolution is crepant, so the point imposes no canonical adjunction condition.

Enriques checks the configuration using a rational cubic surface represented by plane cubics through six points, three collinear and three infinitely near. A quartic de Jonquières transformation of type \((4;3,1^6)\) changes the hyperplane system to \[
D=9h-6p-3(a_1+a_2+a_3)-2(b_1+b_2+b_3)-(o_1+o_2+o_3).
\] All the letters on the right denote total-transform classes on the blown-up plane. Direct calculation gives \[
D^2=81-36-27-12-3=3,\qquad K_XD=-27+6+9+6+3=-3,
\] so \(p_a(D)=1\), as required for a cubic surface.

The central line has class \(E_0=h-a_1-a_2-a_3\); the arms have \(E_i=h-p-b_i-o_i\). Their weighted sum is \[
Z=5h-3p-2\sum a_i-\sum b_i-\sum o_i.
\] Hence \[
D-Z=4h-3p-\sum a_i-\sum b_i.
\] This residual system has square one and arithmetic genus zero. It meets \(E_0\) once and is orthogonal to the three arms. This verifies both the stated quartic residual system and the continued fundamentality of the arms. In the source \(\theta_0\) already denotes the central line counted twice in the relevant cycle; omitting that weight would produce a false degree calculation.

\subsection{The Noether double point with an infinitely near tacnode}

For the next singularity, a convenient local model is \[
z^2+x^3+y^6=0.
\tag{III.8}
\] Its plane branch curve has three smooth branches with common tangent and distinct second-order terms, hence a \([3,3]\) cluster. The weights \((2,1,3)\) give weighted degree six and discrepancy \(-1\). Its minimal resolution has one elliptic exceptional curve \(E\) of square \(-1\). A canonical numerator must lie in the maximal ideal, so the local canonical loss is one.

The arithmetic of pp.~85--87 is most clearly checked on two explicitly different models. On the minimal resolution, a hyperplane pullback \(H\) has \(HE=0\), while \((H-E)E=1\). A degree-one line bundle on an elliptic curve has one section, so this intersection is a fixed point of the residual system. Thus \[
(H-E)^2=H^2-1,
\] but after removing that remaining simple base point the effective moving degree is \(H^2-2\). This is the distinction between virtual and effective degree that Enriques explicitly makes on p.~86.

Blow up that fixed point. Let \(E_0\) be the new rational exceptional curve and \(E_1\) the strict transform of \(E\). Then \[
E_0^2=-1,\quad E_1^2=-2,\quad E_0E_1=1.
\] The reduced total transform \(E_0+E_1\) has genus one and square \(-1\). The maximal-ideal cycle of the original double point is instead \[
Z=2E_0+E_1,\qquad Z^2=-2,\qquad p_a(Z)=0.
\] Removing this cycle accounts for the first genus drop of one; a subsequent removal associated with \(E_1\) accounts for the tacnodal drop of two. Contracting \(E_0\) restores the elliptic curve of square \(-1\).

The genus computation on p.~87 is \(\pi=(\pi-1)+\rho+x-1\). Together with \(x\ge1\) and the known elliptic component, hence \(\rho\ge1\), it gives \(x=\rho=1\); mere nonnegativity would not suffice. The subsequent genus drop of two gives \(C_2E_1=2\) and \(E_1^2=-2\). This bookkeeping explains the intended conclusion without treating every occurrence of \(\theta\) as the same divisor on the same surface. The source passes between a model with a remaining base point and one where it is separated, and between reduced and weighted neighborhoods. Those changes are not made explicit enough for a modern proof. There is also a printed sign inconsistency:\footnote{The p. 87 scan has a minus in ``degree \(-\nu\)'' and in \(-\nu=-\nu_0-\nu_1+2=-1\), while the final formula is \(\nu=-1\). The signed convention adopted in the commentary corrects the left side of the middle identity to \(\nu\).} \(\nu\) is introduced as the signed square \(-1\), but p.~87 occasionally writes ``degree \(-\nu\)'' and \(-\nu=-\nu_0-\nu_1+2=-1\). Retain \(\nu=E^2=-1\), and take \(\nu_i=-E_i^2>0\); then the left side of the latter identity is \(\nu\), not \(-\nu\). The prime after \((\pi-1)\) in the scan has no mathematical role and should be treated as a stray printed mark.

The two rational quartic models cited from Noether are numerically consistent: \[
(7;3,2^9),\qquad (9;3^8,2,1).
\] In either case \(D^2=4\), \(K_XD=0\), and \(p_a(D)=3\). The cubic through the ten specified base points has genus one, square \(-1\), and zero intersection with \(D\). Its intersection with \(D-E\) is one, but that one point is fixed. Thus the statement that it is still fundamental for the residual system means there are no variable intersections; it does not assert that the intersection number remains zero. In the second model, the source itself explains that the tacnodal plane contains the original point and forces splitting of certain plane sections; using a more positive system removes that projective complication.

The source cites Pedoe (1937) for accidental canonical base points. The publication record and abstract were verified, but not the full article; its specific local assertion is therefore retained as a historical dependency, while the calculation above supplies the verification used here.

\subsection{The double-plane verification}

For a double plane branched in degree \(2n\), the pullback of a general conic has genus \(2n-1\), by Riemann--Hurwitz with \(4n\) branch points. A conic through the first point of a \([3,3]\) cluster removes two branch contributions modulo parity, lowering the genus by one; imposing the prescribed infinitely near condition produces a further loss of two. This matches the double point followed by a tacnode. The canonical cluster condition is a single base point, so the virtual canonical loss is one.

The two quartic rational examples project to sextic double planes with, respectively, an ordinary quadruple branch point and a \([3,3]\) branch point. Each imposes one condition on canonical plane constants, leaving no canonical section. These are consistency checks of the already constructed rational models; \(p_g=0\) alone would not prove rationality.

\subsection{Exercise 1: two infinitely near triple points}

Each ordinary triple point gives one canonical condition. For two successive triple points, these are a simple point and the prescribed infinitely near simple point; their local cluster has length two. For a quintic with a double line \(a\), the space of canonical planes through \(a\) has vector dimension two. Assuming the specified cluster is in the stated general position off \(a\), the two conditions remove both sections, giving \(p_g=0\).

Let \(\ell=OO_1\) denote the line in the prescribed tangent direction. If it were not contained in the quintic, its intersection multiplicity at the two successive triple points would be at least six, exceeding five. Hence \(\ell\subset Y_5\). Planes through \(\ell\) have residual quartics. Their plane arithmetic genus is three; the double-line intersection contributes one node and the cluster contributes the other two double-point losses. The general residual quartic is rational. For a complete modern proof of rationality from this pencil, resolve its base points. The smooth generic fiber over \(\mathbb{C}(t)\) has genus zero, hence is a conic via its anticanonical system. Tsen's theorem gives it a \(\mathbb{C}(t)\)-point, so it is \(\mathbb{P}^1_{\mathbb{C}(t)}\) and the surface function field is \(\mathbb{C}(t,u)\). This proves rationality independently of the unspecified spatial transformation (Stacks Project, Theorem 59.67.10, Tag 03RD).

The quadratic spatial transformation cited in the exercise has base the double line and three points of the surface, including \(O,O_1\). Under the genericity assumptions of that transformation, the image is a quartic with a double line, which is a standard rational model. The text does not write coordinates for the transformation; the genus calculation, line containment, and rational pencil above give the requested substantive verification without pretending that its full base scheme was displayed.

\subsection{Exercise 2: a flex imposes one dependent triple-point condition}

The conclusions of Exercise 2 are stated without proof. The natural full-system interpretation contains a substantive dimension error, and its proposed projective model requires a specified subsystem. Let \(S\subset\mathbb{P}^3\) be a smooth quartic containing a line \(a\), let \(H\) be its hyperplane class, and let \(E=H-a\) be the elliptic cubic pencil. Set \(D=H+2E\). Then \[
D^2=16,\qquad p_a(D)=9,\qquad h^0(S,D)=10.
\] These agree with the source. Let \(P\in E\) be a flex of the plane cubic \(E\), and blow it up, with exceptional curve \(e\). Put \[
A=\pi^*D-3e,\qquad E'=\pi^*E-e,\qquad
B=A-E'=\pi^*(H+E)-2e.
\] The numerical assertions also agree: \[
A^2=7,\quad K_XA=3,\quad p_a(A)=6,\quad
E'^2=-1,\quad AE'=0.
\tag{III.9}
\]

The complete imposed subsystem has projective dimension four.\footnote{The p. 89 scan prints \(\infty^3\). Although the final sentence does not repeat ``complete,'' imposing the triple-point condition on the previously specified complete \(\infty^9\) naturally means the full linear subspace. The calculation below gives \(\infty^4\); a specified \(\infty^3\) subsystem recovers the stated septimic presentation.} Here is a proof that also explains the dependence. The pencil \(|E|\) is globally generated and \(H\) is very ample, so \(H+E\) is very ample. It has \(h^0=7\), and its first-jet evaluation at \(P\) is surjective. Therefore \[
h^0(X,B)=7-3=4.
\] Riemann--Roch on \(X\) gives \(\chi(B)=4\), while \(h^2(B)=0\), so \(h^1(B)=0\). Since \(P\) is a flex, \[
A|_{E'}\simeq H|_E(-3P)\simeq\mathcal O_E.
\] The exact sequence \[
0\to B\to A\to\mathcal O_{E'}\to0
\] then gives \[
h^0(X,A)=5,\qquad \dim|A|=4.
\tag{III.10}
\] Equivalently, \(\chi(A)=4\) and \(h^1(A)=1\). The flex is exactly the reason that the six nominal triple-point conditions have rank five. If \(P\) is not a flex, \(A|_{E'}\) is a nontrivial degree-zero bundle, all four remaining sections contain \(E'\), and it is a fixed component rather than the intended contracted curve.

The source does not literally say ``complete'' in its last sentence. Nevertheless, ``the system obtained by imposing a triple point'' on the previously specified complete \(\infty^9\) naturally denotes the full imposed subsystem, and that reading requires \(\infty^4\). A specified projective three-dimensional subsystem does recover the claimed projective model. The following reconstruction preserves the exercise's projective conclusion by making the projection choice and the additional generic hypotheses explicit.

\subsection{A complete construction of the special subsystem}

Assume \(E_0\) is a smooth fiber, \(P\in E_0\setminus a\) is its flex, the tangent plane has contact exactly two with \(S\) at \(P\), and \(S\) contains no line other than \(a\). The condition \(P\notin a\) is printed in the source. The last two are additional general-position hypotheses used for this construction. They hold for general input: imposing another line on a quartic containing \(a\) gives incidence dimensions at most \(24+4=28\) for a disjoint line or \(25+3=28\) for an intersecting one, inside the projective space of dimension 29 of quartics through \(a\). Exact tangent contact two is also available. In the equation \(gC+hT\) used below, keep \(C\) fixed and vary the coefficient \(\beta\) of \(-g^2l\) in \(T\). This changes the quadratic term \(\beta Ux\), while preserving \(a\), \(E_0\), and its flex; choose a smooth general quartic with \(\beta\ne0\). Global generation and the local elliptic calculation below do not require these last two hypotheses.

Choose hyperplane forms \(g,h,l,t\), with \(a=(g=h=0)\), \(h=0\) cutting \(a+E_0\), \(g(P)\ne0\), \(t=0\) the tangent plane, and \(l(P)=0\) a local parameter along \(E_0\). If \(s_a\) cuts out \(a\), put \(u=h/s_a\), \(v=g/s_a\); these generate \(H^0(E)\). Near \(P\) use \[
U=u/v=h/g,\qquad x=l/g,\qquad w=t/g.
\] The flex condition gives \[
w=\alpha U^2+\beta Ux+O(3),\qquad
\operatorname{ord}_x w(0,x)=3.
\] The seven sections of \(H+E\), in this local trivialization, are \[
1,\ U,\ x,\ U^2,\ Ux,\ w,\ Uw.
\] Thus \(B=H+E-2e\) has basis \(U^2,Ux,w,Uw\), represented globally by \(hu,lu,tv,tu\). Away from \(P\), their common vanishing is impossible: if \(u\ne0\), it forces \(h=l=t=0\), whose only ambient point is \(P\); if \(u=0\), it forces \(t=0\) on \(E_0\), whose intersection with the tangent plane is \(3P\). On \(e\), the initial quadrics span \(U^2,Ux\), with the single common zero \(E'\cap e=[0:1]\). Hence \(|B|\) has just that base point. The established surjection \(H^0(A)\to H^0(\mathcal O_{E'})\), together with \(E'+|B|\subset|A|\), removes it and proves that \(|A|\) is globally generated.

Set \(M=A-E=H+E-3e\). The unique tangent-plane section gives \(h^0(H-2e)=1\); Riemann--Roch gives \(h^1(H-2e)=0\). Since \(M=(H-2e)+E'\) and \(M|_{E'}=\mathcal O_{E'}\), one obtains \(h^0(M)=2\). The pencil trick identifies the kernel of \[
H^0(M)\otimes H^0(E)\longrightarrow H^0(A)
\] with \(H^0(H-3e)\). Exact tangent contact two makes this kernel zero. Choose bases \(m_0,m_1\) of \(H^0(M)\) and \(u,v\) of \(H^0(E)\), and a fifth section \(r\). The complete map has image \[
Y\subset Q:\ z_{00}z_{11}-z_{01}z_{10}=0\subset\mathbb{P}^4,
\] where \[
(z_{00},z_{01},z_{10},z_{11},z_4)
=(m_0u,m_0v,m_1u,m_1v,r).
\] The quadric has rank four. Since \(A^2=7\) and the image is nondegenerate, the map is birational and \(\deg Y=7\).

The pencil \(|M|\) has no fixed curve. Its local generators are \(Uw\) and \(w-\alpha U^2-\beta Ux\). The restriction to \(E'\) excludes \(E'\) as a fixed component, and the exact order three of \(Uw\) excludes \(e\). Any other possible fixed curve would be the line \(t=\alpha h+\beta l=0\), excluded by the hypothesis on lines. Because \(M^2=1\), its base scheme consists of one reduced point, disjoint from \(E'\). Blow up that point. The exceptional curve is a section of the resulting pencil, so the general fiber is connected; Bertini gives a smooth irreducible member of genus three. Its image is a plane quartic, since \(AM=4\); birationality on a general member and the equality of its genus with the plane-quartic arithmetic genus three show that this plane image is smooth. A general \(E\)-fiber maps to a smooth plane cubic: \(AE=3\), and \(h^0(A)-h^0(M)=3\) supplies the complete degree-three series.

Choose general members \(M_1=(m_1=0)\), \(E_1=(v=0)\) whose ruling planes avoid \(q=\varphi_A(E')\). Those planes are \[
\Pi_4=(z_{10}=z_{11}=0),\qquad
\Pi_3=(z_{01}=z_{11}=0).
\] Choose \(R=(1,0,0,0,\lambda)\) on their common generator, outside \(Y\) and outside the quadric vertex. The required subsystem is now explicit: \[
V_R=\langle m_0v,\ m_1u,\ m_1v,\ r-\lambda m_0u\rangle.
\] It has vector dimension four and is base-point-free because \(R\notin Y\). In target coordinates \([X:Y:Z:W]\), after shifting the fifth coordinate by \(\lambda z_{00}\), the inverse of projection on \(Q\) is \[
[X:Y:Z:W]\longmapsto[XY:XZ:YZ:Z^2:WZ].
\] Projection is therefore birational. It contracts \(\Pi_4\) to \(L_4=(Y=Z=0)\) and \(\Pi_3\) to \(L_3=(X=Z=0)\). Their curves map onto these lines with degrees four and three.

These are the actual generic sheet multiplicities. Near a general point of \(L_4\), use \(X=1\) and coordinates \((\tau,\rho,w)\) on \(Q\) with \[
z_{00}=\tau,\quad z_{01}=1,\quad z_{11}=\rho,\quad
z_{10}=\tau\rho.
\] Projection is \((\tau,\rho,w)\mapsto(y=\tau\rho,z=\rho,w)\). The plane quartic meets a general projection fiber in four distinct points \(\tau_i\); each local sheet becomes \(y=z\tau_i(z,w)\). Their tangent planes are distinct. Thus \(L_4\) is an ordinary quadruple line at its general point; the same calculation gives the ordinary triple line \(L_3\).

A general plane section has normalization of genus six and already has delta contribution \(6+3=9\) from those two lines, accounting for the whole difference \(15-6\). Hence no further singular curve occurs. Nor does \(A\) contract another curve: for an irreducible \(C\ne E_0\), with \(m=\operatorname{mult}_P C\), one has \(EC\ge m\) and \(HC\ge m\); equality \(AC=0\) forces \(HC=m\), which makes \(C\) a line through \(P\), excluded by hypothesis. Also \(Ae=3\). Away from the two multiple lines and the image of \(E'\), the proper birational map is finite onto a normal hypersurface and hence an isomorphism. The remaining isolated singularity is checked directly next.

\subsection{The double point and its infinitely near tacnode}

Write the original quartic equation as \[
gC(g,l,t)+hT(g,h,l,t)=0,
\] where the smooth cubic \(E_0\) and its flex allow \[
C=g^2t+c_1glt+c_2gt^2+\kappa l^3+c_3l^2t+c_4lt^2+c_5t^3,
\quad\kappa\ne0,
\] and \[
T=-\alpha g^2h-\beta g^2l+\gamma g^2t+gQ_2(h,l,t)+Q_3(h,l,t).
\] Here \(Q_2,Q_3\) are homogeneous forms of degrees two and three. A convenient basis for the five sections of \(A\), expressed as local ratios, is \[
[a_0:b_0:c_0:d_0:f_0]
=[U^3:U^2x:Uw:U^2w:w-\alpha U^2-\beta Ux].
\] These coordinates satisfy \[
d_0f_0=c_0(c_0-\alpha a_0-\beta b_0),\qquad
q=[0:0:0:0:1].
\] On \(f_0=1\), write \(a,b,c\) for the first three coordinates and \(s=c-\alpha a-\beta b\). Then \(d=cs\), and the inverse coordinates on the quartic are \[
[g:h:l:t]=[a:as:bs:cs^2].
\] Substituting into the quartic and dividing out the factor \(as^3\) gives the image equation \[
\begin{aligned}
\Phi={}&a^2+c_1abc+c_2ac^2s+\kappa b^3+c_3b^2cs
+c_4bc^2s^2+c_5c^3s^3\\
&+\gamma a^2c+aQ_2(a,b,cs)+sQ_3(a,b,cs)=0.
\end{aligned}
\] For weights \((a,b,c)=(3,2,1)\), its initial form is \[
\Phi_6=a^2+c_1abc+c_2ac^3+\kappa b^3
+c_3b^2c^2+c_4bc^4+c_5c^6.
\] Every omitted term has weight at least seven. The weighted projective curve \(\Phi_6=0\) is the smooth \(E_0\), via \([g:l:t]=[a:bc:c^3]\). Its cone has an isolated critical point. Consequently \(\Phi=0\) is an isolated normal double hypersurface germ. The first four sections vanish precisely along \(E'\), so its exceptional fiber is the single elliptic curve \(E'\), with square \(-1\).

Blow up this double point in the \(c\)-chart, \(a=cz,b=cy\). Analytic completion of the square gives \[
\begin{aligned}
0={}&Z^2+c\bigl[\kappa y^3+(c_3-c_1^2/4)cy^2\\
&\qquad +(c_4-c_1c_2/2)c^2y+(c_5-c_2^2/4)c^3\bigr]\\
&+\text{terms of degree at least five}.
\end{aligned}
\] Terms of original weight at least seven contribute order at least five to this branch equation: the critical value of \(z\) has order at least two, so an original monomial of weight \(w\) contributes order at least \(w-2\) after division by \(c^2\). The displayed binary quartic has four distinct linear factors: its cubic factor has distinct roots because \(E_0\) is smooth, and \(\kappa\ne0\) separates the factor \(c\). This is the ordinary tacnode of the source. It is the unique singular point on the first exceptional line; the other charts are smooth there. The chosen projection is an isomorphism near \(q\), because \(q\) belongs to neither contracted ruling plane. Thus the septimic has exactly the stated double point with an infinitely near tacnode, under the explicit generic hypotheses.

\subsection{The adjoint cubic}

Let \(\Pi_0=(Z=0)\), the plane containing both multiple lines, and let \(\Pi_q\) be the plane through \(L_4\) and \(q\). On \(X\) their pullbacks are \(M_1+E_1\) and \(M_1+E'+e\). Hence \[
\operatorname{div}(\Pi_0^2\Pi_q)=3M_1+2E_1+E'+e.
\] Removing conductor orders three and two along the ordinary multiple lines and the elliptic adjoint correction \(E'\) leaves \(K_X=e\). This proves the asserted adjoint cubic. Equivalently, the cubic forms in \((Y,Z)^3\cap(X,Z)^2\) are \(Z^2(aY+bZ)\), and passing through \(q\) fixes their ratio. Its residual curve is the plane cubic image of \(e\), whose degree is \(Ae=3\). It is rational, with a cusp at \(q\), as the exceptional-curve parametrization \([U^3:U^2x:f_3(U,x)]\) shows.

If tangent contact at \(P\) is at least three, \(\alpha=\beta=0\) and the quadric drops to rank three. The dimension correction remains valid, but this proof does not assert the same two-line presentation in that special case.

The inverse requested by the exercise has a clean intrinsic answer on the common smooth model: \[
\pi^*H=A+K_X-2E'.
\tag{III.11}
\] Indeed \(K_X=e\), so the right side is \[
(H+2E-3e)+e-2(E-e)=H.
\] Thus the complete system \(|A+K_X-2E'|\) contracts \(e\) and recovers the original quartic embedding. There is also an explicit ambient inverse. In the complete-image coordinates above, put \(s=c_0-\alpha a_0-\beta b_0\). The four quadrics \[
[g:h:l:t]=[a_0f_0:a_0s:b_0s:d_0s]
\tag{III.11a}
\] recover the original quartic coordinates on the common dense open set; this follows immediately from \(d_0f_0=c_0s\). Compose them with the quadratic inverse \([XY:XZ:YZ:Z^2:WZ]\) of the chosen projection, and the intervening linear change from the product basis to \((a_0,b_0,c_0,d_0,f_0)\). This gives four explicit degree-four forms defining the rational inverse from the septimic to \(S\). Their common zero set is the base locus of this chosen presentation; common factors can be cancelled. On the smooth common model the resolved inverse is exactly (III.11). No minimal-degree assertion is made for these quartic representatives. For a fully expanded prescription when \(\beta\ne0\), choose the global sections \(m_0=tu\), \(m_1=tv-\alpha hu-\beta lu+\mu tu\) with general \(\mu\), and complementary section \(r=hu^2=a_0\). Given \([X:Y:Z:W]\), form the quadrics \[
d=XY,\quad c=XZ,\quad s=Y(Z-\mu X),\quad f=Z(Z-\mu X),
\] \[
a=WZ+\lambda XY,\qquad b=(c-s-\alpha a)/\beta.
\] Then the quartic representatives are simply \([af:as:bs:ds]\). In this basis they share the factor \(Z-\mu X\). Cancelling it gives the cubic representatives \([aZ:aY:bY:XY^2]\), without asserting minimality. This illustrates why a chosen presentation degree must be distinguished from the rational map itself.

\subsection{The chapter's explicit limit}

The note on p.~89 expressly says that the preceding examples do not exhaust the singularities or fundamental curves. It asks the reader to examine chains of quadruple points, longer chains of triple points, and mixed clusters. This is not an unproved assertion of a complete classification. The modern branch blowup algorithm (III.5), together with proximity and the actual fiber cycle, is the appropriate extension. Franchetta's Notes I--II are cited on p.~89 as 1945; Calabri and Ferraro's bibliography (2002, reference 13) gives series VIII, volume 1 (1946), pp.~49--57 and 162--168.\footnote{The scan confirms the printed 1945 date. The 1946 metadata come from the modern paper's bibliography, not an independent inspection of the original Lincei scans.}

\section[Curve e superficie biaggiunte]{Curve bicanoniche e superficie biaggiunte (pp.~89--92)}\label{commentary-III-11}

For a degree-\(n\) projective surface, biadjoint forms have degree \(2n-8\), and forms representing \(2K_X+rH\) have degree \(2n-8+r\). The bicanonical space is \(H^0(2K_X)\), not the image of \(\operatorname{Sym}^2H^0(K_X)\). The latter can be zero while the former is nonzero, as the next section shows.

Along an ordinary double curve, a biadjoint vanishes to order two. Along an \(i\)-fold curve for \(i>2\), the intrinsic condition is twice the canonical vanishing on each separated sheet. It need not be represented by ambient order \(2i-2\): in the local ring of \(Y\), forms differing by a multiple of its equation agree. In transverse local coordinates, a useful expression is \[
g\in I^{\,2i-2}+(f),
\] where \(I\) is the ideal of the multiple curve. A representative can have ambient order \(i\), with the leading and subsequent jets chosen to agree with multiples of the equation of \(Y\). For a triple curve, a third-order numerator with the appropriate tangent-cone contact represents the same intrinsic condition as a fourth-order one.

The source proposes combining \(\psi\) with the defining equation \(f\) of \(Y\). This must be understood locally, or after multiplication by forms that equalize their degrees. When \(2n-8<n\), there is no nonzero global degree-\((n-8)\) multiplier available. It would be incorrect to add two homogeneous equations of unequal degree literally. This is a clarification of representatives; it does not invalidate the intrinsic contact condition or the complete biadjoint system.

At an ordinary isolated \(s\)-fold point, the intrinsic numerator condition is \[
g\bmod f\in\mathfrak m^{\,2s-4}\mathcal O_Y.
\tag{III.12}
\] For \(s=3\), the ambient numerator has a double point. For \(s=4\), it has a quadruple point. For \(s\geq5\), it may have order \(s\) together with the extra contact making its class satisfy (III.12). Thus the source's general multiplicity-\(s\) phrase is to be read together with its explicit \(s=3\) exception. In particular, at a quintuple point, matching the tangent cone supplies the next necessary order. The local count, if the tangent cone is ordinary, can be checked from \[
\ell\frac{\mathbb{C}[[x,y,z]]}{(f,\mathfrak m^{2s-4})}
=\binom{2s-2}{3}-\binom{s-2}{3},
\] with the second binomial taken as zero when appropriate. This local refinement is obtained by subtracting the multiples of the degree-\(s\) initial equation; as elsewhere, it is a local condition count, not automatically a global loss.

\subsection{Canonical and bicanonical extension ideals}

For the tacnode \(z^2+f_4(x,y)=0\), the relevant weighted valuation has \(v(x)=v(y)=1\), \(v(z)=2\). A canonical numerator has valuation at least one, but a bicanonical numerator needs valuation at least two. Hence \[
J_1=(x,y,z),\qquad J_2=(x^2,xy,y^2,z).
\] The second ideal has colength three. It is larger than \(J_1^2\), which would impose the unnecessary vanishing of the linear \(z\)-term. This is exactly Enriques's condition: pass through the tacnode and have the prescribed tangent plane. His plane-section ``unloading'' is an elementary way to obtain the same result.

For the Noether point (III.8), \(v(x)=2,v(y)=1,v(z)=3\). The analogous ideals are \[
J_1=(x,y,z),\qquad J_2=(x,z,y^2),
\] of colength one and two respectively. The two bicanonical conditions are the first point and its prescribed infinitely near point. In a double-plane calculation, the formal cluster \((0,2)\) unloads to \((1,1)\). This verifies the stated bigenus loss of two for a \([3,3]\) branch point.

The degree-\(n-4\) plane adjoint named on p.~91 is an index slip: adding a line to the section of a degree-\(n-4\) surface gives degree \(n-3\).\footnote{The scan prints \(\varphi_{n-4}\); the corrected expression is \(\varphi_{n-3}\).}

\subsection{The complete footnote on pp.~91--92}

The long footnote is read continuously across pp.~91--92.\footnote{Both parts are present in the scan of the print. The earlier transcription comments saying that the continuation or its reference mark was outside the packet are incorrect.}

The footnote correctly warns that a complete system downstairs need not pull back to a complete system upstairs. For a smooth double plane branched in degree \(2n\), (III.4) gives \[
P_2=h^0(\mathbb{P}^2,\mathcal O(2n-6))
+h^0(\mathbb{P}^2,\mathcal O(n-6)).
\tag{III.13}
\] The second summand first appears at \(n=6\). For example, a smooth branch curve of degree twelve gives \(28+1=29\) bicanonical sections, whereas pulling back only sextics gives 28. This is a direct computation from the cover algebra, not a failure of completeness of the plane system itself.

The quartic-branch illustration in the footnote concerns pullbacks of conics, not an effective bicanonical system of that rational double plane. Here \(L=\mathcal O(2)\), and \[
h^0(f^*\mathcal O(2))=h^0(\mathcal O(2))+h^0(\mathcal O)=6+1=7.
\] Thus the pulled-back conics form a projective five-dimensional subsystem of a complete projective six-dimensional system. The extra generator is the covering coordinate.

Finally, the footnote uses \(Y\) of degree \(2n\), yet labels its biadjoints \(\psi_{2n-8}\). They must have degree \(2(2n)-8=4n-8\). The printed multiplicity \(2n-2\) at the point at infinity is not the quantity to replace. It is compatible with the corrected degree and the necessary contact. For instance, an invariant representative can be written, in homogeneous coordinates with infinity equation \(w\), \[
w^{2n-2}h_{2n-6}(x,y,w),
\] of degree \(4n-8\). The anti-invariant summand, when present, has representatives of the form \[
z\,w^{3n-3}h_{n-6}(x,y,w).
\] Both degrees check. Multiples of the defining equation can change the ambient multiplicity while preserving the section on \(Y\). The smaller subscript is printed.\footnote{The p. 92 scan has \(\psi_{2n-8}\). With \(\deg Y=2n\), the corrected degree is \(4n-8\); the printed multiplicity \(2n-2\) is not the quantity to replace.}

For the ordinary oscnode \(z^2+f_6(x,y)=0\), the same weighted calculation extends to pluricanonical numerators: \(v(g)\ge2m\), with weights \((1,1,3)\). At \(m=2\), the excluded monomials are the ten monomials in \(x,y\) of degree at most three and the monomial \(z\), so the local bicanonical defect is eleven. The double-cover calculation agrees, giving \(3+2\cdot4=11\), the canonical Euler defect plus the bicanonical contribution of the canonical-square loss. This calculation is restricted to the ordinary weighted model; further infinitely near degenerations require their own valuations.

\section[Genere zero e bigenere positivo]{\texorpdfstring{Superficie di genere \(p_g=0\) e di bigenere \(P_2>0\) (pp.~92--96)}{Superficie di genere p\_g=0 e di bigenere P\_2\textgreater0 (pp.~92--96)}}\label{commentary-III-12}

\subsection{The tetrahedral sextic}

Positive bigenus excludes rationality by birational invariance of plurigenera. The first example is the tetrahedral sextic already announced in Chapter II, Section 8; its complete numerical verification follows here. The displayed equation means \[
Y_6=\varphi(x_2x_3x_4,x_3x_4x_1,x_4x_1x_2,x_1x_2x_3)
+x_1x_2x_3x_4 f(x_1,x_2,x_3,x_4)=0,
\] where both \(f\) and \(\varphi\) are quadratic forms in their respective four arguments. The absence of commas in the last argument list is nineteenth-century notation, not a degree-ten term \(f(x_1x_2x_3x_4)\) in one variable. A general member has double curves along the six edges and the stated ordinary triple and pinch-point behavior.

The source's ``double system'' must be distinguished from the smallest span of products of the four cubics \(t_i=\prod_{j\ne i}x_j\) through the edges. The ten products \(t_it_j\), \(i\leq j\), are distinct monomials, with exponent patterns \((2,2,2,0)\) or \((2,2,1,1)\). The entire degree-six space satisfying the double-edge conditions has dimension fourteen: there are also the four monomials \(x_1x_2x_3x_4x_i^2\), with pattern \((3,1,1,1)\). Indeed an exponent vector \((a_1,\ldots,a_4)\) of sum six must have \(a_i+a_j\geq2\) for every pair corresponding to a coordinate edge ideal; these three patterns exhaust the possibilities. The displayed \(\varphi(t)+x_1x_2x_3x_4f(x)\) spans this full space: its two ten-dimensional summands have a six-dimensional intersection. This is the distinction between the prescribed double-edge system and the minimum product system, not a source transcription error. On a model resolving the entire base ideal, one must additionally track the tetrahedron vertices: each cubic vanishes to order two at a vertex, while the four extra sextics have order three at the corresponding vertex. One should therefore not identify this fourteen-dimensional ambient space with the product span or with a doubled resolved divisor class without specifying those vertex conditions.

No quadric vanishes on all six edges. Indeed a quadratic monomial involves at most two coordinates, so it restricts nontrivially to at least one coordinate edge unless its coefficient is zero. Thus \(p_g=0\). The only quartic vanishing doubly along all six edges is, up to scalar, the product of the four coordinates: the required pairwise exponent inequalities force every exponent to be one. Its divisor on the normalization is twice the conductor divisor. Consequently \[
2K_X\sim0,\qquad K_X\not\sim0,\qquad P_2=1.
\tag{III.14}
\] The normalization is a complex Enriques surface. Dolgachev (2016, Section 2, pp.~3--4, equation (1)) gives the conductor calculation and verifies \(p_g=q=0\) for this model. The genericity qualification is essential: special equations may acquire further singularities or split.

For the hyperplane class, \[
H^2=6,\quad K_XH=0,\quad p_a(H)=4.
\] Because \(q=0\), the characteristic series is complete. It has degree six and projective dimension two and is nonspecial. To see the last assertion directly, Kodaira vanishing gives \(h^0(H)=4\), and the restriction sequence gives \(h^0(H|_C)=3\). A canonical series on a genus-four curve has four sections, so \(H|_C\) is not canonical. Nevertheless, \[
2H|_C\sim2K_C,
\] because \(K_C=(K_X+H)|_C\) and \(2K_X=0\). Thus the example exhibits a nontrivial two-torsion difference between the characteristic and canonical line bundles.

The adjoint system is \(|H+K_X|\), again degree six and genus four. Its double equals \(|2H|\), and applying adjunction twice returns \(|H|\). The source's formulas involving \((2C)''\), \(2C'\), and \(C+C''\) are correct once prime placement is respected. Finally \(K_X^2=0\), so \(p^{(1)}=1\). None of this says that \(K_X\) itself is effective: \(p_g=0\) says the opposite.

\subsection{Castelnuovo's septimic}

The second model is a degree-seven surface with a triple line \(r\), a double conic \(\Gamma\) disjoint from \(r\), and three tacnodes \(A_i\), whose tangent planes \(\alpha_i\) contain \(r\). Enriques constructs a pencil from two reducible surfaces: a quartic double along \(\Gamma\), tangent to the prescribed planes at \(A_i\), multiplied by \(\alpha_1\alpha_2\alpha_3\); and the square of a cubic through \(r,\Gamma\), multiplied by a plane through \(r\). The assertion that a general member is irreducible uses general position and the absence of a forced common component. ``Take a pencil'' alone would not establish irreducibility for arbitrary choices, so this part is an existence sketch rather than a complete parameter-space verification.

A canonical cubic must contain \(\Gamma\) and be double along \(r\). Its restriction to the plane of \(\Gamma\) is forced to contain that plane; the residual quadric is two planes through \(r\). The three tacnodes require three distinct planes through \(r\), which two planes cannot supply. Hence \(p_g=0\).

The exhibited biadjoints are sextics \[
(\text{plane of }\Gamma)^2\alpha_1\alpha_2\alpha_3\beta,
\] where \(\beta\) varies among planes through \(r\). They cut a pencil of plane quartics. The two intersections with the double conic account for two nodes, so the quartics have genus \(3-2=1\). The pencil has square zero and no base points.

The completeness argument needs more than square zero: the general residual quartic is irreducible, so the induced morphism \(f:X\to\mathbb{P}^1\) has connected fibers. Its fiber divisor \(F\) is the exhibited bicanonical divisor. Therefore \[
H^0(X,\mathcal O_X(F))
=H^0(\mathbb{P}^1,f_*\mathcal O_X\otimes\mathcal O(1))
=H^0(\mathbb{P}^1,\mathcal O(1)),
\] and \(P_2=2\). This supplies the missing justification for the source's short completeness sentence. In this model the contact conditions along the triple line force the stronger ambient quadruple vanishing; the intrinsic criterion was not simply ambient fourth order in every situation.

Since \((2K_X)^2=0\) and \(p_a(2K_X)=1\), the source's identities \[
4p^{(1)}-4=0,\qquad 3p^{(1)}-2=1
\] both give \(p^{(1)}=1\). On its minimal model this is a properly elliptic surface, rather than an Enriques surface: \(P_2=2\) and an elliptic bicanonical pencil distinguish it from (III.14). This conclusion uses the actual bicanonical fibration; no unverified list of multiple fibers is inferred.

\subsection{Godeaux and Campedelli: what the historical sentence establishes}

The source refers to the examples of Godeaux (1931) and Campedelli (1932), with respectively \(p_g=0,P_2=2\) and \(p_g=0,P_2=3\). These distinguish \(K_X^2=1\) and \(K_X^2=2\) general-type examples from the preceding \(K_X^2=0\) models. The broad preceding question asks also for a bicanonical system of dimension greater than one. Godeaux's \(P_2=2\) gives only a pencil; Campedelli's \(P_2=3\) supplies the net that meets that additional requirement. One should not turn the combined historical sentence into a claim that both examples have bicanonical projective dimension greater than one.

\subsection{An explicit verification of Campedelli's construction}

Choose nonzero \(a,b\) with \(a^2\ne b^2\), and put \[
q_1=x^2+y^2-a^2z^2,\quad
q_2=x^2+y^2-b^2z^2,\quad
q_3=x^2/a^2+y^2/b^2-z^2.
\] The first two conics are concentric circles in the affine real picture; the third ellipse touches them at the ends of perpendicular axes. Over \(\mathbb{C}\), their six pairwise contacts are \[
I=(1:i:0),\ J=(1:-i:0),\
P_\pm=(\pm a:0:1),\
Q_\pm=(0:\pm b:1).
\] They are distinct. A conic through \(I,J\) has equation \[
\lambda(x^2+y^2)+\ell xz+myz+cz^2=0.
\] The two \(P\)'s force \(\ell=0,c=-\lambda a^2\), and the two \(Q\)'s force \(m=0,c=-\lambda b^2\). Hence all coefficients vanish. There is no conic through the six points. This proves the canonical vanishing without a generic-position slogan.

The quartics having the six prescribed tangent contacts form exactly the vector space \[
\langle q_1q_2,\ q_1q_3,\ q_2q_3\rangle.
\tag{III.15}
\] To prove completeness, restrict any such quartic \(G\) to \(q_1\). Its eight intersections are the four contacts counted twice, precisely the divisor of \(q_2q_3|_{q_1}\). Subtract a scalar multiple of \(q_2q_3\); the remainder is \(q_1R\) for a conic \(R\). At the two contacts of \(q_2,q_3\), the residual conditions force \(R\) into their pencil. This gives (III.15). The three displayed forms are independent. Thus the twelve point-and-tangent conditions on the fifteen-dimensional quartic space are independent.

For a general \[
G=\alpha q_1q_2+\beta q_1q_3+\gamma q_2q_3,
\] the quartic is smooth and its curvature at each contact differs from those of the two conics there. Bertini applies away from the finite base scheme; nonvanishing of the relevant linear terms and distinct second-order terms are open conditions at the six base points. The branch curve \[
B=q_1q_2q_3G=0
\] has degree ten and exactly the six intended \([3,3]\) clusters under these generality assumptions.

Canonical plane numerators would be conics through all six first points. None exists, so \(p_g=0\). The smooth degree-ten double plane has \(\chi=7\); each \([3,3]\) cluster lowers it by one, giving \(\chi=1\), hence \(q=0\). Bicanonical invariant numerators are the quartics in (III.15), with the six \((1,1)\) clusters. The other summand in (III.4) has negative plane degree \(n-6=-1\) and no sections; after resolution it cannot acquire sections, since pushing down to the plane would give an effective divisor of negative degree. Therefore \[
P_2=3.
\tag{III.16}
\]

The source also interprets a virtual canonical double conic as having eight branch points, hence genus \(3\) by Riemann--Hurwitz. This is a virtual computation: no effective conic through the six contacts exists. The following intersection check proves the claimed linear genus without requiring such a conic. On the blown-up plane the quartic system has square \(16-12=4\), and its pullback has square eight. Thus \[
(2K_{\min})^2=8,\qquad K_{\min}^2=2,\qquad p^{(1)}=3.
\] The canonical-resolution calculation provides the same result: start with \(K_X^2=8\), lose two at the second blowup of each of the six clusters, reaching \(-4\), then contract the six disjoint rational \((-1)\)-curves above the first exceptional branch curves, reaching \(2\). Omitting these six contractions would give the wrong answer.

Enriques also proves completeness directly by restricting to the quartic branch component \(G\), of genus three. The residual series has degree four. If it were canonical, the twelve imposed cluster points would lie on a cubic. Such a cubic would meet each \(q_i\) at four prescribed tangencies, at least eight intersections, exceeding its Bézout number six unless it contained \(q_i\). Containing all three conics is impossible for a cubic. The degree-four series is therefore nonspecial and has two sections. There is only one bicanonical divisor containing the ramification curve over \(G\), namely its double. Restriction bounds the total space by \(1+2=3\), matching (III.15). This is a complete argument once the stated intersection interpretation of the cluster is made explicit.

The explicit base ideal also determines the bicanonical map. The ideal of the quartic net at each contact is locally \((u,v^2)\), so the two successive base-point blowups resolve it. It has no other base points and defines a degree-four map of the blown-up plane to \(\mathbb{P}^2\). After lifting to the double cover and contracting the six extraneous \((-1)\)-curves, the complete bicanonical map has degree eight onto \(\mathbb{P}^2\). This follows directly from the explicit base ideal and \((2K_X)^2=8\); it is not needed for Enriques's calculation of \(P_2\).

\section[Criteri d'equivalenza]{Criteri d'equivalenza (pp.~96--100)}\label{commentary-III-13}

\subsection{From adjoints to equivalent restrictions}

Enriques cites Severi's Il teorema d'Abel sulle superficie algebriche (1906, announced in 1905), Osservazioni varie di geometria sopra una superficie algebrica (1906), and Alcune relazioni di equivalenza (1911, abbreviated title in the source) as earlier approaches independent of adjoint surfaces. These citations are historical dependencies, not independently verified proofs of the reconstruction below. The opening observation says that if \(L,M\) cut equivalent divisors on every general member of a sufficiently ample, nonruled system without proper fundamental curves, they are equivalent on the surface. To reduce to the special adjoint form, Enriques chooses a sufficiently high-degree adjoint containing \(L\), with residual curve \(D\), and compares \(L+D\) with \(M+D\). The existence of that adjoint is supplied in modern terms by Serre vanishing for the ideal of \(L\). Proper fundamental divisors remain a possible kernel of restriction; excluding them is necessary.

This is a statement about equivalent line bundles on curves, not equality of numerical degrees. Equal degrees alone are much weaker. Nor does the argument imply that restriction to one arbitrary ample curve is injective on all of \(\operatorname{Pic}(X)\): on \(\mathbb{P}^1\times\mathbb{P}^1\), for example, \(\mathcal O(1,-1)\) restricts trivially to a smooth \((1,1)\)-curve.

\subsection{First criterion: a pencil with a simple base point}

The source assumes a linear pencil of integral curves, no reducible members, and a simple base point \(O\) avoided by \(K_0,K_\infty\). If the restrictions of \(K_0,K_\infty\) are equivalent, a rational function on each member can be normalized to take value one at \(O\). Enriques correctly recognizes that the zero and pole divisors of the resulting function might acquire extra pencil members, and tries to exclude them by the normalization.

A rigorous modern proof can avoid delicate pointwise assertions about a rational function at an indeterminacy point. Resolve the pencil and work over the function field of its parameter line. The line bundle \(M=\mathcal O_X(K_0-K_\infty)\) is trivial on its generic fiber, so choose a rational section trivializing it there. Its divisor is vertical. Push down from the resolved pencil to \(X\): exceptional components disappear, and each remaining vertical prime is a pencil member, because all those members are integral. Consequently \[
K_0-K_\infty\sim sC.
\] The assumed equivalence on \(C\) gives \((K_0-K_\infty)C=0\). A simple base point makes \(C^2>0\), so \(s=0\). This proves the criterion. It also shows why the base-point hypothesis matters and why restrictions must count any intersections at base points rather than silently discarding them.

The transition from fiberwise triviality to generic-fiber triviality uses algebraicity. On a proper flat family of integral curves, a degree-zero line bundle having a section on general fibers has a section on the generic fiber by semicontinuity and base change; a nonzero section of degree zero on an integral proper curve trivializes it. Merely choosing unrelated analytic functions on each complex fiber would not establish rational dependence.

\subsection{Second criterion: equivalence up to vertical divisors}

For a proper pencil \(f:X\to B\) with integral fibers, the same generic-fiber argument yields \[
K_0+\sum a_i C_i\sim K_\infty+\sum b_j C_j,
\tag{III.17}
\] where the sums are effective vertical divisors after separating positive and negative coefficients. This is the criterion on p.~97. If \(B\simeq\mathbb{P}^1\), all fibers are equivalent, so the difference is \(sC\). If the pencil has base points on the original model and those intersections are included, the degree calculation above forces \(s=0\).

The integral-fiber qualification is substantive. With reducible fibers, individual components occur in the vertical divisor; with multiple fibers, primitive components and their multiplicities must be retained. The loose phrase ``a sum of members'' should not erase that information.

\subsection{Trace normalization and extra poles}

The printed proof sets \[
f=\frac{\varphi}{\varphi_1+\cdots+\varphi_m},
\] using the intersection with an auxiliary curve \(L\). The denominator indeed removes the scalar ambiguity when its trace is nonzero. It can, however, vanish on fibers that meet neither \(L\cap K_0\) nor \(L\cap K_\infty\). Replacing \(f\) by \(1/f\) does not replace the trace of \(\varphi\) by the trace of \(1/\varphi\). Thus the asserted identification of the extra pole group on p.~98 is false.

An explicit counterexample is as follows. On \(X=\mathbb{P}^1_x\times\mathbb{P}^1_t\), take fibers \(C_t\), \(K_0=(x=0)\), \(K_\infty=(x=\infty)\), and \(\varphi=x\). Let \[
L:\quad x^2-tx+1=0.
\] This is an irreducible curve, the graph \(t=x+x^{-1}\), with two points over a general \(t\). Their \(\varphi\)-trace is \(t\). Therefore \(f=x/t\), and \[
\operatorname{div}(f)=K_0+C_\infty-K_\infty-C_0.
\tag{III.18}
\] Both \(L\cap K_0\) and \(L\cap K_\infty\) lie over \(t=\infty\), but the additional pole fiber is \(t=0\). It contains neither intersection. This directly contradicts the printed description of \(\sum' C\). It does not contradict (III.17), since the actual correction is still vertical.

The trace need not even be nonzero: with \(L:x^2=t\), the two values are opposite and the trace vanishes identically. A generic auxiliary choice can often avoid this, but the proof must say why and must still handle the resulting trace-zero fibers. The generic-fiber divisor proof supplies both requirements without attempting the false identification.

\subsection{Third criterion: a complete family of index greater than one}

The index \(\nu\) is the number of members through a general point. Enriques states that if a one-dimensional family has index \(\nu>1\), all members are irreducible, and \(K_0,K_\infty\) restrict equivalently to them, then \(K_0\sim K_\infty\). He forms a second trace \[
\psi(P)=f_1(P)+\cdots+f_\nu(P)
\] over the curves through \(P\). He acknowledges possible cancellation of poles, but assumes that all poles are on \(K_\infty\). The first normalization has not proved that assumption: its denominators create other poles. Nor is \(\psi\not\equiv0\) established. The final nonnegative-degree calculation is valid only after these missing premises and the absence of curves orthogonal to the family have been supplied.

Completeness of the family is also essential. Let \(X=\operatorname{Bl}_p\mathbb{P}^2\), with exceptional curve \(E\). Choose a smooth conic in the dual plane, parametrizing a one-dimensional family of lines of index two, and remove the finitely many parameters whose lines pass through \(p\). Their strict transforms are all integral and disjoint from \(E\). Thus \(E\) restricts trivially to every member but is not linearly trivial. To make both compared divisors nonzero and effective, take \(K_0=D+E,K_\infty=D\). Restoring the omitted parameters in the flat closure introduces reducible members containing \(E\). This example shows precisely what an open-family reading of the printed statement would miss.

\subsection{A modern proof with the hypotheses made explicit}

The following modern proof makes the sufficient hypotheses explicit and replaces the incomplete trace argument. Let \(B\) be a smooth complete connected curve parametrizing a flat family of pairwise distinct integral curves \(C_t\subset X\). Let \(I\subset X\times B\) be its incidence surface, \(p:I\to B\), \(q:I\to X\), and assume \(q\) is generically finite of degree \(\nu>1\). Let \(M\) be a line bundle whose restriction to every fiber \(C_t\) is trivial. The same conclusion follows from generic triviality if its extension to all fibers is separately justified; including it in the hypothesis avoids that hidden step.

\begin{enumerate}
\item
  Proper flat base change gives \(p_*\mathcal O_I=\mathcal O_B\) universally. The see-saw lemma gives \(q^*M=p^*N\) for a line bundle \(N\) on \(B\) (Stacks Project, Lemma 37.33.2, Tag 0EX7).
\item
  Resolve \(I\) if necessary. On the resolved incidence surface a divisor pulled back from \(B\) has square zero. The projection formula for the generically finite map therefore gives \[
  \nu M^2=(q^*M)^2=(p^*N)^2=0.
  \] Also \(M C=0\), since its restriction has degree zero.
\item
  Two distinct family members pass through a general point because \(\nu>1\). They have positive local intersection there, so their common numerical class satisfies \(C^2>0\).
\item
  The class \(C\) is ample. Indeed, its intersection with any irreducible curve \(E\) is nonnegative, by taking a different member. If \(CE=0\), completeness makes the proper incidence projection \(q\) surjective onto \(X\). Some member must meet a general point of \(E\); positive local intersection then forces that member to equal \(E\). But its class is \(C\), giving \(CE=C^2>0\), a contradiction. Thus \(CE>0\) for every curve, and the Nakai--Moishezon criterion applies.
\item
  Hodge index, together with \(MC=0,M^2=0,C^2>0\), makes \(M\) numerically trivial. This is numerical, not yet linear, equivalence.
\item
  Since \(C-M\) is ample, Kodaira vanishing and Serre duality give \[
  H^1(X,M-C)^\vee=H^1(X,K_X+C-M)=0.
  \] In the restriction sequence \[
  0\to M(-C)\to M\to\mathcal O_C\to0,
  \] the constant section of \(\mathcal O_C\) therefore lifts to a nonzero global section of \(M\). Its divisor is effective and numerically trivial. An ample divisor has positive intersection with every nonzero effective divisor, so this divisor is zero. Hence \(M\simeq\mathcal O_X\).
\end{enumerate}

This proves the claimed linear equivalence under the displayed complete-family hypotheses. The modern inputs are the Stacks Project (n.d., Tag 0EX7); Brosnan (n.d., Theorems 2.1 and 3.15) for Nakai--Moishezon and Hodge index; and Lazarsfeld (1994, Vanishing Theorem 5.2) for Kodaira vanishing. These ingredients are clearly separated from the incorrect trace pole calculation.

\subsection{The continuous-family observations on p.~100}

A rationally parametrized family of divisors has constant line-bundle class. One modern proof maps the parameter \(\mathbb{P}^1\) to the appropriate component of the Picard variety. Translating to \(\operatorname{Pic}^0(X)\), an abelian variety, makes this map constant: invariant one-forms span its cotangent spaces, whereas \(\mathbb{P}^1\) has no nonzero holomorphic one-forms. Thus the divisors lie in one complete linear system. Enriques attributes the criterion to his 1895 work and refers to Enriques--Chisini for the curve analogue.

The next criterion, attributed to Severi, concerns a connected family \(D_t\) of divisors whose restrictions to one sufficiently general test curve are equivalent. Its modern justification is simpler once the test curve is ample. The differences \(M_t=\mathcal O_X(D_t-D_0)\) lie in \(\operatorname{Pic}^0(X)\), hence are numerically trivial. If \(M_t|_C\) is trivial, the Kodaira-and-restriction argument in step 6 above proves \(M_t\) trivial on \(X\). A complete integral family of test curves with index greater than one supplies ampleness by steps 3--4. This proves the continuous-family version with explicit hypotheses.

The source instead invokes Castelnuovo's enumerative criterion for a series of point groups and a continuity argument that a maximum double-point count persists. It does not write the formula or the conditions preventing specializations of the count. Accordingly it is a sketch relying on the cited curve theory. Baker (1936) revisits a related Severi lemma for point series. Its abstract corroborates that historical discussion; it does not verify a repair of this particular surface trace proof.

The final paragraph attributes to Torelli an inequality for the number of nodal curves in a one-dimensional series, involving the Zeuthen--Segre invariant, with maximum in the linear case. No inequality, normalization of multiplicities, or equality hypotheses are given in these pages. The bibliographic entry can be located as Torelli (1907), volume 42, pp.~86--99, but its original full text was not inspected. It remains an assumed earlier result. No modern formula is invented to stand in for the unspecified historical inequality.

\chapter[Il genere numerico]{Il genere numerico e il teorema di Riemann--Roch per le superficie}\label{enriques-iv}

This chapter constructs the numerical genus from adjoint postulation, proves the surface Riemann--Roch inequality through restriction and subtraction, and then studies regularity, connectedness, and pluricanonical systems. Its twenty-one sections occupy pp.~101--165; p.~166 is blank. The central distinction is between a numerical identity and the geometric hypotheses that make a restriction complete or a system regular. With the notation fixed above,

\[
\begin{aligned}
r&=h^0(X,\mathcal O_X(D))-1\\
 &=p_a(X)+D^2-p_a(D)+1\\
 &\qquad-h^0(X,K_X-D)+h^1(X,\mathcal O_X(D)).
\end{aligned}
\]

Thus the speciality index is \(h^2(D)=h^0(K_X-D)\), while superabundance is \(h^1(D)\). The numerical inequality omits only the nonnegative superabundance term. The empty complete system has dimension \(-1\). Assigned proper and infinitely near base points are incorporated on their resolved smooth model before squares and genera are calculated. The source's ``normal singularities'' remain ordinary projection singularities; its projective ``normality'' means completeness of the hyperplane series. Its ``more ample'' order relation is not automatically positivity in the modern sense.

Enriques supplies complete dimension calculations alongside sketches and imported results. The commentary separates those statuses from the truth of the modern statements. In particular, his explicit working hypothesis for nonreduced connected curves at p.~155 is repaired by the constant-functions theorem for numerically 1-connected divisors (Ramanujam, 1972, Lemma 3; Reid, 1996, Lemma 3.11). The subsequent bicanonical discussion requires the pure, minimal-model interpretation of the canonical system. None of these repairs is retrospectively attributed to the original argument.

\section[Introduzione]{Introduzione (pp.~101--102)}\label{enriques-iv-1}

Enriques starts from results already established in Enriques--Chisini, Book V, volume III: complete linear series on a curve, the birational invariance of the canonical series, completeness of plane adjoints, and curve Riemann--Roch. The opening section imports these foundational theorems rather than proving them. It explains why their method suggests a surface analogue.

Let \(C_0\subset\mathbb P^2\) be an integral plane curve of degree \(m\) with only ordinary nodes, let \(\nu:C\to C_0\) be its normalization, and let \(g=g(C)\). If \(Z\) is the node scheme, its length is

\[
 d=\frac{(m-1)(m-2)}2-g.
\]

Each node has normalization delta-invariant one. With \(H=\nu^*\mathcal O_{C_0}(1)\) and \(G\) the sum of the two branches over every node, adjunction gives

\[
 K_C\sim(m-3)H-G.
\]

Thus degree-\(m-3+s\) plane adjoints, after the conductor contribution is removed, cut \(|K_C+sH|\). ``Complete'' concerns the resulting series on the normalization; it does not mean that arbitrary ambient forms have no kernel. For large \(s\), curve Riemann--Roch and \(H^1(K_C+sH)=0\) give

\[
 \deg(K_C+sH)=2g-2+sm,
 \qquad \dim|K_C+sH|=g-2+sm.
\]

The argument in the introductory paragraph then imposes independent points on a nonspecial complete series. Independence is essential: a special configuration can impose fewer conditions, and that is precisely what the later surface argument is designed to measure. At \(s=0\), the complete series is canonical, of dimension \(g-1\). The nodes impose independent conditions on forms of degree \(m-3\). When \(g=0\), this says the adjoint space is zero-dimensional as a vector space and the projective system is empty; it is not an exception to the formula.

For any divisor \(D\) of degree \(n\), the ``index of speciality'' is the \textbf{vector} dimension

\[
 i(D)=h^0(K_C-D)=h^1(D),
\]

and the precise theorem is

\[
 \dim|D|=n-g+i(D).
\]

This resolves the otherwise easy confusion between ``number of independent canonical groups'' and projective dimension. Vakil's (2002, p.~1) \emph{Complex algebraic surfaces, Class 5} gives a concise modern statement of curve Riemann--Roch, Serre duality, and adjunction on its first page. No additional mathematical gap occurs in §1 once the cited curve theory and the ordinary-node model are acknowledged. The pronoun errors \texttt{contenerle} on pp.~101--102 are printed grammatical slips; the punctuation after the node formula is best read as a semicolon.

\section[Formule di postulazione]{Condizioni imposte ad una superficie dal passaggio per una curva: formule di postulazione (pp.~102--106)}\label{enriques-iv-2}

\subsection{The invariant being counted}

For \(C\subset\mathbb P^3\), the number of independent conditions imposed by containing \(C\) is

\[
 \operatorname{post}_C(N)
 =\operatorname{rank}\bigl[H^0(\mathbb P^3,\mathcal O(N))
 \longrightarrow H^0(C,\mathcal O_C(N))\bigr].
\]

For a smooth degree-\(m\), genus-\(g\) curve, let \(e_N\) be the cokernel dimension of this restriction. Riemann--Roch gives the exact diagnostic identity

\[
 \operatorname{post}_C(N)=mN+1-g+h^1(C,\mathcal O_C(N))-e_N.
\]

Speciality and incompleteness have opposite signs. Enriques correctly distinguishes them. Once \(mN>2g-2\), speciality vanishes and the postulation is at most \(mN+1-g\). Equivalently,

\[
 h^0(\mathcal I_C(N))\geq\binom{N+3}{3}-mN+g-1.
\]

The sentence on p.~103 calling \(C\) a degree-\(n\) curve is printed that way; the curve degree should be \(m\), while \(n\) there denotes the degree of the containing surfaces.

\subsection{\texorpdfstring{Why \(N\geq m-2\) suffices}{Why N\textbackslash geq m-2 suffices}}

The argument is attributed to Castelnuovo's 1893 paper \emph{Sui multipli di una serie lineare di gruppi di punti appartenente ad una curva algebrica}, reprinted at the cited place in \emph{Memorie scelte}. Its structure is a genuine proof, not merely an invocation of ``general position.'' Work over characteristic zero and take a generic projection of the smooth nondegenerate spatial curve to a plane. The projection has

\[
 d=\binom{m-1}{2}-g
\]

ordinary nodes, whose inverse images give a divisor \(G\) of degree \(2d\) on \(C\). Cones over plane adjoints of degree \(N\geq m-2\) give the complete residual series

\[
 |NH-G|=|K_C+(N-m+3)H|.
\]

This is nonspecial because \(N-m+3\geq1\). It remains to show that the \(2d\) points of \(G\) impose independent conditions on spatial forms of degree \(N\). For \(N=m-2\), fix one selected branch point \(P\). By the independence of the \(d\) plane nodes for adjoints of degree \(m-3\), choose a cone of degree \(m-3\) containing the \(d-1\) other secant lines but not the secant containing \(P\). Multiply its equation by that of a plane passing through the other endpoint of this secant but avoiding \(P\). The product vanishes at every point of \(G\) except \(P\). Repeating for each point gives a basis of separators. For larger \(N\), multiply by a form avoiding all of \(G\).

The restriction series therefore contains the complete residual series and acquires exactly \(2d\) independent values along \(G\). It is all of \(H^0(C,NH)\). Consequently

\[
 \operatorname{post}_C(N)=mN-g+1\qquad(N\geq m-2).
\]

This proof assumes the generic nodal-projection theorem and plane-adjoint completeness, already part of the cited curve theory. Those are prerequisites, not newly discovered omissions. Throughout the source argument the surface degree is \(n\), so the intended symbol after ``offre alle'' on p.~103 is \(\Phi_n\).\footnote{The edition's diplomatic text reads \(\Phi_h\), which might be taken for a transcription error. Direct enlargement of p.~103, however, shows the tall left stroke on which the reading \(h\) rests. The small subscript is optically less secure than its mathematical correction to \(n\). The reading \(\Phi_h\) should therefore not be classified categorically as a transcription error: the source most plausibly contains a harmless subscript slip, while the intended \(\Phi_n\) is fixed by the argument.} Repeated equality signs across a printed line break on p.~104 are a typographic continuation convention, not an algebraic equation with an extra equality operator.

\subsection{Nodes, noncoplanar triple points, and component genera}

For singular curves the geometric genus of the normalization is insufficient by itself. Enriques's coefficient \(2t\) is correct because his triple points have three \textbf{noncoplanar} branches. Their completed local ring is modeled by

\[
 k[[x,y,z]]/(xy,xz,yz).
\]

Its normalization is the direct sum of three branch rings. The local functions are triples with one common constant, so the normalization quotient has length two. An ordinary plane triple point, by contrast, has delta-invariant three. Substituting the plane value here would falsely indict the printed formula.

A containing surface must vanish to order at least two at such a triple point: a nonzero linear term cannot vanish on three independent tangent directions. This entails four ambient jet conditions. On the three normalization branches, the forced double-point vanishing removes six intersection conditions. Hence one triple point changes the smooth residual count by \(4-6=-2\). At an ordinary node, the corresponding change is \(1-2=-1\). For an irreducible curve of degree \(m\), normalization genus \(g\), \(t\) such triple points, and \(\tau\) nodes, the eventual result is

\[
 \operatorname{post}_C(N)=mN-2t-\tau-g+1.
\]

Enriques first explains this through residual intersections, then treats several components. The clean modern proof packages all cases into

\[
 0\longrightarrow\mathcal O_C\longrightarrow
 \nu_*\mathcal O_{\widetilde C}\longrightarrow Q\longrightarrow0,
 \qquad\ell(Q)=2t+\tau.
\]

If \(\widetilde C=\coprod_{i=1}^h C_i\), define

\[
 \rho=\sum_{i=1}^h g(C_i)-h+1.
\]

Then \(\chi(\mathcal O_{\widetilde C}(N))=mN+1-\rho\), and

\[
 \chi(\mathcal O_C(N))=mN+1-\rho-2t-\tau.
\]

Serre vanishing makes this the actual restriction rank for \(N\gg0\). Thus the final formula includes all configurations where the stated singularities are ordinary and no other singularities intervene:

\[
 p_a(C)=\rho+2t+\tau,
 \qquad\operatorname{post}_C(N)=mN-2t-\tau-\rho+1.
\]

The integer \(\rho\) is the arithmetic genus of the \textbf{disjoint normalization}, not the sum of its component genera and not the arithmetic genus of the singular curve. It may be negative. This point is what makes the component calculation on pp.~105--106 coherent. The source's repeated minus sign in the broken display \(x=mn-6t-2\tau\) is again a continuation convention.

\section[Dimensione dei sistemi aggiunti]{Dimensione dei sistemi di superficie aggiunte ad una data (pp.~107--110)}\label{enriques-iv-3}

Let \(Y\subset\mathbb P^3\) be a degree-\(n\) projected surface with ordinary singularities, and let its double curve \(\Gamma\) have degree \(d\), normalization-disjoint genus \(\rho\), \(t\) noncoplanar triple points, and \(\tau\) ordinary nodes. The eventual vector dimension of degree-\(m\) adjoint forms is

\[
 N_m=\binom{m+3}{3}-dm+2t+\tau+\rho-1.
\]

Its finite difference is

\[
 N_m-N_{m-1}=\binom{m+2}{2}-d.
\]

Here ``normal singularities'' is the classical term for the prescribed ordinary projection singularities; it does not mean that \(Y\) is normal in the modern scheme-theoretic sense. A surface with a double curve is nonnormal.

The displayed index \(m-1\) on p.~107 is the correct one.\footnote{The prose does not also print \(m-1\), and the edition's \(w-1\) should not be classified as an error confined to the transcription. Direct enlargement instead supports the diplomatic prose reading \(N_m-N_{w-1}\); the following displayed equation plainly reads \(N_m-N_{m-1}\). The prose \(w\) is therefore a source typographical slip. This reading preserves the mathematical finite difference and places the discrepancy in the print rather than in the transcription.} The right side is the vector dimension of the plane adjoints to a general plane section, whenever the \(d\) nodes impose independent conditions. Their \textbf{projective} dimension is \(m(m+3)/2-d\), one less. Enriques keeps those two counts consistent.

The argument descends from high degrees: the kernel of restriction to a plane consists of adjoints divisible by that plane's equation. If the restriction image loses \(\delta\) dimensions, the previous-degree adjoint space gains \(\delta\) relative to the virtual polynomial. Independence of planar nodes in degree at least \(n-3\) lets this continue down to ambient adjoint degree \(n-4\). It gives actual count at least virtual count, and therefore \(p_g\geq p_a\). Below that degree the planar independence guarantee ends. The sentence saying the degree-\(n-4\) plane system \emph{is} superabundant is too categorical; it may be superabundant, and that weaker statement is all the proof needs. A smooth plane curve has no node conditions at all, so cannot exhibit their superabundance.

\subsection{\texorpdfstring{Why the virtual value is \(\chi(\mathcal O_X)-1\), not a convention without content}{Why the virtual value is \textbackslash chi(\textbackslash mathcal O\_X)-1, not a convention without content}}

Let \(\nu:X\to Y\) be the smooth normalization of the ordinary projected model, and write \(H=\nu^*\mathcal O_Y(1)\). Finite duality identifies the conductor ideal on \(Y\) with

\[
 \mathcal{H}om_Y(\nu_*\mathcal O_X,\mathcal O_Y)
 \simeq\nu_*\omega_X\otimes\omega_Y^{-1}.
\]

Since \(\omega_Y\simeq\mathcal O_Y(n-4)\), its preimage in \(\mathcal O_{\mathbb P^3}\) gives the exact sequence

\[
 0\longrightarrow\mathcal O_{\mathbb P^3}(m-n)
 \longrightarrow\mathcal I_\Gamma(m)
 \longrightarrow\nu_*\omega_X(m-n+4)
 \longrightarrow0. \tag{IV.3.1}
\]

For ordinary double curves, including the standard triple/pinch-point local models, this is the classical adjoint ideal. For an arbitrary singular model one must use the actual canonical/resolution adjoint ideal, not assume that vanishing on the reduced singular set alone suffices.

Taking Euler characteristics in (IV.3.1) at \(m=n-4\) gives

\[
 N_{n-4}=\chi(\omega_X)+\chi(\mathcal O_{\mathbb P^3}(-4))
 =\chi(\mathcal O_X)-1=p_a(X).
\]

The contribution \(\chi(\mathcal O_{\mathbb P^3}(-4))=-1\) is essential. The actual canonical forms have dimension \(p_g\), and their excess over \(N_{n-4}\) is precisely \(q\). This modern derivation verifies the intended invariant, while also explaining why the arithmetic genus of the singular hypersurface \(Y\), computed directly from its degree, is generally \textbf{not} Enriques's birational numerical genus.

Substituting into the postulation polynomial gives the printed expression

\[
 p_a=\binom{n-1}{3}-(n-4)d+2t+\tau+\rho-1.
\]

For \(m=n-4+s\), the adjoint divisor is \(K_X+sH\). Let \(\pi=p_a(H)\) and \(\pi_s=p_a(sH)\). Adjunction gives

\[
 \pi_s=1+s(\pi-1)+\frac{s(s-1)}2n,
 \qquad\pi=\binom{n-1}{2}-d.
\]

The dimension of the induced system is obtained by subtracting the scalar dimension and the kernel of restriction to \(Y\): for \(s\gg0\),

\[
 \dim|K_X+sH|=N_{n-4+s}-1-h^0(\mathbb P^3,\mathcal O(s-4))
 =p_a+\pi_s-1.
\]

For \(s=1,2,3\), the kernel is zero. One must not extend the polynomial \(\binom{s-1}{3}\) to \(s=0\) and treat its value \(-1\) as a vector-space dimension. At that endpoint, (IV.3.1) and the actual canonical space must be used. On p.~109 the middle binomial identity has a printed minus in front of \(s(s-1)n/2\); it must be plus. The preceding genus formula, the succeeding binomial expression, and direct expansion all agree on the correction.

The conductor calculation also indicates a further comparison: for explicitly chosen nonordinary singularities, one can track the extra conductor or local geometric-genus contributions to the postulation polynomial. The formulas proved here apply to the stated ordinary class; such an extension is a proposed calculation, not an established generalization.

\section[Il genere numerico]{Il genere numerico (pp.~110--114)}\label{enriques-iv-4}

\subsection{A single exact sequence organizes the whole section}

For a smooth connected effective divisor \(C\subset X\), adjunction gives

\[
 0\longrightarrow\omega_X\longrightarrow\omega_X(C)
 \longrightarrow\omega_C\longrightarrow0. \tag{IV.4.1}
\]

Let \(\delta(C)\) denote the cokernel dimension of the induced map

\[
H^0(\omega_X(C))\to H^0(\omega_C).
\]

Then

\[
 \dim|K_X+C|=p_g+g(C)-1-\delta(C). \tag{IV.4.2}
\]

This includes empty systems without changing conventions. Over a fixed nonzero projective canonical section, the adjoints restricting to it have an affine fibre of dimension \(p_g\), whose projective closure has the same dimension; the kernel itself is \(\mathbb P H^0(K_X)\), dimension \(p_g-1\). This clarifies the source's elementary counting language.

Because \(C\) is connected, \(H^1(K_C)\to H^2(K_X)\) is an isomorphism: its dual is \(H^0(\mathcal O_X)\to H^0(\mathcal O_C)\), the isomorphism on constants. Also \(H^2(K_X+C)=H^0(-C)^*=0\). Therefore the long exact sequence yields

\[
 \delta(C)=q-h^1(K_X+C),\qquad 0\leq\delta(C)\leq q, \tag{IV.4.3}
\]

and hence

\[
 \dim|K_X+C|=p_a+g(C)-1+h^1(K_X+C). \tag{IV.4.4}
\]

Equivalently, by dualizing the relevant restriction sequence,

\[
 \delta(C)=\operatorname{rank}\bigl[H^1(X,\mathcal O_X)
 \longrightarrow H^1(C,\mathcal O_C)\bigr].
\]

The deficiency is thus a rank, not an unexplained geometric defect. Completeness is exactly \(\delta(C)=0\).

\subsection{The comparison argument is valid under its stated hypotheses}

Enriques compares \(C\) with \(C+L\). He assumes moving irreducible systems and an irreducible sum; two pencils must be distinct, so that the relevant intersection \(i=C\cdot L\) is positive. For a smooth \(L\), the restriction

\[
 (K_X+C+L)|_L=K_L+C|_L
\]

has degree \(2g(L)-2+i\), hence is nonspecial and has projective dimension \(g(L)+i-2\). The genus addition formula is

\[
 p_a(C+L)=p_a(C)+p_a(L)+C\cdot L-1.
\]

Subtracting the two instances of (IV.4.2), or using the long exact sequence of restriction to \(L\), gives

\[
 \delta(C)\leq\delta(C+L).
\]

In fact the codimension of the restriction of \(|K_X+C+L|\) inside the complete series on \(L\) is \(\delta(C+L)-\delta(C)\). Thus equality of the deficiencies is equivalent to completeness of this restriction. These are proved claims, not unjustified monotonicity assumptions. For a reducible added divisor, the source carefully asks that its components not be fundamental for the original system and that successive additions retain irreducibility; that positivity is needed for the nonspeciality used in the calculation.

The discussion of base points is most cleanly expressed on a common blowup. If \(\sigma:\widetilde X\to X\) blows up one point of multiplicity \(i\) on \(C\), then

\[
 K_{\widetilde X}=\sigma^*K_X+E,
 \qquad\widetilde C=\sigma^*C-iE,
 \qquad K_{\widetilde X}+\widetilde C
 =\sigma^*(K_X+C)-(i-1)E.
\]

The corresponding adjoint condition has length \(i(i-1)/2\), the number of monomials in two variables of total degree below \(i-1\). For a sufficiently positive system the jets are independent. The source's infinitesimal exceptional curve is the geometric form of this calculation.

For \(L_s=sH\), the large-degree formula of §3 and (IV.4.2) imply \(\delta(L_s)=p_g-p_a=q\). Given a fixed moving \(C\), the difference \(sH-C\) is very ample for \(s\gg0\); comparison then yields \(\delta(C)\leq q\). The p.~112 opening sentence inadvertently calls the adjoint system itself \(|L_s|\), whereas the following formulas require \(L_s=sH\) and \(L'_s=K_X+sH\). This is an antecedent/notation slip, not a change of genus convention.

The maximum-deficiency characterization implies birational invariance of \(p_a\), once canonical forms and moving systems are transferred to smooth models. Modernly, blowing up a smooth point preserves \(H^i(\mathcal O_X)\); birational maps of smooth surfaces factor into such modifications. Thus \(p_g,q,p_a\) are birational invariants. The source defines regularity as \(p_a=p_g\), i.e.~\(q=0\).

There is no contradiction between these preliminary inequalities and the later stronger regularity criterion. In characteristic zero, if \(|C|\) is a moving irreducible system of dimension at least two, resolve its base points. The resulting free class is nef and big: its morphism has two-dimensional image, since a free system composed with a curve would have a general member that is a sum of fibres unless it were a pencil. Ramanujam's vanishing theorem then gives \(H^1(K_X+C)=0\), hence \(\delta(C)=q\). The nef-and-big proof does not address square-zero pencils; moving linear pencils require a separate fibration argument. Isolated or nonmoving divisors can have \(\delta(C)<q\). No strict-deficiency example among complete moving irreducible linear pencils is being claimed. Enriques's earlier estimate is weaker; it is not false merely because later machinery gives equality under stronger hypotheses.

\subsection{The numerical formulas at the end of §4}

Let \(n=C^2\), \(\pi=p_a(C)\), \(C'=K_X+C\), and \(p^{(1)}=K_X^2+1\) on the chosen smooth model. Then \(K_X C=2\pi-2-n\), so

\[
 \begin{aligned}
 \pi'&=p_a(C')=3\pi-3-n+p^{(1)},\\
 n'&=(C')^2=4\pi-5-n+p^{(1)},\\
 \pi-1&=n'-\pi'+1.
 \end{aligned}
\]

The apparent prime after the \textbf{last} \(n\) in the second formula on p.~113 cannot be retained mathematically. The scan has a small high stroke; it is best treated as a printed misprint with moderate visual confidence. The required \(n\), without prime, follows unambiguously from the intersection calculation. The parameter \(p^{(1)}\) is called ``relative'' for good reason: \(K_X^2\) decreases by one under blowup. It is not a birational invariant unless an appropriate minimal-model convention is introduced.

Equation (IV.4.4) gives

\[
 \dim|C'|\geq p_a+\pi-1=p_a+n'-\pi'+1.
\]

Relabelling an adjoint divisor yields the numerical Riemann--Roch inequality printed on p.~114. In full modern form,

\[
 h^0(D)-1=p_a+\frac{D(D-K_X)}2+h^1(D)-h^2(D).
\]

For an adjoint to a nonzero effective divisor, \(h^2(D)=0\); its excess over the numerical value is \(h^1(D)\geq0\). This is the modern content of the regular/superabundant distinction made here. The present section does not yet state the full theorem for all divisors; that is developed later in the chapter.

\section[Complementi]{Complementi (pp.~114--117)}\label{enriques-iv-5}

\subsection{A direct proof of the multiple-curve formula}

Let \(C\subset\mathbb P^3\) be smooth of degree \(m\) and genus \(g\), with ideal \(\mathcal I\). The \(s\)-fold thickening used here is defined by \(\mathcal I^s\). Since \(C\) is a regular embedding of codimension two, \[
 \mathcal I^j/\mathcal I^{j+1}\simeq\operatorname{Sym}^j N_{C/\mathbb P^3}^{\vee}.
\] The normal bundle has rank two and degree \(4m+2g-2\), by adjunction. Its \(j\)th symmetric conormal power has rank \(j+1\) and degree \[
 -\frac{j(j+1)}2(4m+2g-2).
\] Riemann--Roch for vector bundles on a curve therefore gives \[
 \chi\!\left(\operatorname{Sym}^jN^\vee\otimes\mathcal O_C(N)\right)
 =(j+1)(Nm+1-g)-\frac{j(j+1)}2(4m+2g-2).
\] Summing from \(j=0\) to \(s-1\), using \[
 \sum(j+1)=\frac{s(s+1)}2,\qquad
 \sum j(j+1)=\frac{s(s-1)(s+1)}3,
\] produces exactly \[
 \operatorname{post}_{sC}(N)
 =\frac{s(s+1)}6
 \left[(3N-4s+4)m-(2s+1)(g-1)\right]
 \qquad(N\gg0).
\] Thus the displayed Noether--Campedelli formula is correct. The source credits Noether and cites Campedelli's 1931 article; it does not provide this derivation. The footnote saying that singular points of \(C\) require modification is substantive: the conormal filtration need not remain a rank-two vector-bundle calculation, and no correction depending only on degree and normalization genus is supplied. For an isolated assigned \(s\)-fold point in \(\mathbb P^3\), the independent conditions number \[
 \ell\bigl(k[x,y,z]/(x,y,z)^s\bigr)=\binom{s+2}{3}.
\]

\subsection{The indexing errors and the correct telescoping invariant}

Set \(H\) equal to a plane section. On p.~115, degree-\(n-3+s\) adjoints correspond to \[
 K_X+(s+1)H,
\] not to \(K_X+sH\). Since the same paragraph defines \(\pi_s=p_a(sH)\), its printed dimension \(p_a+\pi_s-1\) is an off-by-one error; it must be \(p_a+\pi_{s+1}-1\). The restriction degree \(2\pi-2+sn\) printed immediately afterward is correct, because the restriction to \(H\) is \(K_H+sH_H\).

The eventual finite difference in the number of adjoint forms is \[
 \binom{n+s-1}{2}-d
 =\pi+sn+\frac{s(s-3)}2.
\] This validates the ``arithmetic progression of third order'': it is an eventual cubic Hilbert polynomial whose first difference is this quadratic. Its extrapolated values are virtual counts. The footnote's extension to arbitrary fixed schemes of points and curves is the Hilbert-polynomial theorem; it does not require those base schemes to arise as the singularities of a surface.

Let \(\delta_s\) denote the failure of restriction \[
 H^0(K_X+(s+1)H)\longrightarrow H^0(K_H+sH_H)
\] to be surjective, for \(s\geq0\). The source fixes a threshold \(h\) beyond which all deficiencies vanish. The descent gives \[
 q=\delta_0+\delta_1+\cdots+\delta_{h-1}.
 \tag{IV.5.1}
\] For \(s\geq1\), the curve bundle \(K_H+sH_H\) is nonspecial. The long exact sequence gives \[
 \delta_s=h^1(K_X+sH)-h^1(K_X+(s+1)H).
\] At \(s=0\), the trace isomorphism used in §4 gives \[
 \delta_0=q-h^1(K_X+H).
\] These formulas telescope to (IV.5.1). They explain exactly why the canonical endpoint differs from an ordinary positive-twist step.

The first displayed descending dimension on p.~116 must use \(\pi_{h-1}\), not \(\pi_{h-2}\). The next must use \(\pi_{h-2}\), not \(\pi_{h-3}\), and must include the missing constant \(-1\). When a descent reaches \(K_X\) itself, one uses the separately printed formula \(p_g-1\), rather than forcing the positive-twist notation beyond its range. The final prose formula erroneously begins its sum at \(s=1\); the correct lower bound is \(s=0\). These are errors in the scan, not OCR.

Under the smooth finite-normalization hypothesis and in characteristic zero, \(H\) is ample. Kodaira vanishing gives \(H^1(K_X+sH)=0\) for every \(s\geq1\). Consequently all \(\delta_s\) with \(s\geq1\) vanish and \(\delta_0=q\). Enriques is working toward that stronger picture through classical geometric arguments. His provisional sequence of possible defects is a valid formal estimate, not evidence of actually existing higher defects under these positivity assumptions. If one changes to an arbitrary resolution with exceptional curves, positivity must be checked anew.

\subsection{The multiplication lemma and its special boundary case}

Let \(M,L\) be line bundles on a smooth curve. Suppose \(V\subset H^0(L)\) is a pencil without base points. The exact sequence \[
 0\longrightarrow L^{-1}\longrightarrow V\otimes\mathcal O_C
 \longrightarrow L\longrightarrow0
\] tensored with \(M\) shows that \[
 H^1(M-L)=0
 \quad\Longrightarrow\quad
 V\otimes H^0(M)\twoheadrightarrow H^0(M+L).
 \tag{IV.5.2}
\] This is the modern basepoint-free pencil trick. Enriques's argument with two disjoint divisors is the same proof in projective dimensions. If \(\deg M=m\), \(\deg L=n\), and \(M,M-L\) are nonspecial, the two copies of \(H^0(M)\) meet in \(H^0(M-L)\); their vector-space span has dimension \(m+n-g+1\). The source's ``minimum sum'' means this linear span of products, not the set of pairwise sums of divisors alone. A modern explicit reference is Voight and Zureick-Brown, \emph{The canonical ring of a stacky curve}, §2.4, Lemma 2.4.1.

For \(M=K_C+sL\) with \(s>1\), the residual \(M-L\) has positive twist and (IV.5.2) applies. At \(s=1\), \(M-L=K_C\) has \(h^1=1\). The two-pencil-section image is therefore one dimension short. Its projective dimension is \[
 2(g+n-1)-g-1=g+2n-3,
\] while the complete target \(|K_C+2L|\) has dimension \(g+2n-2\). The scan on p.~117 prints \(g+n-3\) and \(g+n-2\): each occurrence is missing a second \(n\) contribution.

For the plane-line series, a pencil corresponds to lines through a point \(P\) off the curve. The deficient plane adjoints are precisely those through \(P\). Two generic distinct centres give distinct evaluation hyperplanes, whose span is the entire adjoint space. This repairs the boundary case without pretending that \(H^1(K_C)\) vanishes. The argument does not yet prove the \(s=0\) assertion; §6 deals with that case.

\section[Condizioni di regolarità]{Nota sulle condizioni di regolarità di una superficie (pp.~117--120)}\label{enriques-iv-6}

For a smooth surface and a smooth ample hyperplane section \(C\), the decisive exact sequence is (IV.4.1). Kodaira vanishing gives \(H^1(K_X+C)=0\), and hence \[
 \operatorname{coker}\!\left[H^0(K_X+C)\to H^0(K_C)\right]
 \simeq H^1(K_X),\qquad
 \dim\operatorname{coker}=q.
 \tag{IV.6.1}
\] Therefore the adjoints restrict to the \textbf{complete canonical series if and only if \(q=0\)}. In that case \[
 \dim|K_X+C|=p_g+g(C)-1.
\] Kodaira's 1952 paper supplies the later vanishing machinery; this is a modern proof, not a claim that he was the first person to repair the precise 1897 argument reproduced here.

Enriques explicitly calls the ensuing classical demonstration a ``rapid sketch,'' and cites Castelnuovo's \emph{Alcune proprietà fondamentali dei sistemi lineari di curve tracciati sopra una superficie algebrica} (1897). Its first lemma is sound. A divisor \(G\) of degree \(m\) moving in a \(g_m^h\) imposes \(m-h\) canonical conditions, by Riemann--Roch. Under canonical completeness those are the same fixed linear conditions on the global adjoints. Therefore their rank is independent of the smooth plane section through the fixed axis. ``Every plane'' must be read within the smooth/generic-section framework or reformulated using dualizing sheaves at singular members.

The longer construction on pp.~118--119 seeks an adjoint of degree \(n-2\) by sweeping plane adjoints through a pencil of planes. In each plane it imposes: the canonical group cut by \(C'\), \(s-1\) selected axis intersections, and a tangent-plane condition at one selected point. The text asserts a unique curve after these conditions; the swept surface initially has degree \(n-2+h\) and contains the axis with multiplicity \(h\). A tangent-plane argument at one further intersection point is used to force \(h=0\) via the preceding special-series lemma. Varying the data is then said to produce \(s+1\) dimensions of adjoints and \(s\) dimensions of characteristic series.

This is an authentic outline of a geometric proof, but it does not supply the full incidence calculation: independence of the chosen point and tangency conditions, algebraicity of the swept family, control of its degree and fixed components, the limiting tangent-plane argument, and the final dimension variation are compressed. The correct diagnosis is \textbf{sketched, with these details omitted}, rather than ``false.'' Formula (IV.6.1) proves the desired criterion without those auxiliary choices.

Once \(q=0\), characteristic completeness follows directly from \[
 0\longrightarrow\mathcal O_X\longrightarrow\mathcal O_X(C)
 \longrightarrow\mathcal O_C(C)\longrightarrow0.
\] The cokernel of restriction injects into \(H^1(\mathcal O_X)=0\). This proof works for every effective Cartier divisor for which the characteristic-series interpretation is appropriate, and does not need Kodaira vanishing in this direction.

\subsection{The curve multiplication result used at the end}

The additional curve theorem on p.~120 says: if \(L\) is a complete basepoint-free linear series on a positive-genus smooth curve and its map is birational onto its image, then \[
 H^0(K_C)\otimes H^0(L)\twoheadrightarrow H^0(K_C+L).
 \tag{IV.6.2}
\] The phrase ``not composed with an involution'' is the birationality hypothesis. It must not be discarded. For a concrete failure, let \(A\) be the hyperelliptic pencil on a curve of genus \(g\geq2\) and put \(L=2A\). Then \(h^0(L)=3\), while the canonical products have dimension \(g+2\) and \(h^0(K_C+2A)=g+3\). The deficiency is one. This does not describe arbitrarily high multiples, which eventually become very ample.

The source gives a dimension argument using a product-series lemma from Castelnuovo (1893). If \(G\in|L|\) is general, of degree \(n\), and \(h^0(L)=s+1\), the two evaluated spaces have dimensions \(s\) and \(n-s\). The cited lemma gives at least \(n-1\) conditions on their product. Adding the canonical subspace gives \(n+g-2\), exactly the Riemann--Roch dimension of the target. The \(n-1\) lower bound is a substantive lemma about these series; it is not a valid dimension inequality for arbitrary coordinate subspaces.

A modern version of the same proof can use the usual secant lemma for a birational curve series: choose \(s-1\) general points \(D\) so that \(L(-D)\) is a complete basepoint-free pencil. This is the geometric ingredient in the classical multiplication theorem. For \(M=L(-D)\), the pencil trick tensored with \(K_C\) gives a surjective product map, because the dual of \[
 H^1(K_C-M)\longrightarrow H^0(M)\otimes H^1(K_C)
\] is the isomorphism from the dual pencil to \(H^0(M)\). Thus the products already contain \(H^0(K_C+L-D)\). Evaluation of \(H^0(L)\) at \(D\) is surjective. For each point of \(D\), choose a section of \(L\) vanishing at the other points but not there, and a canonical section nonzero there; the products fill every missing coordinate at \(D\). Since \(H^1(K_C+L-D)=H^0(D-L)^*=0\), this fills the entire target. This isolates precisely the secant/general-position input rather than hiding it in a dimension count.

Finally, the note extends the surface criterion to irreducible systems of dimension at least two that need not define birational maps. After resolving base points the moving class is free, nef, and big. Ramanujam's characteristic-zero vanishing gives \(H^1(K_X+C)=0\), so (IV.6.1) remains valid. This proves the extension in the stated dimension range. The source excludes pencils from this formulation, and the nef-and-big proof does not cover square-zero classes. Complete moving irreducible linear pencils can be treated separately by fibration methods in characteristic zero; they are not being offered as counterexamples. Arbitrary isolated effective divisors cannot replace the moving-system hypothesis.

\section[Superficie irregolari]{Esempi di superficie irregolari (pp.~120--123)}\label{enriques-iv-7}

\subsection{Smooth hypersurfaces, rational surfaces, and ruled surfaces}

For a smooth degree-\(n\) hypersurface in \(\mathbb P^3\), \[
 K_X=\mathcal O_X(n-4),\qquad
 p_g=p_a=\binom{n-1}{3},\qquad q=0.
\] The canonical form count follows from the hypersurface sequence and the cohomology of projective space. Smoothness suffices; the surface need not be general in its degree. Planes, smooth quadrics, and smooth cubics give zero. Birational invariance then yields \(p_g=p_a=0\) for every smooth rational surface. The p.~120 footnote formally assigns the empty double curve genus \(\rho=1\), so that its Euler characteristic is zero. That is a useful convention, not the genus of a connected empty curve. Its degree should be denoted \(d=0\), not the printed \(n=0\).

For a ruled surface with base \(B\) of genus \(g\), pass to a smooth \(\mathbb P^1\)-bundle \(f:X\to B\). Then \[
 f_*\mathcal O_X=\mathcal O_B,\qquad R^1f_*\mathcal O_X=0,
\] and Leray gives \(q=g\) and \(p_g=0\), hence \(p_a=-g\). The cone proof in the source is a geometrical counterpart. The first adjoint is empty, yielding a deficiency equal to the plane-section genus; the next adjoint restricts completely, eliminating further deficiencies. Its conclusion is correct. The alternative direct statement about ruled surfaces is only sketched, but the bundle calculation supplies a full proof.

The sentence saying that surfaces with simple double curves are ``usually'' regular is a qualitative observation, not a theorem with specified numerical hypotheses. The historical credits to Cayley, Noether, and Castelnuovo are reported as Enriques's credits. In particular, ``first example'' is not independently certified by the present reconstruction.

\subsection{Why the elementary sextic dimension count needs a generality argument}

A sextic surface has 84 coefficients, hence projective dimension 83. An assigned triple point imposes ten jet conditions. The subtraction \(83-8\cdot10=3\) is only an expected dimension: independence at eight points needs proof. Moreover ``arbitrarily assigned'' cannot mean every configuration, since the next paragraph deliberately supplies a special configuration where independence fails.

Here is an explicit modern repair for eight sufficiently general points. Let \(W=\operatorname{Bl}_{P_1,\ldots,P_8}\mathbb P^3\), let \(H\) be the pullback of a plane, and put \[
 L=2H-\sum E_i.
\] The quadrics through the points form a pencil. On a smooth member \(Q\), blow up the eight points to obtain \(S\) and let \(D=L|_S=-K_S\). The pencil base curve is a smooth elliptic quartic \(E\). The degree-zero bundle \(N=\mathcal O_E(D)\) is not torsion of order 1, 2, or 3 for a sufficiently general configuration. From \[
 0\longrightarrow\mathcal O_S((k-1)D)
 \longrightarrow\mathcal O_S(kD)\longrightarrow N^k\longrightarrow0
\] one obtains \(h^0(S,kD)=1\) for \(k=1,2,3\). Restriction to \(S\in|L|\) on \(W\) then gives \(h^0(W,kL)\leq h^0(W,(k-1)L)+1\). Binary degree-\(k\) polynomials in the two pencil quadrics already give \(k+1\) independent sections. Hence equality holds and \(h^0(W,3L)=4\). Every binary cubic factors over \(\mathbb C\), so every such sextic is a union of three quadrics from the pencil. This is the missing independence argument, with its generality hypothesis explicit.

\subsection{The Cayley-octad net: a complete calculation}

Now let the eight points be the reduced transverse base scheme of a general net of quadrics. On the same blowup \(W\), \(L\) is generated by three sections and defines \[
 f:W\longrightarrow\mathbb P^2.
\] The general fibre is the strict transform of a complete intersection of two quadrics, hence an elliptic quartic. For a general net the morphism is equidimensional with connected fibres; the smooth target and Stein factorization give \(f_*\mathcal O_W=\mathcal O_{\mathbb P^2}\). Consequently \[
 h^0(W,kL)=\binom{k+2}{2}.
\] In particular the sextics triple at the octad form projective dimension 9. This supplies both the lower bound by cubic expressions in the net and the upper bound relegated to the source's footnote.

Let \(B\subset\mathbb P^2\) be a general smooth plane curve of degree \(n\geq3\) and let \(X=f^{-1}(B)\in|nL|\). Generality ensures \(X\) is smooth on the blowup; its image in \(\mathbb P^3\) is a degree-\(2n\) surface with eight ordinary \(n\)-fold points. Since \[
 K_W=-4H+2\sum E_i=-2L,
\] adjunction gives \(K_X=(n-2)L|_X\). With \(f_X:X\to B\) and connected fibres, \[
 p_g(X)=h^0(B,\mathcal O_B(n-2))=\frac{n(n-1)}2.
\] The relative dualizing bundle is \[
 \omega_{X/B}=f_X^*\mathcal O_B(1),
\] because \(\omega_B=\mathcal O_B(n-3)\). Relative duality for this flat genus-one fibration gives \[
 R^1(f_X)_*\mathcal O_X=\mathcal O_B(-1).
\] Leray and \(H^0(B,\mathcal O_B(-1))=0\) now show \[
 q(X)=g(B)=\frac{(n-1)(n-2)}2,\qquad
 \chi(\mathcal O_X)=n,\qquad p_a=n-1.
\] For \(n=3\) these are \(p_g=3,p_a=2,q=1\), exactly as printed. This is a full proof of the general family formulas stated without calculation on p.~122.

Each exceptional \(E_i\simeq\mathbb P^2\) maps isomorphically to the parameter plane of the net, because the three local linear parts of the quadrics are independent at a transverse base point. Therefore \(E_i\cap X\simeq B\). This verifies the source's statement that the tangent-cone section at an ordinary \(n\)-fold point has the same genus as the irrational pencil's base. The exceptional curve also supplies a section of the resolved elliptic fibration. The equality \(q=g(B)\) is a computed fact of this family, not a consequence of having an arbitrary irrational pencil.

\subsection{Campedelli's degree notation}

The p.~123 paragraph cites Campedelli (1936), without title or pages, and says that the construction depends on a genus-one pencil on a quadric ``of degree \(m\),'' with \(m\mid n\), \(n=ms\), and eight ``\(s\)-fold'' base points. It then says that projection gives a Halphen pencil of plane curves of degree \(3m\) with nine \(m\)-fold base points. These data cannot all be correct. They are present in the scan; they are not a transcription error.

The following is a \textbf{modern mathematical reconstruction}, not a claim to have recovered the exact wording or complete proof of Campedelli's original article. The geometric construction is verified below; its identification with a particular emendation of Enriques's letters remains uncertain.

On \(Q\simeq\mathbb P^1\times\mathbb P^1\), let \(E\in|\mathcal O_Q(2,2)|\) be a smooth elliptic quartic through the eight points, and let \(Z=P_1+\cdots+P_8\). Set \[
 N=\mathcal O_E(2,2)(-Z).
\] Its degree is zero. If \(N\) has exact finite order \(m\), then on \(S=\operatorname{Bl}_8Q\) the anticanonical class \(D=-K_S\) has \(D^2=0\) and normal bundle \(N\) along its unique member for \(m>1\). The exact sequences above show \(h^0(kD)=1\) for \(1\leq k<m\) and \(h^1(kD)=0\) in that range, since \(\chi(kD)=1\) and \(h^2(kD)=0\). At \(k=m\), the restriction to \(N^m\simeq\mathcal O_E\) lifts and gives \(h^0(mD)=2\). A new member is disjoint from \(E\), so this pencil is basepoint-free. Its fibres are connected: writing \(S\to B\to\mathbb P^1\) for the Stein factorization, \(H^1(B,\mathcal O_B)\) injects into \(H^1(S,\mathcal O_S)=0\) because \(S\) is rational, so \(B\simeq\mathbb P^1\). If its finite map to the pencil parameter line had degree \(d\), connected-fibre pushforward would give \(h^0(S,mD)=h^0(\mathbb P^1,\mathcal O(d))=d+1\). Thus \(d=1\). Bertini and adjunction then make a general member a smooth connected genus-one curve. This is the quadric form of a Halphen pencil of index \(m\).

Its image on \(Q\) has \textbf{bidegree \((2m,2m)\), spatial degree \(4m\), and eight base points of multiplicity \(m\)}. Projection from one of the base points has plane degree \(4m-m=3m\). The two ruling lines through the projection point contract to two plane base points of multiplicity \(m\); the other seven points give the remaining seven. Thus there are nine \(m\)-fold plane base points, exactly as the final clause says. Neither spatial degree \(m\) nor spatial degree \(2m\) can give that clause.

The other printed condition, that the eight-point group comes by ``division by \(n\)'' from the series cut on \(E\) by degree-\(2n\) surfaces, is coherent: \[
 nZ\sim\mathcal O_E(2n)
 \quad\Longleftrightarrow\quad N^n\simeq\mathcal O_E.
\] It means that the minimal index \(m\) divides \(n\). The degrees are \(8n\) on both sides. Thus the torsion condition itself should be retained while the degree/base-multiplicity wording is corrected or explicitly flagged.

\subsection{Further reconstruction: parameter range and irregularity}

One can push this reconstruction far enough to explain the irregularity and expose a boundary qualification omitted from the short paragraph. This derivation is supplied here, not attributed to an unread historical proof.

On \(W=\operatorname{Bl}_8\mathbb P^3\), again put \(L=2H-\sum E_i\). Suppose the quadric pencil is general subject to the torsion condition of exact index \(m\). For \(k<m\), restriction to \(S\) gives \(h^0(W,kL)=k+1\). The fat-point ideal description gives \[
 \begin{aligned}
 \chi(W,kL)&=\binom{2k+3}{3}-8\binom{k+2}{3}=k+1,\\
 H^2(W,kL)&=H^3(W,kL)=0 \qquad(k\geq0).
 \end{aligned}
\] Hence \(H^1(W,(m-1)L)=0\), and the new section of \(mD\) lifts to a section \(\theta\in H^0(W,mL)\). If \(u,v\) are the pencil quadrics, their common zero scheme on \(W\) is the strict transform of \(E\). The restriction \(\theta|_E\) is a nonzero constant in \(N^m\simeq\mathcal O_E\). Consequently \(u^m,v^m,\theta\) have no common zero, and \(|mL|\) is generated. This supplies the morphism used below and makes Bertini applicable to \(|nL|\) when \(m\mid n\). Moreover, \(u,v,\theta\) are algebraically independent in the graded section ring: \(\theta\) varies along a general quadric, where its restriction supplies the nonconstant Halphen-pencil parameter.

The complete section ring is consequently \[
 \bigoplus_{k\geq0}H^0(W,kL)\simeq\mathbb C[u,v,\theta],
 \qquad\deg u=\deg v=1,\quad\deg\theta=m.
 \tag{IV.7.1}
\] For completeness, the dimension check is not assumed: \(h^0(S,kD)=\lfloor k/m\rfloor+1\) follows from the genus-one pencil, and restriction gives an upper bound whose successive increments are these numbers. The weighted monomials in \(u,v,\theta\) attain precisely that upper bound. They therefore span each graded piece, proving (IV.7.1).

For \(n=ms\), a general \(X\in|nL|\) is the inverse image of a weighted-plane curve \(B\) of degree \(n\) in \(\mathbb P(1,1,m)\). The coefficient of \(\theta^s\) can be chosen nonzero, so \(B\) avoids the weighted-plane singular vertex. The ratio \([u:v]\) gives a degree-\(s\) map \(B\to\mathbb P^1\); the discriminant has degree \(ms(s-1)\) and is simply branched for a general equation. Riemann--Hurwitz yields \[
 g(B)=\frac{(s-1)(ms-2)}2=\frac{(s-1)(n-2)}2.
 \tag{IV.7.2}
\] Adjunction still gives \(K_X=(n-2)L|_X\). Since \(W\) is rational and \(K_W=-2L\), the restriction sequence gives \[
 p_g(X)=h^0(W,(n-2)L),\qquad
 \chi(\mathcal O_X)=n,\qquad p_a(X)=n-1.
\] Using the monomial count in (IV.7.1), \[
 p_g(X)-(n-1)
 =\sum_{j=1}^{\lfloor(n-2)/m\rfloor}(n-jm-1)
 =\frac{(s-1)(n-2)}2.
\] Thus \(q(X)=g(B)\), agreeing with the type of equality asserted by Enriques. The generic fibres of \(X\to B\) are the genus-one members of the quadric Halphen pencil. This supplies a coherent modern explanation for the generalized family under explicit ordinary/genericity hypotheses.

It also makes the parameter qualification visible: with \(n\geq3\), these surfaces are irregular precisely when \(s>1\). The boundary \(s=1\), namely \(n=m\geq3\), has \(q=0\). The source's phrase ``one obtains irregular surfaces'' is safe as an existence assertion after choosing the parameters to make (IV.7.2) positive. Read as an assertion about every divisor \(m\mid n\), it requires this omitted restriction. Because the printed symbols are already inconsistent and the original Campedelli article was not inspected, the present reconstruction does not insist that Campedelli himself made this error; it localizes the issue to the compressed account on p.~123.

\subsection{The general irrational-pencil inequality}

Finally, if \(f:X\to B\) is a proper morphism with connected fibres and \(B\) smooth, so \(f_*\mathcal O_X=\mathcal O_B\), the low-degree Leray sequence begins \[
 0\longrightarrow H^1(B,\mathcal O_B)
 \longrightarrow H^1(X,\mathcal O_X)
 \longrightarrow H^0(B,R^1f_*\mathcal O_X)\longrightarrow0.
\] Therefore \(q(X)\geq g(B)\), and a positive-genus base forces irregularity. The source states the inequality and postpones its proof. This proof only needs the connected-fibre direct-image hypothesis; no assertion about injectivity of differentials under an arbitrary inseparable map is being substituted for it.

The constructions also clarify the scope of the methods. An ambient adjoint of degree \(n-4+t\) represents \(K_X+tH\), which checks all the index shifts in §§3--5. The defect of §4 is a restriction rank on \(H^1(\mathcal O_X)\); the defects of §5 telescope differences of \(h^1(K_X+sH)\). They are not interchangeable. On the quadric the Halphen pencil has spatial degree \(4m\), and projection from an \(m\)-fold point leaves plane degree \(3m\). The boundary \(s=1\) in the weighted construction is regular. Identifying every case and notation in Campedelli's original 1936 paper remains a historical source question, beyond the proved reconstruction.

The Hilbert-polynomial identities, local normalization lengths, blowup formulas, and Leray inequality have broader characteristic validity. The generic nodal-projection arguments, secant lemma for birational curve series, and Kodaira--Ramanujam vanishings have been used in characteristic zero. No unrestricted positive-characteristic version of the regularity equivalences is implied.

\section[Sistemi oltre il canonico]{Il teorema di Riemann--Roch per le superficie: sistemi più ampi del sistema canonico (pp.~123--127)}\label{enriques-iv-8}

The dimension statements concern complete systems. For \(n=C^2\), \(\pi=p_a(C)\), \(r=h^0(C)-1\), \(i(C)=h^0(K_X-C)\), and \(\omega(C)=h^1(C)\), their exact modern form is

\[r=p_a(X)+n-\pi+1-i(C)+\omega(C). \tag{IV.8.1}\]

A regular system has \(\omega(C)=0\), whereas a regular surface has \(q=0\). These conditions are distinct. Arithmetic genus is required for reducible or virtual divisors. The canonical degree denoted \(p^{(2)}\) on p.~144 is \(K_X^2\), and must not be confused with the plurigenus \(P_2\).

\subsection{The main proof and the distinction between containment and positivity}

Enriques recalls the inequality for adjoints \(C=K_X+L\) with \(L\) irreducible. He defines a system ``more ample than the canonical system'' by effective nonzero containment \(C-K_X>0\). This is an order relation among divisor classes, not the modern definition of an ample line bundle. He correctly observes that the residual can have fixed components and need not be irreducible. The subtraction argument removes that restriction.

The proof on pp.~124--125 becomes precise as follows. Fix an ample \(H\), take \(s\) sufficiently large, and set \(A=K_X+sH\), \(D=A-C\). The complete systems can be chosen with the regularity, global generation, and irreducibility required in the argument. These facts follow from Serre vanishing and eventual very ampleness for a fixed line bundle twisted by \(sH\) (Stacks Project, Lemma 30.17.1, Tag 01XO). Choose \(D\) smooth. Since \(C-K_X\) is effective and nonzero, \((K_X-C)D<0\), so adjunction and curve duality give

\[
h^1(D,\mathcal O_D(A))=h^0(D,\mathcal O_D(K_X-C))=0.
\]

Write \(D^2=\nu\), \(p_a(D)=\rho\), and \(CD=t\). The complete restricted line bundle has \(h^0=\nu+t-\rho+1\), and the image of surface sections has no larger dimension. Consequently

\[
r(C)\ge r(A)-(\nu+t-\rho+1).
\]

Together with

\[
A^2=n+\nu+2t,\qquad p_a(A)=\pi+\rho+t-1,
\]

this gives \(r(C)\ge p_a+n-\pi+1\), including reducible \(C\). In (IV.8.1), the reason is \(i(C)=0\) and \(\omega(C)\ge0\). Enriques supplies the subtraction calculation; the high-degree existence assertions are sketches dependent on earlier postulation results.

\subsection{The singular-point illustration on p.~124}

Two displayed errors are in the print. If \(\bar C=C-\theta\), \(C\theta=0\), and \(\theta^2=-\nu\), then

\[
\bar n=n-\nu,\qquad \pi=\bar\pi+\rho+\nu-1.
\]

The source instead prints \(n=\bar n-\nu\) and \(\bar\pi=\bar\pi+\rho+\nu-1\). The scan confirms both the minus sign and the extraneous left overbar. The missing left system bars in \(\bar C=|C-\theta|\) are also printed.

There is a further hypothesis issue. The formulas \(\theta^2=-\nu\) and \(g(\theta)=(\nu-1)(\nu-2)/2\) describe an \textbf{ordinary} isolated hypersurface point with smooth plane tangent cone, not an arbitrary isolated point of multiplicity \(\nu\). Moreover, subtracting one \(\theta\) from \(L=C-K_X\) need not leave an irreducible residual. On the resolution of an ordinary \(\nu\)-fold point of a degree-\(d\) surface,

\[
K_X=(d-4)H-(\nu-2)\theta.
\]

For large \(C=sH\), \(L\) contains fixed contribution \((\nu-2)\theta\), and \(L-\theta\) retains \((\nu-3)\theta\) when \(\nu>3\). The assertion needs a restriction or a different subtraction. This affects an illustration, not the subsequent general proof.

Under the illustration's intended extra assumptions, the corrected calculation does give strict superabundance. Basepoint freedom gives \(h^0(C)-h^0(C-\theta)=1\), while \(\bar n-\bar\pi=n-\pi+\rho-1\). The adjoint lower bound therefore gives \(r(C)\ge p_a+n-\pi+1+\rho\). These corrections belong in commentary, not invisibly in the diplomatic source.

\subsection{The addition assertion on p.~126 is too strong}

For regular nonspecial \(C\), and irreducible smooth \(D\) with genus \(\rho\), square \(\nu\), and \(m=CD>0\), the source correctly obtains

\[
(C+D)^2=n+\nu+2m,\qquad p_a(C+D)=\pi+\rho+m-1.
\]

The condition \(m\ge\rho-\nu\) makes the virtual dimension of \(C+D\) strictly exceed \(r(C)\). With the moving-system and Bertini hypotheses, the intended irreducibility conclusion follows. This does \textbf{not} imply the next assertion that \(C+D\) is regular and its restriction to \(D\) nonspecial. The obstruction is explicit in

\[
0\longrightarrow\mathcal O_X(C)\longrightarrow\mathcal O_X(C+D)
\longrightarrow\mathcal O_D(C+D)\longrightarrow0.
\]

Since \(h^1(C)=h^2(C)=0\), it gives

\[
h^1(C+D)=h^1(D,\mathcal O_D(C+D))
=h^0(D,\mathcal O_D(K_X-C)).
\tag{IV.8.2}
\]

Restriction is complete because \(h^1(C)=0\), but it can be special. A sufficient extra condition is \((C-K_X)D>0\). The source's weaker inequality does not force this.

\textbf{Counterexample.} Let \(X\subset\mathbb P^1\times\mathbb P^2\) be a smooth hypersurface of bidegree \((2,4)\). Put \(F=\mathcal O_X(1,0)\), \(H=\mathcal O_X(0,1)\). Adjunction and intersections give

\[
K_X=H,\quad F^2=0,\quad FH=4,\quad H^2=2,\quad
q=0,\quad p_g=p_a=3.
\]

Take very ample \(C=H+F\) and a smooth fibre \(D\in|F|\), a plane quartic. Then \(\rho=3,\nu=0,m=4\ge3\). Ambient restriction sequences and projective-space cohomology yield:

\textbf{Divisor: \(H+F\).} \(h^0\): 6; \(h^1\): 0; Square: 10; Arithmetic genus: 9; Actual \(r\): 5; RR value with \(i=0\): 5.

\textbf{Divisor: \(H+2F\).} \(h^0\): 9; \(h^1\): 1; Square: 18; Arithmetic genus: 15; Actual \(r\): 8; RR value with \(i=0\): 7.

The ambient exact sequences use \(\mathcal O(-1,-3)\) for the first divisor and \(\mathcal O(0,-3)\) for the second. The latter has one-dimensional second cohomology, producing \(h^1_X(H+2F)=1\); the former has none. Moreover \((H+2F)|_D=\mathcal O_D(1)=K_D\), with speciality one. Every numerical hypothesis in the printed assertion holds, while regularity fails. This is an explicit verification, not a claim of historical novelty.

The asymptotic conclusion on pp.~126--127 survives: for a fixed \(C\), choose \(D=sH_0\) with \(H_0\) ample and \(s\gg0\). Both \(D\) and \(C+D\) are regular and very ample. The false assertion for arbitrary \(D\) is unnecessary. Finally, strict negativity of a fundamental curve's square requires a big moving system. From \(CD=0\) alone one cannot infer \(D^2<0\): fibres of a pencil have square zero. In the big case, the Hodge index theorem supplies the required negativity.

\section[Curve eccezionali]{Eliminazione delle curve eccezionali (pp.~127--128)}\label{enriques-iv-9}

The source distinguishes a rational curve \(E\) with \(E^2=-1\), an exceptional curve of the first kind, from a rational curve with nonnegative square, called the second kind. The latter terminology is not modern exceptional-curve terminology.

Here the addition argument is valid despite the false general predecessor. If \(m=CE>0\) and \(E^2=\nu\ge-1\), then \(\mathcal O_E(C+E)\) on \(\mathbb P^1\) has degree \(m+\nu\ge0\), so \(H^1\) vanishes. The exact sequence proves regularity at every allowed step. Intersections evolve as

\[
(C+kE)E=m+k\nu.
\]

For \(\nu=-1\), the process ends at \(k=m\), when \(E\) has degree zero for \(C+mE\). For \(\nu\ge0\), it continues indefinitely. The additive character is

\[
a(C)=2p_a(C)-2-C^2=K_X C,
\]

and \(a(E)=-2-\nu\le-1\). Indefinite addition therefore produces moving classes with negative canonical intersection. These numerical steps are explicitly proved.

The geometric steps are sketches. Degree zero on \(E\) alone does not establish a birational morphism or a smooth image point. The precise input is Castelnuovo contraction: a smooth rational \((-1)\)-curve contracts to a smooth point. The construction with sufficiently ample \(C\) and \(C+mE\) has a modern proof in Stacks Project, Lemma 54.16.8, Tag 0C2M. This is a modern justification, not a claim that this source was the first historical repair.

Enriques asserts finiteness and excludes intersecting exceptional curves under his nonruled hypothesis; global finiteness is not proved here. If another \((-1)\)-curve \(E'\) meets \(E\) in \(a>0\), its image has square \(-1+a^2\ge0\); for \(a>1\), that image can be singular, so a smooth-rational-curve argument alone does not finish the reasoning. Modern nonruled minimal-model theory supplies the full interpretation. Successive contractions terminate because each lowers Picard number. The reducible configurations discussed are those arising from successive point blowups, contracted in reverse order, not arbitrary negative reducible divisors. Infinitely near points mean centres on successive models.

The closing statement that the transformation ``fails'' for ruled surfaces concerns eliminating both of the source's kinds. Ruled surfaces also have minimal models, meaning models without \((-1)\)-curves, but can retain fibres of square zero. Every positive plurigenus of a ruled surface vanishes because \(K_X\) restricts as \(\mathcal O_{\mathbb P^1}(-2)\) on a general fibre. The classification of the negative-canonical moving-system alternative as ruled is deferred, not proved here.

\section[Restrizione dei sistemi lineari]{Integrità della serie segata da un sistema lineare piccolo sulle curve di un sistema grande (pp.~128--131)}\label{enriques-iv-10}

\subsection{A precise statement and complete modern proof}

For fixed divisors \(C,E_0\), an ample \(H\), and a smooth general \(D_s\in|sH+E_0|\) with \(s\gg0\), the restriction

\[
H^0(X,\mathcal O_X(C))\longrightarrow H^0(D_s,\mathcal O_{D_s}(C))
\]

is surjective, and eventually an isomorphism. In

\[
0\to\mathcal O_X(C-D_s)\to\mathcal O_X(C)
\to\mathcal O_{D_s}(C)\to0,
\]

duality identifies \(H^1(C-D_s)^*\) with \(H^1(K_X-C+D_s)\), which vanishes by Serre vanishing. Negative ample intersection gives \(H^0(C-D_s)=0\). This works for any fixed line bundle \(C\), effective or not, and identifies the required positivity precisely.

\subsection{The geometric proof, its valid dimension step, and its gaps}

Enriques starts from arbitrary \(G\in|\mathcal O_D(C)|\). He seeks \(G'\) such that \(G+G'\) is characteristic and \(G'\) is cut by an irreducible \(B\in|D-C|\). Once available, the pencil through \(G+G'\) has a member containing \(B\); its residual \(C\) cuts \(G\). Resolve the basepoints: the induced pencil has residual degree zero on \(B\), hence is constant there, which gives the claimed reducible member.

The linear-algebra step is sound. Two projective subseries of \(\mathbb P^N\), of dimensions \(a,b\), whose deficiencies are less than their own dimensions satisfy \(N-a<a\) and \(N-b<b\). Thus \(a+b>N\), and they intersect. This is stronger than necessary but valid. No gap should be invented in this criterion.

The subsequent growth calculation has two mistakes and an omitted condition:

\begin{enumerate}
\item
  \(n_s=s^2n\) and \(\pi_s=s\pi+s(s-1)n/2-s+1\) are correct. Substitution gives \[
  r_s\ge p_a+\frac{s^2n}{2}+\frac{sn}{2}-s\pi+s.
  \] The print instead has \(-sn/2\) and constant \(-1\). Its quadratic leading term still supplies the intended asymptotic growth.
\item
  For characteristic defect \(d\) and speciality \(i=h^1(D,\mathcal O_D(D))\), the general identity is \[
  i=d+p_a+\omega(sL).
  \] The printed \(i=d+p_a\) assumes regularity not established at that stage. The weaker \(d\le i-p_a\) is sufficient.
\item
  The strict bound \(i<\deg(K_X|_D)\) is false. Use \[
  h^0(M)\le\max(\deg M+1,0)
  \] on a smooth connected curve. On a K3 surface, \(K_X|_D\) has degree zero and one section. The corrected bound still gives \(i=O(s)\), since \(K_X D=s(2\pi-2-n)\).
\end{enumerate}

Imposing \(G\) of length \(sn\) costs at most \(sn\) conditions, so the relevant series retain leading dimension \(s^2n/2\), and the common subseries exists for large \(s\). The print does not separately ensure that a suitable common divisor comes from an \textbf{irreducible} \(B\), and only sketches passage from the worked special case to general fixed \(C,E_0\). These are missing justifications, not counterexamples to the precise restriction theorem. The cohomological argument above supplies a complete replacement.

\subsection{The literal effective-domination reading is too broad}

The phrase ``more ample than \(sL\)'' on p.~129 cannot safely mean mere effective domination with an arbitrary residual varying with \(s\).

\textbf{Counterexample.} On \(X=\operatorname{Bl}_p\mathbb P^2\), let \(\ell\) be the pulled-back line and \(E\) the exceptional curve. The divisor \(H=C=2\ell-E\) is very ample. Put \(D_s=2s\ell\), choosing a smooth member avoiding \(E\), for \(s\ge2\). Then

\[
D_s-C-(s-2)H=2\ell+(s-1)E
\]

is effective, so \(D_s-C\) contains arbitrarily high multiples of an ample system in the literal order sense. Yet \(D_s\) contracts \(E\), and \(D_s-C=(2s-2)\ell+E\) has fixed \(E\). On \(D_s\), \(C|_{D_s}\simeq\mathcal O_{D_s}(2)\), which has six sections; \(H^0(X,C)\), conics through \(p\), has dimension five. Restriction is incomplete for every \(s\ge2\).

This identifies a needed precision in terminology. It does not refute the ample-twist interpretation or the source's own warning about fundamental curves. A sufficient modern formulation is \(D_s=sH+E_0\) with \(E_0\) fixed, or directly \(H^1(K_X-C+D_s)=0\). The footnote's printed \(L_1\) is also an index slip: surrounding notation requires \(L_s\).

\section[Riemann--Roch per ogni sistema]{Il teorema di Riemann--Roch per i sistemi lineari qualunque (pp.~131--134)}\label{enriques-iv-11}

Adjunction gives \(K_D-D|_D=K_X|_D\). For large \(D\), restriction identifies \(H^0(K_X|_D)\) with \(H^0(K_X)\), so the complete characteristic line bundle has speciality \(p_g\), including \(p_g=0\). Enriques treats the zero case separately: hypothetical speciality forces a sufficiently small \(C\) to lie in its adjoint \(C'\), contradicting absence of canonical divisors. The componentwise intersection comparison works once \(D\) is sufficiently large against every component involved.

For any fixed effective \(C\), the same reasoning gives

\[
h^1(D,\mathcal O_D(C+D))
=h^0(D,(K_X-C)|_D)=h^0(K_X-C)=i(C).
\]

Subtract \(D\) from regular sufficiently positive \(C+D\) to obtain \(r\ge p_a+n-\pi+1-i\). This is the source's general inequality. Equation (IV.8.1) identifies the omitted nonnegative quantity as \(h^1(C)\). The proof is complete once the previous restriction lemma and positivity construction are made precise.

A multiplicity-\(m\) basepoint condition has length \(m(m+1)/2\), and imposes at most that many independent conditions. On the blowup,

\[
\widetilde C=f^*C-mE,\quad
\widetilde C^2=C^2-m^2,\quad
p_a(\widetilde C)=p_a(C)-m(m-1)/2.
\]

Use the complete prescribed-basepoint system on the blowup, not an arbitrary incomplete subsystem. Speciality is compatible because

\[
K_{\widetilde X}-\widetilde C=f^*(K_X-C)+(m+1)E
\]

has the same sections as \(K_X-C\). The calculation iterates for infinitely near conditions. In the same notation, \(\chi(\widetilde C)=\chi(C)-\binom{m+1}{2}\); equality between the length of the assigned jet scheme and its number of independent conditions still requires independence. These formulas follow from the canonical blowup formula and the surface Riemann--Roch identity just obtained.

The residual formulas on pp.~133--134 are correct. For \(C_1+C_2=K_X\), \(n_j=C_j^2\), \(\pi_j=p_a(C_j)\), and \(m=C_1C_2\),

\[
p^{(1)}=\pi_1+\pi_2+m-1,\quad
p^{(1)}-1=n_1+n_2+2m,\quad
m=2\pi_1-2-2n_1.
\]

Hence

\[
\pi_2=p^{(1)}-3\pi_1+2n_1+3,\qquad
n_2=p^{(1)}+3n_1-4\pi_1+3.
\]

The prose once reverses ``genus and degree'' for \(n_1,\pi_1\); the equations consistently use degree and genus. Since \(i(C_1)=r_2+1\), the paired equalities

\[
r_1=p_a+n_1-\pi_1-r_2+\omega_1,\qquad
r_2=p_a+n_2-\pi_2-r_1+\omega_2
\]

give \(\omega_1=\omega_2\). Modern duality is exactly this symmetry, \(h^1(C)=h^1(K_X-C)\). No residual formula needs a sign correction.

\subsection{A complete derivation from curve Riemann--Roch}

Let \(X\) be smooth and projective and put

\[
Q(D)=\chi(\mathcal O_X(D))-\chi(\mathcal O_X)-\tfrac12D(D-K_X).
\]

For a smooth effective curve \(C\subset X\), the exact sequence

\[
0\longrightarrow\mathcal O_X(D-C)\longrightarrow\mathcal O_X(D)
\longrightarrow\mathcal O_C(D)\longrightarrow0
\]

and curve Riemann--Roch give

\[
\chi(D)-\chi(D-C)=D\cdot C+1-g(C).
\]

Adjunction gives \(2g(C)-2=C(C+K_X)\). Substituting shows \(Q(D)=Q(D-C)\). Choose a very ample \(H\) and a sufficiently large \(m\). Choose smooth divisors \(E\in|D+mH|\) and \(T\in|mH|\). Applying the preceding equality to \(E\), with the subtracted curve also \(E\), gives \(Q(E)=Q(0)=0\). Applying it again with the subtracted curve \(T\) gives \(Q(D)=Q(E-T)=Q(E)=0\). This proves

\[
\chi(D)=\chi(\mathcal O_X)+\tfrac12D(D-K_X)
\]

for every divisor class, including classes with no effective representative. Serre duality identifies \(h^2(D)=h^0(K_X-D)\), giving the chapter's exact dimension formula and inequality. This proof makes the extension to ``virtual curves'' automatic; it is not necessary to assume existence of an intermediate effective difference. For the standard theorem and duality statement, see Beauville (1996, I.11--I.12). Curve Riemann--Roch, including the Gorenstein case, is stated in The Stacks Project Authors (2026, Tag 0BS6).

\section[Deficienza della serie caratteristica]{Deficienza della serie caratteristica (pp.~134--137)}\label{enriques-iv-12}

\subsection{Exact meaning and proof of the bound}

For a smooth connected member \(C\) of a complete moving system,

\[
0\to\mathcal O_X\to\mathcal O_X(C)\to\mathcal O_C(C)\to0
\]

identifies the characteristic vector space with the image of \(H^0(C)\) modulo the defining section. Therefore

\[
\boxed{d=\dim\ker[H^1(\mathcal O_X)\to H^1(\mathcal O_X(C))]\le q.}
\tag{IV.12.1}
\]

If \(h^1(C)=0\), equality holds, so sufficiently ample systems attain the maximum \(q\).

The source first assumes \(i(C)=0\) and uses defects \(d,d'\) of the characteristic and canonical restriction series. Curve RR gives

\[
r=n-\pi+p_g+d'-d+1,\qquad
\omega(C)=q+d'-d.
\]

For a regular sufficiently large system, \(d'=0\), hence \(d=q\). The comparison lemma then proves the general bound. Its rank calculation is sound: for \(A=C_1+C_2\), dimensions \(r=\dim|A|\), \(r_j=\dim|C_j|\), and intersection scheme \(G\) imposing \(s-\varepsilon\) conditions, completeness on \(C_1\) gives

\[
d_1=r-r_1-r_2-(s-\varepsilon).
\]

The corresponding image on \(C_2\) has relative defect equal to this expression, which is at most \(d_2\). Thus \(d_1\le d_2\). Scheme lengths replace point counts for nontransverse intersections.

Choose \(D\) sufficiently ample. Then \(C+D\) is regular and restricts completely to \(C\) because \(H^1(D)=0\). The printed condition \((CD)>n-\pi\) reverses the threshold in §8, and the cited arbitrary-addition regularity statement was too strong. Neither defeats this asymptotic proof: the available ample \(D\) supplies all required vanishings.

When \(q=0\), every complete smooth moving system has complete characteristic series. In the projective reformulation, ``normal'' means \textbf{linearly normal}, not normality of local rings or projective normality in all degrees. A linearly normal embedding of a regular surface has linearly normal smooth hyperplane sections. Failure of section linear normality implies irregularity.

\subsection{The scroll illustration needs a correction}

For a ruled surface over a genus-\(g\) curve, \(p_g=0,q=g,p_a=-g\). If \(H\) embeds it as a scroll, then \(H\) has fibre degree one and a smooth general member has genus \(g\). The correct formulas are

\[
r=n-2g+1+h^1(X,H),\qquad
d=g+h^1(C,H|_C)-h^1(X,H).
\tag{IV.12.2}
\]

The source describes the extra term in \(r\) as the speciality of the hyperplane section and concludes that \(d=g\) in every nonconical normal case. This identifies different cohomology groups and is false without another condition.

\textbf{Counterexample.} Let \(B\subset\mathbb P^2\) be a smooth plane quartic, \(X=B\times\mathbb P^1\), and

\[
H=K_B\boxtimes\mathcal O_{\mathbb P^1}(1).
\]

The Segre embedding in \(\mathbb P^5\) is a smooth nonconical linearly normal scroll. One has \(H^2=8\), \(p_a(H)=3\), \(h^0(H)=6\), and \(h^1(H)=2\) by Künneth. A general section is the graph of a degree-four map \(B\to\mathbb P^1\) from a basepoint-free pencil in \(H^0(K_B)\). Hence \(C\simeq B\), \(H|_C\simeq2K_B\), \(h^0(H|_C)=6\), and \(h^1(H|_C)=0\). The characteristic vector image has dimension \(6-1=5\), so \(d=1\), whereas \(q=g=3\). This does not rely on an incomplete embedding or a cone.

The Segre footnote also has a date mismatch. Institutional author and BDIM catalogues identify the relevant \emph{Recherches générales\ldots{}}, part II on ruled surfaces, as \emph{Mathematische Annalen}, 34 (1889), 1--25; the abbreviated printed reference says 1892. This verifies the matching work's date, not the exact theorem or priority of every assertion attributed to it. The original article has not been examined in full here.

\subsection{Irrational pencils and algebraic families}

For a morphism \(f:X\to B\) with connected fibres, Leray injects \(H^1(B,\mathcal O_B)\) into \(H^1(X,\mathcal O_X)\), so \(q(X)\ge g(B)\). Resolve a rational pencil and use Stein factorization before applying this. The source attributes the result to Castelnuovo and relates it to the inequality for families of inequivalent divisors.

Precisely, a family whose image in \(\operatorname{Pic}(X)\) has dimension \(d\) satisfies \(d\le q\). In characteristic zero the Picard scheme is smooth of dimension \(h^1(\mathcal O_X)\), and its tangent space at the identity is \(H^1(\mathcal O_X)\) (Kleiman, 2005, Theorem 5.11 and Corollary 5.14). The source's infinitesimal argument requires a smooth parameter point and the expected tangent rank. It is not a proof that every infinitesimal normal deformation integrates, nor a general completeness theorem for continuous characteristic systems. No stronger assertion is needed. The construction using groups of \(g\) fibres corresponds to the dominant Abel map \(\operatorname{Sym}^gB\to\operatorname{Pic}^gB\).

\subsection{Comparing the three defects}

For a smooth connected member, let \(d\) be the characteristic deficiency and \(d^\prime\) the deficiency of \(H^0(X,K_X)\to H^0(C,K_X|_C)\). Curve and surface Riemann--Roch give

\[d-d^\prime=q-h^1(X,\mathcal O_X(C)).\]

Thus characteristic deficiency, canonical restriction deficiency, and superabundance are distinct. A regular surface has complete characteristic series even for a superabundant complete system. No nonspeciality assumption on the restricted line bundle is needed for the bound \(d\le q\).

\section[Nota storica]{Nota storica (pp.~137--142)}\label{enriques-iv-13}

\subsection{Evidentiary status}

These pages are Enriques's retrospective account as a participant. The following are \textbf{claims made by Enriques}, not independently certified priority findings:

\textbf{Pages: 137--138.} Enriques's report: Cayley's negative arithmetic count for ruled surfaces; Zeuthen's invariance with singularity restrictions; Noether's early comparison of the genera; Assessment: The modern invariants are checked. Exact original restrictions and Noether's reported presumptions have not been reconstructed from all original papers.

\textbf{Pages: 138.} Enriques's report: Castelnuovo 1891 and Picard/Humbert irregular examples; ``hyperelliptic'' surfaces with \(p_g=1,p_a=-1\); Assessment: These invariants fit abelian-type examples. More generally, smooth products of positive-genus curves supply modern irregular examples, so nonnegative arithmetic genus does not force equality with geometric genus. The old term is not today's bielliptic surface, whose \(p_g=0,q=1\). Priority remains attributed.

\textbf{Pages: 138--139.} Enriques's report: Enriques 1893/1896; Picard 1905 adjoint regularity; Severi 1908 algebraic route; Assessment: The modern mathematical content is supplied. Original proof scope has not been exhaustively verified.

\textbf{Pages: 139.} Enriques's report: Noether 1886 implicitly assumed completeness of two restricted series; Assessment: This is Enriques's explicit criticism. The defect identity verifies its mathematical importance; the commentary has not invented it.

\textbf{Pages: 139--140.} Enriques's report: Castelnuovo 1896/1897 maximum defect and Severi 1903/1905 simplifications; Assessment: The resulting theorem and proof route are checked; chronology is not independently adjudicated.

\textbf{Pages: 141--142.} Enriques's report: The 1940 proof is simpler; reducible cases were known early ``at least implicitly,'' and explicit in joint work around 1900; Assessment: Simplicity and implicit knowledge are authorial judgments. The subtraction proof does cover reducible classes. The joint paper's 1900/1901 dating needs care.

The Noether paragraph has a source misprint: the restricted canonical series should have projective dimension \(p_g-i-1\), since \(p_g-i\) is its vector dimension. The printed exponent is \(p-i\). With defects \(d,d'\), the correct equality is

\[
r=p_g+n-\pi+1-i+d'-d.
\]

Completeness of the surface systems alone does not establish completeness of either restriction. Both restriction defects zero suffice for the displayed Noether equality; the defects could also cancel, so the equality alone does not prove individual completeness.

The next historical paragraph displays \(r\ge p_g+n-\pi+1\) without specifying nonspeciality. This cannot apply to every system: on a regular surface with \(p_g>1\), taking \(C=K_X\) gives \(r=p_g-1\), whereas that right side is \(p_g\). The summary therefore needs the restriction \(i(C)=0\), or the missing speciality term. This is a scope qualification of the retrospective paragraph, not a finding about the exact wording of Enriques's 1893 original paper.

Picard's theorem is discussed for simple systems with at least three parameters. A modern translation must retain the actual positivity/birational hypotheses. Fixed ample twists are handled by Serre vanishing; broader big-nef adjoint vanishing is a characteristic-zero theorem. This commentary does not identify an unspecified modern result as the first historical repair by a named author.

\subsection{\texorpdfstring{The \(Af+B\varphi\) route}{The Af+B\textbackslash varphi route}}

Let \(G\) be the plane complete-intersection scheme defined by \(f_n,f_s\) on \(z=0\). If \(f_m\) vanishes on \(G\), then modulo \(z\) its class lies in the ideal generated by \(f_n|_{z=0}\) and \(f_s|_{z=0}\). Choose homogeneous coefficients of degrees \(m-n,m-s\), lift them, and divide the remaining polynomial by \(z\):

\[
f_m=Af_n+Bf_s+zf_{m-1}.
\]

Negative coefficient degree means that coefficient is zero. For general transverse sections, pointwise and scheme-theoretic vanishing coincide. For nonreduced intersections, scheme-theoretic vanishing is required; merely passing through the support is insufficient.

Enriques summarizes Severi's descent from degree \(m\) to \(m-1\) for adjoints through the double points of \(D_s=F_n\cap F_s\). The descent and exact adjoint/conductor conditions are not developed here; they are assumed input, not proved by the polynomial identity alone. For \(s>n-4\), a degree-\((n-4)\) surface through all intersections of \(F_s\) with each double-curve component must contain that component by Bézout, provided the intersections are general and transverse. It is therefore an adjoint of \(F_n\), giving the stated completeness in this setting.

The remaining deductions repeat the preceding route: restriction completeness gives RR; sufficiently high regularity gives characteristic defect \(q\); the comparison lemma gives the universal bound. Eventual regularity is explicitly deferred to §15. The unclosed quotation mark on p.~141 is printed. The later comparison on the same page concerns \(|D_s|\), as does the argument throughout.\footnote{The edition prints \(|D'_s|\); one might reject that prime as a transcription error and suppose that none is printed. Enlargement of p.~141 shows a small high stroke compatible with the prime retained by the supplied diplomatic transcription. Its categorical rejection is therefore unwarranted. If the mark is read as print, the prime is a harmless label slip: no second system has been introduced, and the mathematical system intended here is \(|D_s|\).}

\section[Curve virtuali]{Curve virtuali (pp.~142--147)}\label{enriques-iv-14}

\subsection{Numerical tests and the positive-square criterion}

A virtual curve \(C=A-B\) is a divisor class whose complete system may be empty. Adding a common sufficiently ample divisor makes \(A,B\) regular, smooth in general, basepoint free, and without contracted curves while preserving \(C\). This reduction is crucial; arbitrary nonample representatives cannot be substituted.

If \(C\) is effective,

\[
C(A+B)=A^2-B^2\ge0,\quad
CB=AB-B^2\ge0,\quad CA=A^2-AB\ge0.
\]

Thus \(A^2\ge B^2\) and \(B^2\le AB\le A^2\), also necessary if any positive multiple of \(C\) is effective. For nonzero effective \(C\) and ample \(A,B\), the latter inequalities are strict. The possible zero divisor explains the weak form.

The sufficient criterion is correct: if \(C^2>0\) and \(A^2>AB>B^2\), every sufficiently large \(mC\) is effective. Indeed \(BC>0\), so \(B(K_X-mC)<0\) for \(m\gg0\), and \(h^0(K_X-mC)=0\). Meanwhile

\[
\chi(mC)=\chi(\mathcal O_X)+\frac{m^2C^2-mK_X C}{2}
\]

has positive quadratic leading term. Hence

\[
h^0(mC)=\chi(mC)+h^1(mC)>0.
\]

Enriques instead subtracts \(B\) repeatedly from \(mA\). At the last step the restriction degree is

\[
(mA-(m-1)B)B=AB+(m-1)(AB-B^2),
\]

eventually greater than \(2g(B)-2\). All earlier restriction degrees are at least as large, so they too are nonspecial. This supplies the uniform induction needed to justify the apparently assumed existence of the penultimate residual. The source does address that issue; it should not be called simply circular. The formulas

\[
n_m=m^2n,\qquad \pi_m=m\pi+\frac{m(m-1)}2n-m+1
\]

are correct, including the final sign and constant.

\subsection{The missing hypothesis in the plurigenus bound}

The simplified inequality printed on pp.~142--144 is

\[
P_m\ge p_a+1+\frac{m(m-1)}2K_X^2\qquad(m>1).
\]

The exact modern equality is

\[
\boxed{P_m=\chi(\mathcal O_X)+\frac{m(m-1)}2K_X^2+
h^1(mK_X)-h^0((1-m)K_X).}
\tag{IV.14.1}
\]

The simplified bound follows if \(h^0((1-m)K_X)=0\). A sufficient hypothesis is positive \(K_X H\) for an ample \(H\), as in the intended nontrivial positive canonical class of the nonruled setting. The condition \(K_X^2>0\) alone does not choose the positive component: \(\mathbb P^2\) has \(K_X^2=9\) but every \(P_m=0\). The unrestricted formula also fails on a K3 surface: \(K_X=0,p_a=1,P_m=1\), whereas the claimed lower bound is two. There \((B+mK_X)|_B=\mathcal O_B(B)=K_B\), so the asserted nonspeciality on p.~143 fails by one. These are qualifications of a printed mathematical assertion, not OCR corrections. Retain canonical positivity or the full term in (IV.14.1).

\subsection{Boundary square and torsion}

The source correctly says the positive-square argument is inconclusive for \(C^2=0\) and \(C^2<0\). An explicit square-zero example clarifies why. On \(X=B\times E\), with \(E\) elliptic, take \(b\in B\) and a nontorsion degree-zero \(P\in\operatorname{Pic}^0(E)\). Put

\[
L=\operatorname{pr}_B^*\mathcal O_B(b)\otimes\operatorname{pr}_E^*P.
\]

Then \(L^2=0\) and \(L\) meets an ample divisor positively, yet \(h^0(mL)=0\) for all \(m>0\), because \(P^m\) has no sections. Removing the twist yields a numerically equivalent class with effective multiples. It shows that numerical data do not settle boundary effectivity. For negative square, an effective negative curve and its negative already illustrate opposite effectivity behavior at the same square. Even positive square requires an orientation: on \(\mathbb P^2\), both \(H\) and \(-H\) have square one, while only positive multiples of \(H\) are effective. Intersecting with an ample class selects the relevant component of the positive cone.

Now assume \(A^2=AB=B^2\), retaining positive representatives. Then \(C^2=0\), \(C(A+B)=0\), and \(C\) is numerically trivial. An effective numerically trivial divisor must be zero because every nonzero effective divisor meets an ample divisor positively. Hence

\[
|mC|\ne\varnothing\quad\Longleftrightarrow\quad mC\sim0.
\]

Enriques's \(\equiv\) in this passage means linear equivalence; reading it as today's numerical-equivalence symbol would erase the substantive issue.

On a regular complex surface, numerical triviality implies torsion: some multiple is algebraically trivial, the torsion components of the Picard scheme form a finite group, and \(q=0\) makes \(\operatorname{Pic}^0\) trivial. Therefore some \(mA\sim mB\). Without regularity, some multiples are still algebraically equivalent, but need not be linearly equivalent (Kleiman, 2005, Definition 6.1, Theorem 6.3, Theorem 6.16, Corollary 6.17). Equal zero intersections for arbitrary nonample isotropic divisors do not suffice for this conclusion.

The source's bounded-degree-family argument sketches the existence and boundedness of infinitely many twists; the Picard proof supplies the missing precision. References to Severi, Todd, Hodge, B. Segre, and Bronowski remain Enriques's citations, without an invented claim of historical priority.

\subsection{General virtual RR and speciality}

For \(C=A-B\), define \(i=h^0(K_X-C)\) even when \(|C|\) is empty. Equation (IV.8.1) remains valid. Thus

\[
p_a+n-\pi+1-i\ge0
\]

implies \(h^0(C)\ge1\). This ensures effectivity, not irreducibility or a moving system. The source takes large \(D=sB\), restricts \(K_X-C\) completely to it, and obtains speciality \(i\) for \(\mathcal O_D(C+D)\). Subtraction then gives the inequality. Although written for \(i>0\), the same exact-sequence proof treats \(i=0\).

For effective \(C\), multiplication by its defining section injects \(H^0(K_X-C)\) into \(H^0(K_X)\), giving \(i\le p_g\). For virtual \(C=-mH\), \(i=h^0(K_X+mH)\) grows quadratically and is unbounded. The final assertion that the residual lies inside the canonical system is correct when this means translation by a fixed effective divisor representing \(C\), not literal inclusion of unrelated vector spaces without specifying the multiplier.

\subsection{Limits of the constructions}

Specific limits remain. The isolated-point wording does not state all singularity assumptions. ``Normally large'' admits an overly literal reading, which has been tested separately from a precise sufficient statement. Not all original historical papers were read, so their history is attributed rather than certified. Faint punctuation and ink distinctions remain optically uncertain; those uncertainties do not change a mathematical formula. The uncertainty over the high stroke after \(n\) on p.~113 is recorded where that formula is discussed.

For a concrete comparison, on a product with an elliptic factor let \(M\) be a nontorsion degree-zero line bundle pulled back from that factor and let \(A\) be ample. Then \(A\) and \(A\otimes M\) have identical intersection numbers and are algebraically equivalent, while no positive tensor power of their difference is trivial. This verifies why the regularity assumption in the torsion conclusion is essential.

\section[Sistemi regolari e sovrabbondanti]{Sistemi lineari regolari e sovrabbondanti (pp.~147--150)}\label{enriques-iv-15}

\subsection{Model and cohomological conventions}

Throughout the modern reconstructions, \(X\) is a smooth connected projective surface over \(\mathbb C\). An effective divisor is a nonzero integral combination \(D=\sum m_i\Gamma_i\) of irreducible curves, with \(m_i>0\); it is allowed to be reducible or nonreduced. We write

\[
p_g=h^0(K_X),\qquad q=h^1(\mathcal O_X),\qquad
p_a(X)=\chi(\mathcal O_X)-1=p_g-q,
\]

\[
\begin{aligned}
n&=D^2, &\pi&=p_a(D)=1+\tfrac12(D^2+K_X D),\\
r&=h^0(D)-1, &i&=h^0(K_X-D).
\end{aligned}
\]

The surface Riemann--Roch formula becomes

\[
r=p_a(X)+n-\pi+1-i+h^1(D).
\tag{IV.15.1}
\]

Consequently Enriques's superabundance is precisely \(\omega=h^1(D)\). Empty complete systems have projective dimension \(-1\). Arithmetic genus, rather than the sum of the genera of the reduced components or the genus of a normalization, is required in all addition formulas. Where the classical system has prescribed multiple base points, one must first specify the smooth model and strict-transform system on which its degree and genus are calculated; an adjoint on that model need not be the naive divisor \(K_X+D\) on a model before resolving those base points.

The ``pure'' canonical system in §§18--20 is interpreted on a minimal smooth model. For a surface of general type this gives \(K_X\) nef and \(K_X^2>0\), and the historical absolute linear genus is \(p^{(1)}=K_X^2+1\). On a nonminimal model the actual canonical square is smaller. Purity does \textbf{not} mean that every canonical member is irreducible, reduced, or free of fixed components: §18 exists precisely to disprove that interpretation. An elliptic curve of square \(-1\) is not an exceptional curve of the first kind; that term refers to a smooth rational curve of square \(-1\).

The nefness used below can be justified directly in the situation of the source. Since \(p_g>0\), choose an effective canonical divisor. If an integral curve \(\Gamma\) satisfied \(K_X\Gamma<0\), it would occur in that divisor and would have \(\Gamma^2<0\). Adjunction then gives \(2p_a(\Gamma)-2=K_X\Gamma+\Gamma^2\le-2\). Both negative intersection numbers must equal \(-1\), and \(p_a(\Gamma)=0\), so \(\Gamma\) is a smooth rational \((-1)\)-curve. Removing such curves therefore makes \(K_X\) nef. With \(K_X^2>0\), Riemann--Roch gives quadratic pluricanonical growth, hence bigness and general type. Thus the modern positivity assumptions follow from the stated effective canonical class, absolute linear genus greater than one, and passage to the pure/minimal model; they are not additional unexplained conclusions.

\subsection{Riemann--Roch and the Picard argument}

The definitions at p.~147 follow directly from (IV.15.1). The main adjoint assertion is

\[
h^1(K_X+C)=0,\qquad \dim|K_X+C|=p_a(X)+p_a(C)-1.
\tag{IV.15.2}
\]

For effective \(C>0\), \(h^2(K_X+C)=h^0(-C)=0\), so the dimension follows from the vanishing. The general cohomological diagnostic, valid for every effective divisor, is

\[
h^1(K_X+C)
=h^0(\mathcal O_C)-1+
\dim\ker\bigl[H^1(\mathcal O_X)\longrightarrow H^1(\mathcal O_C)\bigr].
\tag{IV.15.3}
\]

Indeed, take cohomology of \(0\to\mathcal O_X(-C)\to\mathcal O_X\to\mathcal O_C\to0\), use \(H^0(-C)=0\), and apply Serre duality. This is recorded explicitly as Fact 1 in Mendes Lopes (1996, p.~2 of the author manuscript). Formula (IV.15.3) isolates two logically different issues: global functions on a possibly nonreduced curve, and injectivity of the irregularity restriction map.

Enriques states his version with \(\dim|C|\ge q\), but his proof on p.~148 starts with \(\dim|C|>q\). His auxiliary lemma is introduced, ``for simplicity,'' for both systems of dimension greater than one. Thus the displayed proof does not itself cover the equality case or all pencil cases of its statement. This is a difference between theorem and proof scope, not evidence that the theorem is false.

The lemma says that \(|C|\) restricts completely to a general \(L\in|C+D|\) when \(\dim|D|>q\). Its characteristic-series proof uses the earlier bound \(\delta\le q\), a projective dimension count to find a residual group \(G'\), and a pencil from which \(D\) can be detached. It suppresses the exact restriction maps and the assumptions excluding extra fixed intersections. The modern sequence is

\[
0\longrightarrow\mathcal O_X(-D)
\longrightarrow\mathcal O_X(C)
\longrightarrow\mathcal O_L(C)\longrightarrow0.
\tag{IV.15.4}
\]

Thus \(H^1(-D)=0\) is a transparent sufficient condition for complete restriction. When the moving system has a surface as image, its divisor is nef and big on the relevant smooth model; Ramanujam vanishing supplies this condition in characteristic zero. Reid (1996, Theorem 3.13) gives the precise effective nef-and-big statement \(H^1(-D)=0\).

For completeness, the exceptional numerical case of a moving irreducible system with \(C^2=0\) can be treated separately in the complex setting. Its complete irreducible pencil has rational base, since a moving degree-one linear system on its Stein-factor base curve forces that curve to be \(\mathbb P^1\). A holomorphic one-form restricting to zero on a general fiber descends to the base: on the smooth locus of the fibration it has no vertical part, and closedness forces its horizontal coefficient to be constant along fibers; extension across the finitely many exceptional fibers gives a form on the base. Since \(H^0(\mathbb P^1,\Omega^1)=0\), restriction of holomorphic one-forms is injective. Hodge conjugation then gives the required injection on \(H^1(\mathcal O)\). This supplies a modern proof route for the pencil boundary, rather than pretending that the strict inequality in the printed proof covers it. This argument is a modern reconstruction using standard complex Hodge theory, not an identification of Picard's original proof.

The numerical residual calculation on p.~148 is correct. Put \(L=C+D+K_X\), \(\rho=p_a(D+K_X)\), \(m=(D+K_X)^2\), and \(s=C(D+K_X)\). Then

\[
p_a(L)=\pi+\rho+s-1,\qquad
\deg((D+K_X)|_L)=m+s.
\]

The two restricted systems are residual because

\[
(K_X+C)+(K_X+D)=K_X+L.
\]

When their restrictions are complete, curve Riemann--Roch gives

\[
\dim|(K_X+C)|_L|
=p_a+m-\rho+1-\{m+s-p_a(L)\}-1
=p_a+\pi-1.
\]

The kernels \(H^0(-C)\) and \(H^0(-D)\) vanish, so the restricted and surface projective dimensions agree. These are the missing cohomological justifications for the dimension identifications.

\subsection{High multiples and the exact positivity needed}

At pp.~148--149 Enriques proves that high multiples of an irreducible system of dimension at least two are regular when there are no fundamental, that is contracted, curves. On a base point free model, such a divisor \(C\) is nef, has \(C^2>0\), and satisfies \(C\Gamma>0\) for every irreducible curve. Nakai--Moishezon makes it ample; Serre vanishing then gives \(H^1(sC)=0\) for \(s\gg0\). Mere bigness does not give this conclusion.

The source chooses \(D=hC-K_X\ge0\), starts from the regular \(K_X+sC\), and adds the finitely many irreducible components of \(D\), counting repetitions. If \(L^2=\nu\), \(p_a(L)=\rho\), and \(C L=t>0\), the correct identity is

\[
(K_X+sC)L=st+2\rho-2-\nu.
\tag{IV.15.5}
\]

The scan at p.~149 prints \(r\) in place of \(\nu\); this is a source misprint, not OCR. If \(L\) is added again, its self-intersection contributes once more. In the printed negative case \(L^2=-\nu'\), the next intersection is \(st+2\rho-2\), as stated.

This argument invokes the earlier addition criterion of §8. A defective weak numerical criterion in that earlier discussion must not be silently reused here. Fortunately \(s\gg0\) supplies more than the displayed weak inequality: at every one of the finitely many additions, the restriction to the integral component \(L\) has arbitrarily large degree, so its residual to \(\omega_L\) has negative degree and its \(H^1\) vanishes. Explicitly, for any fixed partial sum \(E\le D\),

\[
\deg\bigl((K_X+sC+E+L)|_L\bigr)
=st+(K_X+E+L)L\longrightarrow+\infty.
\]

The finitely many thresholds can be met simultaneously. The restriction exact sequence then preserves \(H^1=0\) at each addition. Thus this conclusion survives independently of the insufficient general addition criterion; it is the unbounded \(st\) term that repairs the proof.

\subsection{Positive-genus contracted curves and a sign slip}

Let a base point free \(C\) contract an integral curve \(\theta\) of genus \(\rho>0\), with \(\theta^2=-\nu<0\). Because the associated morphism contracts \(\theta\), \(\mathcal O_X(sC)|_\theta\simeq\mathcal O_\theta\), and a section not vanishing at the image point shows that restriction has rank one. Consequently

\[
r(sC-\theta)=r(sC)-1.
\]

The weaker numerical assertion \(sC\cdot\theta=0\) alone would not establish triviality or rank one. Adjunction gives

\[
(sC-\theta)^2=(sC)^2-\nu,
\qquad
p_a(sC-\theta)=p_a(sC)-\nu-\rho+1.
\]

For large \(s\), the speciality indices vanish, and the virtual dimension therefore \textbf{increases} by \(\rho-1\) upon removing \(\theta\). The printed p.~150 says that it decreases. The final formula is nevertheless correct:

\[
\omega(sC)=\omega(sC-\theta)+\rho\ge\rho.
\tag{IV.15.6}
\]

The extra term can come from the infinitesimal neighborhoods of the exceptional curve; its reduced genus need not exhaust the defect.

\subsection{Completing the omitted ordinary-multiple-point calculation}

The final example is uniform only with an ordinary isolated hypersurface point. For the displayed total global defect, assume that this point is the only singularity contributing higher direct-image cohomology of the resolution (for example, the projective hypersurface is otherwise smooth). Contributions from other nonrational singularities would have to be added. Assume the degree-\(i\) tangent cone is a smooth plane curve, \(i>3\), and the blowup resolves the point. Its exceptional curve \(E\) then satisfies

\[
g(E)=\frac{(i-1)(i-2)}2,\qquad
\mathcal O_E(-E)=\mathcal O_E(1),\qquad E^2=-i.
\]

Let \(H\) be the pullback of a hyperplane section and \(D_j=sH-jE\). The exact sequences

\[
0\to\mathcal O_X(D_{j+1})\to\mathcal O_X(D_j)
\to\mathcal O_E(j)\to0
\]

have surjective maps on \(H^0\) for \(s\gg0\) and fixed \(j\le i-3\): ambient homogeneous degree-\(j\) initial forms supply every section of \(\mathcal O_E(j)\), and high-degree global sections interpolate the finitely many required jets. Curve duality gives

\[
h^1(\mathcal O_E(j))
=h^0(\mathcal O_E(i-3-j))
=\binom{i-1-j}{2}.
\]

For \(j\ge i-2\), all the subsequent exceptional-neighborhood graded pieces have \(H^1=0\); formal functions and high-degree Serre vanishing give

\[
h^1(D_{i-2})=0.
\]

The recurrence therefore yields the omitted answer

\[
\boxed{h^1(D_j)=\binom{i-j}{3}\quad(0\le j\le i-2),\qquad
\omega(sH)=\binom i3.}
\tag{IV.15.7}
\]

This is a worked calculation under the stated ordinary-point and global-positivity hypotheses. It is not asserted for arbitrary isolated singularities of multiplicity \(i\). For \(i=4\), the exceptional genus is 3 while the defect is 4, concretely explaining the strict inequality in (IV.15.6).

\section[Regolarità degli aggiunti]{Precisazioni sulla regolarità del sistema aggiunto ad una curva sopra una superficie regolare (pp.~150--152)}\label{enriques-iv-16}

\subsection{The reduced argument is sound}

In §16, \(q=0\) and \(p_g=p_a=p>0\). Every system containing the canonical system has an effective residual \(L\), hence is \(|K_X+L|\). For integral \(L\), the adjunction sequence

\[
0\to\mathcal O_X(K_X)\to\mathcal O_X(K_X+L)
\to\omega_L\to0
\]

is surjective on global sections because \(h^1(K_X)=q=0\). Since \(h^0(\mathcal O_L)=1\), duality on the Gorenstein curve \(L\) gives \(h^0(\omega_L)=p_a(L)\), and \(\dim|K_X+L|=p+p_a(L)-1\). Smoothness of the integral curve is unnecessary if dualizing sheaves and arithmetic genus are used.

For a reduced connected divisor, the intersection graph admits an ordering with every initial union connected. Grow a connected set of vertices one vertex at a time; a proper connected set has an edge to its complement. This proves the lemma and corollary at p.~151 without relying on any self-intersection sign.

At each step write \(L=A+B\), where \(B\) is the newly added \textbf{integral} component, \(A\) is the previous connected union, and \(j=AB\ge1\). Then

\[
p_a(L)=p_a(A)+p_a(B)+j-1,
\]

\[
(K_X+L)|_B=\omega_B(A),\quad
\deg\omega_B(A)=2p_a(B)-2+j.
\]

The dual \(H^0(\mathcal O_B(-A))\) vanishes because it has negative degree on the integral curve. This proves nonspeciality and the upper bound \(r_B\le p_a(B)-2+j\) for the image series. Comparing with the surface Riemann--Roch lower bound forces equality, proves regularity, and completes the induction. One must not replace integral \(B\) by an arbitrary connected reducible curve and infer nonspeciality from its total degree alone: degree can be distributed unevenly among its components.

For a reduced curve with \(s\) connected components, (IV.15.3) gives exactly

\[
h^1(K_X+L)=s-1.
\tag{IV.16.1}
\]

The two-disjoint-components calculation at p.~152 is correct: \(p_a(L)=\pi_1+\pi_2-1\), \(\dim|K_X+L|=p+p_a(L)\), and \(\omega=1\). The last sentence of §16 says that analogous exceptional-multiple calculations are easy, but gives no formula. For nonreduced cycles the answer is controlled by \(h^0(\mathcal O_L)\), not merely the count of connected supports.

\section[Aggiunti di curve nonridotte]{Segue: il sistema aggiunto ad una curva connessa formata di componenti multiple sopra una superficie regolare (pp.~152--156)}\label{enriques-iv-17}

\subsection{Numerical 1-connectedness is the correct nonreduced definition}

Enriques's ``arithmetically connected'' means

\[
D=A+B,\quad A,B>0\quad\Longrightarrow\quad AB>0.
\tag{IV.17.1}
\]

Every effective integral subcycle counts; the two pieces may share underlying components. This is modern numerical 1-connectedness. For reduced divisors it agrees with ordinary connectedness, because distinct effective components have nonnegative intersection. For nonreduced divisors, negative self-intersections change the condition.

A simple diagnostic shows why connected support is inadequate. If \(E\simeq\mathbb P^1\) is an exceptional curve with \(E^2=-1\), then \(2E\) has connected support, but

\[
0\to\mathcal O_E(1)\to\mathcal O_{2E}\to\mathcal O_E\to0
\]

gives \(h^0(\mathcal O_{2E})=3\). It is not 1-connected because the split \(E+E\) has intersection \(-1\). These extra sections are nilpotents; no topological path argument can exclude them. For a fully numerical example, blow up a K3 surface at one point. Then \(K_X=E\) and \(q=0\), so \(h^1(K_X+2E)=h^0(\mathcal O_{2E})-1=2\).

The source cites Franchetta for failures of the naive component-addition rule, without displaying his examples. The following worked example makes the phenomenon explicit; it is \textbf{not identified with Franchetta's historical example}. In \(\mathbb P^2\), choose a smooth conic \(A_0\) and a transverse line \(B_0\). Blow up five points of \(A_0\setminus B_0\) and four points of \(B_0\setminus A_0\). Their strict transforms satisfy

\[
A^2=-1,\quad B^2=-3,\quad AB=2.
\]

The cycle \(L=2A+B\) is 1-connected: up to exchanging the two sides, its nontrivial splits have intersections \(A(A+B)=1\) and \(2AB=4\). Moreover \(LB=1>0\). Nevertheless \(L+B=2(A+B)\) is not 1-connected, since the split \((A+B)+(A+B)\) has intersection 0. This directly verifies the phenomenon described at p.~153. The listed splits exhaust its proper subcycles up to exchanging the two sides.

\subsection{Franchetta's two lemmas, as actually proved in the text}

Theorem I is a useful finite descent. If \(D=A+B\) with \(AB\le0\), and \(A\) is not 1-connected, write \(A=A_1+A_2\) with \(A_1A_2\le0\). At least one \(A_iB\le0\), since their sum is \(AB\le0\). For that \(A_i\), its intersection with the entire residual remains nonpositive. Repeating decreases the sum of the coefficients, so the process terminates with a 1-connected subcycle having nonpositive boundary. This proves the contrapositive of the stated criterion. The missing verb on p.~154 is a grammatical source misprint, not a mathematical obstruction.

Theorem II requires \(D\) to have a nontrivial decomposition. Choose a proper 1-connected subcycle \(B\) minimizing its positive boundary \(B(D-B)\). If \(A=D-B\) is not 1-connected, Theorem I provides a 1-connected \(A_1\le A\) with \(A=A_1+A_2\), \(A_1A_2\le0\). Connectedness of \(D\) implies \(BA_1,BA_2>0\). Hence

\[
A_1(B+A_2)\le A_1B<AB,
\]

contradicting minimality. Thus both pieces are 1-connected. The wording must exclude an irreducible reduced divisor, which admits no proper nonzero effective split; this is a harmless scope qualification, not a counterexample to the intended statement.

\subsection{The p.~155 working hypothesis must remain visible}

The geometric proof on pp.~154--155 takes a very positive regular \(C\), sets \(D=C+L\), and aims to prove that \(|C|\) cuts a complete series on \(D\). Restrictions of the fixed \(L\) and \(K_X+L\) to a sufficiently positive \(D\) can be justified by Serre vanishing and duality. The remaining obstruction is specifically \(H^1(-L)\), which does not disappear merely because \(C\) is large.

Enriques constructs a pencil with a prescribed characteristic group and asserts that imposing an additional point on \(L\) forces the entire \(L\), with every multiplicity, to detach. He then explicitly says that he lacks an exhaustive proof and introduces this as a working hypothesis, to which the following result and its consequences are subordinate.\footnote{The scan at p. 155 reads: ``Tale proprietà ci pare vera, ma non ne possediamo una dimostrazione esauriente. Noi la introdurremo come ipotesi di lavoro, restando così inteso che il risultato precedente e le conseguenze che ne trarremo, sono subordinati alla validità di tale ipotesi.''} The result is therefore conditional in the printed argument. Support containment of a nonreduced divisor is insufficient to prove containment of its full scheme structure.

Conditional on complete restriction, the calculation itself is correct. With \(n=C^2\), \(\rho=p_a(C)\), \(m=CL\), and \(\pi=p_a(L)\), the series \(C|_D\) has degree \(n+m\), projective dimension \(p+n-\rho+1\), and ambient curve genus \(\pi+\rho+m-1\). Its speciality index is \(p+\pi\), so its canonical residual has dimension \(p+\pi-1\). The scan prints \((CL')\) for this residual; it should be \((DL')\), since it lies on \(D\).

\subsection{A complete modern proof replacing the hypothesis}

The exact needed theorem is that numerical 1-connectedness implies

\[
H^0(\mathcal O_L)=\mathbb C.
\]

Ramanujam (1972, Lemma 3, reprinted at p.~127 of \emph{C. P. Ramanujam---A tribute}) proves this for effective divisors; Reid (1996, Lemma 3.11(ii)) gives a particularly transparent induction. These are concrete repair references, without an unverified claim about the earliest historical repair.

Here is the short induction in the present notation. Start with a reduced integral component \(A\le L\), so \(h^0(\mathcal O_A)=1\). If \(A<L\), numerical connectedness of the \textbf{whole} \(L\) gives \(A(L-A)>0\), so some integral component \(\Gamma\le L-A\) satisfies \(A\Gamma>0\). The divisor restriction sequence

\[
0\to\mathcal O_\Gamma(-A)\to\mathcal O_{A+\Gamma}
\to\mathcal O_A\to0
\]

has \(H^0(\mathcal O_\Gamma(-A))=0\). Therefore \(h^0(\mathcal O_{A+\Gamma})\le1\), while constants supply a nonzero section. Induct until \(A=L\), counting repeated components. Crucially, the induction preserves \(h^0=1\); it does \textbf{not} claim that every partial sum is 1-connected. This is precisely why Franchetta's obstruction to the naive geometric composition lemma does not invalidate it.

Now \(q=0\), and

\[
0\to\mathcal O_X(-L)\to\mathcal O_X\to\mathcal O_L\to0
\]

gives \(H^1(-L)=0\), because the map on constants is an isomorphism. Duality yields \(H^1(K_X+L)=0\). The assumption \(p_g>0\) is not needed for this adjoint-vanishing argument; \(q=0\) and numerical connectedness suffice. This proves the theorem at p.~156 for all effective 1-connected divisors, including arbitrary negative multiple components, and fills the exact restriction obstruction in the preceding proof. It does not assert that every informal pencil formulation of the detachment principle is valid without its additional base-locus hypotheses.

\section[Riducibilità del sistema canonico]{Sulla riducibilità del sistema canonico\\(pp.~156--158)}\label{enriques-iv-18}

\subsection{What the counterexamples refute}

The alternatives attributed to Noether are exceptional rational curves of square \(-1\), or a pure canonical system of square zero composed with elliptic curves. They are not an exhaustive classification. A minimal general-type surface can have fixed rational \((-2)\)-curves or fixed positive-genus curves. The three examples in the text establish this once their branch curves, resolutions, and minimal models are constructed. The source abbreviates all three constructions and gives only a partial constant-count argument for the second.

The following \textbf{worked reconstructions} verify existence, exact singularity type, and every stated invariant. No claim is made that these are historically new surfaces or the original constructions intended by Enriques. The generic choices and their local consequences are specified below.

\subsection{Interpolation at proper and infinitely near points}

Fix \(r=3\) or 4 distinct points \(p_1,\ldots,p_r\) on the line \(a=(z=0)\subset\mathbb P^2\), with a prescribed tangent direction at each point transverse to \(a\). Blow up each point and then the prescribed point on its first exceptional curve. In the orthogonal total-transform basis write \(H,E_i,E_i'\), with each exceptional basis vector of square \(-1\). Put

\[
A=rH-\sum_{i=1}^r(E_i+E_i').
\tag{IV.18.1}
\]

The system \(|A|\) is base point free. To see every base point explicitly, use an affine coordinate \(x\) on \(a\), write \(p_i=(t_i,0)\), and let the prescribed tangent be \(x-t_i=\alpha_i z\). Choose a homogeneous degree-\(r\) form \(F_0\) whose restriction to \(a\) is \(\prod_i(x-t_i y)\). Its coefficient of \(z\), a degree-\((r-1)\) binary form, can be chosen to have the \(r\) values required to impose the prescribed tangents. These values determine it by interpolation. The system contains

\[
cF_0+z^2R_{r-2}(x,y,z).
\tag{IV.18.2}
\]

Away from \(a\), \(z^r\) is an available nonvanishing section. On \(a\setminus\{p_i\}\), \(F_0\ne 0\). Near \(p_i\), set \(u=x-t_i-\alpha_i z\). The two kinds of sections have leading terms \(u+O(2)\) and \(z^2 R(p_i)+O(3)\), and \(R(p_i)\) can be nonzero. Their common local ideal is the curvilinear length-two ideal \((u+O(2),z^2)\). After the first blowup, using \(u=zt\), and dividing by the prescribed exceptional factor, the leading sections are \(t+O(z)\) and \(z R(p_i)+O(z^2)\). They have exactly one simple base point, the prescribed infinitely near point, and their differentials span its cotangent space. Blowing it up resolves the base ideal; the restrictions to the new exceptional \(E_i'\) span \(H^0(\mathbb P^1,\mathcal O(1))\). There are no other base points on either exceptional curve. This proves the claimed global generation, not merely an expected dimension count.

Also \(A^2=r^2-2 r>0\), so \(A\) is nef and big. The classes \(3 A+H\) for \(r=3\), and \(3A\) for \(r=4\), are therefore free and big. Their general members are smooth irreducible curves on the blown-up surface by Bertini. They avoid the strict transforms \(E_i-E_i'\), on which they have degree zero, and meet each \(E_i'\) transversely in three distinct points. The degree-three restriction is general: products of the full degree-one restrictions of \(|A|\) span all binary cubics on \(E_i'\); multiplying by a section of \(H\) nonzero at \(p_i\) does not change this in the degree-ten case.

Choose these three intersection points away from \(E_i'\cap(E_i-E_i')\). Blowing down the second exceptional curve gives three smooth branches with distinct tangent directions at the first infinitely near point, none tangent to the old exceptional curve. Blowing down the first exceptional curve gives three smooth branches with one common tangent and distinct quadratic terms

\[
u=\lambda_jz^2+O(z^3),\qquad \lambda_1,\lambda_2,\lambda_3\ \text{distinct}.
\]

Thus the downstairs singularity is exactly the ordinary \([3,3]\) cluster intended here: multiplicity 3, followed by an ordinary triple point after one blowup, with no further infinitely near singularity. Smoothness of the strict transform and transversality to the exceptional curves exclude any additional singularities. This is stronger than merely obtaining a nonempty linear system with virtual assigned multiplicities.

\subsection{The double-cover calculation and the exceptional contractions}

For a smooth double cover \(\pi:S\to W\) with branch \(2L\),

\[
\pi_*\mathcal O_S=\mathcal O_W\oplus L^{-1},\qquad
K_S=\pi^*(K_W+L),
\]

\[
\chi(\mathcal O_S)=2\chi(\mathcal O_W)+\tfrac12L(L+K_W),\qquad
K_S^2=2(K_W+L)^2.
\tag{IV.18.3}
\]

The decomposition and Hurwitz formula are the double-cover specialization of the abelian-cover framework described by Pardini (2010, PDF pages 4 and 10); the Euler-characteristic identity follows immediately by applying surface Riemann--Roch to \(L^{-1}\).

At a branch point of multiplicity \(m\), blow up and put \(b=\lfloor m/2\rfloor\). The new half-branch class is \(L'=f^*L-b E\). Direct substitution in (IV.18.3) gives

\[
\Delta\chi=-\tfrac12b(b-1),\qquad
\Delta K_X^2=-2(b-1)^2.
\tag{IV.18.4}
\]

An ordinary \([3,3]\) point has canonical-resolution branch multiplicities 3 and 4, not 3 and 3: the first odd exceptional curve belongs to the branch and passes through the second triple point. Thus \(b=1,2\), and the resolution decreases \(\chi\) by 1 and \(K_X^2\) by 2. After the second blowup the first exceptional curve has square \(-2\), belongs to the branch, and has smooth rational ramified preimage \(R\) with \(R^2=-1\). Contracting \(R\) raises \(K_X^2\) by 1. Therefore each such cluster contributes a net decrease of 1 to \(K_X^2\) on the minimal model. Failing to make this final contraction gives the wrong numbers in all three examples.

The second exceptional curve is not a branch component and meets the branch in four distinct points. Its preimage is elliptic of square \(-2\) before contracting \(R\), and square \(-1\) afterwards. This accounts for the elliptic negative curves in the complete inverse images of the line, while distinguishing them from the rational exceptional curves that were contracted.

\subsection{\texorpdfstring{Example 1: degree \(2n\), one ordinary \((2n-6)\)-fold point, and two \([3,3]\) points}{Example 1: degree 2n, one ordinary (2n-6)-fold point, and two {[}3,3{]} points}}

Take distinct collinear \(O,p_1,p_2\), blow up \(O\) once and the two prescribed transverse length-two clusters at \(p_1,p_2\). Set

\[
A=3H-E_O-\sum_{i=1}^2(E_i+E_i'),\qquad F=H-E_O.
\]

Here \(A^2=4\) and \(|F|\) is the pencil of lines through \(O\). The interpolation proof above works with two tangent conditions and one ordinary point condition. On \(a\), a cubic restricts to the product of the three point factors; the coefficient of \(z\) has two prescribed values at \(p_1,p_2\) and a free value at \(O\). Consequently the two independent first-order directions at \(O\) are both available, so a single blowup resolves its base point. The local argument at each \(p_i\) is unchanged. Thus \(A\) is free, as is \(F\).

For \(n\ge 5\), take a general member of

\[
D_n=3A+(2n-9)F
=2nH-(2n-6)E_O-3\sum_{i=1}^2(E_i+E_i').
\tag{IV.18.5}
\]

This system is free and big. Its restriction to \(E_O\) includes the full space of binary forms of degree \(2n-6\), because the restrictions of both \(A\) and \(F\) span degree-one forms and their products span all the required degree. A general member meets \(E_O\) transversely in \(2n-6\) distinct points. It therefore has an ordinary \((2n-6)\)-fold point downstairs, together with exactly the two ordinary \([3,3]\) points already analyzed and no other singularities. The strict transform of \(a\) has intersection zero with \(D_n\), so a general branch does not contain \(a\).

For a smooth degree-\(2n\) double plane, \(\chi=1+\binom{n-1}{2}\) and \(K_X^2=2(n-3)^2\). The ordinary point has \(b=n-3\); the two clusters have the corrections just computed. Hence

\[
\chi=1+\binom{n-1}{2}-\binom{n-3}{2}-2=2n-6,
\]

\[
K_{\min}^2=2(n-3)^2-2(n-4)^2-2=4n-16.
\]

Canonical plane forms have degree \(n-3\), vanish to order \(n-4\) at \(O\), and pass through \(p_1,p_2\). Their assigned intersection with \(a\) exceeds their degree, so they contain \(a\). The residual forms have degree \(n-4\) and multiplicity \(n-5\) at \(O\), and their vector-space dimension is \(2n-7\). Therefore

\[
p_g=p_a=2n-7,\qquad q=0,\qquad p^{(1)}=4n-15.
\]

The strict line on the cluster surface has square \(-2\); its connected double cover is branched only at the two first exceptional branch curves, so it is rational. Its square on the resolved cover is \(-4\), which becomes \(-2\) after the two rational contractions. Thus the fixed canonical curve \(\theta\) has genus 0, square \(-2\), and

\[
\theta(K_X-\theta)=2p_a(\theta)-2-2\theta^2=2.
\]

The residual moving class is the pullback of \((n-4)H-(n-5)E_O\), which is free and nef. It has intersection 2 with \(\theta\). Consequently \(K_X\theta=0\), and \(K_X\Gamma\ge 0\) for every other irreducible curve \(\Gamma\), since it is a sum of that nef moving class and the effective \(\theta\). This verifies minimality rather than presuming that the contracted rational curves were the only possible \((-1)\)-curves.

\subsection{\texorpdfstring{Example 2: degree 10 and three collinear \([3,3]\) points}{Example 2: degree 10 and three collinear {[}3,3{]} points}}

Use \(D=3 A+H\) with \(r=3\). This is the promised free and big system giving an irreducible branch curve with precisely the required singularities. The constant count can also be completed. On this cluster surface \(K_W=-A\), so \(D-K_W=4 A+H\) is nef and big, and characteristic-zero vanishing gives \(H^1(D)=0\). Since \(D^2=46\) and \(DK_W=-12\), Riemann--Roch gives \(h^0(D)=30\), hence projective dimension 29, as printed.

For the subsystem consisting of twice \(a\) and an octic with three assigned tacnodes, the octic class is

\[
T=8H-2\sum(E_i+E_i')=2A+2H.
\]

Now \(T-K_W=3 A+2 H\) is nef and big, \(T^2=40\), and \(TK_W=-12\), so \(h^0(T)=27\), giving projective dimension 26. In subtracting \(2a\) from the degree-ten cluster class one must unload the old exceptional curves: if \(\widetilde a=H-\sum E_i\), then

\[
D-2\widetilde a=T+\sum(E_i-E_i').
\]

These fixed old exceptional components explain why twice \(a\) plus a tacnodal octic satisfies the virtual \([3,3]\) conditions. Comparing 29 with 26 alone does not rule out all branches containing \(a\) once; the free-system construction resolves that logical omission.

The cover has \(\chi=7-3=4\), \(K_X^2=8-3=5\), and canonical images equal to conics through the three collinear points, namely \(a\) plus arbitrary lines. Thus \(p_g=3\), \(q=0\), and \(p^{(1)}=6\). The line cover has four branch points---three from the first exceptional branch curves and one residual branch intersection---so its fixed component \(\theta\) is elliptic. Its resolved square is \(2(1-3)=-4\), increasing by 3 to \(-1\). Therefore \(\theta(K_X-\theta)=2\), exactly as stated.

The complete inverse image of \(a\) consists of \(\theta\) and three elliptic curves \(E_i^*\), each of square \(-1\), each meeting \(\theta\) once, and pairwise disjoint. Hence

\[
p_a(\theta+\sum E_i^*)=4,\qquad
(\theta+\sum E_i^*)^2=-1-3+6=2.
\]

The residual canonical system is the pullback of lines, so it is nef; \(K_X\theta=-1+2=1>0\), and \(K_X\) is nonnegative on all other curves. This proves that the resulting surface is minimal.

\subsection{\texorpdfstring{Example 3: degree 12 and four collinear \([3,3]\) points}{Example 3: degree 12 and four collinear {[}3,3{]} points}}

Take \(D=3 A\) with \(r=4\). The construction and local verification are identical. Canonical cubics through four collinear points contain \(a\); the residual conics have vector-space dimension 6. The source does not state the surface invariants here, but the same verified computation gives

\[
\chi=11-4=7,\quad p_g=p_a=6,\quad q=0,\quad
K_X^2=18-4=14,\quad p^{(1)}=15.
\tag{IV.18.6}
\]

The line cover is elliptic with four branch points. Its square is initially \(2(1-4)=-6\), and becomes \(-2\) after four contractions. The residual canonical conic system meets it in 4, giving \(\theta(K_X-\theta)=4\). The complete preimage of \(a\) is a star of five elliptic curves, with four edges; its arithmetic genus is 5 and its square is \(-2-4+8=2\). Nefness and minimality follow as in Example 2, now with \(K_X\theta=-2+4=2\).

\subsection{The irregular example in the footnote is also verifiable}

Choose a smooth conic \(Q=0\) and a quartic \(F=0\) transverse to it at eight distinct points, and take three distinct general parameters \(\lambda_j\). The degree-twelve branch

\[
B=\prod_{j=1}^3(F-\lambda_jQ^2)=0
\tag{IV.18.7}
\]

is the union of three smooth quartics in the pencil generated by \(F\) and \(Q^2\). They meet only at the eight base points, with common tangent and pairwise contact order 2. Thus the total branch has eight ordinary \([3,3]\) clusters and no further singularities. Canonical-resolution corrections give \(\chi=11-8=3\) and \(K_X^2=18-8=10\). Canonical cubics through the eight points contain \(Q\), since \(3\cdot 2<8\), and are \(Q\) times arbitrary linear forms. Therefore \(p_g=3\), \(q=1\), \(p_a=2\), and \(p^{(1)}=11\), all as printed.

The conic cover branches at eight points and has genus 3. Its square changes from \(-8\) on the resolved cover to 0 after the eight rational contractions. If \(H\) is the moving pullback of a line, \(K_X=\theta+H\), \(\theta H=4\), \(H^2=2\), and \(K_X^2=10\), confirming both the fixed curve and minimality.

The irrational base can be seen without invoking an unidentified theorem. In the function field of the double plane \(w^2=\prod(F-\lambda_j Q^2)\), put

\[
t=F/Q^2,\qquad u=w/Q^3.
\]

Then \(u^2=\prod(t-\lambda_j)\) defines a genus-one double cover of \(\mathbb P^1\), branched at the three \(\lambda_j\) and infinity. The generic quartic \(F-tQ^2\) is geometrically integral: general members of this pencil are smooth irreducible plane quartics. It remains geometrically integral after the extension of \(\mathbb C(t)\) by \(u^2=\prod(t-\lambda_j)\). Consequently the displayed elliptic function field is algebraically closed in the surface function field, so this is the full Stein-factor base, not merely a curve to which the surface maps. The text's reference to de Franchis is compatible with this calculation, but the exact historical theorem and original citation are not supplied in these pages and have not been independently identified here.

\subsection{Canonical pencils and examples in the final edition}

The main-text claim that examples were not known at p.~158 is superseded by its own footnote naming Pompilj's 1948 examples. Burniat (1952, p.~1030) explicitly describes those regular double planes with canonical genus-two pencils and expands the range of examples. The commentary therefore does not report this as a present open problem or even an unresolved question in the final version of the book.

The numerical pencil bounds are correct on a minimal model. Write \(K_X=Z+sk\), where \(Z\) is fixed, \(p_a(k)=\pi\), and \(k^2=n\ge 0\). Intersecting with \(K_X\) gives

\[
K_X^2=s(2\pi-2-n)+\sum m_i(2p_a(\theta_i)-2-\theta_i^2).
\]

The fixed contributions are \(m_i K_X\theta_i\ge 0\). Moreover \(k(K_X-k)=2\pi-2-2 n\ge 0\). Thus \(n\le\pi-1\) and \(\pi\le K_X^2+1=p^{(1)}\), so \(n\le K_X^2\). The notation \(p^{(2)}=p^{(1)}-1\) on this page denotes canonical degree; it must not be confused with the bigenus \(P_2\).

\subsection{Boundaries of the verified examples}

The branch constructions above permit a finite further investigation: vary the prescribed infinitely near tangent directions and record the boundary cases in which directions coalesce, a line enters the branch, or additional singularities appear. Each branch multiplicity, discrepancy, and subsequent contraction must then be recomputed. This is a proposed extension of the verified generic calculations, not an assertion that their formulas remain valid at every boundary.

\section[Regolarità del sistema bicanonico]{Sulla regolarità del sistema bicanonico (pp.~158--162)}\label{enriques-iv-19}

\subsection{The exact modern theorem and a noncircular proof}

Let \(X\) be a smooth minimal complex surface of general type with \(q=0\) and \(p_g=p>0\). Then every effective canonical divisor is numerically 2-connected, every effective positive multiple of \(K_X\) is numerically 1-connected, and

\[
H^1(mK_X)=H^2(mK_X)=0\qquad(m\ge2).
\tag{IV.19.1}
\]

Therefore

\[
\boxed{P_m=h^0(mK_X)=1+p+\binom m2K_X^2
=1+p+\binom m2\bigl(p^{(1)}-1\bigr)\quad(m\ge2).}
\tag{IV.19.2}
\]

The separate first plurigenus is \(P_1=p\). In particular \(P_2=p+p^{(1)}\), as the source correctly states.

Here is an intersection-theoretic proof of the connectedness step, which avoids first having to prove irreducibility of the bicanonical moving system. An effective nef divisor \(D\) with \(D^2>0\) is 1-connected; this is Reid's Lemma 3.11(i), also Mendes Lopes's Lemma 2.6. Suppose \(D=A+B\) and \(AB\le 0\). Put \(x=DA\ge 0\), \(y=DB\ge 0\), so \(x+y=D^2>0\). If \(x=0\), Hodge index implies \(A^2<0\), because a nonzero effective \(A\) cannot be numerically zero; but \(A^2=x-AB\ge 0\), a contradiction. Similarly \(y>0\). Hodge index now gives

\[
x^2\ge D^2A^2=(x+y)(x-AB)\ge(x+y)x,
\]

contradicting \(xy>0\). Thus \(AB>0\).

Apply this to \(D=m K_X\). For a canonical split \(K_X=A+B\), adjunction additionally gives

\[
AB=A(K_X-A)=2p_a(A)-2-2A^2,
\]

an even positive integer, hence at least 2. This argument is valid for nonreduced subcycles as well as reduced components. By §17, the divisor \((m-1)K_X\) has only constant global functions. Since \(q=0\), it follows that \(H^1(m K_X)=0\). Finally \(H^2(m K_X)=H^0((1-m)K_X)=0\), since a negative multiple of a nef big canonical class cannot be effective. Surface Riemann--Roch proves (IV.19.2).

In the complex setting the more general plurigenus identity is \(P_m=\chi(\mathcal O_X)+\binom m 2 K_X^2\) on every smooth minimal general-type surface for \(m\ge 2\), without the restriction \(q=0\): apply nef-and-big vanishing to \((m-1)K_X\). The elementary effective-divisor proof above uses \(q=0\) and \(p_g>0\); it must not be advertised as a proof of the unrestricted statement. Conversely, its cohomological connectedness part is algebraic and does not intrinsically require characteristic zero once those hypotheses hold. No blanket claim about arbitrary positive-characteristic surfaces is being made.

\subsection{The p.~159 formula is a source error, not an OCR error}

The printed general formula is

\[
P_i=p+\frac{i(i-1)}2p^{(1)}.
\]

The scan of the print confirms every term of this display. It is not corrected by reinterpreting \(p^{(1)}\), because the immediately preceding bigenus formula and the chapter's canonical-degree convention fix \(p^{(1)}=K_X^2+1\). For \(m\ge 3\), the printed expression exceeds the correct answer by \(\binom m 2-1\). For example, the verified degree-ten double plane of §18 has \(p=3\), \(K_X^2=5\), and \(p^{(1)}=6\). Its \(P_2=9\) agrees with the book, whereas \(P_3=19\); the printed general formula gives 21. This independent internal example makes the error concrete.

\subsection{Diagnosing the two bicanonical-pencil arguments}

The source first treats \(p\ge 2\). If a canonical moving pencil has general member \(C\) of genus \(\pi>1\), the adjunction restriction is complete because \(q=0\), and the system \(|K_X+C|\) has variable canonical intersections with \(C\) of degree \(2\pi-2>0\). This is the right geometric idea for ruling out composition of the bicanonical system with that pencil. The passage compresses the standard argument excluding every possible pencil; the independent proof of §20 below avoids this compression. The exclusion of an elliptic pencil is deferred in the source to later classification. On a minimal general-type model it follows immediately: an elliptic fiber \(F\) has \(F^2=K_X F=0\), incompatible with Hodge index and \(K_X^2>0\).

For \(p=1\), the assertion that the unique \(K_X\) decomposes into a fixed part and \(s\) half-fibers is not established. The next dimension estimate also uses the full-fiber genus for a half-fiber. If a resolved fiber is \(F=2 A\), \(F^2=0\), and \(p_a(F)=\pi\), then

\[
A^2=0,\qquad K_X A=\pi-1,\qquad
p_a(A)=\frac{\pi+1}{2},
\]

which is not \(\pi\). This arithmetic diagnosis does not, by itself, validate the earlier half-fiber decomposition. The general theorem is true, but this printed proof is incomplete.

The long footnote supplies a second approach. Its intersection identity

\[
2K_X^2=\sum m_i(2\rho_i-2-\nu_i)
+s(2\pi-2-n)
\]

is correct; nefness makes the fixed terms nonnegative. For a nonspecial complete pencil \(C\), Riemann--Roch gives \(p+n-\pi+1\le 1\), so \(n\le\pi-p\le\pi-1\). The scan prints \(n<\pi-1\), an unjustified strict inequality. Replacing it by \(\le\) preserves the intended contradiction for \(\pi>2\): one gets \(s(\pi-1)\le 2 K_X^2\), hence \(s\le K_X^2\), while \(\dim|2 K_X|\ge p+K_X^2>s\).

When the pencil is special and \(C\) is contained in \(K_X\), the footnote asserts the same weak bound by using \(|2C|\subset|K_X+C|\), but does not fully write the dimension argument. The independent connectedness and dualizing-sheaf proof below supplies the final theorem without requiring that compressed step.

The remaining genus-two case imports a double-plane normal form with specified elementary singularities. That classification is not established in these pages. More seriously, the existence of one doubled canonical member does not by itself show that the fixed plane part of the whole bicanonical system has even multiplicities: multiple-fiber and vertical contributions require separate control. The subsequent deduction \(s=2 h\), \(h\ge 1\Rightarrow p_g>1\), is conditional on that unproved bookkeeping. These are precise limits of these pages, not counterexamples to the bicanonical theorem.

\subsection{Repairing the fixed-subcycle proof of connectedness}

Given an irreducible general moving bicanonical member, any nonpositive split of \(2 K_X=L+C\) has one side supported in the fixed part \(L\). The descent of §17 permits a 1-connected such subcycle \(D\). Let \(\rho=p_a(D)\ge 0\), \(\nu=D^2\). The displayed calculation is correct:

\[
D(2K_X-D)=4\rho-4-3\nu\le0,
\qquad \nu\ge\frac43(\rho-1).
\tag{IV.19.3}
\]

The nonnegativity of \(\rho\) follows from the source's valid induction using Franchetta's Theorem II; alternatively \(h^0(\mathcal O_D)=1\) makes \(\rho=h^1(\mathcal O_D)\ge 0\).

The two low-genus cases admit short rigorous replacements. If \(\rho=0\), integrality gives \(\nu\ge-1\), while nefness gives \(0\le K_X D=-2-\nu\), hence \(\nu\le-2\), contradiction. This avoids calling a reducible or nonreduced genus-zero cycle an irreducible exceptional curve. If \(\rho=1\), nefness and (IV.19.3) force \(\nu=K_X D=0\); Hodge index excludes a nonzero effective cycle of square zero orthogonal to a nef big \(K_X\). The source instead infers an elliptic pencil, but does not establish movement of \(D\), so its inference needs this repair.

For \(\rho>1\), (IV.19.3) gives \(\nu>\rho-1\). The source assigns speciality index \(i(D)=p\). That equality is unnecessary and is not justified for a subcycle of the fixed divisor of \(|2 K_X|\), especially with multiplicities. The always-valid injection \(H^0(K_X-D)\hookrightarrow H^0(K_X)\) gives \(i(D)\le p\). Hence

\[
\dim|D|\ge p+\nu-\rho+1-i(D)
\ge\nu-\rho+1>0.
\]

Replacing \(D\) by a different linearly equivalent effective divisor inside a complete bicanonical member contradicts its being fixed. Thus this part of Enriques's approach can be repaired economically. The direct nef-and-big proof above is shorter and avoids relying on the preceding pencil argument.

\subsection{Why purity and minimality are mathematically essential}

Let \(f:\widetilde X\to X\) blow up \(t\ge 1\) points of a smooth minimal regular general-type surface. Then

\[
K_{\widetilde X}^2=K_X^2-t,\qquad
P_m(\widetilde X)=P_m(X),\qquad
\chi(\mathcal O_{\widetilde X})=\chi(\mathcal O_X).
\]

Pluricanonical forms acquire the prescribed exceptional factors, so the dimensions of their spaces remain birationally invariant. But the Riemann--Roch polynomial on the new surface decreases by \(\binom m 2 t\). Since \(H^2(m K_{\widetilde X})=0\) for \(m\ge 2\), comparison gives the exact defect

\[
\boxed{h^1(mK_{\widetilde X})=\binom m2t\qquad(m\ge2).}
\tag{IV.19.4}
\]

Thus even though the plurigenera are unchanged, the full canonical multiples on the blowup are superabundant. Statements about the pure system cannot be extended to arbitrary blowups by treating minimality as a cosmetic choice. Formula (IV.19.4) is a consistency check derived directly from birational invariance and Riemann--Roch.

\section[Irriducibilità del bicanonico]{Nota sulla irriducibilità del sistema bicanonico (pp.~162--163)}\label{enriques-iv-20}

The author explicitly limits the argument to a reduced canonical divisor. This limitation must be retained. Its first dimension calculation is sound whenever \(K_X=K_1+K_2\), \(K_1\) is integral, \(K_2\) is connected, and \(n=K_1 K_2\ge 2\):

\[
p_a(K_X)=\pi_1+\pi_2+n-1,\qquad
\dim|2K_X|-\dim|K_X+K_2|=\pi_1+n-1>0.
\]

Thus \(K_1\) is not fixed. The restriction to \(K_1\) has degree \(2\pi_1-2+n\ge 2\pi_1\), so its completeness does imply base point freeness on this integral curve; completeness alone, without the degree bound, would not suffice.

Two combinatorial steps are not justified in these pages. First, it chooses an irreducible component of minimum boundary and says the residual remains connected. Theorem II in §17 minimizes over connected subcurves, which need not be irreducible. A general weighted graph can be 2-connected in the numerical edge-cut sense and still have its minimum-degree vertex be a cut vertex: take two copies of a complete four-vertex graph, all edge weights 2, joined to a new central vertex by one weight 2 edge to each copy. The center has degree 4, all other vertices have degree at least 6, and its removal disconnects the graph. \textbf{This is only a diagnostic of the general graph inference; it is not asserted to occur as an effective canonical divisor with nef big \(K_X\), and no canonical counterexample is claimed.} Additional canonical intersection theory might rule out such a configuration, but that argument is absent here. A noncut component does exist by the reduced graph lemma, so minimum boundary is not needed for the first step.

Second, removing rational one-point ``handles'' from a general connected curve need not leave a 2-connected core: separating positive-genus subcurves are possible. The propagation of a fixed handle to a base point on the previously removed component also needs more geometric detail. Again these observations identify missing justification in these pages, not an established counterexample to the special canonical assertion. Mendes Lopes (1996) gives a later systematic treatment of adjoints and rational-chain obstructions; it is not necessary to reconstruct Enriques's peeling procedure to settle the theorem.

A complete modern proof covers nonreduced canonical divisors as well. Catanese, Franciosi, Hulek, and Reid (1999, Theorem 3.3, pp.~194--195) prove that the dualizing sheaf of a numerically 2-connected Gorenstein curve is globally generated unless the curve is \(\mathbb P^1\). Their proof uses duality and factorization through a subcurve: a hypothetical base-point obstruction factors through a generically injective map on a subcurve \(B\); for a proper subcurve, numerical 2-connectedness gives a degree gap at least 2, which cannot be consumed by a length-one base point. The remaining whole-curve case forces the exceptional rational curve. This is exactly the nonreduced case for which the classical argument lacks a proof.

Choose any effective \(D\in|K_X|\). It is a Cartier divisor on a smooth surface, hence Gorenstein; the previous section proves it numerically 2-connected. Its arithmetic genus is \(K_X^2+1\ge 2\), so the rational exception is impossible. Adjunction identifies

\[
\omega_D\simeq\mathcal O_D(2K_X).
\]

Moreover

\[
0\to\mathcal O_X(K_X)\to\mathcal O_X(2K_X)
\to\omega_D\to0
\]

is surjective on global sections because \(H^1(K_X)=q=0\). Thus \(|2 K_X|\) has no base point on \(D\). Outside \(D\), the square of the canonical section defining \(D\) is nonzero, so it has no base point there either. This proves global generation, stronger than absence of fixed components. The same restriction argument appears at the beginning of §6 in Catanese et al.~(1999, p.~206); their later embedding conclusions have stronger hypotheses, which are not being imported into this elementary base-point argument.

Finally \((2 K_X)^2>0\), so the associated morphism has two-dimensional image. Bertini gives a smooth general member, and numerical 1-connectedness makes it connected; a smooth connected curve is irreducible. Thus the general pure bicanonical member is irreducible under the hypotheses of §§19--20, without the source's reducedness restriction and without circularly assuming bicanonical irreducibility to prove canonical connectedness. This conclusion does not assert that the bicanonical map is an embedding or even birational.

\section[Generi di una superficie spezzata]{Generi di una superficie spezzata (pp.~163--165)}\label{enriques-iv-21}

\subsection{The model must be specified before invariants are compared}

The source discusses a projected surface with historical ``normal singularities,'' its limiting apparent double curve, and two components \(Y_1,Y_2\) glued through a further curve \(C\). The degree-\(\delta\) common part of the projected double locus is not automatically the same as the intrinsic gluing/conductor curve. On an arbitrary reducible or nonnormal surface these data can carry multiplicities and embedded structure; numerical formulas cannot be inferred from the reduced supports alone.

The existence of a positive-degree connecting curve is justified in the source by degeneration of plane sections. It relies on the particular connected projective degeneration and the cited degeneration principle. It is not an assertion that every unrelated reducible surface has a positive-dimensional intersection between components.

A clean modern setting is a projective reduced surface \(X_0=X_1\cup_CX_2\), with \(X_i\) smooth and \(C\) a smooth connected curve embedded as a Cartier divisor in both, locally a double normal crossing. Write \(\pi=g(C)\), \(n_i=C^2_{X_i}\), and \(p_i=\chi(\mathcal O_{X_i})-1\). This is an intrinsic gluing statement. To identify these numbers with those of a smooth nearby fiber, additionally assume a proper flat family. To identify canonical squares, require a compatible relative Gorenstein structure, or more generally an integer \(r>0\) for which \(\omega_{\mathcal X/T}^{[r]}\) is invertible and restricts to \(\omega_{X_t}^{[r]}\) on every fiber; flatness of the underlying schemes alone does not supply an arbitrary canonical-divisor comparison. The Stacks Project, Tag 0CBN, supplies the precise normal-crossing definition, and Tag 0B9T gives local constancy of Euler characteristic for proper flat families.

\subsection{Arithmetic genus and virtual dimension}

The exact sequence

\[
0\to\mathcal O_{X_0}\to
\mathcal O_{X_1}\oplus\mathcal O_{X_2}
\to\mathcal O_C\to0
\tag{IV.21.1}
\]

gives

\[
\chi(\mathcal O_{X_0})=\chi(\mathcal O_{X_1})+
\chi(\mathcal O_{X_2})-\chi(\mathcal O_C).
\]

Substituting \(\chi(\mathcal O_C)=1-\pi\) and \(p_a=\chi-1\) yields exactly

\[
\boxed{p_a(X_0)=p_1+p_2+\pi.}
\tag{IV.21.2}
\]

No regularity assumption on \(X_i\) occurs. In particular \(p_i\) cannot be replaced by \(p_g(X_i)\) when the components are irregular. The same Euler-characteristic inclusion-exclusion principle extends to higher dimensions; for unions with more components, higher intersections contribute alternating terms, and the sign convention for arithmetic genus must be respected.

Enriques's virtual count is consistent with this intrinsic proof. Adjoint hypersurfaces of degree \(m_1+m_2-4\) restrict to \(K_{X_1}+C\), with virtual projective dimension \(p_1+\pi-1\). The kernel containing \(Y_1\) consists of residual adjoints to \(Y_2\), with virtual projective dimension \(p_2-1\). Adding image and kernel vector-space dimensions introduces the extra 1 for projective dimensions, giving \(p-1=p_1+p_2+\pi-1\). At p.~165 the author explicitly acknowledges possible failure of actual restriction surjectivity and says that virtual dimensions, or high-degree adjoints, are intended. The Euler-characteristic formulation validates this distinction without presuming surjectivity. A general deformation claim would still need the flatness and conductor hypotheses just stated.

\subsection{Canonical square and the linear-genus formula}

For a double normal crossing, the dualizing sheaf restricts as

\[
\omega_{X_0}|_{X_i}\simeq\omega_{X_i}(C).
\tag{IV.21.3}
\]

Locally this is the residue description on \(xy=0\); the logarithmic poles on the two branches have opposite residues along the common curve. Equivalently it follows from Cartier-divisor adjunction, whose duality formulation is given in the Stacks Project, Tag 0B4A. The canonical square is therefore the sum of the two component squares:

\[
\begin{aligned}
K_{X_0}^2
&=(K_1+C)^2+(K_2+C)^2\\
&=K_1^2+K_2^2+2(K_1C+K_2C)+n_1+n_2\\
&=K_1^2+K_2^2+8\pi-8-n_1-n_2,
\end{aligned}
\tag{IV.21.4}
\]

because \(K_i C=2\pi-2-n_i\). Consequently, if \(p_i^{(1)}=K_i^2+1\) refers to the actual component models being glued,

\[
\boxed{p^{(1)}(X_0)=p_1^{(1)}+p_2^{(1)}+
8\pi-n_1-n_2-9.}
\tag{IV.21.5}
\]

The arithmetic checks in the text agree. On \(X_i\),

\[
p_a(K_i+C)=p_i^{(1)}+3\pi-3-n_i.
\]

Adding these two genera and the common canonical degree \(2\pi-2\), then subtracting 1, yields (IV.21.5). If \(\pi=0\), however, \(2\pi-2=-2\) is not an actual count of points in a canonical group. That part of the geometric language must be read virtually, or replaced by the canonical-square computation, which remains valid. The scan's \(m_1+m_2+ +4\) on p.~165 is a genuine source sign error; the adjoint order is \(m_1+m_2-4\).

The phrase ``linear genus'' must not cause the actual \(K_i^2+1\) in (IV.21.5) to be replaced silently by a birationally invariant minimal-model value: blowing down curves in a component changes both its canonical class and the embedding/normal bundle of the gluing curve. Those adjustments have to be tracked explicitly. Likewise a normalization formula for an arbitrary nonnormal surface needs its complete conductor, not just the reduced intersection curve.

\subsection{A smoothing constraint and an independent numerical test}

If the total space of a semistable smoothing is smooth along \(C\), then

\[
N_{C/X_1}\otimes N_{C/X_2}\simeq\mathcal O_C,
\qquad n_1+n_2=0.
\tag{IV.21.6}
\]

Indeed \(X_1+X_2\) is the principal central fiber, and restricting its line bundle to \(C\) gives the tensor-product relation. This is an additional smooth-total-space condition, not a requirement for the abstract gluing identity (IV.21.4), and not a consequence of flatness alone.

As an explicit consistency test, degenerate a smooth quadric in \(\mathbb P^3\) to two planes meeting in a line. Each plane has \(p_i=0\), \(K_i^2=9\), and \(p_i^{(1)}=10\); the line has \(\pi=0\) and \(n_1=n_2=1\). Formula (IV.21.2) gives 0, and (IV.21.5) gives

\[
10+10-1-1-9=9,
\]

so \(K_X^2=8\), the correct square for a smooth quadric. The family can be written \(xy=tq(z,w)\), with a square-free binary quadratic \(q\). Its total space is singular at the two zeros of \(q\) on the common line, which explains why \(n_1+n_2=2\) does not contradict (IV.21.6). It is nevertheless a flat Gorenstein hypersurface family, so the invariant comparison is legitimate. This is an original consistency check, not an example from the source.

\chapter{Invarianti numerici e piani multipli}
\label{cp1-ch-v}

Chapter V (pp.~167--215) connects the canonical invariants with singular members of pencils, ordinary projections, and branched covers, before turning to the dimension of families of surfaces. The distinctions established in the preceding chapters remain essential: invariants of a specified smooth model must be separated from invariants of its minimal model; an actual family must be separated from its infinitesimal characteristic system; and linear equivalence must be separated from numerical equivalence.

We work over \(\mathbb C\). On a smooth projective surface \(X\) put \(q=h^1(X,\mathcal O_X)\), \(p_g=h^0(X,K_X)\), \(p_a=p_a(X)=\chi(\mathcal O_X)-1=p_g-q\), and \(\lambda=p^{(1)}=K_X^2+1\). Unless minimality is expressly imposed, the last quantity belongs to the chosen model. Write \(e(X)\) for its topological Euler characteristic and \(P_m=h^0(X,mK_X)\) for its plurigenera. When \(H\) is a smooth general member of the system under discussion, \(H^2=n\), \(g(H)=\pi\), and \(K_X H=2\pi-2-n\). A ramification divisor lies on the covering surface; a branch curve lies on the target. Every count below uses intersection or local singularity multiplicities unless genericity makes it an ordinary count of distinct objects. In particular, the book's cusp-count symbol \(\chi\) is replaced by \(c\) when a holomorphic Euler characteristic occurs nearby.

The modern identities used here are adjunction, surface Riemann--Roch, duality, and the blowup formulas already established in the commentary; Beauville (1978, Theorems I.11--I.12 and Sections I.14--I.17) and Reid (1996, Sections A.10--A.12 and B.9) give the general framework. For the exact-sequence derivation of Riemann--Roch see also Brosnan (n.d., Theorem 3.4), referring to Hartshorne (1977, Theorem V.1.6). Enriques provides substantial comparison arguments in the first half of this chapter. The existence theorem for covers has a compressed sufficiency proof, while the stronger classical moduli calculation explicitly depends on an unproved working hypothesis. These different proof statuses are retained below.

\section{L'invariante di Zeuthen--Segre}
\label{cp1-v-1}

\subsection{Pencils and Euler defects}

Resolve the base scheme of a pencil on \(X\) by \(b\) point blowups, producing \(f:\widetilde X\to\mathbb P^1\) with smooth connected generic fibre \(F\) of genus \(g\). For every singular fibre define

\[
\Delta_t=e((F_t)_{\mathrm{red}})-(2-2g).
\]

Then

\[
I(X):=e(X)-4=\sum_t\Delta_t-b-4g. \tag{V.1.1}
\]

For a Lefschetz pencil with distinct simple base points and one node in each singular fibre, \(b=n\) and every \(\Delta_t=1\); this is precisely \(I=\delta-n-4\pi\) on p.~167. ``Relative invariant'' means that \(I\) belongs to a smooth model and changes under blowup. It does not mean a new birational invariant. The introductory claim that it is expressible using \(p_a,p^{(1)}\) is proved later, not here.

\subsection{Segre's comparison and its hypotheses}

On pp.~168--169 he reproduces Segre's comparison of two pencils. Put \(s=C_1C_2\), let \(L\) be their tangency locus, and use the normalizations whenever needed. A general \(C_i\) meets \(L\), outside its own base points, in the ramification divisor of the map to the other pencil; Riemann--Hurwitz gives degree \(2s+2\pi_i-2\). The two induced maps \(L\to\mathbb P^1\) therefore have total ramification

\[
4s+4(\pi_1-1)+2g(L)-2=\delta_1+\tau+n_2,
\] \[
4s+4(\pi_2-1)+2g(L)-2=\delta_2+\tau+n_1.
\]

Here \(\tau\) counts the common third-order tangencies. The terms are justified individually: nodal members of the first pencil, the common third-order contacts, and the base points of the other pencil. Subtraction cancels both \(g(L)\) and \(\tau\). This proves the invariant for pencils in the stated general position. Enriques cites Zeuthen (1871) for the original character and C. Segre (1896) for this comparison.

The argument needs the stated generic relative position, simple ramification in the auxiliary maps, and correct multiplicities at nongeneric contacts. It is not a proof that an unweighted count of singular members is invariant for every pencil. The ``continuity'' passage concerning infinitely near bases on p.~169 is a sketch; resolving the pencil and using (V.1.1) makes its content precise. The number of blowups is not the scheme length \(\sum m_i^2\) of a multiple base scheme. Enriques's convention for ``effective'' and infinitely near bases must be followed together with his corrections to \(\delta\), rather than extracting \(n\) alone and treating it as \(H^2\).

At an ordinary base point of multiplicity \(r\) with varying tangents, Enriques's auxiliary net has Jacobian multiplicity \(3r-1\); subtracting the member of the second pencil leaves multiplicity \(2r-1\) on the two-pencil Jacobian. The relevant intersection with a first-pencil member has multiplicity \(2r\), and after the ramification correction the base point still contributes one to the comparison. These local computations are sketched on pp.~168--169 and require the ordinary moving-tangent hypothesis. They do not identify a multiple base point with \(r^2\) distinct blowups. In the simplest infinitely near example on p.~169, all members have a prescribed common tangent at a simple base point; the two successive nodal degenerations at that point and its first neighbour add two to the relevant count. Resolution explains this convention without appealing only to continuity.

\subsection{An intrinsic proof and singularity weights}

Outside finitely many critical values the map is a differentiable fibre bundle. Additivity of Euler characteristic gives

\[
e(\widetilde X)=2(2-2g)+\sum_t\Delta_t.
\]

Every point blowup increases \(e\) by one; substitution proves (V.1.1). For an isolated reduced plane-curve singularity the local contribution is its Milnor number. An ordinary \(r\)-fold point has \(r\) smooth transverse branches and

\[
\mu=2\delta_{\rm loc}-r+1=r(r-1)-r+1=(r-1)^2.
\]

This proves the multiplicity asserted on p.~169. The adjective ``ordinary,'' or equivalently the generic transversal-branch assumption expressed by Enriques's ``in generale,'' matters: a cusp and a node are both double points but have different Milnor numbers. An arbitrary multiplicity-\(r\) point need not have contribution \((r-1)^2\).

For a generic degree-\(a\) pencil on \(\mathbb P^2\), the base count is \(a^2\), the generic genus is \((a-1)(a-2)/2\), and the nodal contribution is \(3(a-1)^2\). Substitution gives \(I(\mathbb P^2)=-1\) for every \(a\). This provides a check of both the normalization of \(I\) and the distinction between point multiplicity and singular-fibre weight.

\subsection{Reducible and multiple fibres}

The double-component calculation on pp.~170--171 is more explicit than the general formula subsequently quoted. Write a fibre without base points as \(2A+B\), put \(g(A)=\rho\), \(AB=i\), and \(C_2A=\sigma\) for the comparison pencil. The two-pencil Jacobian contains \(A\); write \(L_1\) for its residual curve. The two series on its normalization have degrees \[
m_1=2s+2\pi_1-2,\qquad m_2=2s+2\pi_2-2-\sigma.
\] The first ramification count comprises the ordinary nodal members \(\bar\delta\), the common third-order contacts \(\tau\), the \(n_2\) bases of the comparison pencil, twice the \(i\) intersections with \(B\), and \(2\sigma+2\rho-2\) ramification points on \(A\). Consequently \[
\begin{aligned}
2m_1+2g(L_1)-2&=\bar\delta+\tau+n_2+2i+2\sigma+2\rho-2,\\
2m_2+2g(L_1)-2&=\delta_2+\tau.
\end{aligned}
\] Elimination gives \[
\bar\delta-4\pi_1+2\rho-2+2i=\delta_2-n_2-4\pi_2,
\] so the special fibre replaces \(2\rho-2+2i\) nodal contributions. The method is credited to Castelnuovo and Enriques (1901). The extension from the double component to a component of multiplicity \(h\) is sketched; the fully general formula is quoted from the sources named below.

For a fibre \(F_0=\sum_i m_iD_i\) whose reduced components are smooth and meet transversely, with no triple intersections, write \(g_i=g(D_i)\), \(a_{ij}=D_iD_j\), and sum over unordered pairs \(i<j\). Enriques's Godeaux formula is

\[
\Delta_0=\sum_i(m_i-1)(2g_i-2)
 +\sum_{i<j}a_{ij}(m_i+m_j-1). \tag{V.1.2}
\]

It includes \(2A+B\), for which the contribution is \(2g(A)-2+2i\), and \(hA+B\), for which it is \((h-1)(2g(A)-2)+hi\), provided the residual intersections and components satisfy the same assumptions. These are weighted Euler defects, not literal counts of distinct singular fibres.

Here is a proof not requiring a deformation to nodal fibres. The fibre equations \(F_0D_i=0\), adjunction, and \(K_X F_0=2g-2\) give

\[
2g-2=\sum_i m_i(2g_i-2)+\sum_{i<j}(m_i+m_j)a_{ij}.
\]

On the other hand \(e((F_0)_{\rm red})=\sum_i(2-2g_i)-\sum_{i<j}a_{ij}\). Subtracting \(2-2g\) gives (V.1.2). Enriques provides the two-pencil calculation for a double component on pp.~170--171, sketches the extension to multiplicity \(h\), and cites Godeaux (1920; 1934) and Campedelli (1934) for the general formula. His residual-intersection count on p.~170 must be \(C_2L_1\), not \(C_2L\); the scan preserves the latter, which is a source slip.

For a node on \(D_i\), use either its geometric genus and add \(2m_i-1\), or its arithmetic genus and add one. These two rules agree because a node lowers geometric genus by one. For more general singularities one must use the local delta invariant and the number of branches; one must not extend the ``add one'' rule to all double points. A convenient generalization, derived here, is the correction \(2m_i\delta_p-(r_p-1)\) when the component's geometric genus is used. Blowups first reduce more complicated intersections and base clusters to a configuration where the same Euler-characteristic calculation applies.

A smooth multiple elliptic fibre has Euler defect zero, whereas a nodal elliptic fibre has defect one. Thus a generic genus-one pencil on a K3 surface has total defect \(24\). More generally these weights can include zero contributions; they are not a list of positive numbers of distinct members. This is the uniform interpretation needed in both Sections 1 and 2.

\section{Fascio irrazionale}
\label{cp1-v-2}

For a fibration \(f:X\to B\) over a smooth curve of genus \(b\), with connected generic fibre genus \(g\), the exact formula is

\[
e(X)=4(b-1)(g-1)+\sum_t\Delta_t,
\qquad
I(X)=\sum_t\Delta_t+4(b-1)(g-1)-4. \tag{V.2.1}
\]

This follows from the same Euler-characteristic proof. Enriques alternatively composes \(f\) with a simply branched degree-\(a\) map \(B\to\mathbb P^1\). Such a map can be obtained from a sufficiently ample pencil; his suggested degree \(b+1\) follows from the existence of suitable divisors on a curve. The generic fibre of the composite is a disjoint union of \(a\) fibres of \(f\), and its \textbf{arithmetic} genus is

\[
\Pi=ag-a+1.
\]

It may be disconnected, so this expression must not be called the genus of a connected Riemann surface. The \(2a+2b-2\) simple ramification points of \(B\to\mathbb P^1\) produce fibres containing a double component; each has defect \(2g-2\). Thus

\[
I=\Delta+(2g-2)(2a+2b-2)-4(ag-a+1),
\]

which simplifies to (V.2.1). On p.~173 the printed first factor is \(\pi-2\), whereas both the preceding sentence and the simplification require \(2\pi-2\). The scan confirms a source misprint; the theorem survives the correction.

\section[Zeuthen--Segre e i generi]{Espressione dell'invariante di Zeuthen--Segre per mezzo dei generi}
\label{cp1-v-3}

\subsection{Noether's formula and the choice of model}

A blowup changes \(I\) by \(+1\) and \(\lambda=K_X^2+1\) by \(-1\). Hence \(I+\lambda\) is birationally invariant. Noether's formula is

\[
K_X^2+c_2(X)=12\chi(\mathcal O_X),
\qquad I+\lambda=12p_a+9. \tag{V.3.1}
\]

The modern proof is Hirzebruch--Riemann--Roch applied to \(\mathcal O_X\): the degree-two Todd term is \((c_1^2+c_2)/12\); \(c_1(T_X)^2=K_X^2\), and Chern--Gauss--Bonnet identifies \(c_2\) with \(e(X)\). The same expansion applied to \(\mathcal O_X(D)\) yields

\[
\chi(\mathcal O_X(D))=\chi(\mathcal O_X)+\tfrac12D(D-K_X).
\]

This explains exactly what the historical arithmetic genus measures. It does not replace \(p_a\) by \(p_g\) on irregular surfaces.

\subsection{The ordinary-projection calculation}

The classical proof, credited to Bonnesen (1906), assumes an ordinary projection to a degree-\(n\) surface in \(\mathbb P^3\), with an irreducible ordinary double curve of degree \(d\), normalization genus \(\rho\), ordinary triple points \(t\), and pinch points \(\tau\). Let \(\Delta\subset X\) be the conductor curve, \(\Pi\) its normalization genus, \(H\) a hyperplane section of genus \(\pi\), and \(J=3H+K_X\) the ramification curve of a general projection to the plane. The seven relations being used are

\[
\begin{aligned}
I&=\delta-n-4\pi,\\
\delta+\tau&=3n+12(\pi-1)+\lambda-1,\\
d&=\tfrac12(n-1)(n-2)-\pi,\\
2(\Pi+3t)-2&=2(n-4)d,\\
2\Pi-2&=2(2\rho-2)+\tau,\\
2n^2-7n+15+2\pi(n-7)-\lambda-\tau&=(n-2)d-3t,\\
p_a&=\binom{n-1}{3}-(n-4)d+2t+\rho-1.
\end{aligned} \tag{V.3.2}
\]

Here \(\delta\) is the degree of the dual surface, counted by the nodal members of a general hyperplane pencil. The first relation is the pencil formula. The second is \(J^2-\tau=\delta\). The third is the plane-section genus formula. The fourth uses \(\Delta+K_X\sim(n-4)H\), \(H\Delta=2d\), and adjunction, since \(p_a(\Delta)=\Pi+3t\). The fifth is Riemann--Hurwitz for the double covering of the normalized double curve. For the sixth, \(J\Delta=(n-2)d-3t+\tau\), obtained by the second polar, while direct intersection gives

\[
J\Delta=(3H+K_X)((n-4)H-K_X)
=2n^2-7n+15+2\pi(n-7)-\lambda.
\]

The seventh can be made precise through the conductor exact sequence. Locally at an ordinary triple point the double curve has three spatial branches and delta invariant two; its arithmetic genus is \(\rho+2t\). The conductor curve has three ordinary nodes over each such point, so its arithmetic genus is \(\Pi+3t\). The conductor sequence gives \(\chi(\mathcal O_X)=\chi(\mathcal O_F)+\chi(\mathcal O_\Delta)-\chi(\mathcal O_D)\). Inserting these genera and the fourth equation yields exactly the final equation of (V.3.2). Thus the triple-point corrections have a specific algebraic explanation. Eliminating \(d\) also gives the source's alternative expression \[
p_a=\frac{(n-1)(n-2)(9-2n)}6+\pi(n-4)+2t+\rho-1.
\]

The scan on p.~175 omits \(-\lambda\) from the displayed expansion of \(J\Delta\). Its subsequent equation (6) prints \(2n^2+n+14+2\pi(n-7)-\lambda=(n-2)d-3t\). That expression is inconsistent with the immediately preceding calculations. The sixth line of (V.3.2) is the correction; these are source errors, not transcription errors. The cross-reference for \(d\) should be equation (3), not (2). The ``seven characters'' list on p.~176 omits \(d\).

Solving the corrected equations gives

\[
\begin{aligned}
\delta&=n+4\pi-\lambda+12p_a+9,\\
\rho&=\tfrac12(n^2-7n+48)+(n-12)\pi+9p_a-2\lambda,\\
\Pi&=n^2-6n+2(n-10)\pi+12p_a-3\lambda+36,\\
t&=\tfrac16(n^3-9n^2+26n-78)-(n-8)\pi-4p_a+\lambda,\\
\tau&=2n+8\pi-12p_a+2\lambda-22.
\end{aligned} \tag{V.3.3}
\]

These are the five formulas printed on p.~176; they agree with the corrected system, not with the erroneous equation (6). Consequently \(I=12p_a-\lambda+9\), \(2n+8\pi-\tau=12p_a-2\lambda+22\), and \(2\delta-\tau=36p_a-4\lambda+40\). The last two are relative invariants of the surface expressed using a generic three-dimensional system. Their geometric counting interpretations require the ordinary-projection assumptions. They are not additional independent Chern numbers. In the geometric interpretation on p.~177, \(\tau\) counts the relevant double-base pencils in a generic three-dimensional system; the ordinary projection and the discriminant incidence supply the required multiplicities. This statement is presented without a separate proof there.

\subsection{The plane and ruled surfaces}

A general conic pencil on \(\mathbb P^2\) has four base points, genus zero, and three nodal reducible members; hence \(I=-1\), \(p_a=0\), and \(\lambda=10\), consistently with \(K_X^2=9\). A relatively minimal ruled surface over a genus-\(b\) curve has \(e=4-4b\), \(I=-4b\), \(p_a=-b\), and \(\lambda=9-8b\). ``Absolute linear genus'' for rational surfaces requires particular care: different minimal rational models need not have the same \(K_X^2\). The numerical calculations here are for the specified plane or ruled model.

\section{Curve cuspidate di una rete}
\label{cp1-v-4}

For a sufficiently general net in a sufficiently ample linear system \(H\), whose singular members have the expected stable singularities, the number of cuspidal members is

\[
i=24(p_a+\pi)=24\chi(\mathcal O_X)+12H(H+K_X). \tag{V.4.1}
\]

Enriques works on the ordinary projected surface. Its parabolic curve has class \(P=8H+4K_X\), and the cusp count is \(JP-2\tau\), because a pinch point lies simply on \(J\) and doubly on the parabolic divisor, contributing two to that intersection. Expanding,

\[
JP=24n+20(2\pi-2-n)+4\lambda-4,
\]

and substituting the last formula in (V.3.3) proves (V.4.1). The Hessian has degree \(4n-8\), and subtracting its fourfold double-curve contribution gives \(P\sim4K_X+8H\). The footnote gives a second route: the Jacobian of the polar net is \(J_2=P+H\), whereas its class is \(3J+K_X=9H+4K_X\).

The intermediate sentence on p.~178 is too broad if read as a statement about an arbitrary isolated biplanar double point: its Hessian need not have order four. For example, the homogeneous cubic \(F=uxy+z^3\) has Hessian determinant \(12uxyz\), of local order three at \([0:0:0:1]\). The calculation must retain the surrounding assumption that the point lies on an ordinary double \textbf{curve}. This is a missing local hypothesis, not a counterexample to (V.4.1). The prime attached to one \(F\) in the Hessian transcription is not a different function; the matrix is the symmetric matrix of second derivatives of one polynomial.

Enriques proves the formula for nets embedded in a three-dimensional system, and only sketches extension by adjoining an auxiliary fixed component. A modern intrinsic formulation is the universal \(A_2\)-locus calculation for the first jets of sections: its degree is \(12H^2+12HK_X+2K_X^2+2c_2\), which equals (V.4.1) by Noether. Here is the intrinsic calculation explicitly. On \(X\times\mathbb P^2\), the universal first jet of \(H\) cuts the incidence curve \(Z\) of singular members. With \(h=c_1(\mathcal O_{\mathbb P^2}(1))\), \[
[Z]=c_3(J^1H\otimes\mathcal O(1))
=c_1(J^1H)h^2+c_2(J^1H)h,
\] where \(c_1(J^1H)=K_X+3H\) and \(c_2(J^1H)=c_2(X)+2K_X H+3H^2\). On \(Z\), degeneracy of the Hessian is a section of the line bundle of class \(2K_X+2H+2h\). Its degree is \[
2(K_X+3H)(K_X+H)+2(c_2+2K_X H+3H^2)
=2K_X^2+12K_X H+12H^2+2c_2.
\] For a generic transverse incidence problem, these are precisely the ordinary cuspidal members; Noether gives (V.4.1). This supplies the calculation rather than merely naming a universal formula.

The general-position condition is essential. A deliberately special net, a net composed with a pencil, or a non-simple cover does not have this literal unweighted cusp count.

\section{Notizia storica}
\label{cp1-v-5}

This section supplies historical attributions, not a new proof. Enriques credits Zeuthen (1871) with \(i-24\pi=24p_a\), notes that the early birational-invariance argument had restrictions on surface singularities, and credits Severi (1902) with a direct comparison of nets. Bonnesen (1906), Pannelli (1905), and Caporali are cited for forms of the proof used here. The discussion attributed to Painlevé and Castelnuovo concerns the curve correspondence formula, and the discussion of Severi's equivalence systems, with the bibliographic guide in Severi (1942), concerns an attempted higher-dimensional analogue of divisor linear series. It should not be silently translated into a proved equivalence between every such historical construction and modern Chow groups. The historical citations retained in the references identify the dependencies recorded by Enriques; they do not claim that every original article was independently consulted. Caporali's work is mentioned without a sufficient article locator, so no title or date is supplied by conjecture.

\section{Caratteri di una rete e piani multipli}
\label{cp1-v-6}

Let \(f:X\to\mathbb P^2\) be the finite morphism defined by a generic base-point-free net in \(H\), of degree \(n=H^2\), and suppose it is simply ramified along a smooth ramification curve \(R\sim K_X+3H\). Assume the branch curve \(B\) has only ordinary nodes and cusps and \(R\to B\) is its normalization. Write \(N=\deg B\), \(\rho=g(R)\), \(\delta=\deg B^\vee\), \(k=\#\text{cusps}(B)\), \(d=\#\text{nodes}(B)\), \(i=\#\text{ordinary flexes}(B)\), and \(t=\#\text{ordinary bitangents}(B)\). Then the seven formulas are

\[
\begin{aligned}
N&=2n+2\pi-2,\\
\rho&=\lambda+9\pi-9,\\
\delta&=n+4\pi+12p_a-\lambda+9,\\
k&=3(n+6\pi+\lambda-4p_a-11),\\
i&=24(\pi+p_a),\\
d&=2[(n+\pi)^2-5n-17\pi-2\lambda+6p_a+24],\\
t&=\tfrac12\bigl[(n+4\pi+12p_a-\lambda)^2\\
 &\qquad{}+15n-6\pi+132p_a-17\lambda+74\bigr].
\end{aligned} \tag{V.6.1}
\]

These have all been checked algebraically. The first is \(HR\). The second is adjunction for \(R\). The third is the pencil invariant. The fourth and sixth follow by eliminating \(d,k\) from

\[
\rho=\tfrac12(N-1)(N-2)-d-k,
\qquad \delta=N(N-1)-2d-3k.
\]

The fifth is (V.4.1), equivalently the Plücker flex formula \(i=3N(N-2)-6d-8k\). The last is the genus formula for the dual curve, \(t=\tfrac12(\delta-1)(\delta-2)-\rho-i\). Enriques states these eliminations without writing each line. No numerical inconsistency was found in (V.6.1).

The geometric dictionary is substantive: a triple contact in the net is a cusp of \(B\); a cuspidal member is a flex tangent to \(B\); a two-point tangency in the net gives a node of \(B\); and a two-nodal member gives a bitangent of \(B\). These identifications use stable local models and do not extend unchanged to more singular projections.

Conversely,

\[
\begin{aligned}
\lambda&=\rho-9\pi+9,& p_a&=i/24-\pi,\\
4p_a&=n+\rho-3\pi-k/3-2,& \pi&=N/2-n+1.
\end{aligned}
\]

In the concluding observation the printed \(3(n+6\pi-k)\) must be \(3(n+6\pi)-k\). Its constancy follows from the fourth equation, whereas the printed expression does not. The second invariant \(24\pi-i=-24p_a\) is correctly printed. These branch-curve formulas describe a general net, not every base-point-free net solely on the strength of its being base-point-free; the repeated ``in generale'' is mathematically necessary.

\begin{example}[The branch curve of a quartic]
For a generic projection of a smooth quartic, the tuple
\[
(N,\rho,\delta,k,i,d,t)=(12,19,36,24,96,12,480)
\]
follows by taking $n=4$, $\pi=3$, and $p_a=\lambda=1$.
In particular $t=480$ comes from the genus of the dual curve; it is not a
count obtained by interpreting crossings in a picture. This gives a numerical
check of all seven branch identities simultaneously.
\end{example}

\section[Corrispondenze e piani doppi]{Corrispondenza \texorpdfstring{$(1,n)$}{(1,n)} fra due superficie: piani doppi}
\label{cp1-v-7}

\subsection{Canonical divisors under finite maps}

For a finite generically separable map \(f:X\to Y\) of smooth surfaces,

\[
K_X=f^*K_Y+R_f, \tag{V.7.1}
\]

where \(R_f\) is the divisor of the determinant of \(df\). At a simple ramification point the local model \((u,v)\mapsto(u,v^2)\) gives coefficient one. More generally, a tame ramification index \(e\) gives coefficient \(e-1\). In (V.7.1), \(R_f\) is the ramification divisor, not the whole pullback of the branch divisor: for a general cover \(f^*B\) also contains unramified components. Enriques proves the formula by comparing Jacobians of corresponding nets. That is a valid divisor calculation in the finite, separable situation he assumes.

For a generically finite map with exceptional curves, his instruction to add a ``fundamental curve'' needs the correct discrepancy multiplicities. A single blowup adds its exceptional divisor with coefficient one; a succession of infinitely near blowups can give larger coefficients when expressed on the final model. Thus the one-curve phrase on p.~183 is not a general coefficient-free formula for arbitrary birational correspondences. Resolve the map and apply (V.7.1) and the blowup formula at each stage.

\subsection{Double covers and the complete pluricanonical system}

A flat double cover is specified by a line bundle \(L\) on \(Y\) and a branch section of \(L^{\otimes2}\). Its algebra is

\[
f_*\mathcal O_X=\mathcal O_Y\oplus L^{-1},\qquad
K_X=f^*(K_Y+L).
\]

Projection formula therefore gives the complete plurigenus decomposition

\[
H^0(X,mK_X)=
H^0(Y,mK_Y+mL)\oplus H^0(Y,mK_Y+(m-1)L). \tag{V.7.2}
\]

The two summands give the projective fixed eigenspaces of the involution described on p.~184. For a smooth branch curve of degree \(2n\) on \(\mathbb P^2\), put \(L=\mathcal O_{\mathbb P^2}(n)\). Then

\[
H^0(X,mK_X)=H^0(\mathbb P^2,\mathcal O(m(n-3)))
\oplus H^0(\mathbb P^2,\mathcal O(m(n-3)-n)).
\]

For \(m=1\) the second summand vanishes, so canonical forms correspond to plane curves of degree \(n-3\). For \(m=2\), the eigenspaces are the degree-\(2n-6\) and degree-\(n-6\) systems. A section in the second summand is multiplied by the square-root coordinate; its zero divisor contains the ramification curve. This explains why the whole bicanonical system is invariant under the involution while its general member need not descend to the plane. If the two projective eigenspaces have dimensions \(r\) and \(t\), then \(P_2=r+t+2\), not \(r+t+1\); Enriques correctly writes \(P_2-1=r+t+1\). His \(\pi=t+1\) gives \(P_2=r+\pi+1\). Empty eigenspaces are allowed and have projective dimension \(-1\). The source further identifies \(r+1\) with the genus attached to the first adjoint curves of degree \(2n-3\), and \(\pi=t+1\) with that attached to the degree-\(n-3\) curves. These identifications use the complete adjoint series and the singular-branch conventions developed in Chapter III, rather than just the addition of vector-space dimensions. For any pulled-back plane curve the adjoint is its pulled-back plane adjoint plus ramification; the virtual plane canonical class is the difference of degrees \(m-3\) and \(m\), namely \(-3H\). The canonical transformation formula is therefore valid also when its canonical classes have no effective representatives.

The text permits singular branch curves, importing their multiplicity conventions from Chapter III §7. Those conventions are explained in the commentary on Chapter III. One cannot apply the smooth degree formulas without resolving the branch singularities. A canonical-resolution step at a point of branch multiplicity \(m\) replaces

\[
L\ \text{by}\ L'=\sigma^*L-\lfloor m/2\rfloor E,\qquad
B\ \text{by}\ B'=\sigma^*B-2\lfloor m/2\rfloor E.
\]

The exceptional divisor belongs to \(B'\) when \(m\) is odd. Then \(K_{Y'}+L'=\sigma^*(K_Y+L)+(1-\lfloor m/2\rfloor)E\). Iterating this construction gives the adjoint vanishing conditions and the actual plurigenera. This is a modern repair of the \emph{resolution machinery imported at this point in the book}, not a claim that Enriques invented an incorrect double-cover formula.

The tricanonical and further ``adjoint plus branch'' constructions are instances of (V.7.2); they do not say that every pluricanonical curve belongs to the invariant eigenspace. The final rational/ruled assertion on p.~185 imports the Castelnuovo--Enriques classification of plane curves with vanishing successive adjoints (1900), cited through Conforto (1939, §27, p.~489), based on Enriques's lectures. Enriques does not prove that classification here. The three normal forms he explicitly lists are a degree-\(2n\) curve with a \((2n-2)\)-fold point, a quartic, and a sextic with two infinitely near triple points, together with specializations of these. These are imported Cremona/adjoint classification statements, not three cases derived in the present pages. Over \(\mathbb C\), a surface whose plurigenera all vanish has Kodaira dimension \(-\infty\) and is birationally ruled. In the double-plane case the nonrational base must admit a degree-two map to \(\mathbb P^1\), explaining ``hyperelliptic ruled'' in the historical statement. This implication belongs to surface classification; it does not follow merely by setting a few polynomial degree counts equal to zero.

\section{Formule di corrispondenza}
\label{cp1-v-8}

Let \(f:X\to Y\) have degree \(n\), simple branch curve \(B\), smooth ramification curve \(R\) mapping birationally to \(B\), and no contracted curves. Let \(\rho=g(R)\), \(s=K_YB\), \(\nu=B^2\), \(\varepsilon\) be the number of nodes of \(B\), and \(c\) its number of ordinary cusps. Write \((p_a,\lambda)\) for \(Y\) and \((P_a,\Lambda)\) for \(X\). Pullback of holomorphic two-forms gives \(p_g(X)\ge p_g(Y)\), since the map is dominant and separable. Riemann--Hurwitz on the pullback of a canonical divisor, followed by the genus-addition formula, yields

\[
\Lambda=n(\lambda-1)+\rho+\frac32s. \tag{V.8.1}
\]

More explicitly, the inverse image of a canonical member downstairs has arithmetic genus \(n(\lambda-1)+s/2+1\). Adding the ramification divisor and applying the genus-addition identity gives (V.8.1). This is Enriques's proof on p.~185, with virtual divisors understood through their intersection classes. The term \(s\) can be negative because a canonical divisor may be virtual; the polynomial identity remains an intersection-theoretic statement. The source switches between \(D\) and \(d\) for the branch curve on p.~185; the notation \(B\) downstairs and \(R\) upstairs used here removes that local ambiguity.

\subsection{The pencil calculation on p. 186}

The continuous, legible passage on p.~186 begins ``Si designi dunque'' and supplies a proof of the correspondence formula by comparing pencils. \footnote{Because the book proves the correspondence formula on p. 186, the commentary has not
presented this formula only as a modern reconstruction. The supplementary
evidence of the scan of the print is explicit: the lower passage
does contain the pencil proof, including the negative cusp term and the
final constant $-4$. The upper-left part of this leaf is physically
occluded. The visible “184” belongs to an overlapping portion, so the
underlying folio number is not independently readable; p. 186 is the
edition's positional identification. The beginnings of the preceding
paragraphs and their full linear-genus formula cannot be quoted from this
image and are not reconstructed as recovered text.}

Choose a sufficiently general Lefschetz pencil on \(Y\), in a sufficiently high very ample system, with \(m\) simple base points outside the branch, general member \(C\) of genus \(\pi\), and nodal contribution \(\delta\). Put \(t=CB\). Its pullback pencil has \(M\) base points, smooth general member of genus \(\Pi\), and nodal contribution \(\Delta\). The pullback line bundle is ample and the pullback system is generated; Bertini and connectedness give a smooth connected general member. In the unified notation used here, the printed calculations are \[
i=\delta-m-4\pi,\qquad M=mn,\qquad
\Pi=n(\pi-1)+1+\frac t2,
\] \[
\Delta=n\delta+2t+2\rho-2-c,
\] \[
\begin{aligned}
I&=\Delta-M-4\Pi\\
 &=n\delta+2\rho-2-c-nm-4n(\pi-1)-4.
\end{aligned}
\] Here \(i=I(Y)\) and \(I=I(X)\); the book writes \(\chi\) for the cusp count called \(c\) here. The term \(n\delta\) is explicitly weighted: a lifted singular member has \(n\) nodes and is counted \(n\) times. The phrase ``number of singular members'' cannot be read as an unweighted count.

Every base point is in the étale locus and therefore has \(n\) preimages, which proves \(M=mn\). Riemann--Hurwitz on a general member gives \(2\Pi-2=n(2\pi-2)+t\). Applying the same theorem to the degree-\(t\) map from the normalization of \(B\) to the pencil parameter line gives total ramification \(2t+2\rho-2\). A simple tangency to the smooth branch locus produces an ordinary node upstairs; its local equation is \(aw^2+bv^2+\cdots=0\) with \(ab\ne0\). A branch cusp contributes to the normalization's ramification count but does not produce such an extra node. In the ordinary three-sheet local model \[
f(x,z)=(x,z^3+xz),\qquad
B:\ (x,y)=(-3z^2,-2z^3),
\] a general pencil parameter \(h(x,y)=\alpha x+\beta y+\cdots\), with \(\alpha\ne0\), satisfies \(d(h\circ f)_0=\alpha\,dx\ne0\). Its inverse-image member is smooth there, while on the normalization \(h=-3\alpha z^2-2\beta z^3+\cdots\) has simple ramification. Subtracting one for every cusp gives exactly the displayed \(-c\). At a branch node the two disjoint simple-ramification charts produce no extra singularity for a transverse pencil. Each original Lefschetz node away from \(B\) lifts to \(n\) nodes. These local checks complete the compressed cusp explanation in the printed proof.

The \(2t\) in \(\Delta\) cancels the term \(4(t/2)\) in \(4\Pi\), removing dependence on the chosen pencil. The resulting relations are \[
\begin{aligned}
I(X)&=nI(Y)+4(n-1)+2\rho-2-c,\\
e(X)&=ne(Y)+2\rho-2-c.
\end{aligned}\tag{V.8.2}
\] These are deductions from the visible displays, not quotations from the occluded upper text. For empty branch the convention is \(\rho=1,c=0\), and the formula gives the expected \(e(X)=ne(Y)\).

An independent stratification verifies the same constant and sign. Let \(\varepsilon\) be the number of branch nodes. For an irreducible nodal-cuspidal branch curve, \(e(B)=2-2\rho-\varepsilon\). The fibre cardinalities are \(n\) outside \(B\), \(n-1\) on its smooth locus, and \(n-2\) at a node or cusp, provided node monodromy consists of disjoint transpositions and cusp monodromy is the ordinary three-sheet model. Thus \[
\begin{aligned}
e(X)={}&n\bigl(e(Y)-e(B)\bigr)
 +(n-1)\bigl(e(B)-\varepsilon-c\bigr)\\
 &+(n-2)(\varepsilon+c)\\
={}&ne(Y)-e(B)-\varepsilon-c
 =ne(Y)+2\rho-2-c.
\end{aligned}
\] Nodes cancel in this calculation; cusps do not. For a disconnected branch normalization, replace \(2-2\rho\) by the sum of its components' Euler characteristics. Kulikov and Kulikov (2010, equation (24), p.~138) give the plane-base specialization \(e(X)=3n+2(\rho-1)-c\). Their formula corroborates the calculation and is not being credited as an erratum to Enriques. The historical footnotes on p.~186 cite Enriques (1893/1894) and Severi (1903); the latter's expanded bibliographic title is recorded in the references, without using that identification to supply invisible text.

\subsection{Canonical squares and holomorphic Euler characteristics}

Combining (V.8.1), (V.8.2), and Noether gives the displayed p.~187 formula

\[
24(P_a+1)=24n(p_a+1)+6\rho-6+3s-2c. \tag{V.8.3}
\]

Adjunction on the singular branch curve gives \(s=2(\rho+\varepsilon+c)-2-\nu\). Therefore

\[
\begin{aligned}
\Lambda&=n(\lambda-1)+4\rho+3(\varepsilon+c)-\tfrac32\nu-3,\\
24(P_a+1)&=24n(p_a+1)+12\rho-12+6\varepsilon+4c-3\nu.
\end{aligned} \tag{V.8.4}
\]

The printed second formula on p.~187 ends in \(-3\); it must end in \(-3\nu\). This is confirmed by the substitution, by the next page's \(-3N^2\), and by the double-cover examples. If \(R^2=x\), comparison of \(K_X^2=(f^*K_Y+R)^2\) with (V.8.1) gives

\[
x=\nu/2-(\varepsilon+c).
\]

For a smooth branch component in a double cover, this reduces to \(R^2=B^2/2\), also directly from \(f^*B=2R\) and \((f^*B)^2=2B^2\). With no branch divisor, the conventional \(\rho=1\) makes the terms \(2\rho-2\) vanish; it does not mean that an actual elliptic branch curve exists. For disconnected branch divisors, \(\rho\) must similarly be interpreted by \(2-2\rho=e(\widetilde B)\), not as a positive genus.

Putting \(Y=\mathbb P^2\), \(\lambda=10\), \(p_a=0\), \(\nu=N^2\), \(\rho=g(R)\), \(\varepsilon=d\), and \(c=k\) recovers the two inverse formulas in V.6. Enriques performs these substitutions explicitly on p.~188. This checks consistency on the conventions and the factor \(24\). In an intrinsic form, before introducing the historical genera, the same argument reads \[
12\bigl(\chi(\mathcal O_X)-n\chi(\mathcal O_Y)\bigr)
=K_X^2-nK_Y^2+2\rho-2-c.
\] The canonical squares on both sides refer to these same smooth models. Contracting exceptional curves requires the corresponding square corrections.

\begin{example}[Checks and the necessity of simple branching]
For a smooth sextic double plane, $\rho=10$, and the Euler formula gives
$e(X)=2\cdot3+20-2=24$. The double-cover formula gives
$K_X\sim0$ and $\chi(\mathcal O_X)=2$, confirming the result by Noether's
formula. For a smooth quartic double plane, $\rho=3$ gives $e(X)=10$;
here $K_X^2=2$ and $\chi(\mathcal O_X)=1$, so the relation is not confined
to numerically trivial canonical classes.

A smooth cyclic triple cover branched over a smooth sextic has fibre
cardinalities three off the branch and one on it. Its actual Euler number is
$3e(\mathbb P^2)-2e(B)=9+36=45$, whereas the simple-branch formula would
give $27$. Its covering algebra
$\mathcal O\oplus\mathcal O(-2)\oplus\mathcal O(-4)$ has holomorphic
Euler characteristic $4$, and $K_X=f^*\mathcal O(1)$ has square $3$,
again giving $e(X)=48-3=45$. Thus the generic local ramification hypotheses
cannot be suppressed. For a finite map with constant fibre cardinalities on
strata $T$, the general principle is
$e(X)=\sum_T\#f^{-1}(t)\,e(T)$; higher branch singularities require their
actual local fibre deficits, not a presumed cusp correction.
\end{example}

\subsection{The Kummer double cover (pp. 188--190)}

Enriques assumes the classical \(16_6\) configuration of nodes and trope planes of a Kummer quartic; its proof is referred to Enriques--Chisini, Book VI §40, vol.~IV, p.~246. He chooses a tetrahedron whose four vertices and remaining twelve nodes give odd vanishing orders for its defining product. Three trope planes through a node determine the tetrahedron; the remaining twelve nodes occur three on each face. The construction uses the imported \(16_6\) incidence configuration, rather than establishing that configuration anew. On the smooth Kummer resolution \(Y\), its double cover, which is connected because a split double cover would have no branch, is branched along the sixteen disjoint exceptional \((-2)\)-curves \(E_i\). Thus

\[
B=\sum_{i=1}^{16}E_i,\quad B^2=-32,\quad
\rho=1-16=-15,\quad \varepsilon=c=0.
\]

A generic quartic hyperplane section has genus three; its unramified degree-two inverse image has genus five and self-intersection eight, exactly the \(8=2\cdot5-2\) calculation on p.~189. The cover \(\widetilde X\) consequently has \(K_{\widetilde X}^2=-16\) and \(\Lambda=-15\). This is not a contradiction with the asserted minimal value \(\Lambda=1\): above each \(E_i\) lies a \((-1)\)-curve, and contracting all sixteen increases \(K_X^2\) by sixteen. Equation (V.8.3) gives \(\chi(\mathcal O_X)=0\), hence \(P_a=-1\). The constructed minimal surface has trivial canonical bundle, \(p_g=1\), \(q=2\), and every \(P_m=1\). In modern terminology it is an abelian surface (the Jacobian in the genus-two construction), whereas Enriques's ``hyperelliptic surface'' here uses an older, broader vocabulary.

One can verify the holomorphic Euler characteristic independently from the double-cover formula

\[
\chi(\mathcal O_{\widetilde X})
=2\chi(\mathcal O_Y)+\tfrac12 L(L+K_Y)
=4+\tfrac12(-8)=0,\qquad 2L\sim B.
\]

Since \(K_Y=0\), \(K_{\widetilde X}=f^*L\) is supported on the sixteen exceptional curves, and it disappears after contraction. The genus-two representation is imported from Klein's construction, cited through Enriques and Chisini and Enriques and Severi (1909, as cited). The symmetric square of a genus-two curve is \textbf{birational} to its Jacobian, not isomorphic to it: the Abel map contracts the curve coming from the hyperelliptic pencil. Enriques says ``equivale,'' so the source should not be accused of asserting an isomorphism. More specifically, for a degree-four class \(L\) with complete series \(g^2_4\), the involution on degree-two divisor classes is \(D\mapsto L-D\). Its fixed classes satisfy \(2D=L\); multiplication by two on the genus-two Jacobian has sixteen solutions. For general \(L\), these give the sixteen fixed points of the stated quotient construction.

\subsection{The tetrahedral sextic and its canonical cover (pp. 190--192)}

Let \(Y\) be the smooth model of the sextic double along the six edges of a tetrahedron. The constructed degree-two cover is unramified on the smooth model. The source imports \(p_a=p_g=0\), \(\lambda=1\), and nontrivial bicanonical triviality from Chapter II §8. Equations (V.8.1)--(V.8.3) give \(K_X^2=0\) and \(\chi(\mathcal O_X)=2\). A hyperplane section of the sextic has genus four; its étale double cover has genus seven and self-intersection twelve, so the equality \(12=2\cdot7-2\) on p.~191 is correct. The source's local argument excludes branching not only along the tetrahedral edges but also over each vertex: the normalization separates its three smooth branches and the relevant face-section multiplicity is even. These are the local facts accounted for intrinsically by the canonical-cover identification below.

There is, however, a real gap in the p.~191 implication

\[
|2C'|=|2C|\quad\Longrightarrow\quad |C'|=|C|.
\]

Cancellation by two in \(\operatorname{Pic}\) is not valid: the nontrivial canonical two-torsion of an Enriques surface is precisely the counterexample. The conclusion for this \emph{particular} cover is right, but requires identifying the cover as the canonical cover. A direct repair uses the conductor divisor \(\Delta\) on \(Y\). Normalization adjunction gives \(K_Y\sim2H-\Delta\). The tetrahedron product \(T\), as a section of \(\mathcal O_Y(4H)\), has divisor \(2\Delta\). Consequently adjoining \(\sqrt T\) defines the étale cover associated to the two-torsion line bundle \(\mathcal O_Y(2H-\Delta)=\omega_Y\). Its pullback becomes trivial; therefore \(K_X\sim0\). Together with \(\chi=2\), this gives \(p_g(X)=1\), \(q(X)=0\), all \(P_m=1\), and the modern K3 identification. The numerical multiplication on p.~191 is \(24(P_a+1)=2\cdot24(p_a+1)=48\). \footnote{The edition represents the mark between $2$ and
$24$ as a colon and flags it for correction. It is best to treat the
small lower mark as a possible blemish accompanying a multiplication dot.
The enlarged scan does not make a mathematical division reading defensible:
the degree-two correspondence and the printed result $48$ uniquely require
multiplication. The formula is normalized accordingly without treating the
damaged mark as independent mathematical evidence.} This repair is supplied by the commentary, using the standard cyclic-cover construction, not attributed to an unverified historical author.

The cover statement is compared on p.~192 with Enriques's earlier observation on bigenus one (Enriques, 1908, \emph{Un'osservazione}). The Comessatti splitting criterion in the p.~189 footnote has the following precise content: restricting a double cover to a curve tangent to its branch divisor gives a split cover exactly when the half-contact divisor is linearly equivalent to the restriction of the branch square-root line bundle. Even contact multiplicities alone do not suffice; the square-root class is essential.

\section{Teorema d'esistenza}
\label{cp1-v-9}

\subsection{Discriminants and the line-cover problem}

A polynomial equation \(z^n=f(x,y)\) defines a cyclic cover of the affine plane when its function field has degree \(n\), but its \emph{discriminant divisor} includes \((n-1)(f=0)\) and may branch at infinity. If \(\deg f=m\), the monodromy around infinity is the inverse of an \(m\)th power of an \(n\)-cycle, so it is trivial exactly when \(n\mid m\). More precisely, the coefficient of the infinity component in the discriminant is \(n-\gcd(n,m)\) when it occurs. The latter numerical refinement is derived here from the ramification indices; Enriques only says ``one or more times.'' Irreducibility is automatic for a reduced irreducible branch polynomial over \(\mathbb C\), but not for an arbitrary perfect power.

The polynomial discriminant is obtained by eliminating \(z\) between \(F\) and \(\partial F/\partial z\). For a general projected singular model, it separates as \(\Delta=\Delta_1\Delta_2^2\): \(\Delta_1\) is genuine simple branch, while \(\Delta_2\) comes from the apparent double curve of the chosen equation. This factorization is not a universal factorization for every non-generic polynomial; it describes the ordinary projection under discussion. The finite normalization, not merely the raw polynomial resultant, is the geometric cover.

The degree of a reduced simply branched plane curve is even: Riemann--Hurwitz for a general line cover, or the sign of the product of its transpositions, gives this parity. Choose a line pencil through \(O\notin B\) and a generic fibre meeting a degree-\(2m\) branch curve in \(2m\) points. A connected degree-\(n\) cover of that line is determined by a transitive tuple of \(2m\) transpositions with product one, modulo simultaneous relabelling. Its genus is

\[
g=m-n+1.
\]

The scan on p.~194 really prints \(m/2-n+1\); that is a source error. The statement that there are finitely many possibilities refers to covers with fixed branch locations and fixed degree, not to a finite set of curves independent of those locations. Enriques cites Riemann's existence theorem through Enriques--Chisini and mentions Hurwitz's enumeration without using its number.

\subsection{Local monodromy and the global compatibility conditions}

Moving the line gives braid monodromy. A class of line covers coming from a global surface cover must return to itself around every loop of the pencil parameter. Ordered noncrossing arcs can be transported continuously away from critical values, and a contractible loop in that complement adds no condition. After transporting the vanishing arcs so the colliding punctures are adjacent, with no intervening punctures, the required transpositions are:

At a simple tangent line, the two adjacent meridians become equal, so their transpositions must be identical. At a node, the meridians commute; for a smooth generic covering surface their transpositions are disjoint. At a cusp the relation is \(aba=bab\), and the two transpositions are distinct with one letter in common; their product is a three-cycle.

Equal transpositions also commute at a node but produce a singular double-cover local model; Enriques calls that a ``virtually nonexistent'' node for the generic smooth-cover enumeration. At a cusp, equal transpositions similarly give a singular local covering surface rather than the required generic triple-sheet model. This is why the smoothness/local-model conditions are stronger than an arbitrary representation into \(X_n\).

The reduction to adjacent arcs is the Hurwitz action, not a permission to move arcs through branch points while leaving their transpositions unchanged. In one convention the elementary move is \((a,b)\mapsto(b,b^{-1}ab)\); changing the convention reverses this formula but not the existence condition. Enriques explicitly notes the necessary change of substitutions on p.~197. The three classes of critical values generate the pencil monodromy, together with the projective product relation. Thus checking them is a finite topological formulation of a global condition, not merely three numerical tests on the numbers of nodes and cusps.

\subsection{Existence, algebraization, and the status of the original proof}

A connected degree-\(n\) finite cover of \(\mathbb P^2\) with prescribed simple branch \(B\) and the stated local models exists exactly when there is a transitive monodromy representation

\[
\pi_1(\mathbb P^2\setminus B)\longrightarrow X_n
\]

whose local meridians and local singularity representations have the preceding types. The representation first produces a finite topological cover of the complement. The complex Riemann existence theorem algebraizes it. Normalizing \(\mathbb P^2\) in its function field extends it to a finite normal projective cover, and the prescribed local models show that its total space is smooth. Conversely, any such cover supplies the representation. Grothendieck (Ed., 2003), SGA 1, Exposé XII, Theorem 5.1, is a precise modern reference for the algebraization step; the normalization and local checks are separate steps.

\textbf{What Enriques supplies and what needs completion.} The necessity is developed carefully on pp.~194--197. The sufficiency on p.~198 is a short construction: realize the varying line cover by finitely many sufficiently high-degree plane models, sum the corresponding functions, and vary the line. It assumes that the resulting family is algebraic and descends, and that the sum still generates the intended degree-\(n\) function field. A sum of primitive functions need not itself be primitive. The prescribed distinct points \(O_1,\dots,O_n\) can repair this last issue if they are consistently labelled in every summand: the sum then takes distinct values on that fibre. Nevertheless the algebraization and global gluing remain compressed; the preceding modern proof supplies them without claiming that every trace is automatically primitive. His exclusion of additional branch lines relies on the same distinct-fibre normalization.

\subsection{Deformation and uniqueness}

The existence property is constant along a connected, topologically locally trivial equisingular family, with the local covering types transported along the family. The assertion must not be read as saying that degree, number of nodes, and number of cusps determine the complement's fundamental group. Equisingular strata may have different connected components, and passage through worse singularities can change the answer. Zariski's examples and the modern deformation literature make this distinction indispensable. Likewise existence is different from \textbf{uniqueness}. Enriques's p.~193 appeal to Chisini's 1944 work is a qualified anticipation of a uniqueness problem; it is not a proved general uniqueness theorem in these pages. Modern results for generic linear projections and conjectures for more general generic covers must be distinguished. Kulikov (2008, Theorem, p.~902) proves uniqueness for generic linear projections, with the Veronese exception stated explicitly. That theorem is not a blanket theorem for every generically branched cover of a fixed degree.

\section[Le degenerazioni di Chisini]{Complementi: forme limiti della curva di diramazione secondo Chisini}
\label{cp1-v-10}

Enriques credits the existence theorem to his 1923 paper and Zariski's topological reformulation to papers of 1929 and 1931. The statement that a cusp-free branch curve cannot give a connected simply branched cover of degree \(>2\) uses the standing irreducibility and nodal hypotheses. For an irreducible nodal plane branch curve, the complement group is cyclic; transitive monodromy generated by its one meridian is then cyclic, and a transposition permits only degree two. If multiplicities in the discriminant are allowed, cyclic covers explain the stated exception. With reducible branch divisors and more general ramification there are other abelian covers; the sentence should not be generalized to them. The history of the needed fundamental-group theorem matters: Deligne (1981), Introduction and Theorem 1, explains that Zariski's 1929 argument relied on Severi's then-unproved irreducibility result. Fulton (1980) proved the tame algebraic fundamental-group assertion; Deligne supplied the topological theorem. Thus the rigorous modern nodal-complement justification must not simply be attributed to the 1929 proof.

\subsection{Degeneration of the discriminant}

If a surface equation degenerates into \(F_1F_2=0\), the discriminant acquires the square of the resultant:

\[
\operatorname{disc}(F_1F_2)
=\operatorname{disc}(F_1)\operatorname{disc}(F_2)
 \operatorname{Res}(F_1,F_2)^2
\]

up to a nonzero normalization factor. Thus the limiting branch divisor is \(B_1+B_2+2C\), where \(C\) is the image of the common curve. This algebraic identity explains the construction. Its converse requires a smoothing with the correct monodromy; an arbitrary sum with a doubled curve is not enough.

For a smooth quartic projected from outside, \(n=4,\pi=3,p_a=\lambda=1\) in (V.6.1), so \(B=C_{12}\) has \(24\) cusps and \(12\) nodes. Letting the center move onto a smooth point produces \(C_{10}+2L\), where \(C_{10}\) is a triple-plane branch curve with \(18\) cusps and no nodes, and \(L\) is bitangent. Making the center an ordinary double point produces \(C_6+2Q+2L\), with \(Q\) a six-tangent conic. A triple center gives the further limit \(2C_3+2Q+2L\). The sextic in the second degeneration is \textbf{special enough to possess the six-tangent conic}; the phrase ``affatto generale'' cannot mean a general point of the full space of smooth sextics simultaneously with that condition.

The local counts are consistent: each tangency to an added double component absorbs three cusps; each simple crossing absorbs two nodes. Thus the two tangencies and six further intersections in \(C_{10}+2L\) account for six additional cusps and twelve nodes. The six conic contacts in \(C_6+2Q\) account for eighteen cusps. Enriques explains these through the first and second polars and the bitangents through the moving center, but refers the global smoothing construction to Chisini's papers of 1934, 1938, and 1939.

For the proposed K3 family, put \(n=2\pi-2\), \(p_a=\lambda=1\). The branch data become

\[
\deg B=3n=6\pi-6,\qquad
k(\pi)=24(\pi-2),\qquad
d(\pi)=18\pi^2-78\pi+84.
\]

If such a \textbf{smooth finite generic cover} exists, its pullback \(H\) is ample, \(K_X^2=0\), \(\chi=2\), and \(K_X H=0\). Riemann--Roch gives \(p_g=q+1>0\); a nonzero canonical section has an effective zero divisor of intersection zero with \(H\), hence no zero divisor. Therefore \(K_X\sim0\), \(p_g=1\), and \(q=0\). Kodaira vanishing then gives \(h^0(H)=\pi+1\). If the polarization is very ample, it embeds the surface with degree \(2\pi-2\) in \(\mathbb P^\pi\), with canonical hyperplane sections. Amplitude by itself does not establish this embedding: some K3 polarizations are not very ample, so the generic embedding hypothesis on p.~202 must remain visible. This supplies the content of Enriques's inference that these are surfaces with trivial canonical curve and a \(\pi\)-dimensional hyperplane system. It is stronger than merely saying that their two numerical invariants happen to be one.

The step \(\pi=3\to4\) adds eighteen cusps and twenty-four nodes at the conic stage, then six cusps and twenty-four nodes at the line stage, giving \(48\) cusps and \(60\) nodes for the sextuple plane. The intermediate branch has degree sixteen, \(42\) cusps and \(36\) nodes; the final branch has degree eighteen. These are the separate numerical stages displayed on p.~202. In general the conic stage adds eighteen cusps and \(12n-24\) nodes; the line stage adds six cusps and \(6n\) nodes. Their total node increment is \(18n-24=36\pi-60\), exactly \(d(\pi+1)-d(\pi)\).

The final p.~203 assertion about selecting the line has a more concrete reconstruction than the preceding conic step. The successful intermediate \((n+1)\)-plane has \(p_a=1\), \(K_X^2=-1\), \(H^2=n+1\), and \(K_X H=1\). Since \(\chi=2\), there is an effective canonical divisor. Its ample degree one forces it to be a single reduced irreducible curve \(E\); adjunction gives \(p_a(E)=0\) and \(E^2=-1\). Contracting it produces a K3 surface with a polarization of square \(n+2\). The curve \(E\) maps to a line because \(HE=1\), and its intersection with the ramification divisor is \((3H+E)E=2\), explaining the required two tangencies in the generic model. This fills the source's contraction argument, subject to its ordinary-projection and embedding assumptions. It does not remove the earlier conic-stage obstruction.

\textbf{What is not proved.} Having at least nineteen moduli does not alone establish the existence of a specialization with a prescribed six-tangent conic. Nor does the existence of such a conic ensure that its double can be smoothed with the desired branch monodromy. Enriques explicitly acknowledges the latter obstruction on p.~203 and does not claim to have carried out the induction for all \(\pi\). For the initial \(n=4\) step, Enriques sketches a suitable choice through a Chisini limit possessing a second six-tangent conic to the sextic; this remains a special construction. The low-degree construction is an illustration of Chisini's method; the general K3 existence/classification is deferred to Chapter VII §3, outside Part 1. The historical dependencies on pp.~203--204 are Enriques (1896, 1898), Campedelli (1932a--1932c, 1940a--1940c), Libois (1934), B. Segre (1930), and Pompilj (1939a, 1939b, 1941), for double, biquadratic, triple, and dihedral planes. They are bibliographic pointers; the examples and classifications in those papers are not proved here.

\section{I moduli di una classe di superficie algebriche}
\label{cp1-v-11}

The moduli number means the dimension of a component of birational equivalence classes with specified numerical characters and the polarization needed for the parameter count; it does not imply that every surface with those characters belongs to one smooth family. Enriques restricts first to a minimal regular nonruled surface, \(q=0\), \(p=p_a=p_g\), with some nonzero plurigenus. Choose \(H\) sufficiently positive that the complete linear system is base-point-free, simple, and regular and nonspecial in the cohomological sense \(H^1(H)=H^2(H)=0\). Surface Riemann--Roch gives

\[
r=\dim|H|=p+n-\pi+1.
\]

The inequality \(n\le2\pi-2\) is \(K_X H\ge0\), valid because \(K_X\) is nef on the minimal nonruled complex surface. For branch degree \(N=2n+2\pi-2\), a node imposes at most one condition and an ordinary cusp at most two on the space of degree-\(N\) plane curves when their locations may move. Consequently the equisingular component has dimension at least

\[
s_0=\frac{N(N+3)}2-d-2k
=(n+\pi-1)(2n+2\pi+1)-d-2k.
\]

This dimension lower bound does \textbf{not} assume that the equisingular scheme is smooth or reduced. Write its actual component dimension as \(s_0+\omega\), \(\omega\ge0\). Changing plane coordinates accounts for eight parameters, and choosing a net in the \((r+1)\)-dimensional space of sections accounts for \(\dim\operatorname{Gr}(3,r+1)=3r-6\). The generic degree of the map from covers to branch curves is finite for fixed branch curve and degree, by finite generation of the complement group; uniqueness is unnecessary for this dimension calculation.

For \(q=0\), a continuously varying line-bundle class on a fixed surface is constant, since \(\operatorname{Pic}^0(X)\) is zero-dimensional. The identity component of the automorphism group does not contribute for the regular nonruled surfaces under consideration; its polarized stabilizer is finite. Infinite \textbf{discrete} automorphism groups, such as those of many K3 surfaces, are not ruled out and do not invalidate a dimension calculation. This is the content behind the source's Picard-system argument and its appeal to Enriques--Fano on continuous projective automorphisms.

Subtracting the two redundant dimensions and inserting (V.6.1) gives

\[
M\ge9p-2\lambda+12,\qquad
M=9p-2\lambda+12+\omega. \tag{V.11.1}
\]

The equality expresses \(\omega\) as the actual excess of the relevant branch component, with the preceding fibre-dimension assumptions. Equivalently the first count is \(M\ge s_0-8-(3r-6)=s_0-3r-2\). It does not identify that excess with a tangent-space dimension.

\subsection{The base-point calculation and irregular surfaces}

The text addresses a possible degeneration from systems with base points to the chosen free system. If multiplicities are \(i_j\), moving self-intersection is \(n\), and the resolved generic member has geometric genus \(\pi\), then

\[
L^2=n+\sum i_j^2,\qquad
p_a(L)=\pi+\frac12\sum i_j(i_j-1),
\]

so adjunction gives

\[
K_X L=2\pi-2-n-\sum i_j,
\]

not the printed \(2\pi-2-n-\sum i_j^2\). Equivalently, on the resolved surface, \(K_{\widetilde X}\widetilde L=K_X L+\sum i_j\). The printed square is valid if one holds the \emph{unresolved arithmetic genus} fixed, but that conflicts with the phrase ``genere effettivo'' in this argument. If \(L\) specializes to \(H+\theta\), the corrected formula still gives \(K_X\theta=-\sum i_j<0\). This contradicts nefness. The corresponding residual-square identity is \[
\theta^2=2p_a(\theta)-2+\sum i_j>2p_a(\theta)-2.
\] The source's sum of squares must be replaced here too, and the genus is arithmetic if \(\theta\) is singular or reducible. Thus the intended exclusion survives the quantitative correction.

For an irregular surface, line bundles in \(\operatorname{Pic}^0\) add \(q=p_g-p_a\) parameters to the redundant family. The resulting expression is

\[
M=10p_a-p_g-2\lambda+12+\omega.
\]

If a continuous automorphism group acts, its dimension must be included in the fibre calculation; Enriques's footnote explicitly allows a changed meaning of \(\omega\). One must not identify that redefined \(\omega\) with the original branch excess without an automorphism correction.

The examples \(M=19\) for K3 surfaces and \(M=10\) for Enriques surfaces are asserted here and promised later proof. The first is the dimension of a family with a fixed polarization type. The local analytic deformation space of an unpolarized complex K3 surface has dimension twenty; projectivity with a specified polarization cuts out one condition. Thus nineteen is not an error, but the moduli problem needs to be specified.

\section[La serie caratteristica delle curve piane]{Digressione sulla integrit\`a della serie caratteristica d'un sistema completo di curve piane dotate di nodi e di cuspidi}
\label{cp1-v-12}

Let \(D\) be an irreducible plane curve of degree \(n\) with \(d\) nodes and \(k\) ordinary cusps, and let \(\nu:C\to D\) be its normalization. The line bundle controlling its equisingular first-order deformations is

\[
A=\nu^*\mathcal O_D(n)
 \Bigl(-\sum_{\rm nodes}(P_i+Q_i)-3\sum_{\rm cusps}R_j\Bigr).
\]

The source attributes the absorbed-intersection observation to de Jonquières and Beck (1878, as cited), with a reference to Enriques--Chisini, Book II, Chapter II. At a node, the equisingular Jacobian ideal pulls back to order one on both branches; at a cusp \(x=t^2,y=t^3\), the ideal \((x^2,y)\) pulls back to order three. Hence

\[
m=\deg A=n^2-2d-3k,\qquad
g(C)=\frac{(n-1)(n-2)}2-d-k.
\]

The printed \(m=n-2d-3k\) on p.~208 omits the square. Put \(r=n(n+3)/2-d-2k\); then \(m-g+1=r\), and Riemann--Roch says

\[
h^0(A)=r+h^1(A).
\]

If the actual family dimension is \(r+\omega\), its characteristic \textbf{projective} series has dimension \(r+\omega-1\), whereas the complete projective series has dimension \(h^0(A)-1\). Therefore \(\omega\le h^1(A)\). The source's displayed phrase ``\(\ge r+\omega=m-\pi+\omega\)'' has a further off-by-one: with the corrected \(m\), \(r=m-\pi+1\). The intended characteristic dimension is \(r+\omega-1=m-\pi+\omega\).

Enriques correctly distinguishes a first-order deformation from an actual integrable deformation. He describes the possibility that local equations have dependent differentials while still cutting out the expected smaller-dimensional reduced locus. On p.~209 he explicitly labels the equality \(\omega=h^1(A)\) an \textbf{``Ipotesi di lavoro''} and says he has not proved it. It would be historically wrong to report that he silently assumed the disputed step in this 1949 presentation.

The unrestricted hypothesis is not a theorem that was later proved in full. Nodal/cuspidal equisingular strata can be nonreduced and can have incomplete characteristic series. Wahl (1974), \emph{Deformations of plane curves with nodes and cusps}, is the classical deformation-theoretic treatment; Guffroy (2003), \emph{Sur l'incomplétude de la série linéaire caractéristique d'une famille de courbes planes à nœuds et à cusps}, supplies explicit further examples. Vakil's (2006) Murphy's law, Main Theorem 1.1, parts M5a--M5b, shows the severity of possible singularities for simple covers and nodal/cuspidal plane-curve strata. Consequently a universal replacement must retain an obstruction space; one cannot repair the argument merely by declaring every first-order direction integrable. Singular points of a reduced stratum and generically nonreduced components are logically different phenomena, and the latter are particularly relevant to a hypothesis made at a generic member.

For a precise sheaf-theoretic formulation, Tannenbaum (1984, §1, Definition 1.1 and Remarks 1.2) identifies the equisingular normal sheaf, its tangent sections, and its obstruction cohomology. His Introduction also records the failure of the unrestricted characteristic-series assertion. In the nodal/cuspidal situation above that sheaf is \(\nu_*A\), as the displayed local ideals show.

A precise sufficient condition is \(H^1(C,A)=0\): then the equisingular deformation is unobstructed and of the expected dimension. In the nodal case, \(K_C-A=\nu^*\mathcal O_D(-3)\), so this group vanishes. With cusps, \(K_C-A=\nu^*\mathcal O_D(-3)+\sum R_j\), of degree \(k-3n\); thus \(k<3n\) is an elementary sufficient vanishing condition. This last bound is an explicit deduction in this commentary, not claimed as an optimal theorem from the literature. It also shows why elementary examples can support the hypothesis while more complicated ones invalidate it.

\section[I moduli delle superficie regolari]{I moduli delle superficie regolari di genere \texorpdfstring{$p>3$}{p > 3}}
\label{cp1-v-13}

This section is conditional throughout on V.12's completeness assumption. Enriques says so on p.~209, again on p.~210, and in his ``probabilmente'' conclusion on p.~213. His branch characteristic series pulls back to the ramification curve \(J=3H+K_X\) as the series cut by \(L=4H+K_X\), with the pinch-point divisor \(G\) removed. The geometric explanation on p.~210 is that polars through another point cut the first-polar series of the branch curve. One must distinguish the line bundle of the full infinitesimal series from the potentially smaller series of actual deformations. For a family of ordinary degree-\(n\) projected surfaces, the source also describes the surface characteristic system as cut by degree-\(n\) adjoints with the indicated tangency conditions at the pinch points. That comparison is sketched; it does not replace a proof about actual deformation families.

For sufficiently positive \(H\), the restriction and vanishing hypotheses allow

\[
\dim|4H+K_X|=p+4\pi+6n-4.
\]

If \(\tau=\deg G\) and the relevant speciality is \(\omega\), the characteristic family dimension is

\[
R=13p+4n-4\pi-2\lambda+19+\omega.
\]

Choosing a four-dimensional vector subspace of \(H^0(H)\) and spatial coordinates accounts for \(E_{\mathrm{mod}}=4r-12+15=4r+3=4p+4n-4\pi+7\). Thus \(R-E_{\mathrm{mod}}\) reproduces (V.11.1). At the preceding step the conditional characteristic dimension is \[
R-1=p+4\pi+6n-4-\tau+\omega.
\] The assumed system may be chosen with \(n>\pi-1\), together with the needed vanishing. The source also sketches an equivalence between completeness for ordinary projected surface deformations and the plane-branch hypothesis; it does not prove a full comparison of the two deformation functors. The vanishing and restriction statements are permissible after choosing a sufficiently high positive system; they are not automatic for every divisor merely because a net exists.

\subsection{The canonical calculation and its conflicting symbols}

Assume the canonical map has the generic birational properties required by the source, with \(p>3\). Put \(H=K_X\), so \(n=\lambda-1\), \(\pi=\lambda\), \(J=4K_X\), and by adjunction \(K_J=5K_X|_J\). Let

\[
\theta=\dim|G|_J=h^0(J,\mathcal O_J(G))-1.
\]

The group \(G\) lies in a canonical divisor on \(J\) and is therefore special. Serre duality gives \(h^1(J,K_J-G)=\theta+1\), and \(G\) imposes \(\tau-\theta\) independent conditions on the canonical series. Since \(q=0\), the restriction \(H^0(5K_X)\to H^0(K_J)\) is surjective with kernel \(H^0(K_X)\) of dimension \(p\). Using \(h^0(5K_X)=p+1+10K_X^2\) gives

\[
M=9p-2\lambda+16+\theta. \tag{V.13.1}
\]

This is the coherent meaning of the displayed formula on p.~212. The accompanying symbols are not coherent as printed: the original moduli excess in V.11 equals \(\theta+4\), so \(\theta=\omega_{\rm mod}-4\), whereas the speciality on the canonical ramification curve equals \(\theta+1\), so it cannot be that same \(\omega\). The printed sentence calls \(\theta\) the speciality and then calls it the projective dimension of \(|G|\); these differ by one. In particular, the first displayed characteristic-dimension formula on p.~212 must read

\[
\begin{aligned}
R-1&=p+4\pi+6n-4-\tau+\theta,\\
R&=13p-2\lambda+15+\theta,\qquad E_{\mathrm{mod}}=4p-1.
\end{aligned}
\]

If canonical speciality is denoted by \(\omega_{\rm can}=\theta+1\), the last term in the first equation is \(\omega_{\rm can}-1\), not \(\omega_{\rm can}\). Subtracting the displayed \(E_{\mathrm{mod}}\) gives (V.13.1). The commentary keeps three quantities distinct: branch-component excess \(\omega_{\rm mod}\), canonical speciality \(\theta+1\), and projective divisor-series dimension \(\theta\).

Keeping a canonical net fixed leaves a family, of dimension \(p-4\), of canonical systems of projective dimension three. Enriques argues that their pinch-point divisors on the fixed Jacobian curve move in \(|G|\), giving \(\theta\ge p-4\). The independence and genericity of this family are sketched, not developed as a deformation-theoretic proof. Inserting the inequality in (V.13.1) gives

\[
M\ge10p-2\lambda+12=10\chi(\mathcal O_X)-2K_X^2.
\]

\subsection{The Enriques--Kuranishi inequality}

The inequality survives without the false universal completeness hypothesis. The Kuranishi deformation space of a smooth surface has tangent space \(H^1(X,T_X)\) and an obstruction map into \(H^2(X,T_X)\); hence

\[
\dim\operatorname{Def}(X)
\ge h^1(T_X)-h^2(T_X)
=10\chi(\mathcal O_X)-2K_X^2+h^0(T_X).
\]

Here Hirzebruch--Riemann--Roch gives \(\chi(T_X)=2K_X^2-10\chi(\mathcal O_X)\). For a surface of general type the automorphism group is finite, so \(h^0(T_X)=0\), and the local moduli dimension has the stated lower bound. No claim that the obstruction map vanishes is needed. Catanese (2013), §1.3, arXiv version p.~13, explicitly identifies this as Enriques's inequality proved by Kuranishi. This repairs the numerical conclusion by a different theorem; it does not retroactively prove Enriques's completeness hypothesis.

\section{Nota storica e complementi}
\label{cp1-v-14}

Enriques reports that Noether (1888) anticipated equality \(M=10p-2\lambda+12\) by a negative-postulation parameter count. The count ``number of parameters minus number of conditions'' is a heuristic unless actual local equations, their codimension, and the relevant fibres have been controlled. Equality is stronger than the modern lower bound and is not universal.

The historical admission on p.~214 is unusually explicit. Enriques says his papers of 1908 and 1912 had overlooked the distinction between the differential property and completeness, and that Zariski's \emph{Algebraic surfaces} (1935), V §8, exposed the gap. He compares this with Severi's 1921 criticism of an earlier irregularity argument, to be discussed in Chapter IX §6. The last paragraph of p.~214 and the beginning of p.~215 distinguish the revised \emph{first} moduli bound, which uses actual branch-family dimensions and existence, from the \emph{second} bound still tied to the working hypothesis. Those distinctions must be preserved. The sentence on p.~214 that calls Noether's equality the case \(\omega=0\) needs the prime: it is \(\omega'=0\), the vanishing of the stronger excess in the immediately preceding formula. The scan confirms the missing prime; the transcription preserves the printed reading.

The further displayed expressions are conditional historical proposals:

\[
M=12p_a-2p_g-2\lambda+12+\omega,\qquad \omega\ge0,
\] \[
M\ge10p-2\lambda+11,
\] \[
M=9p_a+p_g-2\lambda+11+\theta.
\]

The last two come from B. Segre's claimed estimates \(\omega\ge p-1\) and \(\omega\ge2p_g-p_a-1\) for the excess in the first formula of V.11. Enriques explicitly says that Segre's approach retained the disputed completeness assumption. These are not independent proofs supplied here, and their individual excess symbols should not be identified across formulas without checking their definitions. In particular, the general modern Kuranishi bound does not automatically establish every sharper irregular estimate simply by renaming an excess term.

The closing Castelnuovo communication asserts \(\omega\ge2p_a+p_g+\lambda-I-3\). If \(I\) retains its chapter-wide meaning as the Zeuthen--Segre invariant, Noether gives

\[
2p_a+p_g+\lambda-I-3=p_g-10p_a+2\lambda-12.
\]

Inserted in \(M=10p_a-p_g-2\lambda+12+\omega\), this yields only \(M\ge0\). Thus, with that symbol identification, the displayed bound is numerically consistent but adds no positive lower bound beyond nonnegativity. The text gives no derivation and no published locator for the private communication. The commentary does not invent one.

\subsection{Excess dimensions and an explicit smooth deformation family}

The difference between an obstruction space and nonzero obstructions is visible in the example of a smooth sextic \(X\subset\mathbb P^3\). Here \(p_g=10\), \(K_X^2=24\), and the family of degree-six equations modulo coordinates has dimension \(\binom93-16=68\), while the Enriques--Kuranishi bound is \(110-48=62\). The normal sequence \[
0\longrightarrow T_X\longrightarrow T_{\mathbb P^3}|_X
\longrightarrow\mathcal O_X(6)\longrightarrow0
\] and the restricted Euler sequence give \[
\begin{gathered}
h^0(\mathcal O_X(6))=83,\qquad h^0(T_{\mathbb P^3}|_X)=15,\\
H^1(T_{\mathbb P^3}|_X)=0,\qquad h^2(T_{\mathbb P^3}|_X)=6.
\end{gathered}
\] The last number is the cokernel dimension for \[
H^2(\mathcal O_X)\longrightarrow H^2(\mathcal O_X(1))^{\oplus4}.
\] Its dual is the kernel of the surjective multiplication map \(H^0(\mathcal O_X(1))^{\oplus4}\to H^0(\mathcal O_X(2))\), from a space of dimension sixteen to one of dimension ten. Hence \(h^1(T_X)=68\) and \(h^2(T_X)=6\). Every first-order abstract deformation comes from an equation, and the smooth open parameter space of equations provides a smooth local deformation family. A finite stabilizer can still give a finite quotient in coarse moduli. Thus the obstruction space is nonzero although the Kuranishi family is smooth; Noether's predicted equality is not universal. The original branch excess of Section 11 is sixteen and the stronger excess is six in this example.

A useful continuing investigation, rather than an assertion of an open theorem, is to fix a nodal-cuspidal branch component and a monodromy representation and compare its branch and cover deformation functors. Actual dimension is bounded by tangent dimension, and higher-order obstructions determine whether equality holds. Likewise Chisini's proposed K3 induction calls for explicit braid and smoothing data: the listed cusp and node counts certify its numerical compatibility, not the existence of the smoothing.

\chapter[Superficie regolari]{Superficie regolari: minimo dei generi e condizioni di razionalit\`a}
\label{cp1-ch-vi}

The opening hypothesis of Chapter VI (pp.~217--235) is \(p_g=p_a=p\), equivalently \(q=0\). Enriques nevertheless points out when an argument does not require regularity. His initial reduction removes exceptional curves of the first kind, using Chapter IV, Section 9; thus \(p^{(1)}=K_X^2+1\) is the absolute linear genus on the nonruled minimal models under discussion. His \(p^{(2)}\) denotes \(K_X^2\), whereas \(P\), translated consistently as \(P_2\), denotes the bigenus. The plane computation later in the chapter concerns the specified plane model, not a model-independent value for every rational surface.

An effective canonical divisor of order zero is the zero divisor, so \(K_X\sim0\). A virtual canonical divisor of order zero has zero intersection with the chosen ample divisor; under the nefness hypotheses here this means \(K_X\equiv0\), and need not imply \(K_X\sim0\). The difference is decisive for the two possibilities \(K_X\sim0\) and \(K_X\not\sim0\), \(2K_X\sim0\). The genus of an arbitrary divisor remains its arithmetic genus \(p_a(D)=1+D(D+K_X)/2\); a singular model must be resolved and its exceptional curves accounted for before this convention is applied. Dolgachev (2016, Section 2) supplies modern context for the numerically trivial cases.

\section{Limite inferiore del genere lineare}
\label{cp1-vi-1}

\textbf{Statement.} If \(X\) is a smooth minimal projective surface and \(P_m(X)>0\) for some \(m\geq1\), then \(K_X^2\geq0\), hence \(p^{(1)}\geq1\). Regularity is unnecessary. Enriques explicitly removes that hypothesis on p.~219. He does not assert the converse: for example, the plane has \(K_X^2=9\) and every positive plurigenus vanishes.

\textbf{The component proof is valid.} First let \(D=K_X=\sum s_iC_i>0\). Put \(\nu_i=C_i^2\), \(\rho_i=p_a(C_i)\), and \(n_{ij}=C_iC_j\geq0\) for \(i\ne j\). Adjunction gives

\[
K_SC_i=2\rho_i-2-\nu_i.
\]

If this were negative, integrality and \(\rho_i\geq0\) would give \(\nu_i\geq-1\). The case \(\nu_i=-1\) forces \(\rho_i=0\); an integral curve of arithmetic genus zero on a smooth surface is a smooth rational curve, so this is precisely an excluded \((-1)\)-curve. Otherwise \(\nu_i\geq0\), and

\[
K_SC_i=s_i\nu_i+\sum_{j\ne i}s_jn_{ij}\geq0,
\]

a contradiction. Hence \(K_SC_i\geq0\) for every component and \(K_X^2=\sum s_iK_SC_i\geq0\). The displayed quadratic expression on p.~218 is

\[
K_X^2=\sum_i s_i^2\nu_i+2\sum_{i<j}s_is_jn_{ij},
\quad
\frac{\partial K_X^2}{\partial s_i}=2K_SC_i.
\]

Euler's homogeneous-function identity gives the same conclusion. It does \textbf{not} show that the intersection matrix is positive semidefinite for arbitrary coefficients; it only evaluates this form at the positive coefficients of the effective canonical divisor. Reading it that way removes an apparent, but nonexistent, gap.

For \(D\in|mK_X|\), the identical argument works: if \(K_SC_i<0\), adjunction and minimality give \(C_i^2\geq0\), whereas \(mK_SC_i=D C_i\geq0\). Curves outside the support of \(D\) also have \(K_SC\geq0\). Thus the argument actually proves that \(K_X\) is nef whenever an effective positive multiple exists on a minimal model. If \(D=0\), then \(mK_X\sim0\), so \(K_X^2=0\). This deals with the separately discussed ``order-zero'' case without regarding the empty divisor as an ordinary curve.

Enriques's calculation with \(C'=C+K_X\) is another correct expression of the latter case. If \(K_X\sim0\), then \(C'=C\), \(n=2\pi-2\), and

\[
p_a(C+K_X)=p_a(C)+p_a(K_X)+C K_X-1
\]

reduces to \(\pi=\pi+p^{(1)}-1\), giving \(p^{(1)}=1\). This is an identity of \emph{virtual arithmetic genera}: \(p_a(0)=1\).

\textbf{The second argument is a sketch, with movable parts requiring care.} On pp.~219--220, Enriques follows the adjoint sequence \(C,C+K_X,C+2K_X,\ldots\). Its intersections with \(K_X\) are

\[
K_X(C+jK_X)=K_X C+jK_X^2.
\]

If \(K_X^2<0\), these strictly decrease. Once a movable divisor \(L\) has \(K_X L<0\), the degrees \((L+jK_X)H\) eventually become negative for a suitable positive test class, so effective adjoints terminate. An effective \(mK_X\), conversely, would give injections \(H^0(C+jK_X)\hookrightarrow H^0(C+(j+m)K_X)\) and prevent termination. These are the mechanisms behind the argument. The prose suppresses the control of fixed components, the choice of a positive test curve, and the construction of an irreducible system after adding sufficiently positive divisors. Those points are not proved in this subsection. The second observation on p.~220 also asserts that an adjoint of a regular irreducible system without fundamental curves should not be reducible in the reductio argument. No unrestricted irreducibility theorem is proved there: sufficient positivity, Bertini, and control of fixed parts are required. Nor does the assertion that a multiple of a reducible system can be added to a suitable irreducible one give, without such control, a general method for eliminating fixed components. These are limitations of the second proof and its observations, not of the component proof.

The first observation on p.~220 must consequently not be read as a general assertion that the mere presence of an exceptional curve produces a \textbf{moving} system with negative canonical intersection and a terminating adjoint process. On the blow-up of a K3 surface, \(K_X=E\): every moving irreducible curve \(C\ne E\) has \(K_X C=EC\geq0\), while the isolated system \(|E|\) has \(E^2=-1>2p_a(E)-2=-2\) but its adjoints \(|(j+1)E|\) never terminate. The distinction between an isolated divisor and the movable systems needed in the proof matters. This is a limitation of the compressed second argument, not a counterexample to the theorem, already proved correctly by components. A modern minimal-model formulation supplies the missing framework; Hassett (2009, Corollary 2.20 and Theorem 2.26) is a verified reference.

\section{Il bigenere e i plurigeneri}
\label{cp1-vi-2}

\subsection{The lower bound and the degree-zero boundary}

Let \(H=C\) be a sufficiently positive smooth irreducible member, \(n=C^2\), \(\pi=g(C)\), and \(d=K_X C=2\pi-2-n\). The second adjoint is \(C''=C+2K_X\), and its restriction to \(C\) has degree \(2\pi-2+d\). If \(d>0\), it is nonspecial and has \(\pi-1+d\) independent sections. The exact sequence

\[
0\longrightarrow\mathcal O_X(2K_X)\longrightarrow
\mathcal O_X(C+2K_X)\longrightarrow
\mathcal O_C(C+2K_X)\longrightarrow0
\]

is the precise modern meaning of ``detaching \(C\).'' It bounds the number of conditions by \(h^0(C,\mathcal O_C(C+2K_X))\), without assuming that the restriction map is surjective.

The exceptional case is visible directly in surface Riemann--Roch and Serre duality:

\[
P_2=\chi(\mathcal O_X)+K_X^2+h^1(2K_X)-h^0(-K_X).
\tag{VI-RR2}
\]

When \(K_X H>0\), \(h^0(-K_X)=0\), so

\[
P_2\geq\chi(\mathcal O_X)+K_X^2=p_a+p^{(1)}.
\]

For a regular surface this is \(P_2\geq p+p^{(1)}\), the formulas on pp.~220--221. The argument on p.~220 initially says that the restricted series is certainly nonspecial even for \(d=0\); degree \(2\pi-2\) alone does not justify that assertion. The printed theorem on p.~221 also retains \(d\geq0\), exactly as transcribed. Enriques then explicitly separates the exceptional \(K_X\sim0\) case on the same page and treats nontrivial degree-zero canonical class on p.~222. The correct final reading is the qualified statement, not an unqualified theorem contradicted by his next paragraph. There is no transcription error in the sign \(\geq\).

\subsection{Equality and elliptic pencils}

If \(q=0\) and \(C_0\in|K_X|\) is an integral canonical divisor, adjunction and

\[
0\to\mathcal O_X(K_X)\to\mathcal O_X(2K_X)
\to\omega_{C_0}\to0
\]

show that restriction is surjective because \(H^1(K_X)\simeq H^1(\mathcal O_X)^*=0\). Since \(h^0(\omega_{C_0})=p_a(C_0)=p^{(1)}\), one obtains \(P_2=p+p^{(1)}\), as on p.~221. Smoothness is convenient but unnecessary if the divisor is integral and its dualizing sheaf is used.

For \(K_X^2=0\) and a nonzero moving pluricanonical system, the associated fibration has genus-one fibres. Its base is \(\mathbb P^1\) when \(q=0\): pullback injects \(H^1(B,\mathcal O_B)\) into \(H^1(X,\mathcal O_X)\). Enriques's assertion that the canonical system decomposes into elliptic members is a structural description, not a proof of semiampleness in all cases; when \(p_g=0\) there is no effective canonical system at all. The statement must concern the effective pluricanonical systems. The modern abundance and genus-one-fibration results cover that distinction (Reid, 1996, §§E.4--E.7).

If \(K_X\sim(p-1)F\) with \(F\) a fibre and no extra fixed part, then

\[
P_2=h^0(\mathbb P^1,\mathcal O(2p-2))=2p-1,
\]

giving Enriques's \(P_2-1=2p-2\). Multiple fibres can make a fixed component become movable after doubling; then \(P_2-1>2p-2\). These are effects of linear equivalence, not failures of the numerical bound.

For later verification, the precise characteristic-zero formula is useful. If \(f:X\to B\) is a relatively minimal genus-one fibration, \(F_j=m_jE_j\) its multiple fibres, and \(\mathcal L=(R^1f_*\mathcal O_X)^\vee\), then

\[
K_X\sim f^*(K_B+\mathcal L)+\sum_j(m_j-1)E_j,
\qquad \deg\mathcal L=\chi(\mathcal O_X).
\tag{VI-CBF}
\]

Thus, for \(B=\mathbb P^1\),

\[
P_m=\max\left\{0,\;1+m(\chi-2)+
\sum_j\left\lfloor\frac{m(m_j-1)}{m_j}\right\rfloor\right\}.
\tag{VI-F}
\]

The floor functions express the fact that only a whole multiple fibre moves. Reid (1996, §E.7.3) states the canonical-bundle formula and explains this divisor calculation; the applications to Enriques's particular examples below are independent computations. In positive characteristic, wild fibres require different coefficients; the characteristic-zero formula must not be transported unchanged.

\subsection{The two numerically trivial cases}

If \(q=0\) and \(K_X\sim0\), then \(p=1\), every \(mK_X\sim0\), and all \(P_m=1\). In (VI-RR2), the term \(h^0(-K_X)=1\) is precisely what defeats the preliminary bound \(P_2\geq p+p^{(1)}=2\). The adjoint identities \(|C|=|C'|=|C''|\), \(n=2\pi-2\), \(p^{(1)}=1\), and \(p-1=P_2-1=0\) on pp.~221--222 are all correct.

If \(q=p_g=0\) and \(K_X\equiv0\), then \(K_X\not\sim0\) and \(h^0(-K_X)=0\): a nonzero effective divisor has positive intersection with an ample divisor. Riemann--Roch for \(2K_X\) now gives \(h^0(2K_X)\geq1\). But an effective numerically trivial divisor must be zero. Hence \(2K_X\sim0\), \(P_2=1\), and \(K_X^2=0\). Therefore

\[
P_{2i}=1,\qquad P_{2i+1}=0.
\]

For sufficiently positive \(C\), the complete systems \(|C+K_X|\) and \(|C|\) differ because \(K_X\) is nontrivial, whereas \(|C+2K_X|=|C|\). This is the cohomological repair of the compressed restriction argument on p.~222: it does not assume, without proof, that two distinct systems have different restrictions to \(C\). Kumar (2008b, Lecture 17, Proposition 1) gives this order-two argument in modern notation. Such a complex surface is an Enriques surface; the \(K_X\sim0,q=0\) case is a K3 surface. These conclusions concern a minimal smooth model, not the branch-plane model before resolving and contracting exceptional curves.

\subsection{The plane calculation and its printed sign error}

On pp.~222--223, take \(C=7H\subset\mathbb P^2\). Then \(g(C)=15\), \(C^2=49\), \(K_X=-3H\), \(C'=4H\), and \(C''=H\). The line series on a septic is \(g^2_7\), with speciality

\[
h^1(C,\mathcal O_C(1))=h^0(C,\mathcal O_C(3))=10.
\]

Thus the nonspecial-series bound cannot be used. Enriques's genus addition calculation is \(0=p^{(1)}+3-12-1\), hence \(p^{(1)}=10=K_X^2+1\), while \(P_2=h^0(\mathbb P^2,-6H)=0\). All those numbers check. The virtual canonical system is \(|C_1-C_4|\), with \(K_X C_1=-3\) and \(K_X C_4=-12\). For negative canonical intersection the second-adjoint restriction has degree \(2\pi-2-d\), which is precisely why it may be special.

There is, however, a source sign misprint on p.~222: the scan prints \(-d=n-(2\pi-2)\) while calling this the negative canonical degree. With \(d>0\) as the following restriction calculation requires, the consistent formulas are

\[
d=n-(2\pi-2)>0,
\qquad K_X C=-d=2\pi-2-n.
\]

The transcription accurately preserves the printed error. The later \(-d=-3\) for a line on p.~223 independently confirms the intended convention.

\subsection{Higher plurigenera and the sixth plurigenus}

The square of an \(m\)-canonical divisor is \(m^2K_X^2=m^2(p^{(1)}-1)\). For \(m\geq2\) and \(K_X H>0\), the same calculation gives

\[
P_m=\chi(\mathcal O_X)+\binom m2K_X^2+h^1(mK_X)
\geq p_a+1+\binom m2(p^{(1)}-1).
\tag{VI-RRm}
\]

The restriction \(m\geq2\) is essential: setting \(m=1\) in the displayed lower bound on p.~223 would give \(p_g\geq p_g+1\) for a regular surface. It is implicit in the word ``plurigenera'' and in the extension from the bigenus. If \(q=0\) and \(C_0\in|(m-1)K_X|\) is integral, the exact sequence \(0\to\mathcal O_X(K_X)\to\mathcal O_X(mK_X)\to\omega_{C_0}\to0\) gives equality, since \(p_a(C_0)=1+\binom m2K_X^2\). This proves the particular equality Enriques states without needing a modern vanishing theorem. On a smooth minimal surface of general type in characteristic zero, nef-and-big vanishing gives the same equality for all \(m\geq2\), without the existence of an integral member at that particular level; this is an extension, not an extra hypothesis silently inserted into the old proof. Apply Esnault and Viehweg (1992, Corollary 5.12(c), pp.~48--49) with zero boundary and then Serre duality to \((m-1)K_X\), which is nef and big.

For \(q=p_g=0,P_2=1\) with the bicanonical divisor of positive order, (VI-RRm) gives \(P_m\geq1\) for every \(m\geq2\), in particular \(P_3\geq1\). Let \(D_2\in|2K_X|\) and \(D_3\in|3K_X|\). The two effective divisors \(3D_2,2D_3\in|6K_X|\) are distinct: equality would imply \(D_2=2E,D_3=3E\), and \(2(E-K_X)\sim3(E-K_X)\sim0\) would imply \(E\sim K_X\), contradicting \(p_g=0\). Thus \(P_6\geq2\), justifying the step only briefly indicated on p.~224. The statement that all plurigenera are positive refers to indices at least two; \(P_1=p_g\) remains zero.

Consequently, within the stated \(q=p_g=0,P_2=1\) class, \(P_3=0\) or \(P_6=1\) detects the order-two, numerically trivial canonical class, whereas positive-order bicanonical divisors give \(P_3\geq1,P_6\geq2\). This is the distinction carried forward into Chapter VII in Part 2.

\section[Esempi di piani doppi]{Esempi: piani doppi di genere lineare \texorpdfstring{$p^{(1)}=1$}{p(1) = 1}}
\label{cp1-vi-3}

\subsection{The imported classification (pp. 224--225)}

Enriques cites his \emph{Sui piani doppi di genere lineare \(p^{(1)}=1\)}, \emph{Rendiconti della R. Accademia dei Lincei}, April and May 1898. That is his reference, not a classification proved in these pages. The precise issue pagination and the completeness of that earlier classification have not been independently verified here. The present excerpt supplies examples and calculations.

A double plane branched over a smooth sextic has \(K_X\sim0\), \(q=0,p_g=1\), and all \(P_m=1\). If \(K_X^2=0\) and some \(P_m>1\), the pluricanonical genus-one fibration is preserved by the covering involution. On an invariant smooth genus-one fibre, the involution either has quotient \(\mathbb P^1\), with four branch points, or has a genus-one quotient, in which case the degree-two map is étale. Alternatively it interchanges fibres, inducing a degree-two map on the base. These are the geometric alternatives behind Enriques's two categories.

In category I, the fibres map to lines through \(O\). A branch curve of degree \(2n\) with an ordinary point of multiplicity \(2n-4\) leaves four branch points on a general such line. With multiplicity \(2n-3\), three residual branch points and one odd contribution over \(O\) again give four. Riemann--Hurwitz gives genus one. The asserted normal forms and allowed additional quadruple or \([3,3]\) singularities are imported from the 1898 analysis; the four-point calculation by itself does not prove their exhaustiveness.

In category II, the quotient pencil is a Halphen pencil: after nine blow-ups, it is \(|-sK_Y|\) on a rational surface \(Y\), with plane members of degree \(3s\) and nine points of multiplicity \(s\). Horizontal branch components would ramify a general fibre and change its genus, so the relevant branch components are fibres and possibly the distinguished anticanonical cubic. Cantat and Dolgachev (2012, Proposition 2.2 and §2.2) verify the modern Halphen description, including infinitely near base points and the unique multiple fibre of multiplicity \(s\) when \(s>1\). Enriques cites Enriques--Conforto, \emph{Le superficie razionali}, for this earlier material. The \(\kappa=0\) Enriques examples below also carry chosen genus-one pencils, although such a pencil is not supplied by their constant pluricanonical ring.

\subsection{Resolution and exceptional contributions}

For a smooth double cover \(\pi:\widetilde X\to Y\) determined by \(B\in|2L|\),

\[
\pi_*\mathcal O_{\widetilde X}=\mathcal O_Y\oplus\mathcal O_Y(-L),\quad
K_{\widetilde X}=\pi^*(K_Y+L),
\] \[
\chi(\mathcal O_{\widetilde X})=2\chi(\mathcal O_Y)+\tfrac12L(L+K_Y),\quad
K_{\widetilde X}^2=2(K_Y+L)^2.
\tag{VI-DC}
\]

These formulas follow directly from the local algebra \(\mathcal O_Y\oplus L^{-1}\), the ramification formula, projection formula, and Riemann--Roch on \(Y\). They also give

\[
H^0(\widetilde X,mK_{\widetilde X})=H^0(Y,m(K_Y+L))
\oplus H^0(Y,m(K_Y+L)-L).
\]

For singular branch curves one first uses the \textbf{canonical resolution}: blowing up a point of branch multiplicity \(b\) replaces \(L\) by \(\sigma^*L-\lfloor b/2\rfloor E\), and adds the exceptional curve to the new reduced branch divisor precisely when \(b\) is odd. This rule follows by removing the largest even power of a local exceptional equation from \(z^2=f\). On the smooth canonical resolution \(\widetilde X\), a branch \((-2)\)-curve pulls back to twice a \((-1)\)-curve; the latter must be contracted before computing the absolute \(p^{(1)}\). Enriques's ``impure'' canonical divisors comprise these exceptional contributions. His ``infinitesimal'' elliptic curves are actual exceptional curves on the resolved double cover, not formal infinitesimals in the modern scheme-theoretic sense.

An aligned \([3,3]\) cluster means a triple point followed by a triple point at the specified first infinitely near point, with the prescribed direction through \(O\). Assume that the stated branch singularities are ordinary in the specified successive blow-ups and that no additional singularity changes the invariants. Write \(H,E_0,E_i,E_i'\) for the line class and the total-transform exceptional classes over \(O,A_i,A_i'\). For branch degree \(2n\), multiplicity \(2n-4\) at \(O\), and \(t\) aligned \([3,3]\) clusters, one obtains

\[
\begin{aligned}
L&=nH-(n-2)E_0-\sum_{i=1}^t(E_i+2E_i'),\\
K_Y+L&=(n-3)(H-E_0)-\sum_{i=1}^tE_i'.
\end{aligned}
\]

Here the coefficient \(2E_i'\) is important: after the first odd triple blow-up the new branch also includes \(E_i\), making the second branch multiplicity four. Consequently

\[
K_{\widetilde X}^2=-2t,\qquad \chi(\mathcal O_{\widetilde X})=n-1-t.
\]

The strict transform \(H-E_0-E_i-E_i'\) of each branch line and the strict transform \(E_i-E_i'\) are disjoint branch \((-2)\)-curves. Their \(2t\) lifts are \((-1)\)-curves. Contracting them gives \(K_X^2=0\). The resulting pencil has \(t\) multiple fibres of multiplicity two. For the examples in the text, the canonical adjoint calculation gives \(q=0\), hence \(p_g=n-2-t\). Formula (VI-F) becomes

\[
P_m=\max\{0,\ 1+m(n-3-t)+t\lfloor m/2\rfloor\}.
\tag{VI-I}
\]

This computation both verifies the numerical examples and makes their genericity assumptions explicit.

\begin{center}
\begin{tabular}{lrrrrr}
\hline
Branch & $p_g$ & Double fibres & $P_2$ & $P_3$ & $P_6$\\
\hline
$D_8$, $O^4$, $t=1$ & $1$ & $1$ & $2$ & $2$ & $4$\\
$D_8$, $O^4$, $t=2$ & $0$ & $2$ & $1$ & $0$ & $1$\\
$D_{10}$, $O^6$, $t=2$ & $1$ & $2$ & $3$ & $3$ & $7$\\
$D_{10}$, $O^6$, $t=3$ & $0$ & $3$ & $2$ & $1$ & $4$\\
\hline
\end{tabular}
\end{center}

In each row \(q=0\). The octics occur on pp.~225--226; the two decics occur on pp.~226--227.

Entries not explicitly computed by Enriques are consequences of (VI-I), included as our extension of his checks.

\textbf{The first octic (pp.~225--226).} Bézout forces the line \(r=OA\) into the branch curve: its required intersection multiplicities exceed eight. The residual septic has \(O^3A^2A_1^2\). The canonical adjoint is the unique line \(r\) through \(O,A\), so \(p_g=1\), and (VI-DC) gives \(\chi=2\), hence \(q=0\). After the two exceptional curves are contracted, the canonical divisor is the reduced half-fibre \(E\), with \(2E\sim F\). Thus \(|2K_X|=|F|\), \(P_2=2\), and \(p^{(1)}=1\). The apparent disconnected rational canonical image on the singular plane model is not the canonical divisor on the smooth minimal model.

For completeness, the conics through \(A,A_1\) have class \(\Gamma=2H-E_A-E_{A_1}\), \(\Gamma^2=2\), and \(\Gamma B=10\). Their double covers have genus \(4\) by Riemann--Hurwitz, self-intersection \(4\), and canonical intersection \(2\); bicanonical intersection is \(4\). They form the stated \(\infty^3\), mapping to a double quadric cone with branch degree ten. These numerical claims check. The further construction projects \(D_{10}\) from a triple point \(P\) to the residual septic \(D_7\), while an infinitely near triple point \(P_1\) maps to \(O\). The direction \(PP_1\) lies in the tangent plane along the relevant ruling, and planes through it cut the bicanonical images. This is a geometric construction sketched on p.~226; it is not independently developed there into a full resolution diagram.

\textbf{The second octic and an actual numerical error (p.~226).} With two clusters, Bézout gives

\[
B=r+s+D_6(O^2A^2A_1^2B^2B_1^2).
\]

There is no canonical adjoint line through \(O,A,B\) when these are not collinear, so \(p_g=0\); \(\chi=1\), \(q=0\), and \(2K_X\sim0\) with \(K_X\not\sim0\). Thus \(P_2=1,P_3=0\), exactly as Enriques says. His bicanonical ``impure'' divisor consists of twice the four exceptional curves before contraction; afterwards it is zero.

The four exceptional curves in the canonical resolution used here include two extra resolution \((-1)\)-curves over the clusters. Contracting these first leaves the two exceptional branch-line curves counted in Enriques's description; the difference in the number of curves is a difference of intermediate models, not a further numerical error.

The next sentence is internally inconsistent. The scan really says that the \(\infty^4\) of cubics through \(O,A,A_1,B,B_1\) gives genus \(4\), degree \(6\), and a \textbf{double cubic} surface, normal in \(\mathbb P^4\), with branch degree \(10\). The indicated full cubic system has instead

\[
\begin{aligned}
\Gamma&=3H-E_0-E_A-E_{A_1}-E_B-E_{B_1}=-K_Y,\\
\Gamma^2&=4,\qquad p_a(\Gamma)=1,\qquad \Gamma B=8.
\end{aligned}
\]

Its pullback descends after contracting the four exceptional curves and has square \(8\) and genus \(5\). It has five sections, so its projective dimension is indeed four. The image of \(Y\) is a four-nodal \textbf{quartic} del Pezzo surface in \(\mathbb P^4\), and the one-dimensional branch curve has degree \(8\), with the nodes also involved in the double-cover construction. Thus, for the construction actually named, the consistent tuple is

\[
(\text{genus},\text{degree},\text{image degree},\text{branch degree})
=(5,8,4,8),
\]

not \((4,6,3,10)\). This repair is independently confirmed by Dolgachev (2016, §4, p.~13 in the author PDF): he describes exactly the anticanonical map from the blow-up of the five singular points of an Enriques octic, the four-nodal quartic, and branching in \(|\mathcal O_D(2)|\). The printed tuple is recorded above; the correction concerns the construction explicitly named in the text. Which different construction, if any, Enriques had in mind for the erroneous tuple is not established.

The promised sextic model with double tetrahedral edges is a forward reference, not proved on p.~226. Its birational construction is deferred beyond the present chapter.

\textbf{Two decic examples (pp.~226--227).} For \(D_{10}\) with \(O^6\) and two clusters, the forced lines \(r,s\) leave \(D_8(O^4A^2A_1^2B^2B_1^2)\). The canonical adjoint is the conic \(r+s\); after removing exceptional contributions the canonical divisor consists of the two reduced double fibres. With \(2E_A\sim2E_B\sim F\),

\[
K_X\sim E_A+E_B,\qquad 2K_X\sim2F.
\]

Hence \(p_g=1\), \(P_2=3>p+p^{(1)}=2\), and the bicanonical net corresponds to pairs of lines through \(O\). Linear equivalence must be retained: although \(K_X\equiv F\), it does not follow that \(K_X\sim F\); the latter would wrongly give \(p_g=2\).

Adding a third cluster forces the third branch line. There is no conic with a double point at \(O\) passing through three distinct directions \(A,B,E\): such a conic is a union of two lines through \(O\). Thus \(p_g=0\) and, by \(\chi=1\), \(q=0\). With three double fibres, the canonical formula is

\[
K_X\sim-F+E_A+E_B+E_E.
\]

It follows that \(2K_X\sim F\), \(3K_X\sim E_A+E_B+E_E\), and \(6K_X\sim3F\). These identities establish \(P_2=2,P_3=1,P_6=4\), and explain the three double members of the bicanonical pencil and the image by triples of lines. They verify every plurigenus asserted in this paragraph.

\subsection{Halphen covers and an endpoint missed in the text (pp. 227–229)}

Let \(Y\) be a Halphen surface of index \(h\), with \(J=-K_Y\) represented by its distinguished cubic. Then \(J^2=0\), the pencil is \(|hJ|\), and for \(a\geq0\),

\[
h^0(Y,aJ)=\lfloor a/h\rfloor+1;
\qquad h^0(Y,aJ)=0\quad(a<0).
\tag{VI-H}
\]

One proof is to use the unique multiple fibre \(hJ\), strip off the fixed remainder of \(aJ\), and identify the movable part with the pullback of the line bundle of degree \(\lfloor a/h\rfloor\) on \(\mathbb P^1\). This is exactly the feature recorded by the normal bundle of \(J\), of order \(h\) in \(\operatorname{Pic}^0(J)\); merely knowing the branch degree and nine multiplicities without the irreducibility and pencil information is not enough to fix all plurigenera. In the source's examples, however, the stated irreducibility actually forces the required index, so this is not an additional hypothesis missing from these particular examples. Indeed, write \(a=qh+r\) with \(0\leq r<h\). Every member of \(|aJ|\) is the sum of \(q\) fibres and the fixed divisor \(rJ\). An irreducible reduced member with \(a>1\) therefore forces \(q=1,r=0\), hence \(a=h\). Applied to the irreducible \(D_{6n}\) and \(D_{6n-3}\), this gives \(h=2n\) and \(h=2n-1\), respectively.

\textbf{First family: an irreducible \(D_{6n}\), \(n>1\).} The nine prescribed points have ordinary multiplicity \(2n\). Here \(h=2n\), \(B\sim2nJ\), and the covering line bundle is \(nJ\). Therefore

\[
K_X=(n-1)\pi^*J,\qquad \chi(\mathcal O_X)=2,
\] \[
P_m=h^0(Y,m(n-1)J)+h^0(Y,(m(n-1)-n)J)
=\left\lfloor\frac{m(n-1)}n\right\rfloor+1.
\tag{VI-H1}
\]

In particular \(p_g=1,q=0,K_X^2=0,P_2=2\). The inverse image \(E=\pi^*J\) is the connected étale double cover of the distinguished cubic: the restriction of \(nJ\) to \(J\) is its nontrivial order-two line bundle. A general fibre of the genus-one pencil on \(X\) has class \(F=nE\). General fibres on \(Y\) split into pairs on \(X\), which is the geometric point of the Comessatti remarks in the footnotes. This splitting follows from the degree-two Stein factor on the base, or equivalently from the square criterion for the restricted cover; parity of intersection multiplicities alone is insufficient.

Thus \(2K_X\sim F+(n-2)E\), exactly explaining the moving elliptic component and fixed part of multiplicity \(n-2\) described on p.~228. But the printed statement that the fixed component is a \(1/(n-1)\)-part of \(L\) is wrong: it is a \(1/n\)-part. The source prints \((n-1)\)-ma, as does the transcription. Moreover,

\[
P_3=\begin{cases}2,&n=2,\\3,&n\geq3.\end{cases}
\]

The displayed \(P_3=3\) on p.~228 overlooks the allowed endpoint \(n=2\). More precisely, \(3K_X=2F+(n-3)E\) only when \(n\geq3\); at \(n=2\) one has \(3K_X=F+E\). Even for \(n>3\), ``pairs of curves \(L\)'' describes the movable part and leaves a fixed \((n-3)E\) unmentioned. These corrections are our explicit double-cover computation, not a claim that a named later author published an erratum to this page.

\textbf{Second family: \(D_{6n-3}+D_3\).} The irreducible component has multiplicity \(2n-1\) at the nine Halphen points. Now the Halphen index is \(2n-1\), the branch divisor is \(G+J\sim2nJ\), and \(J\) itself is branched. If \(E\) is its reduced inverse image, then \(\pi^*J=2E\), a general fibre is \(F=(2n-1)E\), and

\[
K_X=2(n-1)E,\qquad
P_m=\left\lfloor\frac{m(2n-2)}{2n-1}\right\rfloor+1.
\tag{VI-H2}
\]

For \(n>1\), this gives \(p_g=q+p_a=1\) with \(q=0\), \(K_X^2=0\), and \(P_2=2\), verifying the claims actually printed at the bottom of p.~228. Here \(2K_X=F+(2n-3)E\). The final sentence of the printed paragraph describes only the moving fibre and the doubled image \(D_3^{n-2}\), accounting for \(F+2(n-2)E\); it omits one reduced ramification cubic \(E\). This is visible already in the preceding construction: the ramification divisor is \(C_G+E\), and adding \(\pi^*((n-2)J)\) gives \(C_G+(2n-3)E\). At \(n=2\), the fixed part is still \(E\), whereas the final printed exponent would leave no fixed part. This is a source omission in the conclusion, not a transcription error; the preceding ramification construction and the asserted \(P_2=2\) are correct.

\textbf{The irregular extensions.} The sentence on p.~229 allows branch divisors composed of several Halphen fibres and says that in general \(p_a=1,p_g>1,p^{(1)}=1\). Under the stated vertical-cover assumptions, (VI-DC) still gives \(\chi=2\), and the Stein base of the lifted pencil can have genus \(b>0\). Then \(q=b\), \(p_g=1+b\), and \(K_X^2=0\). The source makes a broad structural remark, not a classification of every possible singular fibre configuration. Its first subscript is \(a\); it must not be changed to \(i\) or to \(g\).

\subsection{The decic with a nonic and a line (pp. 229--230)}

Enriques cites two papers of L. Campedelli in the \emph{R. Accademia dei Lincei} (1932), whose titles are recorded in the references. The earlier existence and classification are assumed here. The stated branch is a Halphen nonic with nine triple points, three pairs infinitely near and the other three on \(r\), together with \(r\). The points \(A,B,E\) are not collinear. Enriques describes this model as falling outside the stated second-category form because of its exceptional branch line; that component must be removed before comparing the minimal forms. Any candidate canonical conic through the six specified proper points would contain \(r\), leaving an impossible line through \(A,B,E\); hence \(p_g=0\). The branch line has self-intersection \(-2\) on the resolved base and lifts to an exceptional \((-1)\)-curve. Removing it is necessary for the absolute genus calculation. The hypothetical line through \(A,B,E\) used in the virtual canonical discussion would have four branch points, the three odd contributions and the intersection with \(r\), so its double would have genus one by Riemann--Hurwitz; its actual nonexistence is the point of the noncollinearity argument. A direct invariant calculation supplies the other needed information: a smooth decic double plane starts with \(\chi=7,K_X^2=8\). The three proper triple points of the nonic on \(r\) are quadruple points of the total branch, and the second point in each of the three \([3,3]\) clusters also has total branch multiplicity four after the first odd blow-up. These six multiplicity-four steps each lower \(\chi\) by one and \(K_X^2\) by two. The intervening multiplicity-three steps change neither. Hence the smooth canonical resolution has \(\chi=1,K_{\widetilde X}^2=-4\). The branch line and the three first exceptional branch \((-2)\)-curves lift to four \((-1)\)-curves; contracting them gives \(K_X^2=0\). With \(p_g=0\), this also proves \(q=0\).

On the minimal surface, let \(C_2\) be the irreducible bicanonical elliptic curve represented by the cubic \(D_3\), and \(C_3\) the tricanonical elliptic curve represented by \(D_9\). The geometric description yields

\[
2K_X\sim C_2,\qquad3K_X\sim C_3,\qquad
6K_X\sim3C_2\sim2C_3=F.
\]

Thus the fibration has a triple fibre \(3C_2\) and a double fibre \(2C_3\), with \(\chi=1\). Formula (VI-F) gives

\[
P_m=\max\{0,\;1-m+\lfloor m/2\rfloor+\lfloor2m/3\rfloor\},
\]

so \(P_2=P_3=1,P_6=2\), in agreement with the stated invariants. These numbers provide a concrete example of the positive-order case distinguished in §2. The isolation of \(C_3\) can also be justified without invoking the whole formula: if \(|3K_X|\) moved and a member passed through a point of \(C_2\), the zero intersection \(C_2(3K_X)=0\) would force containment of \(C_2\), leaving an effective canonical divisor, a contradiction.

The sentence near the end of p.~230 nevertheless calls the unique tricanonical image \(D_3\). The scan confirms that subscript. It is a source slip: \(D_3\) has just been identified as bicanonical, while \(D_9\) is tricanonical. The claimed \(P_3=1,P_6=2\) is unaffected. The statements about irreducibility of the lifted cubic and nonics rely on Comessatti's criterion, cited through Enriques--Chisini, Book V, Chapter IV, §38, vol.~III, p.~445. The commentary checks their compatibility with the divisor and fibration calculation; it does not supply new equations for the branch curve or independently re-prove the existence of the required nine-point configuration.

\section{Condizioni di razionalit\`a di una superficie}
\label{cp1-vi-4}

\subsection{Exact statement and what the excerpt establishes}

For a smooth projective surface over an algebraically closed field of characteristic zero,

\[
X\text{ is rational}\quad\Longleftrightarrow\quad q(X)=0,\ P_2(X)=0.
\]

Equivalently, in Enriques's notation, \(p_a=P_2=0\). Indeed a nonzero canonical section has a nonzero square, so \(P_2=0\) implies \(p_g=0\); hence \(p_a=p_g-q=0\) then implies \(q=0\). Rationality gives the converse vanishing by birational invariance and the direct calculation on \(\mathbb P^2\). Vanishing of \(p_g,q\) alone is insufficient, as the preceding Enriques examples show.

The source credits Castelnuovo's \emph{Sulle superficie di genere zero} (1896), citing \emph{Memorie scelte}, p.~307. It explicitly uses earlier rationality theorems of Noether and Clebsch--Noether, referring to Enriques--Conforto. The present argument is substantial: it divides into two cases, analyses the last adjoints, and treats a fixed-component case. It is not merely an assertion of the criterion. Nevertheless some residual-system and irreducibility steps are compressed into assertions, so the historical proof should not be labelled fully expanded in modern foundations.

\subsection{Case I: a system of negative canonical intersection (pp. 231–234)}

Take an irreducible moving system \(C\) with \(n=C^2>2\pi-2\); equivalently \(K_X C<0\). After making it a suitable hyperplane system, put \(d=n-(2\pi-2)>0\). The adjoints \(C+jK_X\) have degrees against \(C\)

\[
2\pi-2,\quad2\pi-2-d,\quad2\pi-2-2d,\ldots.
\]

A positive ample degree cannot become negative for an effective divisor, so this process terminates. The proof assumes that the last system of positive dimension can be tracked through removing fixed components and, if necessary, replacing a system by its moving part.

If the last irreducible system has no positive-dimensional adjoint, its genus is zero or one. The relevant cohomological fact is

\[
0\longrightarrow\mathcal O_X(K_X)\longrightarrow
\mathcal O_X(K_X+D)\longrightarrow\omega_D\longrightarrow0.
\]

For \(p_g=q=0\) and a smooth irreducible \(D\), it yields \(h^0(K_X+D)=g(D)\). Thus an adjoint of projective dimension at most zero means \(g(D)\leq1\). This fills in the numerical step on p.~231.

If the genus is zero, a pencil of rational curves makes the surface rational. On resolving its base points, the generic fibre is a genus-zero curve over \(\mathbb C(t)\); the Noether--Enriques theorem, or the existence of a rational section over this field, gives a birational ruled surface over \(\mathbb P^1\). Enriques assumes that result from earlier work.

For an elliptic complete system of dimension \(r\geq2\), the sequence \(0\to\mathcal O_X\to\mathcal O_X(D)\to\mathcal O_D(D)\to0\) and \(q=0\) show that its characteristic series is complete. Since \(D^2>0\), curve Riemann--Roch gives \(r=D^2\). Imposing \(r-2\) general simple base points leaves a net of degree two on the blown-up surface. Its general member has genus one, so the resulting double-plane branch has degree four. This is the classical Clebsch case, whose rationality is imported. ``Containing a net of degree two'' includes the imposed base points; merely taking an arbitrary three-dimensional subspace would not lower \(D^2\).

If the final elliptic system is a pencil, Enriques passes to the preceding genus-two system, of dimension at least two. He argues that the hyperelliptic conjugate of a point remains fixed as the curve varies through that point; otherwise it traces a rational component of the elliptic pencil. This is a geometric sketch of the existence of a surface involution. For an adjoint \(D=K_X+C\), \(p_g=q=0\) and smooth connected \(C\), the same cohomology sequence gives \(H^1(X,D)=0\): the map \(H^1(\omega_C)\to H^2(K_X)\) is the trace isomorphism. Hence its characteristic series is not only complete but nonspecial.

On a hyperelliptic curve of genus \(g\), a complete series composed with the hyperelliptic \(g^1_2\) is a multiple \(a g^1_2\) only while no extra anti-invariant sections appear. If it is also nonspecial, the borderline value is \(a=g\), giving degree \(2g\) and projective dimension \(g\). For a surface system of dimension \(r\), the characteristic projective dimension is \(r-1\), so

\[
g=r-1,\qquad D^2=2(r-1).
\]

For \(g=2\), this gives \(r=3,D^2=4\), exactly the \(\infty^3\) printed on p.~232. The map has degree two onto a quadric. In the branch configuration retained after excluding rational pencils, the quadric is a cone; its ruling has three moving branch points and the vertex supplies the fourth. A smooth quadric would have branch bidegree \((a,b)\) with even entries and \(a+b=6\); one ruling then has at most two branch points, returning to a rational pencil already disposed of. This gives the missing qualification behind the short parity argument in the text: degree six alone does not force bidegree \((3,3)\). The double-cone model is birational to Noether's double plane with a sextic having an infinitely near pair of triple points. Its rationality is again assumed, not established anew in this subsection. The source also assumes the earlier existence of the particular double-plane models; the divisor calculation is not an independent construction of their defining equations.

\subsection{Fixed parts and composite adjoints (pp. 232–234)}

If an adjoint has fixed part \(\theta\) and an irreducible moving part of dimension at least two, Enriques says that one can continue after stripping \(\theta\). This requires control of the cohomology and arithmetic genus of the new system; a fixed-part removal does not automatically preserve all the vanishing statements just proved for an actual adjoint. The excerpt does not supply that control in full generality. Under the required completeness and nonspeciality hypotheses, the moving rational and elliptic systems have degree \(r-1\) and \(r\), respectively, by curve Riemann--Roch. These numerical assertions survive; the missing point is their applicability after arbitrary fixed parts have been removed.

For a moving part composed with a pencil \(|L|\), regularity makes the base rational. Rational members settle rationality immediately. Otherwise, let their genus be \(\rho>0\). The fixed part is fundamental for the preceding irreducible system: the restricted canonical series has no base points, so \(\theta\) cannot meet a general preceding member. Restriction of the complete canonical series to that member is composed with the pencil. A degree-one restriction would return to the rational case already treated; in the remaining situation its degree \(\nu>1\) is necessarily two, since a canonical map of a nonhyperelliptic curve is birational. The preceding curves are therefore hyperelliptic, and the pencil consists of their bisecants.

Using completeness and nonspeciality again, the text obtains

\[
n_{h-1}=2(r_{h-1}-1),\qquad
\pi_{h-1}=r_{h-1}-1.
\]

Subtracting \(L\) is said to impose two conditions and to lower the genus by one. If the residual curve has the required connectedness and arithmetic genus, the addition formula

\[
\pi_{h-1}=(\pi_{h-1}-1)+\rho+i-1
\]

with \(i=L\cdot(\text{residual})\geq1\) gives \(\rho+i=2\), hence \(\rho=i=1\). These numerical deductions are correct. The omitted issue is establishing those residual hypotheses and the asserted behaviour after all fixed parts have been removed. Connectedness of the original member alone should not be substituted for an argument about every residual system.

On an elliptic curve a degree-one complete linear series is fixed. Identifying that fixed point with the pencil base point is part of the compressed incidence argument. Enriques uses it to conclude that residual curves split into members of the pencil and that the degree-two image is a cone of degree \(m=r_{h-1}-1\) in \(\mathbb P^{m+1}\). The hyperelliptic branch degree on a general hyperplane section is \(2+2\pi_{h-1}=2m+2\). Each ruling must contain at least three moving branch points if its double cover has positive genus. Therefore

\[
3m\leq2m+2,
\]

so \(m\leq2\). The excluded plane/rational cases leave \(m=2\), the quadric-cone case already treated. This verifies the arithmetic of pp.~233--234 while identifying the geometric assertions that remain sketched.

\subsection{Case II and the minimal-model proof}

If every relevant irreducible moving system satisfies \(K_X C\geq0\), Enriques uses §2 to deduce \(P_2\geq p^{(1)}\). Since \(P_2=0\), \(p^{(1)}\leq0\), so \(K_X^2\leq-1\). Successive adjoints decrease intersection with \(K_X\) by \(K_X^2\); either an adjoint violates the assumed nonnegativity, or the process terminates and returns to Case I. His final assertion that an appropriate multiple of a reducible system can be combined with a regular system to produce an irreducible system is another indicated construction rather than a proof in the excerpt.

A modern proof avoids these residual-system choices. Contract \((-1)\)-curves. The surface minimal-model theorem gives a minimal model with \(K_X\) nef, or a geometrically ruled/rational model. Suppose \(K_X\) were nef. Because \(P_2=0\), \(p_g=0\); together with \(q=0\), this gives \(\chi=1\). Riemann--Roch for \(-K_X\) and Serre duality yield

\[
h^0(-K_X)+h^0(2K_X)=1+K_X^2+h^1(-K_X)\geq1.
\]

Thus \(-K_X\) is effective. A nef \(K_X\) has nonnegative intersection with an ample class, whereas a nonzero effective \(-K_X\) has positive ample degree, a contradiction. If the divisor is zero instead, then \(K_X\sim0\), contradicting \(P_2=0\). Hence the minimal model is ruled or rational. For a ruled model over \(B\), \(q=g(B)\); therefore \(B\simeq\mathbb P^1\), and the surface is rational.

This is a complete proof \textbf{relative to the surface minimal-model theorem}, not a claim to re-prove that theorem in a paragraph. Hassett (2009, Theorem 2.26, p.~178, with Corollary 2.20) supplies this modern formulation; Kumar (2008a, Lecture 9, Theorem 1 and its reductions) supplies an alternative proof closer to the classical adjunction method. Zariski's \emph{On Castelnuovo's criterion of rationality \(p_a=P_2=0\) of an algebraic surface} (1958) is the precise later reference for the characteristic-\(p\) rationality theorem, as explicitly recorded by Hassett (2009, Remark 2.29). It is an extension and foundational treatment, not evidence that the characteristic-zero theorem was false. The accompanying \emph{The problem of minimal models in the theory of algebraic surfaces} (1958) addresses the minimal-model foundations. These references are not being used to attribute the local misprints on pp.~222, 226, or 228 to a documented historical erratum.

The field must be algebraically closed. For example, a smooth anisotropic quadric over \(\mathbb R\) can have \(q=P_2=0\) and no real point, hence cannot be \(\mathbb R\)-rational. The classical criterion concerns geometric rationality, not arbitrary-field rationality. In positive characteristic the cohomological \(q=h^1(\mathcal O)\) must also be distinguished from the dimension of a possibly nonreduced Picard scheme; none of the complex-analytic identifications should be silently transferred.

The contrast between the two principal arguments is precise. The component proof of \(K_X^2\ge0\) and the distinct-products proof of \(P_6\ge2\) are complete once their stated hypotheses are imposed. The adjoint-system proof of rationality additionally needs the geometric and cohomological control of moving and residual systems identified above. The original attribution remains Castelnuovo's method with Enriques's simplified details, rather than a claim that this commentary has independently compared every step with the 1896 paper.

\begin{remark}[Further calculations already determined by the examples]
The four category-I configurations, both Halphen families, and the final
$(2,3)$-multiple-fibre example determine every $P_m$ through the displayed
floor-function formulas. They give periodic corrections to linear growth,
and identify the exceptional $n=2$ value in the first Halphen family. A
resolution graph recording branch parity at each blowup and the contraction
of each lifted branch $(-2)$-curve would make the different intermediate
exceptional-curve counts explicit; the existence and genericity of any
chosen branch equation would remain a separate task. The origin of the
incorrect tuple on p. 226 is a textual-history question: the named cubic
system determines the corrected quartic construction, but these pages do
not show whether the printed tuple came from another construction or from
several printing errors.
\end{remark}

\section*{References}
\markboth{References}{References}

The list consolidates the works cited throughout Part 1. Statements about external texts being checked, inspected, or unavailable record how far those texts were consulted for this commentary; the book's disputed readings are checked against the published edition and the scan of the print. Historical citations that the source gives only in abbreviated form remain explicitly qualified. An unverified title, date, volume, or page range has not been supplied merely to complete a bibliographic field. Distinct editions, offprints, and similarly titled works are identified where they affect the discussion.

\begin{sloppypar}
\begin{itemize}
\item Albanese, G. (1924). Trasformazione birazionale di una superficie algebrica qualunque in un'altra priva di punti multipli. \emph{Rendiconti del Circolo Matematico di Palermo}. Historical reference printed on p.~6; bibliographic fields not established here are left unspecified.

\item Alzati, A., \& Tortora, A. (2002). On connected divisors. \emph{Advances in Geometry, 2}(3), 243--258. \href{https://doi.org/10.1515/advg.2002.012}{DOI: \nolinkurl{10.1515/advg.2002.012}}. Background terminology; the exceptional-divisor proof in the commentary is given explicitly.

\item Baker, H. F. (1936). On the proof of a lemma enunciated by Severi. \emph{Mathematical Proceedings of the Cambridge Philosophical Society, 32}(2), 253--259. \href{https://doi.org/10.1017/S0305004100001808}{DOI: \nolinkurl{10.1017/S0305004100001808}}. The publication record and abstract were checked, not the full original proof.

\item Barber, S. F., \& Zariski, O. (1935). Reducible exceptional curves of the first kind. \emph{American Journal of Mathematics, 57}(1), 119--141. \href{https://doi.org/10.2307/2372025}{DOI: \nolinkurl{10.2307/2372025}}. Bibliographic identity checked against the journal issue; the whole paper was not re-proved or fully read.

\item Bauer, I., \& Catanese, F. (2021). \(\mathfrak S_5\)-equivariant syzygies for the Del Pezzo surface of degree 5. \emph{Rendiconti del Circolo Matematico di Palermo, Series 2, 70}, 97--120. \href{https://doi.org/10.1007/s12215-020-00483-9}{DOI: \nolinkurl{10.1007/s12215-020-00483-9}}. Published online in 2020; volume year 2021. Section 2 and the abstract verify the anticanonical model and its automorphism group.

\item Beauville, A. (1978). \emph{Surfaces alg\'ebriques complexes}. Ast\'erisque, 54. Soci\'et\'e Math\'ematique de France. Especially I.12--I.16, III.20, VIII.2--VIII.5, and VIII, Exercise 11. \href{https://www.numdam.org/item/AST_1978__54__R1_0/}{NUMDAM record}. Additional retained locators are Theorems I.11--I.12, Sections I.14--I.17, and Chapter II.

\item Beauville, A. (1996). \emph{Complex algebraic surfaces} (2nd ed.). Cambridge University Press. \href{https://doi.org/10.1017/CBO9780511623936}{DOI: \nolinkurl{10.1017/CBO9780511623936}}. Evidence: publisher's book and chapter metadata; companion author's lecture text below inspected for formulas, blowups, and contraction statements. The commentary on Chapter IV uses Theorems I.11--I.12 and Remarks I.13, available in the \href{https://assets.cambridge.org/97805214/95103/excerpt/9780521495103_excerpt.pdf}{publisher excerpt}.

\item Beauville, A. (2020, July). \emph{Algebraic surfaces} {[}Lecture slides{]}. Université Côte d'Azur. \href{https://math.univ-cotedazur.fr/~beauvill/conf/Shanghai.pdf}{Full text}. Locators: Lecture I, slides 2--12, intersection and genus formulas; Lecture II, slides 18--27, blowups and minimal surfaces. This is an auxiliary accessible primary exposition, not a substitute for exact source-page citations to Enriques.

\item Bertini, E. (1882). Sui sistemi lineari. \emph{Rendiconti del Reale Istituto Lombardo, Series 2, 15}, 24--28. The December 1880 date used by Enriques is distinguished from the 1882 publication date by Kleiman (1998).

\item Bertini, E. (1894). La geometria delle serie lineari sopra una curva piana secondo il metodo algebrico. \emph{Annali di Matematica Pura ed Applicata}, series II, \emph{22}. Enriques cites Section 21(b), (g), on pp. 64--65. The original priority comparison was not independently verified.

\item Bonnesen, T. (1906). Sur les s\'eries lin\'eaires triplement infinies de courbes alg\'ebriques sur une surface alg\'ebrique. \emph{Bulletin de l'Acad\'emie Royale des Sciences et des Lettres de Danemark}, no. 4. As cited on p. 173, whose footnote prints initial P. although the main text gives T.

\item Brill, A., \& Noether, M. (n.d.). Work on the curve-theoretic Restsatz cited in \emph{Mathematische Annalen}, volume 7. Historical locator as printed on p.~28; the abbreviated reference is retained without an unverified expansion.

\item Bronowski, J. (1938). {[}Title not supplied.{]} \emph{Journal of the London Mathematical Society}. Cited by Enriques at p.~146 in connection with Hodge's base theorem.

\item Brosnan, P. (n.d.). \emph{UMD 808J: Algebraic surfaces} {[}Course notes{]}. University of Maryland. \href{https://math.umd.edu/~pbrosnan/notes/Surfaces/index.html}{Online record}. Locators include Section 3.1, Proposition 3.1; Section 3.2, Theorem 3.4; and Section 4.2 for adjunction, Riemann--Roch, and blowups. Additional locators retained from the later chapters are Sections 4.6 (elimination of indeterminacy) and 4.10 (Castelnuovo contraction), Theorems 2.1 and 3.15 (Nakai--Moishezon and Hodge index), and Theorem 3.4 in \href{https://math.umd.edu/~pbrosnan/notes/Surfaces/sect0017.html}{sect0017.html}. Also cited are the author's pages \href{https://math.umd.edu/~pbrosnan/notes/Surfaces/sect0016.html}{sect0016.html}, \href{https://math.umd.edu/~pbrosnan/notes/Surfaces/sect0028.html}{sect0028.html}, \href{https://math.umd.edu/~pbrosnan/notes/Surfaces/sect0007.html}{sect0007.html}, and \href{https://math.umd.edu/~pbrosnan/notes/Surfaces/shit.html}{shit.html}.

\item Burniat, P. (1952). Surfaces algébriques à système canonique pur dégénéré. \emph{Bulletin de la Classe des sciences, Académie Royale de Belgique, 38}, 1030--1043. \href{https://doi.org/10.3406/barb.1952.69760}{DOI: \nolinkurl{10.3406/barb.1952.69760}}. The discussion of Pompilj's 1948 examples on p.~1030 was checked; the entire article was not examined.

\item Calabri, A., \& Ciliberto, C. (2010). Birational classification of curves on rational surfaces. \emph{Nagoya Mathematical Journal, 199}, 43--93. \href{https://doi.org/10.1215/00277630-2010-003}{DOI: \nolinkurl{10.1215/00277630-2010-003}}. Evidence: journal metadata and its Barber--Zariski bibliography; the full text was not consulted, so no page-specific proof is attributed to it.

\item Calabri, A., \& Ferraro, R. (2002). Explicit resolutions of double point singularities of surfaces. \emph{Collectanea Mathematica, 53}(2), 99--131. \href{https://arxiv.org/abs/math/9911183}{arXiv: 9911183}. The cited version is version 3, with two authors. Locators: Section 3, equations (3.2)--(3.3); Theorem 7.4; Section 9 and Theorem 9.7; Sections 10--11; Section 14, especially Theorem 14.13. Preprint and journal pagination differ.

\item Campedelli, L. (1931). Sulla postulazione di una curva \(i\)-pla. \emph{Rendiconti del Circolo Matematico di Palermo, 55}. Cited by Enriques at p.~114; original proof not independently inspected.

\item Campedelli, L. (1932). Sui piani doppi con curve di diramazione del decimo ordine. \emph{Rendiconti dell'Accademia Nazionale dei Lincei}, series VI, \emph{15}, 358--362. Metadata recorded in the university bibliography \href{https://www.mat.uniroma3.it/users/sernesi/BIBLIOGRAFIA/campedelli.htm}{Online record}; the commentary verifies the construction from Enriques's account, without claiming an independent reading of the original paper.

\item Campedelli, L. (1932a). Sopra alcuni piani doppi notevoli con curva di diramazione del $10^\circ$ ordine. \emph{Reale Accademia dei Lincei}. As cited on pp. 204 and 229.

\item Campedelli, L. (1932b). Sui piani doppi con curva di diramazione del $10^\circ$ ordine. \emph{Reale Accademia dei Lincei}. As cited on pp. 204 and 229.

\item Campedelli, L. (1932c). Sui piani doppi con curva di diramazione dell'ottavo ordine. \emph{Reale Accademia dei Lincei}. As cited on p. 204.

\item Campedelli, L. (1934a). Sul computo dell'invariante di Zeuthen--Segre per una superficie algebrica. \emph{Rendiconti dell'Accademia dei Lincei}, series 6, \emph{XIX}. As cited on p. 172.

\item Campedelli, L. (1934b). \emph{Ancora sul computo ecc.} \emph{Rendiconti dell'Accademia dei Lincei}, series 6, \emph{XIX}. Abbreviated title as cited on p. 172.

\item Campedelli, L. (1936). {[}Work on surfaces of degree \(2n\) with eight \(n\)-fold points.{]} \emph{Atti della Reale Accademia di Torino}. Enriques gives no title or pages at p.~123. The author-maintained institutional bibliography identifies the likely title \emph{Sulle superficie di ordine 2n con otto punti n-pli}; this remains a bibliographic lead, and the original article was not independently read.

\item Campedelli, L. (1940a). La classificazione dei piani doppi con tutti i generi uguali all'unit\`a. \emph{Atti del 2$^\circ$ Congresso dell'Unione Matematica Italiana}. Bologna. As cited on p. 204.

\item Campedelli, L. (1940b). Le superficie con i generi uguali all'unit\`a, rappresentabili in infiniti modi sul piano doppio. \emph{Rendiconti del Seminario Matematico della Reale Universit\`a di Roma}, series V, \emph{I}. As cited on p. 204.

\item Campedelli, L. (1940c). Sopra i piani doppi con tutti i generi uguali all'unit\`a. \emph{Rendiconti del Seminario Matematico di Padova, XI}. As cited on p. 204.

\item Cantat, S., \& Dolgachev, I. (2012). Rational surfaces with a large group of automorphisms. \emph{Journal of the American Mathematical Society, 25}(3), 863--905. \href{https://doi.org/10.1090/S0894-0347-2012-00732-2}{DOI: \nolinkurl{10.1090/S0894-0347-2012-00732-2}}. Proposition 2.2 and Section 2.2.

\item Caporali, E. (n.d.). [Work on the Jacobian of a net of Jacobians; title and publication not identified]. The attribution on p. 180 gives no further locator; no bibliographic expansion is asserted.

\item Casas-Alvero, E. (2000). \emph{Singularities of plane curves}. Cambridge University Press. \href{https://doi.org/10.1017/CBO9780511569326}{DOI: \nolinkurl{10.1017/CBO9780511569326}}. Chapter 3, ``Infinitely near points,'' pp.~67--118, \href{https://doi.org/10.1017/CBO9780511569326.005}{DOI: \nolinkurl{10.1017/CBO9780511569326.005}}; chapter 4, ``Virtual multiplicities,'' pp.~119--158, \href{https://doi.org/10.1017/CBO9780511569326.006}{DOI: \nolinkurl{10.1017/CBO9780511569326.006}}. Evidence: publisher metadata and chapter summaries; no uninspected theorem number is claimed.

\item Castelnuovo, G., \& Enriques, F. (1900/1901). Sopra alcune questioni fondamentali nella teoria delle superficie algebriche. \emph{Annali di Matematica, Serie 3, 6}, §4. Enriques's p.~142 references give variant titles and a 1900 date; the 1900/1901 dating is flagged rather than resolved from an uninspected original. The same work is cited without its title for exceptional-curve elimination on p.~58 and in abbreviated form on pp.~170 and 172. These pointers identify the historical dependency; the modern smooth-point contraction proof is given separately.

\item Castelnuovo, G. (1891). Osservazioni intorno alla geometria sopra una superficie. Nota I. \emph{Rendiconti dell'Istituto Lombardo}. Reprinted in \emph{Memorie scelte}, item XVII. Cited by Enriques at p.~122.

\item Castelnuovo, G. (1891). Ricerche generali sopra i sistemi lineari di curve piane. \emph{Memorie della Reale Accademia delle Scienze di Torino, Series 2, 42}. Reprinted in \emph{Memorie scelte di geometria algebrica}, pp.~137 ff., Zanichelli, 1937. Historical attribution of the characteristic series printed on p.~19.

\item Castelnuovo, G. (1893). Sui multipli di una serie lineare di gruppi di punti appartenente ad una curva algebrica. \emph{Rendiconti del Circolo Matematico di Palermo, 7}. \href{https://doi.org/10.1007/BF03012436}{DOI: \nolinkurl{10.1007/BF03012436}}. Enriques cites \emph{Memorie scelte}, p.~95 and item VIII (pp.~103 and 120). Bibliographic metadata was verified; the original article was not independently examined.

\item Castelnuovo, G. (1896). Alcuni risultati sui sistemi lineari di curve appartenenti ad una superficie algebrica. \emph{Memorie della Società Italiana delle Scienze (detta dei XL), Serie 3, 10}, 82--123. \href{https://media.accademiaxl.it/memorie/S3-VX-1896/Castelnuovo82-123.pdf}{Academy scan}. Enriques cites \emph{Memorie scelte}, p.~335 and item XXII (pp.~115 and 140). Metadata was checked; the original proof was not read end to end.

\item Castelnuovo, G. (1896). Sulle superficie di genere zero. \emph{Memorie della Societ\`a Italiana delle Scienze}, series III, \emph{10}. Historical reference printed on p. 93; original full text not independently inspected. The Chapter VI locator is the reprint in \emph{Memorie scelte}, p.~307 (cited on p.~230).

\item Castelnuovo, G. (1897). Alcune proprietà fondamentali dei sistemi lineari di curve tracciati sopra una superficie algebrica. \emph{Annali di Matematica, Serie 2, 25}. Reprinted in \emph{Memorie scelte}, item XXIII. Cited by Enriques at pp.~118 and 140; historical proof attribution is retained.

\item Castelnuovo, G. (1897). Sul genere lineare di una superficie e sulla classificazione a cui esso d\`a luogo. \emph{Rendiconti della Reale Accademia dei Lincei}, Vol. VI. Enriques's p. 55 footnote also points to \emph{Memorie scelte}, p. 443. Exact original pagination was not verified. The Chapter IV footnote on p.~142 also identifies \emph{Memorie scelte}, item XXIV.

\item Castelnuovo, G. (1898, December 25). Communication to F. Enriques, no. 4. \emph{Rendiconti del Circolo Matematico di Palermo}. Historical locator as printed on p.~40; the letter itself was not independently located.

\item Catanese, F. (2013). A superficial working guide to deformations and moduli. In \emph{Handbook of moduli} (Vol. I, pp. 161--215). International Press. \href{https://arxiv.org/abs/1106.1368}{arXiv: 1106.1368}. Section 1.3, p. 13 of the preprint, for the Enriques--Kuranishi inequality and its attribution.

\item Catanese, F., Franciosi, M., Hulek, K., \& Reid, M. (1999). Embeddings of curves and surfaces. \emph{Nagoya Mathematical Journal, 154}, 185--220. \href{https://doi.org/10.1017/S0027763000025381}{DOI: \nolinkurl{10.1017/S0027763000025381}}. Theorem 3.3, pp.~194--195; §6, p.~206. The numerically 2-connected dualizing-sheaf theorem is distinct from the constant-functions lemma for 1-connected divisors.

\item Cayley, A. (1870). {[}Spatial-curve postulation formulas; title and publication not supplied at the cited passage.{]} Attributed by Enriques at p.~102.

\item Cayley, A. (1871). On the deficiency of certain surfaces. \emph{Mathematische Annalen, 3}. Cited by Enriques at pp.~121 and 137. The preliminary discussion at p.~102 separately attributes spatial-curve postulation to Cayley (1870), without supplying an article title.

\item Chisini, O. (1921). La risoluzione delle singolarità delle superficie mediante trasformazioni birazionali dello spazio. \emph{Memorie dell'Accademia delle Scienze di Bologna}. Historical citation printed on p.~6; no claim of full-text verification is made.

\item Chisini, O. (1934). Un teorema di esistenza dei piani multipli (Notes I and II). \emph{Atti dell'Accademia dei Lincei}, series VI, \emph{19}. As cited on p. 199.

\item Chisini, O. (1938). Un pi\`u generale teorema di esistenza dei piani multipli. \emph{Atti dell'Accademia dei Lincei}, series VI, \emph{27}. As cited on p. 199.

\item Chisini, O. (1939). Altre curve di diramazione dei piani $n$-pli. \emph{Atti dell'Accademia dei Lincei}, series VI, \emph{29}. As cited on p. 199.

\item Chisini, O. (1944). Sulla identit\`a birazionale delle funzioni algebriche di due variabili dotate di una medesima curva di diramazione. \emph{Rendiconti dell'Istituto Lombardo di Scienze e Lettere}. As cited on p. 193.

\item Ciliberto, C., \& Flamini, F. (2011). On the branch curve of a general projection of a surface to a plane. \emph{Transactions of the American Mathematical Society, 363}(7), 3457--3471. \href{https://doi.org/10.1090/S0002-9947-2011-05401-2}{DOI: \nolinkurl{10.1090/S0002-9947-2011-05401-2}} . The full author preprint \href{https://arxiv.org/abs/0811.0467}{arXiv: 0811.0467} (v2, 2010) was checked, especially §2, Definitions 2.1 and 2.3, Remark 2.2, and Theorems 2.4--2.6. Its newly proved plane-branch theorem is beyond the present extract; it must not be confused with the \(\mathbb P^3\) double-curve theorem quoted there.

\item Clebsch, A. (1868). [Note on everywhere finite double integrals of the first kind, cited by Enriques on p. 62.] \emph{Comptes rendus}. The title and complete bibliographic details are not given in that citation and were not independently verified.

\item Clebsch, A. (1870; frequently cited as 1871 by bound volume). Ueber den Zusammenhang einer Klasse von Fl\"achen-abbildungen mit der Zweitheilung der Abel'schen Functionen. \emph{Mathematische Annalen, 3}, 45--75. \href{https://doi.org/10.1007/BF01443905}{DOI: \nolinkurl{10.1007/BF01443905}}. The publisher dates the issue to March 1870; the source's 1861 is erroneous. The original article body was not independently inspected.

\item Conforto, F. (1939). \emph{Le superficie razionali}. Zanichelli. Based on Enriques's lectures; the source also cites the 1945 edition on p.~26. This is the historical Enriques--Conforto reference, not a second modern proof of the reconstructed quintic. Further locators are Book II, chapter II, Section 7, p.~273; Book I, chapter VIII; pp.~196, 204, and 206--207; Section 27, p.~489; and chapter VI, Section 37, p.~496. These are historical dependencies, not independently reverified classifications.

\item De Franchis, M. (n.d.). {[}General theorem on irregular double planes; title, date, and publication not given.{]} Attribution in Enriques's footnote at p.~157. The explicit irregular example is reconstructed independently in the commentary.

\item De Jonqui\`eres, E., and Beck (1878, as cited). [Observation on infinitesimal intersections at nodes and cusps]. \emph{Mathematische Annalen, 14}. Page 207 associates this observation with both names and gives no article title; it also refers to Enriques--Chisini, Book II, Chapter II. This records the printed attribution, not an assertion of joint authorship of an identified article.

\item Deligne, P. (1981). Le groupe fondamental du compl\'ement d'une courbe plane n'ayant que des points doubles ordinaires est ab\'elien (d'apr\`es W. Fulton). \emph{S\'eminaire Bourbaki, 1979--1980}, Expos\'e 543, 1--10. \href{https://www.numdam.org/item/SB_1979-1980__22__1_0/}{NUMDAM record}. Introduction and Theorem 1.

\item Dolgachev, I. V. (2012). \emph{Classical algebraic geometry: A modern view}. Cambridge University Press. \href{https://doi.org/10.1017/CBO9781139084437}{DOI: \nolinkurl{10.1017/CBO9781139084437}}. Its bibliography is cited for the 1871 bound-volume dating of Clebsch's volume-3 paper; this is a bibliographic comparison, not verification of Clebsch's original proof.

\item Dolgachev, I. V. (2016). A brief introduction to Enriques surfaces. In \emph{Development of moduli theory---Kyoto 2013} (Advanced Studies in Pure Mathematics, Vol. 69, pp. 1--32). Mathematical Society of Japan. \href{https://arxiv.org/abs/1412.7744}{arXiv: 1412.7744}. The relevant conductor and tetrahedral-sextic discussion is Section 2, pp. 3--4, equation (1), in the version-2 preprint. Chapter VI also uses Section 4, especially p.~13 of the author manuscript, for the corrected quartic del Pezzo model.

\item Du Val, P. (1936). Reducible exceptional curves. \emph{American Journal of Mathematics, 58}(2), 285--289. \href{https://doi.org/10.2307/2371038}{DOI: \nolinkurl{10.2307/2371038}}. Bibliographic identity checked against the journal issue.

\item Ein, L., \& Lazarsfeld, R. (1997). Singularities of theta divisors and the birational geometry of irregular varieties. \emph{Journal of the American Mathematical Society, 10}(1), 243--258. Proposition 3.1 and equation (3.1.1), p. 254. \href{https://www.math.stonybrook.edu/~cschnell/pdf/references/EinLazarsfeld.pdf}{Full text}.

\item Enriques, F., \& Chisini, O. (1915, 1918, 1924, 1934). \emph{Lezioni sulla teoria geometrica delle equazioni e delle funzioni algebriche} (4 vols.; also cited in abbreviated form without \emph{Lezioni sulla}). Zanichelli. Historical prerequisite throughout the book. Retained opening locators include Book I, Section 23; Book II, chapter I, Sections 3 and 5, volume I, pp.~167 and 180; Books III--IV, volume II, including p.~289 and pp.~327 ff.; and Book V, volume III, including pp.~78, 399, 405, and 485. The anomalous Book I/volume II pointer on p.~39 is reported as printed. Chapter II cites volume II (1918), Book III, chapter I, Section 5, and the printed ``Libro XV, cap. IV, Section 39, vol. II, pag. 634 e seg.''; the unusual Book XV locator is not silently emended. Chapter III cites Book IV; Book V, Section 12, volume III, pp.~90 and 93; Book V, chapter I, note to Section 10, volume III, pp.~73 and 485; and Book V, chapter IV, Section 42, volume III, p.~482. Chapter IV chiefly invokes Book V, volume III, for curve theory, and cites the work again at p.~164. Chapters V--VI add Book II, chapter II, volume I, p.~331 and volume II, p.~286; Book V, Section 3, volume III, p.~18; Book V, chapter III, Section 33, volume III, p.~355; Book V, chapter IV, Section 38, volume III, p.~445; and Book VI, Section 40, volume IV, p.~246.

\item Enriques, F., \& Fano, G. (n.d.). [Result concerning continuous projective automorphisms.] Enriques's p.~206 footnote cites the result through Conforto (1939), chapter VI, Section 37, p.~496, without identifying an original paper. The indirect attribution is retained; no unverified publication data are supplied.

\item Enriques, F., \& Severi, F. (1909, as cited). M\'emoire sur les surfaces hyperelliptiques. \emph{Acta Mathematica}. As cited on p. 190; no article pagination is supplied by that citation.

\item Enriques, F. (1893). Una questione sulla linearità dei sistemi di curve appartenenti ad una superficie algebrica. \emph{Rendiconti della Reale Accademia dei Lincei, Series 5, 2}, 3--8. Official collected-works scan: \href{https://www.federigoenriques.org/wp-content/uploads/opere/pdf/GM_Enriques_1893_3.pdf}{National-edition scan}. The statement, counterexample, and proof on original pp.~4--7 were checked directly.

\item Enriques, F. (1893/1894). \emph{Ricerche di geometria sulle superficie algebriche}. Clausen (1893 offprint); \emph{Memorie della Reale Accademia delle Scienze di Torino}, series II, \emph{44}, 171--232 (1894 journal publication). The official \href{https://www.federigoenriques.org/indice-cronologico-delle-opere/}{chronological archive} distinguishes the two versions. The book cites the 1893 work at pp.~21, 142, 150, and 186, including IV, Section 4; the date of the offprint must not be silently replaced by the journal year. This is a different work from Enriques's 1893 paper on the linearity criterion.

\item Enriques, F. (1895). [Rational divisor-family criterion, attributed without a work title on p. 100.] The attribution is retained from the source; the modern Picard-variety proof is supplied in the commentary.

\item Enriques, F. (1896). Introduzione alla geometria sopra le superficie algebriche. \emph{Memorie della Società Italiana delle Scienze detta dei XL}, Series 3, \emph{10}, 1--81. \href{https://www.federigoenriques.org/indice-cronologico-delle-opere/}{National-edition record}. Evidence: national-edition bibliographic record; the book cites §§15--16, but the article's scan was not consulted in full.

\item Enriques, F. (1896). Sui piani doppi di genere uno. \emph{Memorie della Societ\`a Italiana delle Scienze}. Historical reference printed on p. 79; the original proof was not independently inspected. The same title appears on p.~203, with the numeral 1 in place of \emph{uno}.

\item Enriques, F. (1898, April--May). Sui piani doppi di genere lineare $p^{(1)}=1$. \emph{Rendiconti della Reale Accademia dei Lincei}. As cited on pp. 203 and 224; the original issue pagination and the completeness of that classification were not independently verified.

\item Enriques, F. (1899). Una proprietà delle serie continue di curve appartenenti ad una superficie algebrica regolare. \emph{Rendiconti del Circolo Matematico di Palermo, 13}. Cited at p.~137; original work not independently examined.

\item Enriques, F. (1899). Universit\`a di Bologna. Programma del corso di Geometria superiore per l'anno 1897--98. \emph{Bollettino di Bibliografia e Storia delle Scienze Matematiche, 2}, 76--78.

\item Enriques, F. (1901). Intorno ai fondamenti della geometria sopra le superficie algebriche. \emph{Atti della Reale Accademia delle Scienze di Torino, 37}, 19--40. The lecture-year recollection in the opening differs from that in the 1949 book; the programme title supplies the independent 1897--98 locator.

\item Enriques, F. (1904). [Algebraic proof of the irregularity assertion, referred to in Castelnuovo's preface without a full work identification.] The historical claim and its criticism are retained as reported by the preface; no exact title or original-paper locator has been established.

\item Enriques, F. (1908, June). Sui moduli delle superficie algebriche. \emph{Rendiconti dei Lincei}. As cited on p. 214.

\item Enriques, F. (1908). Un'osservazione relativa alle superficie di bigenere uno. \emph{Rendiconti dell'Accademia di Bologna}. As cited on p. 192.

\item Enriques, F. (1912). Sui moduli di una classe di superficie algebriche e sul teorema d'esistenza per le funzioni algebriche di due variabili. \emph{Atti della Reale Accademia di Torino}. As cited on p. 214.

\item Enriques, F. (1914). Sul teorema d'invarianza della serie canonica \(g_{2p-2}^{p-1}\) ecc. \emph{Memorie della Reale Accademia di Bologna}. Abbreviated title retained from p. 65; an unverified full title and pagination are not supplied.

\item Enriques, F. (1923). Sulla costruzione delle funzioni algebriche di due variabili possedenti una data curva di diramazione. \emph{Annali di Matematica}, series IV, \emph{I}. As cited on p. 199; p. 214 gives the Annali year range 1923--1924.

\item Enriques, F. (1932). \emph{Lezioni sulla teoria delle superficie algebriche, raccolte dal Dott. Luigi Campedelli}. CEDAM. Historical precursor identified in the prefaces.

\item Enriques, F. (1934). Sulla classificazione delle superficie algebriche, particolarmente di genere zero: Lezioni raccolte dal Dott. Luigi Campedelli. \emph{Rendiconti del Seminario Matematico della R. Università di Roma}. Historical precursor identified in the prefaces.

\item Enriques, F. (1936). Sulle singolarità che nascono per proiezione di una superficie o varietà algebrica. In \emph{Scritti matematici offerti a Luigi Berzolari} (p.~351 ff.). Pavia. Historical projection reference printed on p.~7; no full-paper inspection is asserted.

\item Enriques, F. (1940). Sur l'extension du théorème de Riemann-Roch aux systèmes linéaires de courbes appartenant à une surface algébrique. \emph{Bulletin des Sciences Mathématiques}. Cited at pp.~141 and 147, especially §5.

\item Enriques, F. (1949). \emph{Le superficie algebriche}. Nicola Zanichelli. Primary source: the published edition and the scan of the print, front matter and pp.~1--236. Page citations throughout this commentary refer to the printed book. The scanned blank leaves at pp.~14, 66, 166, 216, and 236 contain no additional mathematics.

\item Esnault, H., \& Viehweg, E. (1992). \emph{Lectures on vanishing theorems} (DMV Seminar, Vol. 20). Birkh\"auser. \href{https://page.mi.fu-berlin.de/esnault/books/esvibuch.pdf}{Full text}. Corollary 5.12(c), pp. 48--49.

\item FederigoEnriques.org. (n.d.). \emph{Indice cronologico delle opere}. \href{https://www.federigoenriques.org/indice-cronologico-delle-opere/}{National-edition record}. Used to verify the offprint/journal distinction and the course-programme and memoir metadata.

\item Franchetta, A. (1940). Sulle curve eccezionali riducibili di prima specie. \emph{Bollettino dell'Unione Matematica Italiana}, 332--341. Cited on p. 58. Conflicting secondary volume numbers are not resolved here, and the original article was not inspected; the discrepancy formula is proved by blowup induction in the commentary.

\item Franchetta, A. (1941). Sulla caratterizzazione delle curve eccezionali riducibili di prima specie. \emph{Bollettino dell'Unione Matematica Italiana}, 372--375. Cited on p. 60. The exact broad \(p_g=0\) qualification and conflicting secondary volume numbers were not independently resolved. The Chapter IV citation on p.~153 refers to the same paper; it supplies no independent resolution of the volume-number question.

\item Franchetta, A. (1941). Sulla curva doppia della proiezione di una superficie generale dell'\(S_4\), da un punto generico su un \(S_3\). \emph{Atti della Reale Accademia d'Italia. Rendiconti della Classe di Scienze Fisiche, Matematiche e Naturali}, ser. 7, \emph{2}, 282--288. The year 1941 agrees with Enriques and Lopez--Ran; some later bibliographies give 1940. The original journal issue was not independently inspected, so the year discrepancy remains bibliographic, not silently ``corrected.''

\item Franchetta, A. (1946). Sui punti doppi isolati delle superficie algebriche. Nota I; Nota II. \emph{Atti dell'Accademia Nazionale dei Lincei, Rendiconti, Classe di Scienze Fisiche, Matematiche e Naturali}, series VIII, \emph{1}, 49--57; 162--168. Calabri and Ferraro (2002), bibliography item 13, supply this metadata; Enriques prints 1945 on p. 89. The original Lincei scans were not independently inspected.

\item Franchetta, A. (1947). Sulla curva doppia della proiezione di una superficie generale dell'\(S_4\), da un punto generico su un \(S_3\). \emph{Rendiconti dell'Accademia Nazionale dei Lincei}, ser. 8, \emph{2}, 276--279. Title and pages corroborated by the checked modern references; the original article was not independently read.

\item Friedman, M., Lehman, R., Leyenson, M., \& Teicher, M. (2012). On ramified covers of the projective plane II: Generalizing Segre's theory. \emph{Journal of the European Mathematical Society, 14}, 971--996. \href{https://doi.org/10.4171/JEMS/324}{DOI: \nolinkurl{10.4171/JEMS/324}}. Section 2 was consulted for ordinary projection singularities and the double-curve theorem; the quintic's numerical reconstruction is the commentary's own calculation.

\item Fulton, W. (1980). On the fundamental group of the complement of a node curve. \emph{Annals of Mathematics, 111}(2), 407--409. \href{https://doi.org/10.2307/1971204}{DOI: \nolinkurl{10.2307/1971204}}.

\item Fulton, W. (1998). \emph{Intersection theory} (2nd ed.). Springer. \href{https://doi.org/10.1007/978-1-4612-1700-8}{DOI: \nolinkurl{10.1007/978-1-4612-1700-8}}. Chapter 14, §14.4, Thom--Porteous formula; publisher chapter preview: \href{https://doi.org/10.1007/978-1-4612-1700-8_15}{DOI: \nolinkurl{10.1007/978-1-4612-1700-8_15}}.

\item Godeaux, L. (1920). Sur le calcul de l'invariant de Zeuthen--Segre. \emph{M\'emoires de la Soci\'et\'e des Sciences du Hainaut, 66}. As cited on p. 172.

\item Godeaux, L. (1931). Sur une surface alg\'ebrique de genre z\'ero et de bigenre deux. \emph{Rendiconti dei Lincei}, Vol. XIV. Historical reference printed on p. 95; complete pagination and the original proof were not verified.

\item Godeaux, L. (1934, October 13). [Communication cited without title]. \emph{Bulletin de l'Acad\'emie Royale de Belgique}. As cited on p. 172.

\item Grothendieck, A. (1961). Techniques de construction et théorèmes d'existence en géométrie algébrique IV : Les schémas de Hilbert. \emph{Séminaire Bourbaki, 6}, Exposé 221, 249--276. \href{https://www.numdam.org/item/SB_1960-1961__6__249_0.pdf}{Original text}. Corollary 5.4 and §§5.5--5.6 were inspected for tangent spaces, the characteristic homomorphism, and the vanishing cases. This is the relevant Exposé for the completeness discussion; it must not be confused with Exposé 236 merely because the latter begins on printed p.~221.

\item Grothendieck, A. (1962a). Technique de descente et théorèmes d'existence en géométrie algébrique. V. Les schémas de Picard : Théorèmes d'existence. \emph{Séminaire Bourbaki, 7}, Exposé 232, 143--161. \href{https://numdam.org/book-part/SB_1961-1962__7__143_0/}{Bibliographic record and original text}. Used for the construction framework; the detailed modern presentation actually checked is Kleiman (2005).

\item Grothendieck, A. (1962b). Technique de descente et théorèmes d'existence en géométrie algébrique. VI. Les schémas de Picard : Propriétés générales. \emph{Séminaire Bourbaki, 7}, Exposé 236, 221--243. \href{https://numdam.org/item/SB_1961-1962__7__221_0.pdf}{Original text}. Proposition 2.10(iii), concerning tangent dimension and characteristic-zero smoothness, was checked directly. A published erratum exists; none of the more general relative-family claims subject to additional qualifications is needed here.

\item Grothendieck, A. (Ed.). (2003). \emph{Rev\^etements \'etales et groupe fondamental (SGA 1)} (Documents Math\'ematiques, Vol. 3; revised reissue of Lecture Notes in Mathematics, Vol. 224, 1971, including two expos\'es by M. Raynaud). Soci\'et\'e Math\'ematique de France. \href{https://arxiv.org/abs/math/0206203}{arXiv: 0206203}. Expos\'e XII, Theorem 5.1, original pp. 332--333.

\item Guffroy, S. (2003). Sur l'incompl\'etude de la s\'erie lin\'eaire caract\'eristique d'une famille de courbes planes \`a n\oe uds et \`a cusps. \emph{Nagoya Mathematical Journal, 171}, 51--83. \href{https://doi.org/10.1017/S0027763000025514}{DOI: \nolinkurl{10.1017/S0027763000025514}}. The publisher's abstract distinguishes the generically nonreduced example of Wahl from examples whose reduction is singular.

\item Hartshorne, R. (1977). \emph{Algebraic geometry} (Graduate Texts in Mathematics, Vol. 52). Springer. \href{https://doi.org/10.1007/978-1-4757-3849-0}{Publisher record}. Chapter III supplies coherent cohomology and duality; Chapter V, §§1, 3, and 5 supplies intersection theory, blowups, and birational geometry of surfaces, including surface Riemann--Roch at Theorem V.1.6. Publisher and chapter metadata were verified. The full book was not available for a fresh page-by-page inspection; the actual formulas and elementary proofs used here are written out in the commentary and cross-checked against the accessible references. The commentary on Chapters II--III also cites II, Proposition 8.20 and V, Propositions 3.1--3.4, checked through Brosnan's authored notes; it does not claim inspection of every original textbook page.

\item Hartshorne, R. (2010). \emph{Deformation theory}. Springer. \href{https://doi.org/10.1007/978-1-4419-1596-2}{DOI: \nolinkurl{10.1007/978-1-4419-1596-2}}. Chapters 1--2 provide the normal-sheaf interpretation of embedded first-order deformations and their obstruction theory.

\item Hassett, B. (2009). Rational surfaces over nonclosed fields. In H. Darmon, D. A. Ellwood, B. Hassett, \& Y. Tschinkel (Eds.), \emph{Arithmetic geometry} (Clay Mathematics Proceedings, Vol. 8, pp. 155--209). American Mathematical Society. \href{https://www.claymath.org/wp-content/uploads/2022/03/cmip08c.pdf}{Full text}. Corollary 2.20, Theorem 2.26, and Remark 2.29.

\item Hodge, W. V. D. (1937). {[}Theorem on a base; article title not supplied.{]} \emph{Journal of the London Mathematical Society, 12}. Cited at p.~146.

\item Humbert, G. (n.d.). {[}Irregular-surface examples; exact title and date not supplied in the Chapter IV discussion.{]} Enriques's historical attribution at p.~138 is retained without an independent priority determination.

\item Kleiman, S. L. (1973). Completeness of the characteristic system. \emph{Advances in Mathematics, 11}(3), 304--310. The theorem's location on p.~307 is verified through Kleiman (2005, Remark 5.18); the corrected page range 304--310 is corroborated by the author's 2014 bibliography, reference 37. The 2005 preprint bibliography starts the range at 305. The original 1973 article was not directly inspected in full, and no DOI is guessed.

\item Kleiman, S. L. (1998). Bertini and his two fundamental theorems. \emph{Rendiconti del Circolo Matematico di Palermo, Series 2, Supplement 55}, 9--37. \href{https://arxiv.org/abs/alg-geom/9704018}{arXiv: 9704018}. The 1997 preprint was checked at §§4--5, especially Theorems 4.1 and 5.3 and the discussion of van der Waerden. The publication year and pages are verified by the Stacks Project bibliography: \href{https://stacks.math.columbia.edu/bibliography/Kleiman-Bertini}{Online record}.

\item Kleiman, S. L. (2005). The Picard scheme. In B. Fantechi, L. Göttsche, L. Illusie, S. L. Kleiman, N. Nitsure, \& A. Vistoli, \emph{Fundamental algebraic geometry: Grothendieck's FGA explained} (Mathematical Surveys and Monographs, Vol. 123, pp.~235--321). American Mathematical Society. \href{https://arxiv.org/abs/math/0504020}{Author preprint}. Checked: the historical introduction; Theorems 5.4 and 5.11; Corollaries 5.13--5.14; and Remark 5.18. The theorem/remark labels refer to that checked preprint. Its stray phrase equating a zero semiregularity map with injectivity is not used: the definition and the exact sequence require \textbf{injectivity}. The commentary on Chapter IV additionally uses Definition 6.1, Theorems 6.3 and 6.16, and Corollary 6.17 of the same arXiv text.

\item Kleiman, S. L. (2014). \emph{The Picard scheme} {[}Author preprint{]}. arXiv. \href{https://arxiv.org/abs/1402.0409}{Text}. This is a separate historical exposition from the 2005 chapter, despite the identical title. Used narrowly to verify the 1973 bibliography and distinguish the two publications.

\item Kodaira, K. (1952). On the theorem of Riemann-Roch for adjoint systems on Kählerian varieties. \emph{Proceedings of the National Academy of Sciences, 38}(6), 522--527. \href{https://doi.org/10.1073/pnas.38.6.522}{DOI: \nolinkurl{10.1073/pnas.38.6.522}}.

\item Kulikov, V. S., \& Kulikov, Vik. S. (2010). On complete degenerations of surfaces with ordinary singularities in $\mathbb P^3$. \emph{Sbornik: Mathematics, 201}(1), 129--158. \href{https://doi.org/10.1070/SM2010v201n01ABEH004068}{DOI: \nolinkurl{10.1070/SM2010v201n01ABEH004068}}. Equation (24), p. 138. The authors' different abbreviated first names are retained.

\item Kulikov, V. S. (2008). On Chisini's conjecture. II. \emph{Izvestiya: Mathematics, 72}(5), 901--913. \href{https://doi.org/10.1070/IM2008v072n05ABEH002423}{DOI: \nolinkurl{10.1070/IM2008v072n05ABEH002423}}. Theorem, p. 902, concerning generic linear projections and its Veronese exception.

\item Kumar, A. (2008a). \emph{Algebraic surfaces, Lecture 9: Castelnuovo's criterion for rationality}. MIT OpenCourseWare, 18.727. Notes prepared by K. Venkatram in collaboration with A. Kumar. \href{https://ocw.mit.edu/courses/18-727-topics-in-algebraic-geometry-algebraic-surfaces-spring-2008/01d543f81d0b743f06deed68d9eec77d_lect9.pdf}{MIT lecture notes}. Theorem 1 and its reductions.

\item Kumar, A. (2008b). \emph{Algebraic surfaces, Lecture 17: Enriques surfaces}. MIT OpenCourseWare, 18.727. Notes prepared by K. Venkatram in collaboration with A. Kumar. \href{https://ocw.mit.edu/courses/18-727-topics-in-algebraic-geometry-algebraic-surfaces-spring-2008/9fa718160b9c9c0ecc0dceeaa339e644_lect17.pdf}{MIT lecture notes}. Proposition 1.

\item Lazarsfeld, R., with the assistance of G. Fern\'andez del Busto. (1994). \emph{Lectures on linear series} (arXiv:alg-geom/9408011). Vanishing Theorem 5.2, p. 26. \href{https://www.math.stonybrook.edu/~roblaz/Reprints/Lazarsfeld.LLS.pdf}{Full text}.

\item Lefschetz, S. (n.d.). Work on the topological interpretation of virtual curves, cited historically on p.~30. Enriques does not specify a work at that point; no exact passage or missing bibliographic fields are invented.

\item Levi, B. (1897). Risoluzione delle singolarità puntuali delle superficie algebriche. \emph{Atti dell'Accademia delle Scienze di Torino}; compare \emph{Annali di Matematica} (1897). Historical citation as printed on p.~6; the priority assertion is Enriques's.

\item Libois, P. (1934). [Work on biquadratic degree-four covers; title not supplied]. \emph{Rendiconti dell'Accademia dei Lincei}. As cited on p. 204.

\item Liedtke, C. (2008). \emph{A note on non-reduced Picard schemes} {[}Preprint, version 2{]}. arXiv. \href{https://arxiv.org/abs/0805.0848}{arXiv: 0805.0848}. The Introduction and Section 2 were consulted. Published in 2009 in \emph{Journal of Pure and Applied Algebra, 213}(5), 737--741, \href{https://doi.org/10.1016/j.jpaa.2008.09.001}{DOI: \nolinkurl{10.1016/j.jpaa.2008.09.001}}; citations to 2008 identify the consulted preprint.

\item Lipman, J. (1978). Desingularization of two-dimensional schemes. \emph{Annals of Mathematics, 107}, 151--207. \href{https://doi.org/10.2307/1971141}{DOI: \nolinkurl{10.2307/1971141}} . The theorem on p.~151 and the introductory account of its dependencies were checked in the author's PDF at \href{https://www.math.purdue.edu/~jlipman/papers-older/\%5B1978\%5D\%20Desingularization\%20of\%20two-dimensional\%20schemes.pdf}{Full text} .

\item Lopez, A., \& Ran, Z. (2003). On the irreducibility of secant cones, and an application to linear normality. \emph{Duke Mathematical Journal, 117}, 389--401. \href{https://doi.org/10.1215/S0012-7094-03-11731-7}{DOI: \nolinkurl{10.1215/S0012-7094-03-11731-7}} . The proof cited here was checked in the author preprint \href{https://arxiv.org/abs/math/0111147}{arXiv: 0111147} (2001), pp.~2 and 7; no assertion that every theorem number is unchanged between versions is needed.

\item Mendes Lopes, M. (1996). Adjoint systems on surfaces. \emph{Bollettino dell'Unione Matematica Italiana, Serie 7, 10-A}, 169--179. \href{https://www.math.tecnico.ulisboa.pt/~mmlopes/ADJOINT.pdf}{Author manuscript}. Fact 1 on manuscript p.~2, Lemma 2.6, and Theorem 4.1; manuscript and journal pagination differ.

\item Mezzetti, E., \& Portelli, D. (1997). A tour through some classical theorems on algebraic surfaces. \emph{Analele Ştiinţifice ale Universităţii Ovidius Constanţa, Seria Matematică, 5}(2), 51--78. Bibliographic identity verified; theorem pointers 3, 5, and 8 are reported by Ciliberto--Flamini, not independently read.

\item Milne, J. S. (2008). \emph{Abelian varieties} (Version 2.00) {[}Course notes{]}. \href{https://www.jmilne.org/math/CourseNotes/AV.pdf}{Full text}. Part I, Propositions 3.9--3.10, supplies the constancy of rational maps from rational curves and unirational varieties to abelian varieties.

\item Moishezon, B. (1977). \emph{Complex surfaces and connected sums of complex projective planes} (Lecture Notes in Mathematics, Vol. 603). Springer. \href{https://doi.org/10.1007/BFb0063355}{DOI: \nolinkurl{10.1007/BFb0063355}} . The publisher verifies the chapter ``Generic projections of algebraic surfaces into \(\mathbf{CP}^3\),'' pp.~72--115; the specific p.~115 attribution to Mumford is supplied by the checked Lopez--Ran text. The chapter itself was not independently read in full.

\item Mumford, D. (1966). \emph{Lectures on curves on an algebraic surface} (with a section by G. M. Bergman; Annals of Mathematics Studies, No.~59). Princeton University Press. \href{https://stacks.math.columbia.edu/bibliography/Mum}{Bibliographic record}. The tangent-space and obstruction pointers, pp.~154 and 157--159, are independently identified in the inspected Kleiman (2005, Remark 5.18); the full book was not independently read.

\item Noether, M. (1870). Zur Theorie des eindeutigen Entsprechens algebraischer Gebilde. \emph{Mathematische Annalen, 2}. The title is supplied by Enriques's footnote on p.~62 for both volumes 2 and 8. The book prints 1869 for volume 2; the standard publication year is 1870, but the original journal pages were not inspected. The volume-8 continuation is listed separately under 1875.

\item Noether, M. (1871). {[}Spatial-curve postulation formulas; title and publication not supplied at the cited passage.{]} Attributed by Enriques at p.~102. The modern proof of the relevant formula is given in Chapter IV.

\item Noether, M. (1875). Zur Theorie des eindeutigen Entsprechens algebraischer Gebilde. \emph{Mathematische Annalen, 8}, 495--533. Publisher issue record: \href{https://link.springer.com/journal/208/volumes-and-issues/8-4}{Online record}. This verifies the date correction; no claim is made that the entire original article was read. The Chapter II footnote on p.~62 supplies the title, while its year 1874 differs from the 1875 date elsewhere in the book. The publisher issue record resolves the publication year as 1875. The p.~138 discussion of Noether's view of the genera remains an attributed historical report, not a verified quotation of the original paper; p.~173 cites only volume VIII and \emph{op. cit.}

\item Noether, M. (1878). Ueber die Einzweideutigen Ebenentransformationen. \emph{Sitzungsberichte der Physikalisch-medizinischen Societ\"at zu Erlangen}, Vol. 10. Historical title and locator retained from p. 79; original full proof not independently inspected.

\item Noether, M. (1886). Extension du théorème de Riemann-Roch aux surfaces algébriques. \emph{Comptes Rendus de l'Académie des Sciences, Paris, 103}. Enriques's criticism at p.~139 is retained as attributed history.

\item Noether, M. (1888). Announcement on ordinary models of algebraic surfaces, cited without a title on p.~8. The historical attribution is retained; the supplied evidence does not identify an exact original item.

\item Noether, M. (1888). Anzahl der Moduln einer Classe algebraischer Fl\"achen. \emph{Sitzungsberichte der Akademie}, Berlin. As cited on p. 213.

\item Pannelli, M. (1905). Sui sistemi lineari triplamente infiniti di curve tracciati sopra una superficie algebrica. \emph{Rendiconti del Circolo Matematico di Palermo, XX}. As cited on p. 180.

\item Pardini, R. (2010, October 21--24). \emph{Abelian covers and complete moduli of varieties of general type} {[}Lecture slides; joint work with V. Alexeev{]}. University of Georgia. \href{https://pagine.dm.unipi.it/pardini/files/UGA10B.pdf}{Slides}. PDF pp.~4 and 10 give the eigensheaf decomposition and Hurwitz formula; the explicit branch constructions in Chapter IV are separate calculations.

\item Pedoe, D. (1937). On the canonical systems of certain algebraic surfaces. \emph{Mathematical Proceedings of the Cambridge Philosophical Society, 33}(3), 311--314. \href{https://doi.org/10.1017/S030500410001968X}{DOI: \nolinkurl{10.1017/S030500410001968X}}. Publication record and abstract checked; the precise local claim is not attributed to an independently inspected passage.

\item Picard, E., \& Simart, G. (n.d.). \emph{Traité des fonctions algébriques de deux variables indépendantes} (Vol. II), p.~437. Locator supplied by Enriques at p.~139; no independently verified edition date is assigned.

\item Picard, E. (1905). Sur quelques questions se rattachant à la connexion linéaire dans la théorie des fonctions algébriques de deux variables indépendantes. \emph{Journal für die reine und angewandte Mathematik, 129}. Cited at p.~139; original proof scope is not independently reconstructed.

\item Piene, R. (1978). Some formulas for a surface in \(\mathbb P^3\). In \emph{Algebraic geometry, Tromsø 1977} (Lecture Notes in Mathematics, Vol. 687, pp.~196--235). Springer. This historical projection reference is taken from Friedman et al.~(2012); it was not independently inspected.

\item Pompilj, G. (1939a). Sulla rappresentazione algebrica dei piani tripli. \emph{Rendiconti del Seminario Matematico della Reale Universit\`a di Roma}. As cited on p. 204.

\item Pompilj, G. (1939b). Sui piani multipli diedrici. \emph{Rendiconti del Seminario Matematico della Reale Universit\`a di Roma}. As cited on p. 204.

\item Pompilj, G. (1941). Osservazioni sui piani tripli. \emph{Reale Istituto Lombardo}. As cited on p. 204.

\item Pompilj, G. (1948). Alcuni esempi di superficie algebriche a sistema canonico puro degenere. \emph{Rendiconti dell'Accademia dei Lincei, Serie 8, 4}. Cited in Enriques's footnote at p.~158 and corroborated by Burniat (1952, p.~1030); the complete original article was not inspected.

\item Ramanujam, C. P. (1972). Remarks on the Kodaira vanishing theorem. \emph{Journal of the Indian Mathematical Society, 36}, 41--51. Reprinted in \emph{C. P. Ramanujam---A tribute} (1978, pp.~123--134), Tata Institute of Fundamental Research/Springer. \href{https://mathweb.tifr.res.in/Documents/Publications/Studies/CP-Ramanujam-A-Tribute.pdf}{Institutional reprint}. Lemma 3, reprint p.~127; the introduction explicitly discusses the missing lemma in Franchetta's approach.

\item Ramanujam, C. P. (1974). Supplement to the article ``Remarks on the Kodaira vanishing theorem.'' \emph{Journal of the Indian Mathematical Society, 38}, 121--124. Reprinted in \emph{C. P. Ramanujam---A tribute} (1978), beginning at p.~159, Tata Institute of Fundamental Research/Springer. \href{https://mathweb.tifr.res.in/Documents/Publications/Studies/CP-Ramanujam-A-Tribute.pdf}{Institutional reprint}. The theorem on reprint p.~160 gives characteristic-zero vanishing for the inverse of a nef line bundle of positive square.

\item Reid, M. (1996). \emph{Chapters on algebraic surfaces} {[}Preprint{]}. arXiv:alg-geom/9602006. \href{https://arxiv.org/pdf/alg-geom/9602006}{Full text}. Lemmas 3.10--3.11, pp.~72--73, and Theorem 3.13, p.~74, in this version. The effectivity and characteristic hypotheses of the cited vanishing theorem are retained. The commentary on Chapters V--VI additionally uses Sections A.10--A.12, B.9, D.2.2--D.2.3, and E.4--E.7, especially E.7.3 and E.7.7, in the same preprint.

\item Segre, B. (1930). Sulla caratterizzazione delle curve di diramazione dei piani multipli generali. \emph{Reale Accademia d'Italia, I}. As cited on p. 204.

\item Segre, B. (1934). Sui moduli delle superficie algebriche irregolari. \emph{Rendiconti dell'Accademia dei Lincei, XIX}. As cited on p. 215.

\item Segre, B. (1937). {[}Title not supplied.{]} \emph{Annali di Matematica}. Cited by Enriques at p.~146 in connection with Hodge's base theorem.

\item Segre, C. (1889). Recherches générales sur les courbes et les surfaces réglées algébriques (II partie, Surfaces réglées algébriques). \emph{Mathematische Annalen, 34}, 1--25. \href{https://bdim.dma.unina.it/item?id=GM_Segre_1889_1}{Institutional record}. The matching work is dated 1889 in the institutional bibliography; Enriques's abbreviated reference at p.~136 prints 1892. This resolves the date, without claiming full verification of the original theorem.

\item Segre, C. (1894). Introduzione alla geometria sopra un ente algebrico semplicemente infinito. \emph{Annali di Matematica Pura ed Applicata, Series 2, 22}. Historical reference printed on p.~19; exact additional fields have not been inferred. The Chapter II citation specifies Section 75 (p.~64). The complete original priority comparison was not independently verified.

\item Segre, C. (1896). \emph{Intorno ad un carattere delle superficie ecc.} \emph{Atti dell'Accademia delle Scienze di Torino}. Title abbreviated as cited on p. 167.

\item Serre, J.-P. (1955). Faisceaux algébriques cohérents. \emph{Annals of Mathematics, 61}(2), 197--278. \href{https://doi.org/10.2307/1969915}{DOI: \nolinkurl{10.2307/1969915}}.

\item Severi, F. (1902). Il genere aritmetico ed il genere lineare in relazione alle reti di curve tracciate sopra una superficie algebrica. \emph{Atti dell'Accademia di Torino, XXXVII}. As cited on p. 180; the print reads \emph{tracciate}.

\item Severi, F. (1903). Sulla serie caratteristica di un sistema lineare di curve appartenente ad una superficie algebrica. \emph{Rendiconti della Reale Accademia dei Lincei, Serie 5, 12}. Cited by Enriques at p.~140.

\item Severi, F. (1903). Sulle relazioni che legano i caratteri invarianti di due superficie in corrispondenza algebrica. \emph{Rendiconti del Reale Istituto Lombardo di Scienze e Lettere}, series 2, \emph{36}, 495--511. The title is abbreviated on p. 186. Bibliographic identification: \href{https://www.mat.uniroma3.it/users/sernesi/BIBLIOGRAFIA/severi.htm}{Online record}, entry 17; the original full text was not independently inspected.

\item Severi, F. (1904). Osservazioni sui sistemi continui di curve appartenenti ad una superficie algebrica. \emph{Atti della Reale Accademia delle Scienze di Torino, 39}, 376--392. Bibliographic identity corroborated in Tannenbaum (1984), reference 9; the original paper's text was not inspected. Enriques cites the work again at p.~137.

\item Severi, F. (1905). Sul teorema di Riemann-Roch e sulle serie continue di curve appartenenti ad una superficie algebrica. \emph{Atti dell'Accademia delle Scienze di Torino, 40}. Cited by Enriques at p.~140.

\item Severi, F. (1905). Sulle curve algebriche virtuali appartenenti ad una superficie algebrica. \emph{Rendiconti dell'Istituto Lombardo, Serie 2, 38}. Cited by Enriques at p.~146. The Chapter IV citation on p.~146 supplies the printed journal locator. The original text was not independently inspected.

\item Severi, F. (1906; announcement 1905). Il teorema d'Abel sulle superficie algebriche. \emph{Annali di Matematica Pura ed Applicata}, series III, \emph{12}, 55--79. \href{https://doi.org/10.1007/BF02419496}{DOI: \nolinkurl{10.1007/BF02419496}}. Enriques gives 1905; the publication record dates the issue to December 1906 and identifies the earlier announcement. Original full proof not independently inspected.

\item Severi, F. (1906). Osservazioni varie di geometria sopra una superficie algebrica. \emph{Atti dell'Istituto Veneto}, \emph{65}, part II. Citation retained from p. 96; original proof not independently inspected. The source dates this work to 1906 in the p.~96 footnote, but the p.~100 passage gives Severi (1905) and refers back to ``Osservazioni ecc., l.c. That internal date variation is retained without an unsupported historical resolution; it is distinct from the announcement/publication dates of the Abel paper.

\item Severi, F. (1906). {[}Title not supplied in Enriques's citation.{]} \emph{Mathematische Annalen, 62}, 194. Cited at p.~146 for the equivalence discussion.

\item Severi, F. (1908, November 8). Sulla regolarità del sistema aggiunto. \emph{Rendiconti della Reale Accademia dei Lincei}. Enriques's p.~139 citation explicitly notes dependence on the completeness theorem for continuous characteristic systems.

\item Severi, F. (1909). Fondamenti per la geometria sulle varietà algebriche. \emph{Rendiconti del Circolo Matematico di Palermo}. Cited by Enriques at p.~164 for higher-dimensional inclusion--exclusion.

\item Severi, F. (1911). Alcune relazioni di equivalenza ecc. \emph{Atti dell'Istituto Veneto}, \emph{70}. Abbreviated title retained from p. 96; no unverified expansion or proof locator is supplied.

\item Severi, F. (1921). [Criticism of the completeness step in an earlier irregularity argument; title not identified in the pages treated here]. Historical attribution on p. 214, with a forward reference to Chapter IX, Section 6.

\item Severi, F. (1942). \emph{Serie, sistemi di equivalenza e corrispondenze algebriche sulle variet\`a algebriche}. Cremonese. As cited on p. 180.

\item Tannenbaum, A. (1984). On the classical characteristic linear series of plane curves with nodes and cuspidal points: Two examples of Beniamino Segre. \emph{Compositio Mathematica, 51}(2), 169--183. \href{https://www.numdam.org/item/CM_1984__51_2_169_0/}{NUMDAM record}. Reference 9 is used to identify Severi's bibliography, without attributing the elliptic-product example to this paper. Chapter V additionally uses the Introduction, Definition 1.1, and Remarks 1.2 for characteristic systems of nodal and cuspidal plane curves.

\item The Stacks Project Authors. (2026). \emph{The Stacks Project}. \href{https://stacks.math.columbia.edu/}{Online record}. The cited tags were checked on November 7, 2026; stable tags, rather than potentially changing section numbers alone, identify the cited statements. Curve genus, Riemann--Roch, and duality: \href{https://stacks.math.columbia.edu/tag/0BY6}{0BY6}, \href{https://stacks.math.columbia.edu/tag/0B5B}{0B5B}, and \href{https://stacks.math.columbia.edu/tag/0BS6}{0BS6} (Lemma 53.5.2). Ample twists and vanishing: \href{https://stacks.math.columbia.edu/tag/0B5U}{0B5U} and \href{https://stacks.math.columbia.edu/tag/01XO}{01XO} (Lemma 30.17.1). Projective-space cohomology and linear series: \href{https://stacks.math.columbia.edu/tag/01XS}{01XS}, \href{https://stacks.math.columbia.edu/tag/01ND}{01ND}, and \href{https://stacks.math.columbia.edu/tag/0CCM}{0CCM}. Bertini: \href{https://stacks.math.columbia.edu/tag/0FD4}{0FD4} and \href{https://stacks.math.columbia.edu/tag/0G4C}{0G4C}. Surface resolution: \href{https://stacks.math.columbia.edu/tag/0ADX}{0ADX} and \href{https://stacks.math.columbia.edu/tag/0BGP}{0BGP} (Theorem 54.14.5). Exceptional-curve contraction: \href{https://stacks.math.columbia.edu/tag/0C2I}{0C2I}, \href{https://stacks.math.columbia.edu/tag/0C2L}{0C2L}, and \href{https://stacks.math.columbia.edu/tag/0C2M}{0C2M} (Lemma 54.16.8). Normal bundles and effective divisors: \href{https://stacks.math.columbia.edu/tag/0C4S}{0C4S} and \href{https://stacks.math.columbia.edu/tag/0B3P}{0B3P}. Finite and Cartier-divisor duality: \href{https://stacks.math.columbia.edu/tag/0AU5}{0AU5} and \href{https://stacks.math.columbia.edu/tag/0B4A}{0B4A}. Principal parts: \href{https://stacks.math.columbia.edu/tag/09CQ}{09CQ}. Descent of a fiberwise trivial bundle: \href{https://stacks.math.columbia.edu/tag/0EX7}{0EX7}. Tsen's theorem: \href{https://stacks.math.columbia.edu/tag/03RD}{03RD}. Euler characteristic in proper flat families: \href{https://stacks.math.columbia.edu/tag/0B9T}{0B9T}. Normal crossings: \href{https://stacks.math.columbia.edu/tag/0CBN}{0CBN}. These checks concern the relevant statements and displayed proofs or dependencies; this commentary does not claim a new inspection of the entire project.

\item Todd, {[}printed as ``A. Todd''{]}. (1939, March 23). {[}Title not supplied.{]} \emph{Proceedings of the Cambridge Philosophical Society}. Citation retained as printed at p.~146; the initials and precise bibliographic identification are not silently emended.

\item Torelli, R. (1907). Sui sistemi algebrici di curve appartenenti ad una superficie algebrica. \emph{Atti della Reale Accademia delle Scienze di Torino, 42}, 86--99. Metadata recorded in \href{https://www.mat.uniroma3.it/users/sernesi/BIBLIOGRAFIA/torelli.htm}{Online record}. The source does not print the nodal-curve inequality on p. 100, and its full original statement was not inspected.

\item Vakil, R. (2002). \emph{Complex algebraic surfaces, Class 5} {[}Lecture notes{]}. Stanford University. \href{https://math.stanford.edu/~vakil/02-245/sclass5A.pdf}{Notes}. PDF p.~1: curve Riemann--Roch, Serre duality, and adjunction.

\item Vakil, R. (2002). \emph{Complex algebraic surfaces: Class 6} {[}Course notes{]}. Stanford University. \href{https://math.stanford.edu/~vakil/02-245/sclass6A.pdf}{Full text}. Locator: §2, surface Riemann--Roch and the genus formula.

\item Vakil, R. (2006). Murphy's law in algebraic geometry: Badly-behaved deformation spaces. \emph{Inventiones Mathematicae, 164}, 569--590. \href{https://arxiv.org/abs/math/0411469}{arXiv: 0411469}. Main Theorem 1.1, parts M5a--M5b.

\item van der Waerden, B. L. (1937). Zur algebraischen Geometrie. X. Über lineare Scharen von reduziblen Mannigfaltigkeiten. \emph{Mathematische Annalen, 113}, 705--712. Reprinted in \emph{Zur algebraischen Geometrie: Selected papers} (1983), pp.~225--232, \href{https://doi.org/10.1007/978-3-642-61782-9_17}{DOI: \nolinkurl{10.1007/978-3-642-61782-9_17}}. Bibliography and the reprint's statement were checked; the specific historical repair is documented in Kleiman's analysis, not claimed from an independent reading of every page of the 1937 paper.

\item Voight, J., \& Zureick-Brown, D. (n.d.). \emph{The canonical ring of a stacky curve} {[}Author manuscript with errata incorporated{]}. \href{https://jvoight.github.io/articles/stacky-canonical-rings-final-fixederrata.pdf}{Manuscript}. §2.4, Lemma 2.4.1, the basepoint-free pencil trick. The inspected version is cited without assigning an unverified publication year.

\item Wahl, J. M. (1974). Deformations of plane curves with nodes and cusps. \emph{American Journal of Mathematics, 96}(4), 529--577. \href{https://doi.org/10.2307/2373706}{DOI: \nolinkurl{10.2307/2373706}}. The example descriptions used here are supported by Guffroy (2003) and the later deformation literature; the original full article was not independently inspected.

\item Walker, R. J. (1935). Reduction of the singularities of an algebraic surface. \emph{Annals of Mathematics, 36}(2), 336--365. \href{https://doi.org/10.2307/1968575}{DOI: \nolinkurl{10.2307/1968575}} . Cited by Enriques on p.~6; bibliographic identity checked, full proof not independently examined.

\item Zariski, O. (1929a). On the problem of existence of algebraic functions of two variables possessing a given branch curve. \emph{American Journal of Mathematics, 51}. As cited on p. 199.

\item Zariski, O. (1929b). On the linear connection index of the algebraic surfaces $z^n=f(xy)$. \emph{Proceedings of the National Academy of Sciences of the U.S.A., 15}. Title and locator as cited on p. 199.

\item Zariski, O. (1931). On the irregularity of cyclic multiple planes. \emph{Annals of Mathematics}, series II, \emph{32}. As cited on p. 199.

\item Zariski, O. (1935). \emph{Algebraic surfaces}. Springer. Chapter V, Section 8, is the locator given on p. 214 for the criticism of characteristic-system completeness.

\item Zariski, O. (1939). The reduction of the singularities of an algebraic surface. \emph{Annals of Mathematics, 40}(3), 639--689. \href{https://doi.org/10.2307/1968949}{DOI: \nolinkurl{10.2307/1968949}} . Cited by Enriques on p.~6; bibliographic identity checked, full proof not independently examined.

\item Zariski, O. (1941). [Additional reference on resolution, title unspecified.] \emph{Annals of Mathematics}. Enriques's p.~6 footnote adds ``Cfr. Annals of Math., 1941.'' This item is not identified more precisely here. It must not be conflated with Zariski's 1941 \emph{Transactions of the American Mathematical Society} paper on Bertini's theorem.

\item Zariski, O. (1941). Pencils on an algebraic variety and a new proof of a theorem of Bertini. \emph{Transactions of the American Mathematical Society, 50}(1), 48--70. \href{https://doi.org/10.1090/S0002-9947-1941-0004241-4}{DOI: \nolinkurl{10.1090/S0002-9947-1941-0004241-4}}. Bibliographic identity verified against AMS/JSTOR records; direct AMS PDF retrieval was unavailable.

\item Zariski, O. (1944). The theorem of Bertini on the variable singular points of a linear system of varieties. \emph{Transactions of the American Mathematical Society, 56}, 130--140. Bibliography and the intrinsic argument are verified through Kleiman, §4 and reference 46. Direct original full-text access was not obtained, so no unverified internal theorem number is supplied.

\item Zariski, O. (1958a). On Castelnuovo's criterion of rationality $p_a=P_2=0$ of an algebraic surface. \emph{Illinois Journal of Mathematics, 2}(3), 303--315. \href{https://doi.org/10.1215/ijm/1255454536}{DOI: \nolinkurl{10.1215/ijm/1255454536}}. Its role in positive characteristic is documented by Hassett (2009, Remark 2.29); the full original proof was not independently consulted.

\item Zariski, O. (1958b). The problem of minimal models in the theory of algebraic surfaces. \emph{American Journal of Mathematics, 80}, 146--184. The bibliographic identification and role used here follow Hassett (2009).

\item Zariski, O. (1971). \emph{Algebraic surfaces} (2nd supplemented ed.). Springer. \href{https://doi.org/10.1007/978-3-642-92450-7}{DOI: \nolinkurl{10.1007/978-3-642-92450-7}}. Publisher metadata was verified, and the work is used as classical context, without a claim of fresh full-text inspection.

\item Zeuthen, H. G. (1871). \emph{\'Etudes g\'eom\'etriques etc.} \emph{Mathematische Annalen, IV}. Title abbreviated as cited on pp. 167 and 180; no unverified expansion is asserted here.

\item Zeuthen, H. G. (n.d.). {[}Work on invariance of numerical genus; exact title and date not supplied in the Chapter IV historical discussion.{]} Enriques's attribution at pp.~137--138 is reported without an independently certified original proof scope.
\end{itemize}
\end{sloppypar}
